\documentclass[10pt,a4paper]{report}

\usepackage{blindtext}

\usepackage[utf8]{inputenc}

\usepackage{amsmath}
\usepackage{amsfonts}
\usepackage{amssymb}
\usepackage{graphicx}
\usepackage{mathrsfs}
\author{Eric Yuta Otomo}

\setlength\topmargin{0in}
\setlength\headheight{0in}
\setlength\headsep{0in}
\setlength\textheight{8.9in} %8.5
\setlength\textwidth{6.5in}
\setlength\oddsidemargin{0in}
\setlength\evensidemargin{0in}

\renewcommand{\contentsname}{Conteúdo}

\begin{document}
\title{Meta-Teoria Principal}
\maketitle

\tableofcontents

\hspace{\baselineskip}

%\input{includeDia.tex}

\newpage

%%%%%%%%%%%%%%%%%%%%%%%%%%%%%%%%%%%%%%%%%%%%%%%%%%%%%%%%%%%%%%%%%%%%%%%%%%%%%%%%%%%%%%%%%%%%%%%%%%%%%%%%%%%%
%%%%%%%%%%%%%%%%%%%%%%%%%%%%%%%%%%%%%%%%%%%%%%%%%%%%%%%%%%%%%%%%%%%%%%%%%%%%%%%%%%%%%%%%%%%%%%%%%%%%%%%%%%%%
%%%%%%%%%%%%%%%%%%%%%%%%%%%%%%%%%%%%%%%%%%%%%%%%%%%%%%%%%%%%%%%%%%%%%%%%%%%%%%%%%%%%%%%%%%%%%%%%%%%%%%%%%%%%
%%%%%%%%%%%%%%%%%%%%%%%%%%%%%%%%%%%%%%%%%%%%%%%%%%%%%%%%%%%%%%%%%%%%%%%%%%%%%%%%%%%%%%%%%%%%%%%%%%%%%%%%%%%%
%%%%%%%%%%%%%%%%%%%%%%%%%%%%%%%%%%%%%%%%%%%%%%%%%%%%%%%%%%%%%%%%%%%%%%%%%%%%%%%%%%%%%%%%%%%%%%%%%%%%%%%%%%%%

\section*{Introdução}

%%%%%%%%%%%%%%%%%%%%%%%%%%%%%%%%%%%%%%%%%%%%%%%%%%%%%%%%%%%%%%%%%%%%%%%%%%%%%%%%%%%%%%%%%%%%%%%%%%%%%%%%%%%%
% As três teorias principais
\hspace{\baselineskip}

Este texto é uma coleção de três teorias distintas, a \textbf{Teoria de Conceitos Transcendentais}, a \textbf{Teoria de Projeções Dialéticas},  e a \textbf{Teoria Lógica Formal}. Todas estas são meta-teorias, são teorias cuja finalidade é auxiliar na criação de teorias. A Teoria de Conceitos Transcendentais busca investigar os conceitos transcendentais, conceitos que o ser humano tem consciência de sua existência mas é incapaz de definir com precisão sendo necessário um processo indutivo, a Teoria de Projeções Dialéticas é essencialmente um método de modelagem com base em estruturas chamadas de bases dialéticas, e por último a Teoria Lógica Formal que busca definir o método de teorização axiomática de forma ágil.

%%%%%%%%%%%%%%%%%%%%%%%%%%%%%%%%%%%%%%%%%%%%%%%%%%%%%%%%%%%%%%%%%%%%%%%%%%%%%%%%%%%%%%%%%%%%%%%%%%%%%%%%%%%%
% Por que o Nome é Teoria Principal

%%%%%%%%%%%%%%%%%%%%%%%%%%%%%%%%%%%%%%%%%%%%%%%%%%%%%%%%%%%%%%%%%%%%%%%%%%%%%%%%%%%%%%%%%%%%%%%%%%%%%%%%%%%%
% Por que esta ordem
\hspace{\baselineskip}

Apesar de cada filosofia poder ser utilizada de maneira independente, a ordem que cada teoria é exposta não é ao acaso, eu busco primeiro investigar os conceitos que são difíceis de definir, busco entender como tais conceitos estão relacionados usando modelos e posteriormente busco formalizar tais conceitos para construção de teorias axiomáticas.

%%%%%%%%%%%%%%%%%%%%%%%%%%%%%%%%%%%%%%%%%%%%%%%%%%%%%%%%%%%%%%%%%%%%%%%%%%%%%%%%%%%%%%%%%%%%%%%%%%%%%%%%%%%%
% Pretensão ideal
\hspace{\baselineskip}

Este texto tem como objetivo definir o método de criação de teorias e expor de uma maneira eficiente para um Ser Humano, este com vida limitada e portanto com tempo precioso, portanto não pretendo fazer uma investigação minuciosa, mas uma exposição ágil para que se torne uma linguagem comum entre as diferentes atividades intelectuais que o adepto desta filosofia possa ter. Investigações mais cuidadosas e mais refinadas já existem para a lógica, como a teoria de conjuntos, a lógica proposicional, mas seu conteúdo é extenso para finalidade de propósito geral, portanto fora do escopo deste documento.

%%%%%%%%%%%%%%%%%%%%%%%%%%%%%%%%%%%%%%%%%%%%%%%%%%%%%%%%%%%%%%%%%%%%%%%%%%%%%%%%%%%%%%%%%%%%%%%%%%%%%%%%%%%%
% 
\hspace{\baselineskip}

Não perca seu tempo relendo este livro, tente reescreve-lo às suas palavras, use os títulos e subtítulos como uma estrutura. A verdade é esta, a estrutura desta teoria é bem definida, mas a forma do discurso não. A filosofia dá forma a algo que pode ter múltiplas formas, portanto a maneira como ela é expressa não é a essência, e sim a capacidade de expressa-la de diferentes formas, sob diferentes situações. Sedimentar seu conteúdo é sedimentar sua estrutura e ser capaz de expor seu conteúdo de formas diferentes.

%%%%%%%%%%%%%%%%%%%%%%%%%%%%%%%%%%%%%%%%%%%%%%%%%%%%%%%%%%%%%%%%%%%%%%%%%%%%%%%%%%%%%%%%%%%%%%%%%%%%%%%%%%%%
%%%%%%%%%%%%%%%%%%%%%%%%%%%%%%%%%%%%%%%%%%%%%%%%%%%%%%%%%%%%%%%%%%%%%%%%%%%%%%%%%%%%%%%%%%%%%%%%%%%%%%%%%%%%
%%%%%%%%%%%%%%%%%%%%%%%%%%%%%%%%%%%%%%%%%%%%%%%%%%%%%%%%%%%%%%%%%%%%%%%%%%%%%%%%%%%%%%%%%%%%%%%%%%%%%%%%%%%%
%%%%%%%%%%%%%%%%%%%%%%%%%%%%%%%%%%%%%%%%%%%%%%%%%%%%%%%%%%%%%%%%%%%%%%%%%%%%%%%%%%%%%%%%%%%%%%%%%%%%%%%%%%%%
%%%%%%%%%%%%%%%%%%%%%%%%%%%%%%%%%%%%%%%%%%%%%%%%%%%%%%%%%%%%%%%%%%%%%%%%%%%%%%%%%%%%%%%%%%%%%%%%%%%%%%%%%%%%

\chapter{Teoria de Conceitos Transcendentais}

%%%%%%%%%%%%%%%%%%%%%%%%%%%%%%%%%%%%%%%%%%%%%%%%%%%%%%%%%%%%%%%%%%%%%%%%%%%%%%%%%%%%%%%%%%%%%%%%%%%%%%%%%%%%
% O que é ser transcendental ?
\hspace{\baselineskip}

O recurso mais elementar, talvez o único que faça sentido para o início de qualquer investigação intelectual é a linguagem. Não poderia utilizar recursos mais avançados sem antes introduzi-los por meio da linguagem, portanto é natural que o começo de minha investigação seja através de um discurso. Seria possível utilizar recursos mais elementares, por exemplo o desenho, a geometria, entretanto o desenho pode ser ambíguo, a linguagem consegue definir conceitos com suficiente rigor mesmo fazendo uso de outros conceitos indefinidos.

%%%%%%%%%%%%%%%%%%%%%%%%%%%%%%%%%%%%%%%%%%%%%%%%%%%%%%%%%%%%%%%%%%%%%%%%%%%%%%%%%%%%%%%%%%%%%%%%%%%%%%%%%%%%
%%%%%%%%%%%%%%%%%%%%%%%%%%%%%%%%%%%%%%%%%%%%%%%%%%%%%%%%%%%%%%%%%%%%%%%%%%%%%%%%%%%%%%%%%%%%%%%%%%%%%%%%%%%%
%%%%%%%%%%%%%%%%%%%%%%%%%%%%%%%%%%%%%%%%%%%%%%%%%%%%%%%%%%%%%%%%%%%%%%%%%%%%%%%%%%%%%%%%%%%%%%%%%%%%%%%%%%%%

\section{A Diferença}

\hspace{\baselineskip}

Sentimos a diferença, ela é uma sensação, e como sensação é difícil de explicar, apenas acreditamos nesta capacidade, sem a percepção da diferença todos os outros conceitos não fariam o menor sentido. Sentimos as cores azul e vermelho (pelo menos a grande maioria dos seres humanos), mas como explicar a diferença de cores para um cego de nascença, poderíamos explicar fazendo uma analogia com texturas diferentes, mas essa explicação é insuficiente, falta informação, nunca irá abranger a sensação do porquê vemos o azul como azul ou o vermelho como vermelho, apenas sentimos, sabemos que sentimos da maneira que sentimos, mas nem sabemos se sentimos da mesma maneira. 

%Vamos exemplificar o que eu disse através da sensação de cores diferentes, 

\hspace{\baselineskip}

A diferença é um destes conceitos que também não dá para explicar direito, apenas sentimos, entretanto mesmo que não consigamos explicar, podemos aceitar sua existência, como se fosse algo inerente de nós mesmos. A percepção da diferença é transcendental neste sentido, a percepção da diferença é uma capacidade a priori do ser humano. Não dá para fazer absolutamente nada sem o conceito de diferença e sem a capacidade de perceber diferenças, sem a diferença nem dá para dizer que existimos. A diferença é essência da dialética que será discutida na seção posterior. % cabe a filosofia das projeções dialéticas definir dialética, nesta etapa vou me ater a comentar sem muito rigor.

\hspace{\baselineskip}

Somente é possível definir conceitos por causa da habilidade de perceber diferenças, toda discussão intelectual requer assumir essa capacidade.% Há uma profundidade muito maior quando se diz em respeitar diferenças.

%%%%%%%%%%%%%%%%%%%%%%%%%%%%%%%%%%%%%%%%%%%%%%%%%%%%%%%%%%%%%%%%%%%%%%%%%%%%%%%%%%%%%%%%%%%%%%%%%%%%%%%%%%%%
%%%%%%%%%%%%%%%%%%%%%%%%%%%%%%%%%%%%%%%%%%%%%%%%%%%%%%%%%%%%%%%%%%%%%%%%%%%%%%%%%%%%%%%%%%%%%%%%%%%%%%%%%%%%
%%%%%%%%%%%%%%%%%%%%%%%%%%%%%%%%%%%%%%%%%%%%%%%%%%%%%%%%%%%%%%%%%%%%%%%%%%%%%%%%%%%%%%%%%%%%%%%%%%%%%%%%%%%%

\section{Realidade e Representação}

\hspace{\baselineskip}

%%%%%%%%%%%%%%%%%%%%%%%%%%%%%%%%%%%%%%%%%%%%%%%%%%%%%%%%%%%%%%%%%%%%%%%%%%%%%%%%%%%%%%%%%%%%%%%%%%%%%%%%%%%%
Uma diferença elementar para qualquer atividade intelectual é a percepção da realidade com a representação da realidade, o conceito de realidade pode parecer razoável para muitos, mas a realidade em si é transcendental para o ser humano, nossos sentidos foram bem "construídos" mas não são perfeitos. Essa percepção falha da realidade pode ser explorada através das ilusões, porém irei convencer por algo bem mais simples, o ser humano erra, o ser humano se confunde, o erro é uma pequena ilusão que nós cometemos ao acreditar que estamos certos, mas ao acreditar sentimos como se fosse real. Portanto o que nos referimos como realidade na verdade é uma representação mais fiel possível para o ser humano de algo transcendental a nós, este sim chamado de realidade.

Apesar de ser assimilado de maneira transcendental, este conceito não é natural, um ser humano que não tenha vivenciado e aceitado a limitação humana em sentir a realidade não é capaz de compreender este conceito, e tais experiências não se limitam apenas a ilusões, mas dependem também do desenvolvimento cultural científico, pois é a ciência capaz de fornecer instrumentos que permitem uma observação mais precisa da realidade ao mesmo tempo que escancara a limitação humana.\\

Uma representação é uma tentativa de simplificar a compreensão da realidade usando outros recursos. Representar é reconhecer a limitação humana sobre a percepção da realidade, uma boa teoria deve levar em conta uma boa notação para representar suas ideias. Entretanto devemos nos ater a não adicionar informações a realidade, uma boa teoria deve tentar representar a realidade e ser fiel a ela, tentar modificar de qualquer maneira é criar uma obra de ficção, porém ao longo deste capítulo busco convencer que conceitos transcendentais podem ser interpretados como adições a realidade, e tais conceitos são necessários, eles refletem uma das limitações humanas na compreensão da realidade, portanto a adição de tais conceitos devem ser feitos com cautela, devemos nos ater a dar pulinhos de gatos.\\

É importante comentar que aqui temos um paradoxo, enquanto a lógica se atém a não mudar a realidade, ela está considerando a realidade de maneira estática e que qualquer raciocínio posterior a análise não irá afetar esta realidade que foi congelada no tempo, ela funciona bem na matemática, pois medidas são estáticas, conjuntos são feitos para serem estáticos, quantidade são estáticas quando medidas, porém o homem se insere na realidade, então se considerarmos o fluxo do tempo, o futuro, não se pode dizer que uma informação não irá afetar a humanidade, pois neste contexto a informação é parte da realidade, uma teoria afeta nosso comportamento, desinformação afeta nosso comportamento. Considerar o futuro como parte da realidade é transcendental a considerar a realidade como seções estáticas.





%%%%%%%%%%%%%%%%%%%%%%%%%%%%%%%%%%%%%%%%%%%%%%%%%%%%%%%%%%%%%%%%%%%%%%%%%%%%%%%%%%%%%%%%%%%%%%%%%%%%%%%%%%%%
%%%%%%%%%%%%%%%%%%%%%%%%%%%%%%%%%%%%%%%%%%%%%%%%%%%%%%%%%%%%%%%%%%%%%%%%%%%%%%%%%%%%%%%%%%%%%%%%%%%%%%%%%%%%
%%%%%%%%%%%%%%%%%%%%%%%%%%%%%%%%%%%%%%%%%%%%%%%%%%%%%%%%%%%%%%%%%%%%%%%%%%%%%%%%%%%%%%%%%%%%%%%%%%%%%%%%%%%%

\section{Sentidos e Memória}

%%%%%%%%%%%%%%%%%%%%%%%%%%%%%%%%%%%%%%%%%%%%%%%%%%%%%%%%%%%%%%%%%%%%%%%%%%%%%%%%%%%%%%%%%%%%%%%%%%%%%%%%%%%%
% Sentidos e Memória ~ Espaço e Tempo
\hspace{\baselineskip}

O indutor capaz de definir todo conceito é a experiência da realidade, a sua percepção e a sua lembrança através da memória, a linguagem é apenas um recurso para mimetizar tais experiências. São os sentidos que induzem todos os conceitos, portanto para entender os conceitos transcendentais é interessante investigar quais conceitos são naturalmente induzidos através dos sentidos. 

%Antes vamos enumerar os sentidos, tais que considero importantes para percepção da realidade. A visão que fornece de forma mais precisa as informações da realidade e é responsável pela noção de simultaneidade do espaço, a audição que necessariamente precisa ser interpretada no tempo mas fornece também informações espaciais, o tato que fornece informações espaciais mas necessita de proximidade. Irei desconsiderar o paladar e o olfato no presente momento. A visão permite que seja possível enxergar limites, contornos e também criar a concepção de vazio, ela permite criar a noção de simultaneidade, a noção de espaço, a noção de objeto.

A visão fornece de forma mais precisa as informações da realidade e é responsável pela noção de simultaneidade do espaço, permite que seja possível enxergar limites, contornos e também criar a concepção de vazio, ela permite criar a noção de simultaneidade, a noção de espaço, a noção de objeto. É por enxergar objetos a distância que a representação pontual de objetos faz sentido. Estes são conceitos próximos ao conceito de espaço. 

A noção de tempo "somente" pode ser entendida e induzida com a memória, o fluxo de tempo depende da percepção da realidade, da capacidade de lembrar de eventos passados e da capacidade de imaginar eventos futuros. A percepção do tempo também induz o conceito de infinito através de processos.

%%%%%%%%%%%%%%%%%%%%%%%%%%%%%%%%%%%%%%%%%%%%%%%%%%%%%%%%%%%%%%%%%%%%%%%%%%%%%%%%%%%%%%%%%%%%%%%%%%%%%%%%%%%%
%%%%%%%%%%%%%%%%%%%%%%%%%%%%%%%%%%%%%%%%%%%%%%%%%%%%%%%%%%%%%%%%%%%%%%%%%%%%%%%%%%%%%%%%%%%%%%%%%%%%%%%%%%%%
%%%%%%%%%%%%%%%%%%%%%%%%%%%%%%%%%%%%%%%%%%%%%%%%%%%%%%%%%%%%%%%%%%%%%%%%%%%%%%%%%%%%%%%%%%%%%%%%%%%%%%%%%%%%

\section{A liberdade}

\hspace{\baselineskip}

A liberdade também é algo a priori, e portanto transcendental, sem o conceito de liberdade seria difícil entender a nossa capacidade de agir. A mesma define esta capacidade, ação e liberdade estão juntas como se estivessem em uma estrutura, não se pode agir sem liberdade e não se pode ter liberdade sem poder agir. \\

Entretanto liberdade também induz ao conceito de perceber a diferença entre o que se pode e o que não se pode fazer, falar em liberdade pode ser considerado como definir o que é possível fazer do que é errado fazer dentro de um contexto, e isso é importante para definição da lógica, cujo intuito é definir o que se pode e o que não se pode fazer para que possamos preservar a verdade em um raciocínio.\\

É papel da lógica formal definir as liberdades que permitem a preservação da verdade, entretanto vamos afrouxar este requisito para filosofia geral, a liberdade de raciocínio permite a preservação de algo, não necessariamente a verdade da lógica formal, um exemplo deste caso é a modelagem matemática usada na física, cujo o objetivo é achar um bom modelo mais do que garantir que cada etapa do raciocínio seja rigorosamente verdadeiro do ponto de vista matemático.\\

A liberdade dentro de um raciocínio define o que se pode ou não fazer, isso define um tipo de verdade diferente da realidade, é como estabelecer regras para um jogo, tais regras são artificiais e nada tem a ver com a realidade, mas dentro do contexto eles permitem separar o que faz ou não sentido dentro de um raciocínio.

%%%%%%%%%%%%%%%%%%%%%%%%%%%%%%%%%%%%%%%%%%%%%%%%%%%%%%%%%%%%%%%%%%%%%%%%%%%%%%%%%%%%%%%%%%%%%%%%%%%%%%%%%%%%
%%%%%%%%%%%%%%%%%%%%%%%%%%%%%%%%%%%%%%%%%%%%%%%%%%%%%%%%%%%%%%%%%%%%%%%%%%%%%%%%%%%%%%%%%%%%%%%%%%%%%%%%%%%%
%%%%%%%%%%%%%%%%%%%%%%%%%%%%%%%%%%%%%%%%%%%%%%%%%%%%%%%%%%%%%%%%%%%%%%%%%%%%%%%%%%%%%%%%%%%%%%%%%%%%%%%%%%%%

\section{Objeto e Identidade}

%%%%%%%%%%%%%%%%%%%%%%%%%%%%%%%%%%%%%%%%%%%%%%%%%%%%%%%%%%%%%%%%%%%%%%%%%%%%%%%%%%%%%%%%%%%%%%%%%%%%%%%%%%%%
% O lugar comum chamado objeto
\hspace{\baselineskip}

Um objeto é algo que possui identidade e possui limites que isolam o mesmo dos demais objetos. A coleção de experiências de um objeto cria um lugar comum na memória para o mesmo objeto, dando a sensação de identidade, o que permite representa-lo na linguagem através de um nome. Um objeto concreto, por exemplo um livro, é compreendido de forma dinâmica, é necessário ter consciência de sua relação com o espaço e de sua relação com o tempo. A consciência de um objeto é o lugar comum à sua percepção real, imaginações de eventos futuros e lembranças de eventos passados do objeto. 

%%%%%%%%%%%%%%%%%%%%%%%%%%%%%%%%%%%%%%%%%%%%%%%%%%%%%%%%%%%%%%%%%%%%%%%%%%%%%%%%%%%%%%%%%%%%%%%%%%%%%%%%%%%%
% Objeto ~ Abstração
\hspace{\baselineskip}

Para ser um objeto é necessário ter identidade, mas a consciência da identidade de um objeto não pode ser compreendida por uma entidade consciente se não tiver uma coleção de experiências que criem um lugar comum do objeto na memória. Este processo que acontece na nossa mente de criar lugares comuns para coleções de experiências é chamado de abstração (pode ser chamado, lembre-se de que no presente momento não há rigor nas definições), sentimos a existência de um objeto quando as coleções de experiências forem suficientes para criar um lugar comum, uma vez formado este lugar comum na consciência podemos nomear com um nome específico. Contudo, por meio da linguagem, podemos fazer o processo inverso, podemos ter um nome para algo desconhecido e posteriormente adquirir experiência que definam a entidade em questão.

%%%%%%%%%%%%%%%%%%%%%%%%%%%%%%%%%%%%%%%%%%%%%%%%%%%%%%%%%%%%%%%%%%%%%%%%%%%%%%%%%%%%%%%%%%%%%%%%%%%%%%%%%%%%
% Linguagem como substrato para explorar os lugares comuns
\hspace{\baselineskip}

A linguagem permite que nossa mente possa explorar estes lugares comuns com alguma precisão. Provavelmente nos tempos anteriores a linguagem nossa mente assim como a mente dos animais sentia a presença de objetos, entidades com identidade na realidade, entretanto tinham capacidade reduzida de reflexão e mesmo imaginação de situações futuras sobre estas entidades. Porém é uma tese que não pretendo investigar no momento.

%%%%%%%%%%%%%%%%%%%%%%%%%%%%%%%%%%%%%%%%%%%%%%%%%%%%%%%%%%%%%%%%%%%%%%%%%%%%%%%%%%%%%%%%%%%%%%%%%%%%%%%%%%%%
% Abstração como fundamento

\hspace{\baselineskip}

O processo de abstração, formação de lugares comuns na memória, criação de nomes, percepção de identidade, todos estes são fundamentais para o que chamamos de raciocínio, o nome que descreve o lugar como da nossa própria mente e da dinâmica das ideias. 

%%%%%%%%%%%%%%%%%%%%%%%%%%%%%%%%%%%%%%%%%%%%%%%%%%%%%%%%%%%%%%%%%%%%%%%%%%%%%%%%%%%%%%%%%%%%%%%%%%%%%%%%%%%%
% Identidade como conceito transcendental
\hspace{\baselineskip}

É importante perceber que o conceito de identidade é transcendental, nós sentimos que algo tem identidade, por isso nomeamos, a identidade é um conceito resultante do funcionamento de nossa consciência, pela formação destes lugares comuns. Por ser transcendental ele precisa estar a priori na nossa mente, toda forma definição apenas resgata algo que já está em nossa mente. A identidade está também relacionado a diferença, perceber a identidade de algo é perceber a diferença do mesmo em relação a outros objetos, portanto existe uma relação conceitual entre identidade e diferença.

%%%%%%%%%%%%%%%%%%%%%%%%%%%%%%%%%%%%%%%%%%%%%%%%%%%%%%%%%%%%%%%%%%%%%%%%%%%%%%%%%%%%%%%%%%%%%%%%%%%%%%%%%%%%
% Conceito de Classe e Processo de Abstração

%%%%%%%%%%%%%%%%%%%%%%%%%%%%%%%%%%%%%%%%%%%%%%%%%%%%%%%%%%%%%%%%%%%%%%%%%%%%%%%%%%%%%%%%%%%%%%%%%%%%%%%%%%%%
% identidade é um conceito transcendental, nós sentimos a identidade como lugar comum na nossa memória

%%%%%%%%%%%%%%%%%%%%%%%%%%%%%%%%%%%%%%%%%%%%%%%%%%%%%%%%%%%%%%%%%%%%%%%%%%%%%%%%%%%%%%%%%%%%%%%%%%%%%%%%%%%%
%%%%%%%%%%%%%%%%%%%%%%%%%%%%%%%%%%%%%%%%%%%%%%%%%%%%%%%%%%%%%%%%%%%%%%%%%%%%%%%%%%%%%%%%%%%%%%%%%%%%%%%%%%%%
%%%%%%%%%%%%%%%%%%%%%%%%%%%%%%%%%%%%%%%%%%%%%%%%%%%%%%%%%%%%%%%%%%%%%%%%%%%%%%%%%%%%%%%%%%%%%%%%%%%%%%%%%%%%

%\newpage

\section{Realidade e Estabilidade}

\hspace{\baselineskip}

A realidade diferente de um sonho é estável, ela permanece ou se altera de maneira mais lenta, estável, uniforme, ela é capaz de armazenar informações por conta desta estabilidade, em um sonho todas as percepções são voláteis, parecem reais mas não são tão estáveis quanto a realidade. Perceba que é uma ideia relativa, a realidade também pode ter fenômenos caóticos. \\

Essa estabilidade permite que possamos criar estes lugares comuns, criar conceitos, se o universo fosse absolutamente caótico nunca poderíamos ter a chance de ter consciência de nossa existência. Porém a estabilidade absoluta compartilha do mesmo mal, suponha que um ser humano fosse criado no escuro e não pudesse se mover, ou se fosse criado em um ambiente totalmente ilumidado e também não pudesse se mover, a estabilidade aqui não permitiria o conceito de diferença, requerido como disse anteriormente para o desenvolvimento da consciência. Neste sentido o caos absoluto e a estabilidade absoluta são conceitos correlatos, talvez da mesma maneira que o vazio e o infinito, estes também transcendentais.\\

Estabilidade ainda é mais interessante pelo fato de estar relacionado com o conceito de identidade, sentimos a estabilidade e percebemos a identidade de um objeto, são conceitos que estão relacionados, também faz sentido dizer que algo somente possui identidade se ele tiver alguma estabilidade, o que aproxima o conceito de caos com o conceito de estabilidade. Portanto identidade e estabilidade estão relacionados, posteriormente veremos que na lógica formal todas as suas construções são consideradas absolutas, não é permitido que suas propriedades estejam mal definidas, o que abre o caminho para estender a lógica formal admitindo construções que possuem propriedades parcialmente estáveis.


%%%%%%%%%%%%%%%%%%%%%%%%%%%%%%%%%%%%%%%%%%%%%%%%%%%%%%%%%%%%%%%%%%%%%%%%%%%%%%%%%%%%%%%%%%%%%%%%%%%%%%%%%%%%
%%%%%%%%%%%%%%%%%%%%%%%%%%%%%%%%%%%%%%%%%%%%%%%%%%%%%%%%%%%%%%%%%%%%%%%%%%%%%%%%%%%%%%%%%%%%%%%%%%%%%%%%%%%%
%%%%%%%%%%%%%%%%%%%%%%%%%%%%%%%%%%%%%%%%%%%%%%%%%%%%%%%%%%%%%%%%%%%%%%%%%%%%%%%%%%%%%%%%%%%%%%%%%%%%%%%%%%%%



\section{A existência}

%%%%%%%%%%%%%%%%%%%%%%%%%%%%%%%%%%%%%%%%%%%%%%%%%%%%%%%%%%%%%%%%%%%%%%%%%%%%%%%%%%%%%%%%%%%%%%%%%%%%%%%%%%%%
% 
\hspace{\baselineskip}

Existência defino como um conceito similar ao conceito de objeto, na verdade, existência é um tipo de objeto, entretanto ele possui uma qualidade que é muito importante no raciocínio lógico, uma existência possui relação com pelo menos outra existência diferente dela, algo existe em um universo se tal tem capacidade de interferir, ou seja, é intelectualmente relevante para este universo.\\

Dizemos que algo efetivamente existe se sua presença é relevante para um determinado contexto, algo pode existir e não ser efetivo se ele tem capacidade de interferência desprezível no sistema, o que não significa que ele não exista ou que ele não pode interferir caso outros eventos permitam em tempo futuro. Existência é um conceito relativo que está assentado em algo chamado de relação de interferência e seu uso requer boa dose de arbitrariedade, portanto convém coloca-lo no hall de conceitos transcendentais pela sua inerente imprecisão e por seu caráter priorístico.\\

Naturalmente uma existência, assim como um objeto necessita de alguma estabilidade, pois deve ser reconhecido para ser nomeado, e portanto identificável. Explorar existências que possui instabilidades inerentes é muito interessante e é explorado por exemplo na teoria da evolução, em que as diferentes formas de vida, são de fato mal definidas se considerarmos que sua definição reside na informação genética, entretanto por sua estabilidade permite até certo ponto sua identificação e análise.

%%%%%%%%%%%%%%%%%%%%%%%%%%%%%%%%%%%%%%%%%%%%%%%%%%%%%%%%%%%%%%%%%%%%%%%%%%%%%%%%%%%%%%%%%%%%%%%%%%%%%%%%%%%%
%%%%%%%%%%%%%%%%%%%%%%%%%%%%%%%%%%%%%%%%%%%%%%%%%%%%%%%%%%%%%%%%%%%%%%%%%%%%%%%%%%%%%%%%%%%%%%%%%%%%%%%%%%%%
%%%%%%%%%%%%%%%%%%%%%%%%%%%%%%%%%%%%%%%%%%%%%%%%%%%%%%%%%%%%%%%%%%%%%%%%%%%%%%%%%%%%%%%%%%%%%%%%%%%%%%%%%%%%

\section{Nomes Particulares e Nomes Genéricos}

%%%%%%%%%%%%%%%%%%%%%%%%%%%%%%%%%%%%%%%%%%%%%%%%%%%%%%%%%%%%%%%%%%%%%%%%%%%%%%%%%%%%%%%%%%%%%%%%%%%%%%%%%%%%
% identidade e unicidade e substituição
\hspace{\baselineskip}

Nós sabemos que possuímos identidade, mas também sabemos que nossa identidade é única. O conceito de unicidade parece ser mais forte que o conceito de identidade, apesar de correlatos. A maneira que utilizo para estabelecer a diferença entre estes conceitos é o conceito de substituição, apesar deste conceito parecer bem mais complicado é algo que é natural para um ser humano. Algo é único se é insubstituível, algo que não é único pode ser substituível, a nossa consciência é insubstituível, não podemos substituir nossa consciência por outra, nem faz sentido, por outro lado alguém para conversar é substituível, podemos conversar com um amigo, podemos conversar com um parente, podemos conversar com qualquer outra pessoa.

%%%%%%%%%%%%%%%%%%%%%%%%%%%%%%%%%%%%%%%%%%%%%%%%%%%%%%%%%%%%%%%%%%%%%%%%%%%%%%%%%%%%%%%%%%%%%%%%%%%%%%%%%%%%
% nome para coisas únicas e nome para coisas genéricas ~ o genérico pode ser substituído
\hspace{\baselineskip}

Existem nomes para coisas com unicidade e identidade e existem nomes para coisas que podem ter identidade mas não são únicas, um ser humano é único, o eu é único e possui identidade, entretanto a palavra (Ser Humano) designa o lugar comum entre todas as pessoas, o conceito designado pela expressão (Ser humano) expressa o lugar comum entre várias possíveis entidades. Por exemplo na expressão (o ser humano tem dois pés e duas mãos) podemos substituir Ser Humano por (eu possuo dois pés e duas mãos), ou podemos substituir por (seu amigo possui dois pés e duas mãos), Ser Humano é um nome genérico, ele pode representar uma coleção de entidades.

%%%%%%%%%%%%%%%%%%%%%%%%%%%%%%%%%%%%%%%%%%%%%%%%%%%%%%%%%%%%%%%%%%%%%%%%%%%%%%%%%%%%%%%%%%%%%%%%%%%%%%%%%%%%
% Inferência Primitiva pela Relação entre Genérico e Particular, indução e dedução
\hspace{\baselineskip}

A relação entre nomes particulares e genéricos induz um mecanismo de inferência, um raciocínio, sempre que em um determinado contexto um nome genérico represente uma afirmação universal a respeito de todos os particulares que o genérico representa então podemos substituir este nome genérico por seu particular, gerando uma outra informação "verdadeira" se a primeira for, pois a "verdade" da segunda está na primeira, entretanto essa regra de substituição não é válida caso queiramos substituir um particular por um genérico, pois a "verdade" do particular não contém a "verdade" de um genérico. Também não é verdade a regra de substituição quando dentro da informação o nome genérico não represente uma afirmação universal. Por exemplo quando dizemos "Um veículo é uma máquina" podemos substituir "veículo" por carro, pois na expressão "veículo" se refere de forma universal, entretanto não podemos substituir "máquina" por "furadeira", pois "máquina" dentro deste contexto não se refere a uma afirmação universal. As regras de substituição ficarão bem definidas quando os operadores universais ($\forall$, para todo) e ($\exists$, Existe) forem definidos mais adiante, entretanto convém que seja introduzido fórmulas antes de discutir substituições.\\

Chamamos de dedução a substituição de nomes genéricos por particulares e indução a substituição de nomes particulares por genéricos, a primeira uma regra de inferência, a segunda um mecanismo de teorização. É importante frisar que em um discurso o genérico é tradado como se fosse único que sua identidade é desconhecida, um nome genérico não é sinônimo de conjunto ou coleção, pois se comporta como se fosse único.

\hspace{\baselineskip}

Por hora caracterizamos conceitos transcendentais como todos os conceitos e informações gerados pelo processo indutivo, seja lá o que isso for. Posteriormente com mais recursos construídos poderemos sair de uma definição ingênua para uma definição mais explícita do que é transcendental.

%\newpage

%%%%%%%%%%%%%%%%%%%%%%%%%%%%%%%%%%%%%%%%%%%%%%%%%%%%%%%%%%%%%%%%%%%%%%%%%%%%%%%%%%%%%%%%%%%%%%%%%%%%%%%%%%%%
% Confusão entre conceito e instância e nome
\hspace{\baselineskip} 

Durante este discurso pode ficar confuso quando quisermos nos referir a um conceito e quando quisermos nos referir ao próprio objeto genérico que quisermos representar, por isso irei usar uma notação, quando eu quiser me referir ao conceito, ou seja, a um nome para algo real ou transcendental eu coloco um asterisco, por exemplo: Todo objeto real existe, mas objeto* é transcendental. Essa notação foi roubada da computação, variáveis com asteriscos são chamados de ponteiros, para o leitor que desejar uma reflexão sobre o porquê da escolha desta notação. \\

Ainda mais, talvez seja necessário diferenciar entre o nome (sequência de caracteres de uma palavra ou expressão) do conceito e de sua referência, uso a notação objeto$^\circ$, para fazer referência as letras (o,b,j,e,t,o) ligadas para formar a palavra e não ao conceito, tão pouco a um determinado objeto. Colocar ($\circ$) é querer se referir a representação em si, não ao conceito nem a um particular.\\

%%%%%%%%%%%%%%%%%%%%%%%%%%%%%%%%%%%%%%%%%%%%%%%%%%%%%%%%%%%%%%%%%%%%%%%%%%%%%%%%%%%%%%%%%%%%%%%%%%%%%%%%%%%%
% 
Podemos usar esta notação para explicitar quem pode ser substituído em uma inferência dedutiva. Por exemplo (Veículo$^\circ$ é um nome* para uma máquina*, veículo é uma máquina*), podemos substituir para (veículo$^\circ$ é um nome* para uma máquina*, carro é uma máquina*), pois na substituição se faz referência a um genérico, enquanto em qualquer outra ocorrência não podemos substituir, pois (máquina*) faz referência a um conceito, que é algo único na realidade e (veículo$^\circ$) faz referência a um nome que também designa algo único bem definido em uma linguagem. Em uma linguagem sabemos diferenciar a quê se refere uma palavra, pois existem mecanismos definidos na sintaxe, por isso em uma linguagem não necessitamos de notações diferentes para tal finalidade.


%%%%%%%%%%%%%%%%%%%%%%%%%%%%%%%%%%%%%%%%%%%%%%%%%%%%%%%%%%%%%%%%%%%%%%%%%%%%%%%%%%%%%%%%%%%%%%%%%%%%%%%%%%%%
%%%%%%%%%%%%%%%%%%%%%%%%%%%%%%%%%%%%%%%%%%%%%%%%%%%%%%%%%%%%%%%%%%%%%%%%%%%%%%%%%%%%%%%%%%%%%%%%%%%%%%%%%%%%
%%%%%%%%%%%%%%%%%%%%%%%%%%%%%%%%%%%%%%%%%%%%%%%%%%%%%%%%%%%%%%%%%%%%%%%%%%%%%%%%%%%%%%%%%%%%%%%%%%%%%%%%%%%%

%\section{Classes e Conjuntos}

%%%%%%%%%%%%%%%%%%%%%%%%%%%%%%%%%%%%%%%%%%%%%%%%%%%%%%%%%%%%%%%%%%%%%%%%%%%%%%%%%%%%%%%%%%%%%%%%%%%%%%%%%%%%
% 


%%%%%%%%%%%%%%%%%%%%%%%%%%%%%%%%%%%%%%%%%%%%%%%%%%%%%%%%%%%%%%%%%%%%%%%%%%%%%%%%%%%%%%%%%%%%%%%%%%%%%%%%%%%%
%%%%%%%%%%%%%%%%%%%%%%%%%%%%%%%%%%%%%%%%%%%%%%%%%%%%%%%%%%%%%%%%%%%%%%%%%%%%%%%%%%%%%%%%%%%%%%%%%%%%%%%%%%%%
%%%%%%%%%%%%%%%%%%%%%%%%%%%%%%%%%%%%%%%%%%%%%%%%%%%%%%%%%%%%%%%%%%%%%%%%%%%%%%%%%%%%%%%%%%%%%%%%%%%%%%%%%%%%

\section{O objeto e o Processo}

%%%%%%%%%%%%%%%%%%%%%%%%%%%%%%%%%%%%%%%%%%%%%%%%%%%%%%%%%%%%%%%%%%%%%%%%%%%%%%%%%%%%%%%%%%%%%%%%%%%%%%%%%%%%
% O objeto como o mais genérico dos lugares comuns, qualquer coisa é um objeto
\hspace{\baselineskip}

O objeto é o lugar comum mais genérico, ele pode representar qualquer coisa com identidade, qualquer coisa que possua contorno definido ou possa ser sentido como lugar comum de uma coleção de experiências. Eu poderia chama-lo de coisa, mas soaria deselegante, por isso uso o termo objeto, mas é importante notar que eu precisei usar um sinônimo para definir, o que a rigor não é definição alguma, este é o desafio da filosofia, definir sabendo que a linguagem é um recurso inadequado. O objeto é o mais genérico de todos os nomes, qualquer outro nome pode ser substituído em um processo de dedução.

%%%%%%%%%%%%%%%%%%%%%%%%%%%%%%%%%%%%%%%%%%%%%%%%%%%%%%%%%%%%%%%%%%%%%%%%%%%%%%%%%%%%%%%%%%%%%%%%%%%%%%%%%%%%
% O processo
\hspace{\baselineskip}

O conceito de objeto é algo estático, queremos também um conceito genérico que represente o fluxo de tempo, ele mesmo sendo também um objeto, mas que externalize uma natureza diferente, represente um lugar comum diferente. Chamo de processo este lugar comum, este objeto que representa evolução e fluxo de tempo, um processo conecta dois outros objetos de maneira anti simétrica.

% dentro de um contexto de processo podemos chamar os objetos que o processo conecta de eventos.

%%%%%%%%%%%%%%%%%%%%%%%%%%%%%%%%%%%%%%%%%%%%%%%%%%%%%%%%%%%%%%%%%%%%%%%%%%%%%%%%%%%%%%%%%%%%%%%%%%%%%%%%%%%%
% Com o processo e com o objeto temos a abstração da própria arquitetura cerebral
\hspace{\baselineskip}

É interessante notar que o Objeto e o Processo abstrai a própria arquitetura do nosso cérebro, muito provavelmente não é coincidência. Com este recurso em mãos podemos tentar entender o transcendental.


%%%%%%%%%%%%%%%%%%%%%%%%%%%%%%%%%%%%%%%%%%%%%%%%%%%%%%%%%%%%%%%%%%%%%%%%%%%%%%%%%%%%%%%%%%%%%%%%%%%%%%%%%%%%
%%%%%%%%%%%%%%%%%%%%%%%%%%%%%%%%%%%%%%%%%%%%%%%%%%%%%%%%%%%%%%%%%%%%%%%%%%%%%%%%%%%%%%%%%%%%%%%%%%%%%%%%%%%%
%%%%%%%%%%%%%%%%%%%%%%%%%%%%%%%%%%%%%%%%%%%%%%%%%%%%%%%%%%%%%%%%%%%%%%%%%%%%%%%%%%%%%%%%%%%%%%%%%%%%%%%%%%%%


\section{Abstração de Ponto e Seta}

%%%%%%%%%%%%%%%%%%%%%%%%%%%%%%%%%%%%%%%%%%%%%%%%%%%%%%%%%%%%%%%%%%%%%%%%%%%%%%%%%%%%%%%%%%%%%%%%%%%%%%%%%%%%
% Abstração de Ponto e Seta
\hspace{\baselineskip}

Visto os poucos recursos e a dificuldade inerente de caracterizar o transcendental irei utilizar a imaginação, o desenho para tentar construir o conceito de transcendentalidade. Para construir irei usar dois objetos da geometria: o ponto e a seta; o ponto é algo com dimensões desprezíveis que representa algo isolado, distinto dos demais, o ponto representa um objeto, e a seta representa a abstração de processo. Vamos literalmente construir conceitos através de desenhos.

$$ \bullet : \text{Objeto} $$
$$ \rightarrow : \text{Processo} $$
$$ \bullet \rightarrow \bullet : \text{Processo} : \text{Relação Antisimétrica} $$
$$ \bullet \sim \bullet : \text{Relação} : \text{Relação Simétrica} $$
$$ \bullet \sim \bullet \sim \bullet : \text{Topologia} $$

Uma relação generaliza um processo, abstraindo a direção, como se direção não fosse importante para conectar dois objetos, também pode representar dois processos, um de ida e outro de volta. O conceito de topologia é o conceito associado a processos intermediários, objetos intermediários, conceito associado a proximidade. O conceito de topologia será fundamental para caracterização dos conceitos transcendentais.

\hspace{\baselineskip}

Modéstia à parte, é elegante construir uma teoria a partir da abstração do cérebro, é o mesmo que dizer que para estudar as coisas inerentes de nossa consciência precisamos compreender a essência de nossa mente e a essência é obtida pela abstração. Não deve existir nada mais a priori para o intelecto do que a própria arquitetura do nosso cérebro.

%%%%%%%%%%%%%%%%%%%%%%%%%%%%%%%%%%%%%%%%%%%%%%%%%%%%%%%%%%%%%%%%%%%%%%%%%%%%%%%%%%%%%%%%%%%%%%%%%%%%%%%%%%%%
% Ambiente de Raciocínio
\hspace{\baselineskip}

Vamos agora criar regras ao representar ideias através destas formas abstratas (ponto e seta), irei chamar este conjunto de regras de abstração de ponto e seta.

%%%%%%%%%%%%%%%%%%%%%%%%%%%%%%%%%%%%%%%%%%%%%%%%%%%%%%%%%%%%%%%%%%%%%%%%%%%%%%%%%%%%%%%%%%%%%%%%%%%%%%%%%%%%
% Contexto
\hspace{\baselineskip}

\textbf{[ Contexto ]} Conjunto de regras em cima de um substrato que definem um ambiente separado de raciocínio, um ambiente separado de representação de ideias. Em um contexto não quero que haja regras de inferência, pois é interessante analisar com cautela.

%%%%%%%%%%%%%%%%%%%%%%%%%%%%%%%%%%%%%%%%%%%%%%%%%%%%%%%%%%%%%%%%%%%%%%%%%%%%%%%%%%%%%%%%%%%%%%%%%%%%%%%%%%%%
% Contexto Raciocínio de Redes
\hspace{\baselineskip}

\textbf{Contexto [ de Redes ]}

\begin{itemize}
	\item[ (i) ] \textbf{[ Substrato ]} O substrato é onde os objetos são desenhados, os pontos e setas podem ser desenhados em um papel em branco ou imaginados como se estivessem desenhados. A representação na realidade por meio de um desenho permite que possamos colocar muitos objetos e muitas setas, o que pode se tornar laborioso se for utilizado apenas a imaginação. O substrato em geral representa um contexto exterior ao desenho, por exemplo a realidade.
	\item[ (ii) ] \textbf{[ Ponto ]} Representa qualquer objeto com identidade distinto dos demais pontos em um mesmo desenho. Em relação aos demais pontos em um desenho o ponto é único, entretanto ele pode representar um nome genérico. Por exemplo podemos representar o conceito carro como um ponto e o conceito casa como outro ponto em um mesmo desenho, eles são conceitos distintos entre si, mas representam nomes genéricos, ou podemos representar meu carro e seu carro como dois pontos, eles são distintos entre si, mas representam objetos concretos.
	\item[ (iii) ] \textbf{[ Seta ]} Uma seta é um objeto e pode ligar um ponto e outro ponto desde que tais pontos estejam no desenho, sendo que uma seta também pode ser representada por um objeto, ou seja, ela possui identidade e pode ser nomeada. É assumido que tenha consciência do que uma seta rotacionada, portanto $A \rightarrow B$ é o mesmo que $B \leftarrow A$, esta afirmação em geral não é assumida em lógica formal, é preciso dizer explicitamente que possui tal propriedade.
	\item[ (iv)  ] \textbf{[ Identificação ]} Podemos nomear cada objeto, e o nome dentro de um desenho fica próximo do objeto seja um ponto ou uma seta. Se queremos dizer que dois objetos são iguais no sentido de representar o mesmo objeto fora do contexto do desenho então identificamos pelo mesmo nome, se quisermos representar objetos diferentes em um contexto exterior então devemos colocar nomes diferentes. Entretanto podemos colocar nomes diferentes em objetos iguais e isso não implica que são diferentes, mas no contexto pode ou não ser tratado como diferentes.
\end{itemize}

A linguagem possui um contexto, que é uma folha com uma sequência de caracteres organizado em linhas. Este é um exemplo que também expõe a dificuldade inerente da filosofia em tentar entender algo cuja única ferramenta para tal é ele próprio. Voltaremos a falar desta construção em lógica formal.\\

Uma teoria se constrói através de unidades chamadas de informação, podemos utilizar a abstração de ponto e seta para interpretar uma informação como sendo uma conexão ordenada entre dois pontos, um ponto que abstrai a premissa, outro ponto que abstrai a conclusão. Um raciocínio seria a viagem que a premissa estimulada pela memória e imaginação conduz naturalmente a uma conclusão.

$$ \texttt{Premissa} \rightarrow \texttt{Conclusão} $$
$$ \texttt{Causa} \rightarrow \texttt{Efeito} $$
$$ \texttt{Evento Inicial} \rightarrow \texttt{Evento Final} $$

%% Todos os processos intelectuais de raciocínio podem ser abstraídos e projetados no ponto e na seta ??

Suponha agora uma experiência que é possível reter na memória, por exemplo o evento de se sentar, vamos imaginar experiências similares entre si, podemos nos sentar no chão, podemos nos sentar em uma cadeira, podemos nos sentar em uma mesa, a ação sentar pode ser construída como um lugar comum entre as diferentes experiências que temos e que atribuímos o verbo sentar, isso exemplifica o processo de abstração e a construção de conceitos através de exemplos. Vamos considerar que cada experiência estimula nossa consciência como uma sensação, e que podemos diferenciar diferentes experiências por conta da capacidade de sentir diferenças, também vamos considerar que existe uma conexão natural entre experiências similares, quando sentimos algo na realidade a mente busca as experiências parecidas, se agora associarmos um nome as experiências parecidas, estamos criando uma abstração das experiências, esta abstração podemos chamar de evento.

\section{Construção do Conceito de Transcendentalidade}

%%%%%%%%%%%%%%%%%%%%%%%%%%%%%%%%%%%%%%%%%%%%%%%%%%%%%%%%%%%%%%%%%%%%%%%%%%%%%%%%%%%%%%%%%%%%%%%%%%%%%%%%%%%%
% Uso da abstração de ponto e seta para construção mental do conceito a priori
\hspace{\baselineskip}

Vamos agora utilizar esta abstração para construir os conceitos transcendentais. Vamos partir do conceito de topologia, o conceito de topologia é criado com pelo menos mais de dois objetos, com dois objetos temos o conceito de interação, ou relação, mas com três objetos ou mais temos o conceito de intermediário e proximidade que auxiliam na construção do conceito de topologia. Imagine agora que o (eu) uma consciência humana é um objeto e existe um outro objeto que está circundado por outros objetos e somente podemos acessar este objeto interior através dos objetos que o circundam, este objeto por ser inacessível é transcendental ao ser humano, precisamos dos outros objetos exteriores para tentar entender este objeto interior. Note que neste conceito de transcendentalidade ainda não foi necessário inserir o infinito.

%%%%%%%%%%%%%%%%%%%%%%%%%%%%%%%%%%%%%%%%%%%%%%%%%%%%%%%%%%%%%%%%%%%%%%%%%%%%%%%%%%%%%%%%%%%%%%%%%%%%%%%%%%%%
% Existência e Infinitude Compreensão , Abstração e Finitude
\hspace{\baselineskip}

Suponha agora que não fosse necessário para um conceito transcendental o infinito, mais especificamente, suponha que o objeto interior da discussão acima fosse completamente determinado pelo comportamento dos objetos exteriores e existissem apenas finitos destes objetos, ora podemos usar o conceito de existência para dizer que este objeto interior não é necessário em uma representação mental, bastando conhecer o comportamento dos objetos exteriores e assim podemos até duvidar de sua existência real (no sentido usual). Na prática isso não é verdade, conceitos construídos de maneira finita abstraem ideias, e abstrações simplificam raciocínios, portanto são úteis, e reconhecer a utilidade da abstração é reconhecer a limitação da consciência humana, entretanto não é seguro dizer que tais conceitos são transcendentais, por depender de finitos elementos.

%%%%%%%%%%%%%%%%%%%%%%%%%%%%%%%%%%%%%%%%%%%%%%%%%%%%%%%%%%%%%%%%%%%%%%%%%%%%%%%%%%%%%%%%%%%%%%%%%%%%%%%%%%%%
% A transcendentalidade de uma informação ou da definição
\hspace{\baselineskip}

Podemos construir a transcendentalidade destes objetos interiores mesmo que com objetos exteriores finitos. Para isso vamos usar a informação* como exemplo. Uma informação bem formulada pode ser vista como um objeto único que interage com outras informações, dizer que se conhece uma informação é essencialmente dizer que se sabe como esta informação se comporta em relação as outras informações ou saber explicitar sua definição, a primeira define o conhecimento* de uma informação de maneira transcendental, pois uma informação pode ter impacto imprevisível na mente humana e assim podem existir infinitas outras informações que são utilizadas ou interagem com a informação em questão, a segunda definição de conhecer é meramente saber enumerar suas propriedades, ou expressar uma formulação, esta segunda não garante que a informação possa ser utilizada na prática. Portanto inserindo o infinito o conhecer uma informação se torna transcendental.

%%%%%%%%%%%%%%%%%%%%%%%%%%%%%%%%%%%%%%%%%%%%%%%%%%%%%%%%%%%%%%%%%%%%%%%%%%%%%%%%%%%%%%%%%%%%%%%%%%%%%%%%%%%%
% construção de Transcendentalidade por Interação estável
\hspace{\baselineskip}

Podemos construir outro conceito de transcendentalidade a partir de dois objetos e uma relação, suponha dois objetos e eles possuem uma interação estável, ora uma interação é um conceito virtual, quando imaginamos dois objetos interagindo e formulamos a ideia de que existe uma relação entre eles estável e chamamos esta relação de informação, logo a informação é como se fosse a identificação de algo transcendental e este identificado pela própria relação. É um conceito similar ao conceito acima entretanto o infinito se expressa no caráter recorrente que experiências com estes dois objetos possam ter e não é construída com infinitos outros pontos.\\

A primeira transcendentalidade é dita construída pelo infinito espacialmente, a segunda é dita construída pelo infinito temporalmente. Não deve ser mera coincidência que existam estas duas interpretações, pois a consciência de espaço e tempo é elementar e portanto priorística para o intelecto humano. Estas estruturas recorrentes serão a base da formulação da teoria de projeções dialéticas, adiantando o assunto, a base espaço e tempo é uma base dialética do qual outros conceitos podem ser projetados nela para serem compreendidos.

%%%%%%%%%%%%%%%%%%%%%%%%%%%%%%%%%%%%%%%%%%%%%%%%%%%%%%%%%%%%%%%%%%%%%%%%%%%%%%%%%%%%%%%%%%%%%%%%%%%%%%%%%%%%
% O que é conhecer algo? Conhecer algo é fundamentalmente transcendental

%%%%%%%%%%%%%%%%%%%%%%%%%%%%%%%%%%%%%%%%%%%%%%%%%%%%%%%%%%%%%%%%%%%%%%%%%%%%%%%%%%%%%%%%%%%%%%%%%%%%%%%%%%%%
% Perfeição e Infinito e conceitos a priori garantia de perfeição e racicínio defensivo



%%%%%%%%%%%%%%%%%%%%%%%%%%%%%%%%%%%%%%%%%%%%%%%%%%%%%%%%%%%%%%%%%%%%%%%%%%%%%%%%%%%%%%%%%%%%%%%%%%%%%%%%%%%%
% Diferenciação de conceitos transcendentais a priori e conceitos transcendentais baseados no infinito
% Existe diferença ou se trata de mesmos conceitos?







%%%%%%%%%%%%%%%%%%%%%%%%%%%%%%%%%%%%%%%%%%%%%%%%%%%%%%%%%%%%%%%%%%%%%%%%%%%%%%%%%%%%%%%%%%%%%%%%%%%%%%%%%%%%
%%%%%%%%%%%%%%%%%%%%%%%%%%%%%%%%%%%%%%%%%%%%%%%%%%%%%%%%%%%%%%%%%%%%%%%%%%%%%%%%%%%%%%%%%%%%%%%%%%%%%%%%%%%%
%%%%%%%%%%%%%%%%%%%%%%%%%%%%%%%%%%%%%%%%%%%%%%%%%%%%%%%%%%%%%%%%%%%%%%%%%%%%%%%%%%%%%%%%%%%%%%%%%%%%%%%%%%%%

\section{Consciência de Objetos Finitos}

%%%%%%%%%%%%%%%%%%%%%%%%%%%%%%%%%%%%%%%%%%%%%%%%%%%%%%%%%%%%%%%%%%%%%%%%%%%%%%%%%%%%%%%%%%%%%%%%%%%%%%%%%%%%
% Imediato e Consciência de Objetos Finitos
\hspace{\baselineskip}

A consciência de objetos finitos cuja quantidade de objetos seja muito grande não é algo que temos consciência imediata, para formalizar esta ideia um tanto vaga vamos supor que ter consciência finita de objetos é a capacidade de imaginar tais objetos por meio de pontos sem que se perca a posição de tais objetos distintos, é imediato que a tentativa de imaginar 100 objetos por meio de pontos e nomeá-los é uma tarefa impossível para um ser humano, iremos perder a informação dos objetos na memória com bem menos que estes tantos objetos, portanto esta ideia parece fazer sentido se definido o que é ter consciência de objetos finitos.

%%%%%%%%%%%%%%%%%%%%%%%%%%%%%%%%%%%%%%%%%%%%%%%%%%%%%%%%%%%%%%%%%%%%%%%%%%%%%%%%%%%%%%%%%%%%%%%%%%%%%%%%%%%%
% Consciência de Objetos Finitos
\hspace{\baselineskip}

\textbf{[ Consciência Espacial Estática de Objetos Finitos ]} Temos consciência espacial estática de uma quantidade de objetos finitos se na imaginação pudermos representar os objetos por pontos nomeados e conseguirmos localizar cada objeto com seu devido nome sem perdemos tanto nomes quanto objetos sempre que nossa atenção se voltar para um ou outro objeto ou subgrupo de objetos.

%%%%%%%%%%%%%%%%%%%%%%%%%%%%%%%%%%%%%%%%%%%%%%%%%%%%%%%%%%%%%%%%%%%%%%%%%%%%%%%%%%%%%%%%%%%%%%%%%%%%%%%%%%%%
% Consciência Espacial Dinâmica de Objetos Finitos
\hspace{\baselineskip}

Temos claramente consciência espacial estática de dois objetos distintos, na verdade podemos até imaginar rotacionar a imaginação e mesmo assim não perdermos a consciência destes objetos. Temos consciência espacial estática de três objetos distintos, entretanto se rotacionarmos os objetos ou trocarmos de posição os objetos começamos a sentir certa dificuldade de resgatar a identidade de cada objeto (aqui podemos tentar convencer com uma brincadeira muito tradicional, pegue três copos que não sejam transparentes e um objeto que possa ser ocultado, a brincadeira é mostrar onde se escondeu o objeto uma vez embaralhado as posições dos objetos, claramente ter três copos torna o problema bem mais difícil de acompanhar que ter apenas dois copos), portanto ter consciência de três objetos sugere uma definição mais refinada deste conceito de consciência. 

%%%%%%%%%%%%%%%%%%%%%%%%%%%%%%%%%%%%%%%%%%%%%%%%%%%%%%%%%%%%%%%%%%%%%%%%%%%%%%%%%%%%%%%%%%%%%%%%%%%%%%%%%%%%
% Consciência Espacial Dinâmica de Objetos Finitos ~ Definição
\hspace{\baselineskip}

\textbf{[ Consciência Espacial Dinâmica de Objetos Finitos ]} É consciência espacial estática de uma quantidade de objetos e a capacidade de acompanhar tais objetos mesmo que cada objeto mude de posição à sua maneira, sem perder a identidade e a quantidade de objetos na imaginação.\\

Provavelmente a consciência (espacial dinâmica) de dois objetos é auxiliada pelo fato de que saber a posição relativa de um objeto nos torna capazes de inferir a posição relativa de outro objeto. Quando temos três objetos fica mais difícil de acompanhar pois a consciência de um objeto não permite inferir a posição relativa dos outros dois, entretanto a consciência estática ainda não é perdida, e é auxiliada pelas noções de cima e embaixo, frente e atrás, direito e esquerdo, estes todos frutos de pares. A última parte sugere que estruturas e noções já conhecidas auxiliam na consciência de objetos finitos, por exemplo podemos usar os membros do nosso corpo para termos consciência espacial de cinco objetos distintos, basta atrelarmos os cinco objetos a cada membro do nosso corpo (cabeça, mãos, pés), entretanto é necessário que cada membro represente efetivamente cada objeto em questão (ou seja, devemos pela imagem de cada membro resgatar o objeto).

%%%%%%%%%%%%%%%%%%%%%%%%%%%%%%%%%%%%%%%%%%%%%%%%%%%%%%%%%%%%%%%%%%%%%%%%%%%%%%%%%%%%%%%%%%%%%%%%%%%%%%%%%%%%
% Ganho que se obtem com a Analogia 
\hspace{\baselineskip}

\textbf{[ Analogia ]} É o processo de representar de forma efetiva determinado número de objetos fazendo uma conexão para cada objeto que queremos representar com cada objeto de alguma estrutura conhecida, esta estrutura chamada de analogia. No exemplo motivador acima a analogia é o corpo humano. \\

No processo de analogia podemos facilitar o trabalho de representação dos objetos usando analogias que representem coisas próximas, por exemplo, podemos ter a consciência dos membros de um animal fazendo a analogia com os membros de nosso próprio corpo, pois são semelhantes, entretanto representar números para cada membro do nosso corpo pode ser uma tarefa mais laboriosa, pois por não serem semelhantes, não terem estruturas comuns. Este conceito de analogia será fundamental para a filosofia de projeções dialéticas, definido no capítulo seguinte, ter consciência do ganho para o ser humano de ter acesso a boas analogias é essencial para entender o valor das projeções dialéticas, se considerarmos de forma utilitária o conhecimento.

%%%%%%%%%%%%%%%%%%%%%%%%%%%%%%%%%%%%%%%%%%%%%%%%%%%%%%%%%%%%%%%%%%%%%%%%%%%%%%%%%%%%%%%%%%%%%%%%%%%%%%%%%%%%
% Transcendência de Quantidades Grandes de Objetos
\hspace{\baselineskip}

Estas definições servem para ilustrar que a consciência de quantidades grandes de objetos é essencialmente fictícia, e portanto, não somente o infinito é transcendental para o ser humano, mas também quantidades grandes de números. A transcendência é o que se chama de pulo do gato que o ser humano faz diariamente ou o salto de fé (leap of faith). 

\hspace{\baselineskip}

Vamos assumir aqui que temos consciência de objetos finitos, mesmo que em grande quantidade, pois não precisamos imaginar os objetos, podemos desenha-los, este é o grande ganho que se obtém por representar fora da nossa mente uma ideia. A matemática depende do lápis e papel, ou de algum recurso similar para compreender os números. Portanto na nossa definição de transcendental vamos desconsiderar objetos finitos, entretanto devemos ter em mente que a consciência de grandes números transcende a consciência de dois objetos.

%%%%%%%%%%%%%%%%%%%%%%%%%%%%%%%%%%%%%%%%%%%%%%%%%%%%%%%%%%%%%%%%%%%%%%%%%%%%%%%%%%%%%%%%%%%%%%%%%%%%%%%%%%%%
%%%%%%%%%%%%%%%%%%%%%%%%%%%%%%%%%%%%%%%%%%%%%%%%%%%%%%%%%%%%%%%%%%%%%%%%%%%%%%%%%%%%%%%%%%%%%%%%%%%%%%%%%%%%
%%%%%%%%%%%%%%%%%%%%%%%%%%%%%%%%%%%%%%%%%%%%%%%%%%%%%%%%%%%%%%%%%%%%%%%%%%%%%%%%%%%%%%%%%%%%%%%%%%%%%%%%%%%%

\section{Indução Natural e Definição de Filosofia} 

%%%%%%%%%%%%%%%%%%%%%%%%%%%%%%%%%%%%%%%%%%%%%%%%%%%%%%%%%%%%%%%%%%%%%%%%%%%%%%%%%%%%%%%%%%%%%%%%%%%%%%%%%%%%
% 
\hspace{\baselineskip}

Para finalizar este capítulo irei definir indução natural como o menor passo transcendental, vulgo menor pulo de gato, conectando ideias em um processo construtivo de ideias. É papel da filosofia geral garantir que este passo seja suficientemente pequeno para convencer um ser humano das ideias ali contidas, diferente da lógica formal que utiliza a demonstração formal como guia entre as ideias abstratas, para filosofia geral este conceito de indução natural é o único mecanismo para garantir que o propósito da investigação tenha sido alcançado. Me atrevo a definir filosofia como a construção natural de ideias, embora esta ideia seja imprecisa, pois depende de cada ser humano julgar se algo é natural ou é obscuro do ponto de vista intelectual.\\

No próximo capítulo irei criar um mecanismo intelectual mais concreto para construção de ideias, entretanto tal mecanismo ainda conserva uma inerente imprecisão interiorizada através de conceitos transcendentais. Busco obviamente introduzir tais ideias da maneira que julgo suficientemente natural, fazendo uso dos conceitos aqui discutidos.

 

%%%%%%%%%%%%%%%%%%%%%%%%%%%%%%%%%%%%%%%%%%%%%%%%%%%%%%%%%%%%%%%%%%%%%%%%%%%%%%%%%%%%%%%%%%%%%%%%%%%%%%%%%%%%
%%%%%%%%%%%%%%%%%%%%%%%%%%%%%%%%%%%%%%%%%%%%%%%%%%%%%%%%%%%%%%%%%%%%%%%%%%%%%%%%%%%%%%%%%%%%%%%%%%%%%%%%%%%%
%%%%%%%%%%%%%%%%%%%%%%%%%%%%%%%%%%%%%%%%%%%%%%%%%%%%%%%%%%%%%%%%%%%%%%%%%%%%%%%%%%%%%%%%%%%%%%%%%%%%%%%%%%%%
%%%%%%%%%%%%%%%%%%%%%%%%%%%%%%%%%%%%%%%%%%%%%%%%%%%%%%%%%%%%%%%%%%%%%%%%%%%%%%%%%%%%%%%%%%%%%%%%%%%%%%%%%%%%
%%%%%%%%%%%%%%%%%%%%%%%%%%%%%%%%%%%%%%%%%%%%%%%%%%%%%%%%%%%%%%%%%%%%%%%%%%%%%%%%%%%%%%%%%%%%%%%%%%%%%%%%%%%%

\chapter{Teoria de Projeções Dialéticas}

%%%%%%%%%%%%%%%%%%%%%%%%%%%%%%%%%%%%%%%%%%%%%%%%%%%%%%%%%%%%%%%%%%%%%%%%%%%%%%%%%%%%%%%%%%%%%%%%%%%%%%%%%%%%
% Introdução
\hspace{\baselineskip}

A teoria de projeções dialéticas consiste de três elementos, o objeto que se quer estudar, uma entidade abstrata chamada de base dialética e um processo de raciocínio chamado de projeção dialética. A base dialética é uma estrutura que serve como uma analogia generalizada e projeção dialética é a ação de enxergar um objeto através desta analogia. Diferente de uma analogia convencional a base dialética é abstrata e seus conceitos são gerais. É importante lembrar que não se trata de uma ferramenta rigorosa, mas um método de modelagem, portanto passível de interpretações e arbitrariedade.\\

Antes convém que se comente a cerca do que se trata dialética socrática neste texto, não sendo exatamente algo que se encaixe como projeção dialética ela é uma ferramenta importante e muitas vezes suficiente para uma análise cuidadosa, assim é interessante ser o primeiro tópico neste capítulo.


%\section{Filosofia e Projeções Dialéticas}

\hspace{\baselineskip}

%\textbf{Questão [ Externalização Trialética ]} Por que o nome deste documento é Externalização Trialética? O que diabos vem a ser Externalização e o que diabos vem a ser Trialética?\\

%Caso você ainda não tenha percebido, o método trialético não é justificado cientificamente, tão pouco é uma teoria lógica formal, a única justificação que tem é com base na filosofia, o que é fraco para os padrões de objetividade atuais, mas paciência, porque é o que tem para hoje.

%Uma resposta rudimentar para a pergunta é que o documento é resultado da aplicação do método trialético, é uma consequência da aplicação do método trialético para uma determinada teoria (coleção de informações), é usando um jargão próprio uma externalização do método trialético.

\section{A Dialética}

\subsection{Dialética Socrática}

%%%%%%%%%%%%%%%%%%%%%%%%%%%%%%%%%%%%%%%%%%%%%%%%%%%%%%%%%%%%%%%%%%%%%%%%%%%%%%%%%%%%%%%%%%%%%%%%%%%%%%%%%%%%
% 
\hspace{\baselineskip}

%\textbf{Questão [ Dialética ]} O que é dialética? O que é trialética?\\

%\textbf{Paradigma [ Dialética Socrática Clássica ]} 

Sócrates foi um filósofo grego que criou um método para analisar problemas chamado de diálogo ou dialética que consiste de duas etapas tradicionalmente conhecidas como ironia e maiêutica. Existem alguns elementos que precisam estar presentes na prática da filosofia socrática, primeiro deve haver pelo menos duas pessoas, o inquiridor e o interlocutor, o primeiro tem como papel fazer questionamentos sobre determinado assunto próprio do interlocutor (que o interlocutor supostamente domina), Sócrates fazia o papel de inquiridor, o segundo é uma pessoa que acha que domina determinado assunto. O inquiridor deve fazer perguntas detalhadas o suficiente para que a segunda pessoa se declare ignorante a respeito do que ele sabia, este é o objetivo da ironia, desmascarar a ignorância, a segunda etapa da maiêutica é auxiliar de forma propositiva o interlocutor para que ele mesmo possa achar a resposta, mas a resposta deve vir do próprio interlocutor, Sócrates apenas auxilia a encontrar o caminho. A etapa da ironia é encapsulado na frase célebre atribuído a Sócrates: "Só sei que nada sei, e o fato de saber isso, me coloca em vantagem sobre aqueles que acham que sabem alguma coisa".\\

\textbf{Questão [ Arrogância x Ignorância ]} No sentido atribuído a Sócrates, qual a diferença entre ignorância e arrogância? \\

\textbf{Paradigma [ Dialética Socrática ]} Fazendo uma releitura da dialética Socrática, defino como Dialética Socrática o procedimento de enfrentar um problema através de dois processos, o primeiro processo é o processo de questionamento, que requer uma sensibilidade, uma capacidade de perceber inconsistências, de perceber insuficiências, de aumentar a resolução das perguntas para ser possível fazer questionamentos mais profundos, e o segundo processo é o processo de proposição, em que possíveis soluções são apresentadas para responder as questões, tais proposições não são definitivas, toda resposta tem uma pergunta inerente, a resposta é verdadeira ou não? Fazendo um paralelo temos que o questionamento seria a etapa de ironia, e a proposição seria a etapa de maiêutica, entretanto existem diferenças pontuais, primeiro não é necessário duas pessoas, segundo a divisão em duas etapas não implica que as duas etapas precisam estar separadas, mas sim que existam dois processos que podem ser chamados de forma arbitrária de acordo com a necessidade do problema.\\

\textbf{Paradigma [ Dividir para Conquistar ]} Também conhecido como dividir e governar, atribuído a história-política como uma estratégia de governo em dividir o poder em poderes menores subordinados para conseguir controlar um reino/império. Em ciência da computação é um paradigma de dividir um determinado problema em subproblemas menores e mais fáceis de resolver.\\

O diálogo socrático no tempo de Sócrates era de fato um diálogo, duas pessoas conversando verbalmente, e existe uma dificuldade inerente em um diálogo preservar a estruturação e o conteúdo da discussão e das perguntas, algo que é ruim se for o objetivo preservar todo conteúdo para uma análise posterior e também se o assunto em questão for muito complicado. Um problema consegue ser estruturado naturalmente em uma estrutura recursiva, uma árvore composta de uma pergunta principal vaga, sub-perguntas mais refinadas e possíveis respostas, mais sub-perguntas e assim por diante, portanto é uma aplicação do paradigma de dividir para conquistar. Unir dois paradigmas milenares é extremamente interessante. Esta estruturação recursiva chamo simplesmente de dialética ou caso requeira uma distinção de dialética socrática. No caso específico de lógica formal eu defino uma estruturação chamada de dialética demonstrativa que essencialmente é uma estrutura dialética, algo que será visto adiante no capítulo 3. \\

%\textbf{Questão [ Dialética Socrática e Dividir para Conquistar ]} Qual a relação entre dividir para conquistar e a dialética socrática?\\

%A dialética socrática também é um processo de decomposição de um problema em subproblemas, assim a aplicação do método socrático herda os mesmos benefícios do paradigma de Dividir para Conquistar, ele pode simplificar uma questão complicada em questões menores. Entretanto a dialética socrática não consiste em apenas simplificar, mas também em aumentar a resolução, ou seja, refinar os problemas encontrando mais problemas que não se tinha consciência de existirem. Dividir para Conquistar tem um propósito incutido, quero resolver o problema, quero encontrar a solução, enquanto a dialética socrática, pode ser usada para encontrar soluções ou para aumentar os problemas. Dialética Socrática se encaixa na academia enquanto que Dividir para Conquistar se encaixa na sociedade em geral. Do meu ponto de vista a Dialética Socrática estende o paradigma de Dividir para Conquistar.

Ter consciência do método socrático talvez seja o mais importante que você possa aprender neste capítulo, caso não tenha dificuldades, de posse do método socrático você já pode fazer muita coisa, apenas usando a linguagem. Caso você não esteja satisfeito com seu desempenho então é interessante refinar o seus métodos, utilizar outros paradigmas, que é o objetivo central deste capítulo.\\

Para resumir o conteúdo da dialética socrática irei representa-lo por meio de uma notação em redes.\\

$$ \texttt{Problema} \overset{\text{ação}}{\rightarrow} \{ \texttt{Questionar} , \texttt{Responder} \} $$ 
$$ \texttt{Problema} \overset{\text{questionar}}{\rightarrow} \texttt{Questão} $$ 
$$ \texttt{Problema} \overset{\text{responder}}{\rightarrow} \texttt{Proposição} $$ 
$$ \{ \texttt{Questão} , \texttt{Proposição} \} \overset{\text{generalização}}{\rightarrow} \texttt{Problema} $$

\hspace{\baselineskip}

A primeira abstração diz que para atacar um problema você tem duas operações, o ato de questionar e o ato de formular proposições. A última abstração diz que tanto a questão gerada quanto uma possível resposta (chamada aqui de proposição) são igualmente tratadas como novos problemas. Todas as bases dialéticas e a lógica formal servem como catalizadores para as ações de questionar e de responder, elas servem para permitir um maior refinamento das perguntas e respostas, portanto a dialética socrática é essencialmente a base de todos os paradigmas aqui investigados.

\subsection{Bases Dialéticas Elementares}

\hspace{\baselineskip}

%%%%%%%%%%%%%%%%%%%%%%%%%%%%%%%%%%%%%%%%%%%%%%%%%%%%%%%%%%%%%%%%%%%%%%%%%%%%%%%%%%%%%%%%%%%%%%%%%%%%%%%%%%%%
% 
%\textbf{Questão [ Dialética Materialista e Projeções Dialéticas ]} O que tem a ver a dialética materialista com a dialética socrática?\\

%Bem, este é um ponto interessante, porque a dialética materialista por si só foi capaz de gerar a forma mais terrível de governo com base no Ideário do Socialismo Marxista, e continua a dar muita dor de cabeça para a política atual. Mesmo com essa bagagem negativa, vejo que existe uma ideia interessante sob o ponto de vista filosófico, mas que é preciso aparar alguns aspectos para que não gere problemas para o futuro. 

%Os socialistas devem estar contorcendo o nariz com o que disse acima, eu peço que não olhe o socialismo como exclusividade do Marxismo, por exemplo o catolicismo, a igreja católica sempre foi uma instituição socialista, sempre buscou ajudar os outros, a se solidarizar com os outros, o Marxismo sequestrou o socialismo (ou pelo menos as pessoas que são atraídas pelo socialismo) e agora os socialistas sofrem de uma síndrome de estocolmo, tem empatia pelo seu sequestrador, eu sugiro que se você ficou em qualquer nível incomodado com minhas palavras que pense em criar um socialismo do zero, não busque o que não deu certo, se liberte.

%Voltando ao que interessa, a dialética materialista é como se fosse um modelo que analisa um determinado processo através da interação entre dois objetos, o objeto inicial chamado de tese, um segundo objeto chamado de antítese, e um terceiro elemento resultante que na verdade é a absorção da tese pela antítese ou da antítese pela tese, esse terceiro elemento é chamado de síntese. Podemos entender a dialética socrática através da dialética materialista, a tese seria uma afirmação, e a antítese seria uma afirmação contrária, na dialética materialista o conflito tem um papel importante, assim as duas afirmações contrárias gerariam um conflito, que é o questionamento, e como são opostos, um ou outro estaria certo, assim um ou outro seria absorvido ou eliminado para gerar a síntese que seria o ganhador do processo. A primeira falha da dialética materialista é sua mistificação, esse processo se tornou mistificado, como se fosse uma fórmula mágica, um segredo do universo, e isso fez com que o socialismo marxista se tornasse uma religião, a segunda falha é a ênfase no conflito, as interações duais não são sempre conflituosas, nem sempre há necessidade da destruição do outro, a terceira falha é que mesmo no conflito deve existir um propósito comum, um guia que unifique os processos, o objetivo não é a destruição, na dialética socrática é chegar em uma solução.

%Olhando de forma geral o que faz a dialética materialista, ela representa um processo por meio de componentes, na matemática existe um processo de representação de entidades matemáticas através de uma sequência de outros objetos, por exemplo uma função pode ser representada por uma soma de tipos especiais de funções multiplicadas por constantes, esses tipos especiais são chamados de base e as constantes associadas a cada elemento da base são chamados de componentes ou projeções. Unificando as ideias, na dialética materialista você busca representar um processo por meio da base (tese, síntese, antítese), por que não podemos ter bases com números arbitrários de elementos?\\

%\textbf{Paradigma [ Projeções Dialéticas ]} Uma base dialética é uma sequência de componentes (a,b,c,...) que são usadas para modelar algo, chamamos o processo de modelagem de projeção dialética, nela buscamos os componentes específicos de cada entrada da base. Por exemplo: Na ironia da dialética socrática sob a ótica da dialética materialista (base dialética materialista) o problema (questão) é visto como o conflito entre duas ideias contrárias, a primeira ideia seria correspondente a tese e a segunda ideia seria correspondente a antítese que geraria o conflito.\\

\textbf{[ Estrutura ]} Uma coleção de objetos que possuem significado entre si, ou seja, que existem propriedades entre estes objetos que caracterizam a estrutura. Dizemos que uma estrutura é preservada entre duas coleções se cada uma delas possuem estas mesmas propriedades, ou seja, o que caracteriza uma estrutura é seu significado e não os seus componentes em si. Uma estrutura conceitual é uma estrutura formada de conceitos. Toda estrutura será representada por meio de uma lista ordenada e um nome. \\

\textbf{[ Base Dialética ]} Uma coleção de componentes que formam uma estrutura.\\

Esta definição de base dialética é bem vaga, mas é interessante assim ser, pois a analogia é algo que se pode fazer com qualquer coisa, logo qualquer estrutura poderia ser uma boa analogia para algo.\\


\textbf{Relação [ de Analogia Estrutural ]} Para denotar que $A$ e $B$ são estruturas equivalentes e é possível estabelecer uma analogia usamos o símbolo ($\sim$), neste caso $A \sim B$. Para o caso de listas ordenadas dizer por exemplo que $(A,B) \sim (C,D)$ é o mesmo que dizer $A \sim C$ e $B \sim D$,  portanto a ordem importa. Lembre-se de que é uma definição imprecisa, portanto é interessante sempre vir acompanhada de uma justificativa, uma explicação do porquê existir tal relação.\\

Uma relação de analogia não é determinística, de acordo com a circunstância podemos considerar $(A,B) \sim (C,D)$ ou $(A,B) \sim (D,C)$, tudo depende do argumento usado para estabelecer a relação. Portanto estamos fora do escopo da lógica formal, mas no escopo da modelagem e da filosofia geral.

%%%%%%%%%%%%%%%%%%%%%%%%%%%%%%%%%%%%%%%%%%%%%%%%%%%%%%%%%%%%%%%%%%%%%%%%%%%%%%%%%%%%%%%%%%%%%%%%%%%%%%%%%%%%
% Resgate ao que foi dito na teoria transcendental de ponto e rede
\hspace{\baselineskip}

No capítulo anterior foi discutido sobre a abstração de ponto e seta, ou abstração de redes, e para entender o conceito de transcendentalidade foi utilizado duas abordagens, a abordagem dita espacial que é considerar o transcendental de conhecer uma informação como conhecer todas as informações que é possível se associar com a mesma, e a abordagem dita temporal, que é considerar o transcendental como a própria identidade de uma interação recíproca de pelo menos dois objetos. Podemos estabelecer que:

$$ (\texttt{Ponto}, \texttt{Seta}) \sim ( \texttt{Espaço} , \texttt{Tempo} ) $$

Apesar da relação de analogia ser simétrica é importante notar que o ponto (pontos) abstrai o conceito de espaço (existências e vazio) e a seta abstrai a ideia de processo e portanto tempo, existe uma ordem entre eles. Entretanto em nossa discussão consideramos (Espaço, Tempo) como uma base dialética, pois queremos estruturas conceituais como ferramentas de modelagem.

%%%%%%%%%%%%%%%%%%%%%%%%%%%%%%%%%%%%%%%%%%%%%%%%%%%%%%%%%%%%%%%%%%%%%%%%%%%%%%%%%%%%%%%%%%%%%%%%%%%%%%%%%%%%
% 
\hspace{\baselineskip}

\textbf{Base Dialética [ Espaço e Tempo ]} É a estrutura conceitual $(\texttt{Espaço}, \texttt{Tempo})$. Para simplificar a notação podemos chamar apenas de base $\texttt{E} \times \texttt{T}$. Ela é a primeira base conceitual estabelecida pois a consciência de espaço e tempo deve ser a priori, mais priorística que a consciência da linguagem.\\

A segunda base dialética mais fundamental é chamada de base existencial ou base topológica, ela é construída em cima do conceito de transcendentalidade via interpretação espacial, portanto pode se dizer derivada do espaço. Para construir sua representação e interpretar sua ideia imagine que todos os infinitos objetos exteriores formam uma linha contínua, e portanto são vistas como um contorno, para simplificar, um círculo no entorno do objeto, assim um objeto visto sob essa base topológica se constitui de um contorno chamado de forma externa e de uma identidade representada por um ponto em seu interior. \\

Como algo que possui uma forma externa só faz sentido se vista ou afetada por outra, então resgatamos o conceito transcendental de existência discutido no capítulo anterior para construir a dialética existencial.

%%%%%%%%%%%%%%%%%%%%%%%%%%%%%%%%%%%%%%%%%%%%%%%%%%%%%%%%%%%%%%%%%%%%%%%%%%%%%%%%%%%%%%%%%%%%%%%%%%%%%%%%%%%%
% Base Dialética Existencial
\hspace{\baselineskip}

\textbf{Base Dialética [ Dialética Existencial ]} A dialética existencial olha a realidade como uma coleção de coisas, que irei chamar de existências, cada existência possui dois componentes, um componente tangível chamado de forma externa e um componente intangível chamado de caráter interno. A acessibilidade é a propriedade que caracteriza essa estrutura, o caráter interno somente é acessível através da forma externa para nossa consciência (ou para um objeto distinto qualquer do universo). Para denotar a base existencial podemos escrever ($\texttt{Forma Externa}$,$\texttt{Caráter Interno}$) ou de forma simplificada como $\bigodot$. \\



% Indução 
A base dialética existencial é construída de maneira transcendental (processo infinito) por (ponto, seta), assim podemos dizer que a base dialética de espaço e tempo induz através da abstração de rede a base dialética existencial.

$$ \texttt{E} \times \texttt{T} \overset{ind.nat.}{\rightarrow} (\texttt{ponto,seta}) \overset{ind.nat.}{\rightarrow} \bigodot $$

% Colocar elementos da teoria transcendental e da abstração de seta e ponto

A dialética existencial modela o universo como um conjunto discreto (enumerável, contável) de objetos que chamamos de existências. Por conta disso a dialética existencial é elementar, ele está lidando com o que é mais elementar na consciência humana que é a representação do universo através de sua delimitação em partes menores distinguíveis entre si. A representação natural de uma existência é um ponto, ou uma bola, ou uma bola com um ponto dentro representando uma forma externa e possivelmente algo que não enxergamos como seu centro que é o caráter interno.

%\textbf{Estrutura [ Universo Existencial ]} (Universo ,  Existência)$_{U.E.}$ formam 

%\begin{itemize}
%	\item[ (i) ] \textbf{[ Existência ]} Uma existência existe no universo se ela é capaz de interagir/influenciar outra existência do universo.
%	\item[ (ii) ] \textbf{[ Universo ]} Podemos criar uma cadeia de relações de interação/influência com quaisquer duas outras existências de um Universo. Um universo é composto de existências.
%\end{itemize}

%\textbf{Observação :} Neste contexto, quando dizemos que algo não existe, significa que ele não tem capacidade de influenciar qualquer outra existência. E assim na linguagem comum é como se não existisse. Mas não significa que não exista !!! \\

%\textbf{Paradigma [ de Definição Estrutural ]} Observe que a definição de existência e universo não podem ser dissociadas, uma definição depende da outra e seu entendimento se passa por perceber que seu aprendizado não deve ser feito de forma linear (deve ser relido pelo menos outra vez). Quando um conceito/definição precisa ser construído (ou é vantajoso ser construído) de maneira não linear é interessante optar por uma construção estrutural, evidenciando que ela não é linear. \\


%\textbf{Paradigma [ Discretização ]} Este paradigma diz que precisamos discretizar a realidade em objetos finitos, distinguíveis para poder entender o universo, o contínuo é visto como uma coisa só, ou como uma coleção de coisas delimitadas por fronteiras finitas, não conseguimos trabalhar bem com o contínuo apenas usando a imaginação (por exemplo não conseguimos medir distâncias precisamente, não conseguimos medir velocidades precisamente, entretanto enxergamos a diferença, assim podemos dizer o que é maior do que é menor, o que é mais rápido do que é lento). Isso justifica a modelagem em bases, isso justifica a representação do universo em existências, isso também justifica a matemática para trabalhar com medidas.\\

%%%%%%%%%%%%%%%%%%%%%%%%%%%%%%%%%%%%%%%%%%%%%%%%%%%%%%%%%%%%%%%%%%%%%%%%%%%%%%%%%%%%%%%%%%%%%%%%%%%%%%%%%%%%%%%%%%%%%%%%%%%%%%%%%%%%%%%%%%%%%%%%%%%%%%%%%%%%%%%%%%%%%%%% Dialética Contrapositiva
\hspace{\baselineskip}

\textbf{Base Dialética [ de Contraste ]} A essência do contraste é a diferença, não faz muito sentido tentar entender coisas completamente diferentes, portanto o contraste entre dois objetos requer que eles tenham alguma relação em comum. Representamos então $A \times B$ quando $A$ e $B$ estão relacionados (são existenciais), porém queremos evidenciar sua diferença. Quando quisermos fazer referência a base em si denotamos por $\bigotimes$, por exemplo $\succ A \bullet \bigotimes$ ? para representar o problema: quero encontrar em $A$ elementos de contraste.\\

A base dialética de contraste está assentada em uma percepção elementar, que é a percepção da diferença, em termos de formulação de perguntas ela é interessante entretanto ela não é suficiente para resolver problemas, apenas para evidencia-los, o que não diminui seu valor, pois permite que teorias sejam refinadas e aumentem a resolução de suas perguntas.



%%%%%%%%%%%%%%%%%%%%%%%%%%%%%%%%%%%%%%%%%%%%%%%%%%%%%%%%%%%%%%%%%%%%%%%%%%%%%%%%%%%%%%%%%%%%%%%%%%%%%%%%%%%%%%%%%%%%%%%%%%%%%%%%%%%%%%%%%%%%%%%%%%%%%%%%%%%%%%%%%%%%%%
%%%%%%%%%%%%%%%%%%%%%%%%%%%%%%%%%%%%%%%%%%%%%%%%%%%%%%%%%%%%%%%%%%%%%%%%%%%%%%%%%%%%%%%%%%%%%%%%%%%%%%%%%%%%%%%%%%%%%%%%%%%%%%%%%%%%%%%%%%%%%%%%%%%%%%%%%%%%%%%%%%%%%%
%%%%%%%%%%%%%%%%%%%%%%%%%%%%%%%%%%%%%%%%%%%%%%%%%%%%%%%%%%%%%%%%%%%%%%%%%%%%%%%%%%%%%%%%%%%%%%%%%%%%%%%%%%%%%%%%%%%%%%%%%%%%%%%%%%%%%%%%%%%%%%%%%%%%%%%%%%%%%%%%%%%%%%

\subsection{Base Dialética e Projeção Dialética}

%%%%%%%%%%%%%%%%%%%%%%%%%%%%%%%%%%%%%%%%%%%%%%%%%%%%%%%%%%%%%%%%%%%%%%%%%%%%%%%%%%%%%%%%%%%%%%%%%%%%%%%%%%%%%%%%%%%%%%%%%%%%%%%%%%%%%%%%%%%%%%%%%%%%%%%%%%%%%%%%%%%%%%
% Base Dialética e Projeção
\hspace{\baselineskip}

Uma das coisas que de forma recorrente é utilizada na matemática é a representação de um determinado objeto através de componentes de um sistema de representação, tal que o objeto possua um determinado valor associado a um determinado componente. Chamamos os diferentes componentes de base, e a sequência ordenada dos valores em cada componente da base de representação do objeto. Em geral na matemática se utiliza sistemas em que todo objeto de um conjunto possui uma representação e que tal representação é única, vamos considerar casos menos rígidos, vamos considerar que nem todos possuem representação perfeita e que as representações podem não ser estritamente iguais (consideramos na dialética que o raciocínio não é estritamente determinístico, mas subjetivo e possivelmente arbitrário) sendo necessário constatar a relação e não apenas aplicar um operador.\\

\textbf{[ Projeção Dialética ]} Seja uma base dialética $ \beta $ e um objeto $A$, chamamos de projeção dialética de $A$ na base $\beta$ sua possível representação $A_{\beta}$ nesta base. A motivação de dar nome de projeção é implicitamente dizer que a representação pode não ser perfeita e podemos estar perdendo informação. Para codificar essa informação uso a notação: $ A \prec A_{\beta} \bullet \beta $ que significa ver $A$ através de $\beta$. Para dizer que quero achar a representação uso: $ \succ A \bullet \beta $ ?. O símbolo $\bullet$ representa a ideia de que dois objetos possuem uma interação, neste caso entre base e representação.\\

\textbf{Exemplos [ de Projeções Dialéticas ]}

\begin{itemize}
	\item[ (i) ] Lógica $\prec$ ( Teoria de Conjuntos, Lógica Proposicional ) $\bullet$ (Espaço, Tempo)
	\item[ (ii) ] Matemática $\prec$ (Cálculo Diferencial, Probabilidade) $\bullet$ (Espaço, Tempo)
	\item[ (iii) ] Programação $\prec$ (Orientação a Objeto, Programação Imperativa) $\bullet$ (Espaço, Tempo)
\end{itemize}

O símbolo $\bullet$ tem função similar a $\sim$ definido anteriormente para equivalência estrutural, entretanto $\bullet$ é usado no contexto de projeção (temos o símbolo $\succ$ ou $\prec$), portanto possuem propósitos diferentes.

%%%%%%%%%%%%%%%%%%%%%%%%%%%%%%%%%%%%%%%%%%%%%%%%%%%%%%%%%%%%%%%%%%%%%%%%%%%%%%%%%%%%%%%%%%%%%%%%%%%%%%%%%%%%%%%%%%%%%%%%%%%%%%%%%%%%%%%%%%%%%%%%%%%%%%%%%%%%%%%%%%%%%%
%%%%%%%%%%%%%%%%%%%%%%%%%%%%%%%%%%%%%%%%%%%%%%%%%%%%%%%%%%%%%%%%%%%%%%%%%%%%%%%%%%%%%%%%%%%%%%%%%%%%%%%%%%%%%%%%%%%%%%%%%%%%%%%%%%%%%%%%%%%%%%%%%%%%%%%%%%%%%%%%%%%%%%
%%%%%%%%%%%%%%%%%%%%%%%%%%%%%%%%%%%%%%%%%%%%%%%%%%%%%%%%%%%%%%%%%%%%%%%%%%%%%%%%%%%%%%%%%%%%%%%%%%%%%%%%%%%%%%%%%%%%%%%%%%%%%%%%%%%%%%%%%%%%%%%%%%%%%%%%%%%%%%%%%%%%%%

\subsection{Paradigmas de Análise Dialética}

\hspace{\baselineskip}

\textbf{Paradigma [ Primeiro Paradigma da Análise Dialética ]} Suponha que queiramos analisar uma existência, precisamos ter acesso a suas propriedades através de uma "forma externa" (forma externa em relação a mente) e que possivelmente sem ela não poderíamos concluir nada (seria o mesmo que não ter consciência dela), essa forma externa pode ser interpretada como uma projeção da existência em alguma base ou alguma estrutura que permite importar propriedades estabelecidas pela estrutura (sem a dialética poderia ser apenas uma imagem do conceito ou uma analogia). Assim o primeiro paradigma da Análise Dialética é que toda análise dialética de um determinado objeto deve começar com um nome para existência e achar projeções da mesma para que seja possível "manipula-la na mente". \\

\textbf{Paradigma [ Segundo Paradigma da Análise Dialética ]} Depois de ter estabelecido uma forma externa para seu objeto, agora é a hora de tentar entende-lo, assim este paradigma estabelece que você comece a tentar estabelecer relações com outros objetos e comece a elaborar perguntas com quaisquer meios disponíveis, usando quaisquer bases dialéticas definidas anteriormente e até mesmo mudando a representação do seu objeto problema. \\

\textbf{Paradigma [ Terceiro Paradigma da Análise Dialética ]} Depois de estabelecido os resultados do seu estudo, é hora de avaliar a presença de artefatos e discutir possíveis má interpretações, organizar e simplificar conteúdos e estabelecer conclusões.

%%%%%%%%%%%%%%%%%%%%%%%%%%%%%%%%%%%%%%%%%%%%%%%%%%%%%%%%%%%%%%%%%%%%%%%%%%%%%%%%%%%%%%%%%%%%%%%%%%%%%%%%%%%%%%%%%%%%%%%%%%%%%%%%%%%%%%%%%%%%%%%%%%%%%%%%%%%%%%%%%%%%%%
%%%%%%%%%%%%%%%%%%%%%%%%%%%%%%%%%%%%%%%%%%%%%%%%%%%%%%%%%%%%%%%%%%%%%%%%%%%%%%%%%%%%%%%%%%%%%%%%%%%%%%%%%%%%%%%%%%%%%%%%%%%%%%%%%%%%%%%%%%%%%%%%%%%%%%%%%%%%%%%%%%%%%%
%%%%%%%%%%%%%%%%%%%%%%%%%%%%%%%%%%%%%%%%%%%%%%%%%%%%%%%%%%%%%%%%%%%%%%%%%%%%%%%%%%%%%%%%%%%%%%%%%%%%%%%%%%%%%%%%%%%%%%%%%%%%%%%%%%%%%%%%%%%%%%%%%%%%%%%%%%%%%%%%%%%%%%

\subsection{Bases Trialéticas e Tetraéticas}

% Base Trialética
\hspace{\baselineskip}

\textbf{Base [ Trialética ]} A trialética tem relação com ter três componentes. Existem dois tipos e preservam uma relação análoga à abstração de redes e a dialética espaço e tempo, elas são definidas pelas estruturas conceituais (Início, Meio, Fim) $\sim$ (Emoção, Intelecto, Vontade). A primeira mais abstrata e geral e a segunda mais próxima do ser humano (trialética humana), sendo que a segunda pode ser considerada como a projeção da consciência humana na primeira. \\


De forma a simplificar seu uso a base trialética será denotada pelas letras gregas $\alpha \sim \texttt{Emoção}$, $\gamma \sim \texttt{Intelecto}$ e $\beta \sim \texttt{Vontade}$. Para auxiliar ainda mais é possível utilizar cores, por exemplo, imaginar o $\alpha$ como representado pela cor vermelha, $\beta$ representado pela cor verde e $\gamma$ pela cor azul, apesar de isso ser um pouco arbitrário. \\

É justo evidenciar que a filosofia da matemática está mais próxima do intelecto e da vontade, se projetarmos a filosofia nesta base poderíamos definir uma filosofia mais próxima da emoção que lida com questões existenciais da vida humana, filosofia intelectual que lida com a lógica e o raciocínio e a filosofia volitiva que lida com a vontade e a motivação das ações, a matemática apesar de ser uma restrição da filosofia por não lidar com questões existenciais diretamente não deixa de ser filosofia por lidar com lógica e com a busca por soluções de problemas. \\

\textbf{Exemplo [ Auto Exemplo ]}
$$ \texttt{TP} := \texttt{Teoria Principal} $$
$$ \texttt{TCT} := \texttt{Teoria de Conceitos Transcendentais} $$
$$ \texttt{TPD} := \texttt{Teoria de Projeções Dialéticas} $$
$$ \texttt{TLF} := \texttt{Teoria de Lógica Formal} $$
$$ \texttt{TP} \prec ( \texttt{TCT} ,  \texttt{TPD} , \texttt{TLF} ) \bullet (\alpha, \beta, \gamma) $$

\hspace{\baselineskip}

A própria organização deste documento segue uma projeção na trialética humana. Perceba também que na seção anterior os paradigmas de análise dialética podem ser projetados na trialética humana. Como existe uma íntima relação da trialética com a consciência é natural estar quase onipresente nas diferentes atividades intelectuais e culturais.

\hspace{\baselineskip}

A princípio não existe restrição na aplicação das bases dialéticas, assim podemos aplicar a base trialética em um dos componentes dela mesma, em específico ao intelecto.

\hspace{\baselineskip}

\textbf{[ Projeção Trialética do Intelecto ]} Temos três direções principais de aprendizado. Ter consciência destas três direções irá poupar seu tempo e tornar o seu estudo mais eficiente, é importante que se tenha bons métodos a sua disposição que otimizem o aprendizado para estes três caminhos, não se iluda com o aprendizado natural, não existe campo do intelecto humano que já não tenha superado a capacidade de aprendizagem natural do ser humano e você certamente tem outros projetos para sua vida.

\begin{itemize}
	\item[ (i) ] \textbf{[ Intelecto Alfa : Conhecimento ]} Esta direção de aprendizado está relacionado com a memorização, com a velocidade de absorção de conhecimento e com a sua estabilidade, é uma parte muito importante do aprendizado.
	\item[ (ii) ] \textbf{[ Intelecto Gama : Análise ]} Esta associado a busca efetiva pela verdade, e pela preservação da mesma ao longo de um raciocínio. A lógica e a matemática são exemplos de atividade intelectual predominantemente do tipo gama.
	\item[ (iii) ] \textbf{[ Intelecto Beta : Aplicação ]} O intelecto tipo gama e beta formam a dialética tradicional de ensino que se constitui em uma exposição através de um orador e a etapa de aplicação pela resolução de problemas. 
\end{itemize}

Em geral rejeito o estudo tradicional por conta desta projeção. Para efetivamente entender teorias mais elaboradas é necessário dar atenção ao componente alfa e não somente aos dois componentes gama e beta do método de ensino tradicional. Para teorias simples talvez o estudo tradicional cumpra bem seu papel.

\hspace{\baselineskip}

\textbf{Base Dialética [ Trialética Existencial ]} A trialética existencial é uma extensão da dialética existencial, a trialética vê uma existência como um processo de interação, temos a recepção, um estimulo que chega a forma externa e alcança seu interior, e o interior responde com uma mudança na forma externa, essa mudança é chamada de resposta. Vamos chamar o terceiro elemento de processamento na falta de outro termo, assim temos a trialética existencial como (Recepção, Processamento, Resposta). A propriedade estrutural é que os padrões de recepção e resposta caracterizam a existência, que os padrões de recepção e resposta é uma forma de tangenciar o caráter interno da existência, e que o processamento é a parte mais próxima do caráter interno. Portanto os padrões de recepção e resposta caracterizam a existência.

$$ \texttt{Trialética Existencial} \sim \texttt{Trialética Humana} $$

Note que a trialética existencial pode ser vista pela dialética existencial, a (recepção, resposta) corresponderia a forma externa do processamento, o processamento seria uma entidade mais interna, não exatamente o caráter interno, mas algo que está comunicando diretamente com o caráter interno.

$$ \texttt{Trialética Existencial} \prec ( \texttt{(Recepção, Resposta)}, \texttt{(Processamento) } ) \bullet \bigodot $$

%\textbf{Questão [ Existencial ]} Algo existe se não pode interagir? Qual a diferença entre algo que não existe e algo que não tem efeito nenhum sobre nosso universo? \\

\textbf{Base Dialética [ Trialética Humana ]} A trialética humana é uma projeção do ser humano como existência na trialética existencial, entretanto ela mesma é uma base dialética, vamos chamar (Emoção, Intelecto, Vontade) como a projeção da mente humana na (Recepção, Processamento, Resposta), é importante notar que estamos projetando a mente humana e não o corpo humano, a mente humana, a consciência humana ainda é interior ao nosso corpo. A emoção lida com a recepção de nossa consciência, o intelecto lida com o processamento, e a vontade lida com a resposta voluntária de nossa consciência que promove uma resposta (ou não) do corpo.\\

\textbf{Base Dialética [ Fundamento de Quatro Posições ]} O fundamento de quatro posições tem origem religiosa (encontrada no livro Princípio Divino, na religião conhecida como Igreja da Unificação ou Unificacionismo), mas assim como a dialética materialista ela será abordada de uma forma mais objetiva e menos metafísica. O fundamento de quatro posições busca corrigir a base dialética materialista adicionando um quarto elemento, este na religião é conhecida como Deus, mas que pode ser abordado mais objetivamente como a origem do propósito, ou propósito. Todo processo acontecendo no tempo e no espaço é uma interação entre dois elementos distribuídos no espaço, um elemento chamado de sujeito, mais ativo, possui mais iniciativa, e um elemento chamado de objeto, mais passivo e temos dois elementos associados a transformação temporal, o propósito e o resultado, ficamos então com (propósito, sujeito, objeto, resultado), entretanto a transformação temporal pode ser vista trialeticamente como (Origem, Divisão, União), este é a projeção temporal do fundamento de quatro posições, a origem corresponde ao propósito, a divisão corresponde a interação entre sujeito e objeto e a união é o resultado. O sujeito tem uma conexão mais íntima com o propósito enquanto o objeto tem conexão mais íntima com o resultado. Como o propósito é subjetivo o método de estudo será trialético, pois definir um propósito é particular do leitor, entretanto é interessante ter consciência de que é preciso de um propósito em suas ações, o ser humano deve agir de forma consciente para fazer jus a sua existência.\\

\textbf{[ Ação de Dar e Receber ]} A ação de dar e receber é um processo intrínseco da estrutura do fundamento de quatro posições, ela descreve o processo de interação entre o sujeito e o objeto, quando a interação se torna estável então dizemos que é formado uma nova existência, a existência resultante, ela pode ser de fato uma nova existência ou pode ser a internalização da interação entre o sujeito e objeto formando um novo centro, onde a própria interação se torna uma existência.\\

O fundamento de quatro posições é um exemplo de base tetraética, tem quatro elementos, entretanto o fundamento de quatro posições incorpora tanto a dialética quanto a trialética, a dialética é incorporada na relação entre sujeito e objeto (espacial) e a trialética é incorporada na origem, divisão e união (temporal). Portanto o fundamento de quatro posições estende ambos os conceitos e serve como um catalizador para ambos. Neste documento não será explorado o fundamento de quatro posições, pois é interessante analisar com mais cuidado as bases que a compõe, além disso existe uma dificuldade inerente na utilização desta base por conta do critério para escolha das entidades sujeito e objeto, como a modelagem dialética é subjetiva pode levar a interpretações equivocadas.\\ 

É interessante observar que as noções de tempo e espaço estão incutidas no fundamento de quatro posições, o espaço descrito como sujeito e objeto e o tempo descrito como o processo de origem, divisão e união. O caráter temporal apesar de ser um atributo da estrutura do fundamento de quatro posições, pode também ser atribuído a qualquer trialética, uma trialética tem uma característica temporal inerente. Quanto ao caráter espacial, ele pode ser visto como essencialmente dialético, dual, ou pode ser visto como uma coleção de elementos que podem estabelecer uma relação dialética. Essa característica do Fundamento de Quatro posições conter tanto dialética quanto trialética vem de um dos objetivos que justificam o propósito de sua criação, de que o fundamento de quatro posições seria assim como a dialética materialista, uma verdade universal, sua ambição é ser uma fórmula geral para as projeções dialéticas.\\

Porém me abstenho de utilizar o fundamento de quatro posições desta forma, para garantir a objetividade deste documento irei me resguardar das suas implicações esotéricas e me ater as suas implicações racionais.

%%%%%%%%%%%%%%%%%%%%%%%%%%%%%%%%%%%%%%%%%%%%%%%%%%%%%%%%%%%%%%%%%%%%%%%%%%%%%%%%%%%%%%%%%%%%%%%%%%%%%%%%%%%%%%%%%%%%%%%%%%%%%%%%%%%%%%%%%%%%%%%%%%%%%%%%%%%%%%%%%%%%%%
%%%%%%%%%%%%%%%%%%%%%%%%%%%%%%%%%%%%%%%%%%%%%%%%%%%%%%%%%%%%%%%%%%%%%%%%%%%%%%%%%%%%%%%%%%%%%%%%%%%%%%%%%%%%%%%%%%%%%%%%%%%%%%%%%%%%%%%%%%%%%%%%%%%%%%%%%%%%%%%%%%%%%%
%%%%%%%%%%%%%%%%%%%%%%%%%%%%%%%%%%%%%%%%%%%%%%%%%%%%%%%%%%%%%%%%%%%%%%%%%%%%%%%%%%%%%%%%%%%%%%%%%%%%%%%%%%%%%%%%%%%%%%%%%%%%%%%%%%%%%%%%%%%%%%%%%%%%%%%%%%%%%%%%%%%%%%

\subsection{Construção de Dialéticas Socráticas}

\hspace{\baselineskip}

Nesta subseção irei discutir sobre as implicações das bases dialéticas na dialética socrática. Para isso irei construir e nomear dialéticas socráticas como modelos para induzir perguntas. Para manter uma organização durante a construção irei definir dialética socrática como uma estrutura composta de um motivador da pergunta, chamado de indutor, e uma simulação de perguntas chamado de dialética. Todas as definições podem ser refinadas, elas não são fixas, porém para simplificar irei considerar apenas dialéticas simplificada, as dialéticas aqui são apenas sugestões, por isso um bom exercício é ser capaz de criar dialéticas mais refinadas. Caso exista uma relação intrínseca entre o indutor e uma finalidade irei também explicitar no corpo da definição da dialética.

%%%%%%%%%%%%%%%%%%%%%%%%%%%%%%%%%%%%%%%%%%%%%%%%%%%%%%%%%%%%%%%%%%%%%%%%%%%%%%%%%%%%%%%%%%%%%%%%%%%%%%%%%%%%%%%%%%%%%%%%%%%%%%%%%%%%%%%%%%%%%%%%%%%%%%%%%%%
% Dialética Socrática Contrapositiva
\hspace{\baselineskip}

\textbf{Dialética Socrática [ Trivial ]}

\begin{itemize}
	\item[ ] \textbf{Indutor :} Duas existências $A$ e $B$
	\item[ ] \textbf{Finalidade :} Relação, Interação, Comparação entre duas existências.	
	\item[ ] \textbf{Dialética :}
	\begin{itemize}
		\item[ (i) ] \textbf{Pergunta [ Positiva ]} ?` $A \sim B$ ? $\equiv$ ?` Existe alguma relação entre $A$ e $B$ ?
		\item[ (ii) ] \textbf{Pergunta [ Negativa ]} ?` $ A \times B $ ? $\equiv$ ?` Qual a diferença entre $A$ e $B$ ?
	\end{itemize}
\end{itemize}

Esta dialética é chamada trivial, pois pode ser considerado um dos primeiros tipos de pergunta a ser feito em qualquer tentativa de análise. É trivial mas é essencial que seja feito estas perguntas.




%%%%%%%%%%%%%%%%%%%%%%%%%%%%%%%%%%%%%%%%%%%%%%%%%%%%%%%%%%%%%%%%%%%%%%%%%%%%%%%%%%%%%%%%%%%%%%%%%%%%%%%%%%%%%%%%%%%%%%%%%%%%%%%%%%%%%%%%%%%%%%%%%%%%%%%%%%%
% Dialética Socrática [ de Modelagem ]
\hspace{\baselineskip}

\textbf{Dialética Socrática [ Modelagem ]} Suponha um problema $P$

\begin{itemize}
	\item[ ] \textbf{Indutor :} Seja um problema $P$ e um modelo ou estrutura $E$
	%\item[ ] $\succ P$ $\bullet$ \{Modelagem\}  ? $\equiv$ ?` Seja um problema $P$, existe possibilidade de modelar $P$ com alguma estrutura, existe alguma estrutura análoga a $P$ ?
	\item[ ] \textbf{Finalidade :} Compreensão, Entendimento, Consciência do problema $P$	
	\item[ ] \textbf{Dialética :}
	\begin{itemize}
		\item[ ] ?` Quais argumentos justificam $E$ como um bom modelo para $P$ ?
		\item[ ] ?` Quais argumentos justificam $E$ como sendo um mal modelo para $P$ ?
		\item[ ] ?` Existem outros modelos além de $E$ ?
	\end{itemize}
\end{itemize}

%%%%%%%%%%%%%%%%%%%%%%%%%%%%%%%%%%%%%%%%%%%%%%%%%%%%%%%%%%%%%%%%%%%%%%%%%%%%%%%%%%%%%%%%%%%%%%%%%%%%%%%%%%%%%%%%%%%%%%%%%%%%%%%%%%%%%%%%%%%%%%%%%%%%%%%%%%%
% Dialética Socrática [ Multiplicidade de Modelos ]
\hspace{\baselineskip}

\textbf{Dialética Socrática [ Multiplicidade de Modelos ]} Suponha um problema $P$ e modelos $E_n$

\begin{itemize}
	\item[ ] \textbf{Indutor :} Existência de vários modelos para um mesmo problema.
	\item[ ] \textbf{Finalidade :} Comparação, Modelagem
	\item[ ] \textbf{Dialética :}
	\begin{itemize}
		\item[ ] ?` Existe melhor modelo ?
		\item[ ] ?` Qual o melhor modelo ?
		\item[ ] ?` Existem casos particulares em que um modelo é melhor que outro ?
	\end{itemize}
\end{itemize}

%%%%%%%%%%%%%%%%%%%%%%%%%%%%%%%%%%%%%%%%%%%%%%%%%%%%%%%%%%%%%%%%%%%%%%%%%%%%%%%%%%%%%%%%%%%%%%%%%%%%%%%%%%%%%%%%%%%%%%%%%%%%%%%%%%%%%%%%%%%%%%%%%%%%%%%%%%%
% Dialética Socrática
\hspace{\baselineskip}

\textbf{Dialética Socrática [ Definição Transcendental ]}

\begin{itemize}
	\item[ ] \textbf{Indutor :} Falha na tentativa de construir um conceito de maneira finita. Relação do conceito com o infinito.
	\item[ ] \textbf{Finalidade :} Definição, Modelagem
	\item[ ] \textbf{Dialética :}
	\begin{itemize}
		\item[ ] ?` Existe possibilidade de construir o conceito de maneira transcendental usando abstração de redes via interpretação espacial ?
		\item[ ] ?` Existe possibilidade de construir o conceito de maneira transcendental usando abstração de redes via interpretação temporal ?
	\end{itemize}
\end{itemize}

Para finalizar esta subseção irei definir a dialética socrática associada a uma base dialética genérica, ela é similar a dialética de modelagem por ser de fato um processo de modelagem. Como a ideia da projeção dialética é compartilhada entre as diferentes bases então basta que haja apenas uma dialética socrática para uma base genérica. \\

\textbf{Dialética Socrática [ Projeções Dialéticas ]}

\begin{itemize}
	\item[ ] \textbf{Indutor :} Possuir alguma característica estrutural em comum com a base dialética, por exemplo número de elementos, relações análogas entre componentes, ou existir a possibilidade de agrupar componentes e compara-los como se fossem parte de um elemento só.
	\item[ ] \textbf{Finalidade :} Modelagem, Indutor da Base Dialética
	\item[ ] \textbf{Dialética :}
	\begin{itemize}
		\item[ ] ?` É possível dividir o conceito em componentes ?
		\item[ ] ?` É possível agrupar os componentes do conceito ?
		\item[ ] ?` Existe alguma base dialética em que os componentes ou grupos de componentes possam estar associados ? $\rightarrow$ D.S. [ de Modelagem ]
	\end{itemize}
\end{itemize}

A seguir os indutores de base que fornecem propósitos específicos associados a base dialética:\\

\textbf{Indutor de Base [ Ponto e Seta ]} Construir de maneira finita conceitos ou de maneira transcendental. \\

\textbf{Indutor de Base [ Espaço e Tempo ]} Entender como um processo evoluindo no tempo com conceitos de simultaneidade e transição.\\

\textbf{Indutor de Base [ Trialética ]} Entender como um processo evoluindo no tempo.\\

\textbf{Indutor de Base [ Trialética Existencial ]} Entender como um processo evoluindo no tempo tal que o processo em si é visto em um contexto diferente dos objetos que o processo conecta.\\

\textbf{Indutor de Base [ Trialética Humana ]} Herdando características da trialética existencial e da trialética, busca entender como o ser humano vê e deseja os componentes.\\

\textbf{Indutor de Base [ Fundamento de Quatro Posições ]} Diferenciar os papeis de sujeito e objeto para alcançar um objetivo ao longo de um processo no tempo através de uma interação estável e harmônica.

%%%%%%%%%%%%%%%%%%%%%%%%%%%%%%%%%%%%%%%%%%%%%%%%%%%%%%%%%%%%%%%%%%%%%%%%%%%%%%%%%%%%%%%%%%%%%%%%%%%%%%%%%%%%%%%%%%%%%%%%%%%%%%%%%%%%%%%%%%%%%%%%%%%%%%%%%%%
%%%%%%%%%%%%%%%%%%%%%%%%%%%%%%%%%%%%%%%%%%%%%%%%%%%%%%%%%%%%%%%%%%%%%%%%%%%%%%%%%%%%%%%%%%%%%%%%%%%%%%%%%%%%%%%%%%%%%%%%%%%%%%%%%%%%%%%%%%%%%%%%%%%%%%%%%%%
%%%%%%%%%%%%%%%%%%%%%%%%%%%%%%%%%%%%%%%%%%%%%%%%%%%%%%%%%%%%%%%%%%%%%%%%%%%%%%%%%%%%%%%%%%%%%%%%%%%%%%%%%%%%%%%%%%%%%%%%%%%%%%%%%%%%%%%%%%%%%%%%%%%%%%%%%%%


\section{Teoria de Evolução Existencial}

%%%%%%%%%%%%%%%%%%%%%%%%%%%%%%%%%%%%%%%%%%%%%%%%%%%%%%%%%%%%%%%%%%%%%%%%%%%%%%%%%%%%%%%%%%%%%%%%%%%%%%%%%%%%%%%%%%%%%%%%%%%%%%%%%%%%%%%%%%%%%%%%%%%%%%%%%%%
% Introdução
\hspace{\baselineskip}

A teoria de evolução dialética se propõe a tratar a teoria de evolução biológica em um contexto geral e neutro, portanto sem assumir a existência de entidades metafísicas. Entretanto a teoria de evolução servirá como base para análise da consciência humana, logo terá caráter transcendental e pode levar a induções que vão além da objetividade do mundo físico. 

\subsection{Definição de Existência Evolutiva}

%%%%%%%%%%%%%%%%%%%%%%%%%%%%%%%%%%%%%%%%%%%%%%%%%%%%%%%%%%%%%%%%%%%%%%%%%%%%%%%%%%%%%%%%%%%%%%%%%%%%%%%%%%%%%%%%%%%%%%%%%%%%%%%%%%%%%%%%%%%%%%%%%%%%%%%%%%%
% Caracterização de Existência
\hspace{\baselineskip}

\textbf{[ Existência Evolutiva ]} 

\begin{itemize}
	\item[ (i) ] \textbf{[ Identidade e Estabilidade ]} Ambos os conceitos estão associados, para ter identidade é necessário ser reconhecível e portanto é necessário estabilidade. A estabilidade permite que algo seja reconhecido.
	\item[ (ii) ] \textbf{[ Mutabilidade ]} Uma existência evolutiva não possui estabilidade absoluta, ela mantêm a identidade (sob a ótica humana) ao mesmo tempo que se permite alterar seus componentes, por exemplo, seu comportamento, ou mesmo os estados de algum componente externo.
\end{itemize}

$$ E.E. \prec \texttt{Parte Estável} \times \texttt{Parte Instável} \bullet \bigotimes $$

$$ (\texttt{Parte Estável}, \texttt{Parte Instável}) \sim \bigodot \sim ( \texttt{Identidade}, \texttt{Comportamento} ) $$

%%%%%%%%%%%%%%%%%%%%%%%%%%%%%%%%%%%%%%%%%%%%%%%%%%%%%%%%%%%%%%%%%%%%%%%%%%%%%%%%%%%%%%%%%%%%%%%%%%%%%%%%%%%%%%%%%%%%%%%%%%%%%%%%%%%%%%%%%%%%%%%%%%%%%%%%%%%
% Explicação da Definição
Ambos os conceitos são importantes para a evolução, não existe evolução de imutáveis e não existe como o ser humano enxergar a evolução de irreconhecíveis. Não obstante, por ser derivado de existência, pressupõe-se que ela possa interagir com outras existências.

%%%%%%%%%%%%%%%%%%%%%%%%%%%%%%%%%%%%%%%%%%%%%%%%%%%%%%%%%%%%%%%%%%%%%%%%%%%%%%%%%%%%%%%%%%%%%%%%%%%%%%%%%%%%%%%%%%%%%%%%%%%%%%%%%%%%%%%%%%%%%%%%%%%%%%%%%%%
% Evento como unidade de simultaneidade
\hspace{\baselineskip}

\textbf{[ Evento ]} Um evento é identificado pela simultaneidade, se contrapondo ao tempo. Pode se dizer como uma unidade identificável de ocorrências simultâneas ou quase simultâneas. \\

É possível afrouxar a definição de evento, por exemplo a queda de um meteoro, uma explosão, amanhecer podem ser considerados eventos na linguagem comum, entretanto eles a rigor são processos e não eventos simultâneos. Neste texto, porém, iremos desconsiderar processos mantendo a definição de evento como escrito acima.\\

\textbf{[ Evolução ]} A evolução é a transição de eventos. \\

\textbf{[ Sistema ]} É difícil falar acerca de eventos e evolução sem delimitar as existências do universo para um conjunto finito de existências, portanto definimos um sistema para ser considerado de maneira independente do universo. Um sistema se define pelas suas existências (não necessariamente finitas).\\

Naturalmente temos (Evento, Evolução) $\sim$ (Ponto, Seta), sendo que ambos associados ao conceito de existência evolutiva.\\

\textbf{Rede [ de Eventos de um Sistema ]} É a representação da evolução de um sistema através de eventos representado como uma rede com os pontos representando eventos e com as setas representando a probabilidade ou possibilidade de transição entre os eventos.

%%%%%%%%%%%%%%%%%%%%%%%%%%%%%%%%%%%%%%%%%%%%%%%%%%%%%%%%%%%%%%%%%%%%%%%%%%%%%%%%%%%%%%%%%%%%%%%%%%%%%%%%%%%%%%%%%%%%%%%%%%%%%%%%%%%%%%%%%%%%%%%%%%%%%%%%%%%
% Aleatório e Determinístico
\hspace{\baselineskip}

\textbf{[ Aleatório x Determinístico ]} Um evento futuro é determinístico se ele sempre ocorre através de um processo e um evento passado. Um evento aleatório é um evento que não é determinado completamente por nenhuma condição passada.\\

A existência do aleatório é uma pergunta em aberto, com a computação e a geração de números aleatórios surgiu o conceito de caos determinístico, algo é caótico determinístico se ele aparenta ser aleatório mas fundamentalmente é determinístico. Ambos o caos determinístico e o aleatório conservam o conceito de caos, algo é caótico se é errático, se aparenta desordem, ruído, flutuações sem padrão definido. Ao contrário do aleatório, o caos determinístico existe por conta do desenvolvimento tecnológico humano, pois a geração de números aleatórios feita pelo computador é determinística, o computador não consegue gerar números aleatórios reais, o que justifica o termo pseudo-aleatório usado para referir a estes.

\hspace{\baselineskip}

\textbf{Axioma [ Aleatório ]} É assumido a existência de fenômenos aleatórios na teoria de evolução, a justificativa é que é estranho existir liberdade e consciência em um contexto absolutamente determinístico, assumir que o aleatório não existe é assumir também que não existe liberdade, que a liberdade e desejo são ilusões e que estamos assistindo a um filme muito interessante sobre nós mesmos acreditando que somos conscientes e escrevendo sobre isso.

%%%%%%%%%%%%%%%%%%%%%%%%%%%%%%%%%%%%%%%%%%%%%%%%%%%%%%%%%%%%%%%%%%%%%%%%%%%%%%%%%%%%%%%%%%%%%%%%%%%%%%%%%%%%%%%%%%%%%%%%%%%%%%%%%%%%%%%%%%%%%%%%%%%%%%%%%%%
% Evento Estável Aleatório
\hspace{\baselineskip}

\textbf{[ Evento Estável Aleatório Local ]} Seja um sistema que pode transitar entre eventos de forma aleatória, um destes eventos tem estabilidade aleatória local se ela tem a probabilidade de permanecer no evento maior que a probabilidade de mudar para outros eventos adjacentes (ligados por setas em uma representação em rede).\\

O exemplo elementar de estabilidade aleatória local é a hereditariedade genética, se a informação gerada por mutações não deletérias (que não provocam morte imediata ou defeito grave ao organismo) junto com informações antigas não pudesse se manter com suficiente estabilidade ao longo das gerações, não haveria aumento de complexidade dos organismos. \\

\textbf{[ Evento Estável Aleatório Global ]} Seja um sistema que pode transitar entre eventos de forma aleatória, um evento aleatório é estável globalmente se a probabilidade de encontrar o sistema neste evento é maior que a probabilidade de encontrar o sistema nos outros eventos.\\

Herdar informações não deletérias não garante maior sobrevivência se tal informação não conferir alguma vantagem em relação a outras possíveis mutações, portanto a estabilidade para manter informações é diferente da estabilidade associado a vantagem que a informação confere no sistema como um todo. A análise global de estabilidade é complicada por conta da incapacidade humana de analisar todas as possibilidades, tal incapacidade não se deve somente a capacidade de analisar grandes quantidades de informação, mas pela própria natureza da evolução de estar analisando o futuro. Como prever a ação de uma entidade antes inexistente de forma efetiva no universo? Se ela era inefetiva, então não seria possível obter informações previamente sobre seu comportamento, nem sobre suas características, nem sobre o porquê se tornou efetivo, logo sua repentina capacidade de interação não pode ser prevista por um ser humano, que se insere como existência no universo. Em nosso universo, novas existências são geradas, a complexidade de análise é comparável a adição de existências antes inefetivas.\\

Por conta deste motivo a análise da evolução deve se dar com base em análises locais, o que faz perder o caráter determinístico das previsões da teoria. Porém é possível estender a análise local com uma análise relativamente global, por exemplo a adição do conceito de vantagem no exemplo de hereditariedade é uma extensão da estabilidade da informação durante a passagem da informação genética.



%%%%%%%%%%%%%%%%%%%%%%%%%%%%%%%%%%%%%%%%%%%%%%%%%%%%%%%%%%%%%%%%%%%%%%%%%%%%%%%%%%%%%%%%%%%%%%%%%%%%%%%%%%%%%%%%%%%%%%%%%%%%%%%%%%%%%%%%%%%%%%%%%%%%%%%%%%%
% Tendência
\hspace{\baselineskip}

\textbf{Argumento [ Tendência ]} Seja um evento futuro, uma tendência é um tipo de argumento utilizado para mostrar que, em determinado caso hipotético, um evento futuro irá provavelmente acontecer. Uma tendência não diz de forma determinística que algo irá acontecer, pois a interferência de outras existências fora da hipótese pode alterar o resultado final.\\

\textbf{Tendência [ Fundamental Aleatória ]} Suponha que a transição entre eventos de um sistema se dá de maneira aleatória, se existir um estado estável o sistema tende a ficar nele. Tudo tende a estabilidade.

%%%%%%%%%%%%%%%%%%%%%%%%%%%%%%%%%%%%%%%%%%%%%%%%%%%%%%%%%%%%%%%%%%%%%%%%%%%%%%%%%%%%%%%%%%%%%%%%%%%%%%%%%%%%%%%%%%%%%%%%%%%%%%%%%%%%%%%%%%%%%%%%%%%%%%%%%%%
% Estabilidade é um conceito relativo

\subsection{Evolução Dialética e Trialética}

\hspace{\baselineskip}

%%%%%%%%%%%%%%%%%%%%%%%%%%%%%%%%%%%%%%%%%%%%%%%%%%%%%%%%%%%%%%%%%%%%%%%%%%%%%%%%%%%%%%%%%%%%%%%%%%%%%%%%%%%%%%%%%%%%%%%%%%%%%%%%%%%%%%%%%%%%%%%%%%%%%%%%%%%
% Tendência Fundamental Aleatória e Evolução de Comportamento a Tendência de Existir

%%%%%%%%%%%%%%%%%%%%%%%%%%%%%%%%%%%%%%%%%%%%%%%%%%%%%%%%%%%%%%%%%%%%%%%%%%%%%%%%%%%%%%%%%%%%%%%%%%%%%%%%%%%%%%%%%%%%%%%%%%%%%%%%%%%%%%%%%%%%%%%%%%%%%%%%%%%
% Análise Local
Devemos portanto nos ater a análises locais e estender caso necessário a análises mais globais. É natural portanto considerar primeiramente a análise de interações entre existências, visto que não faz sentido analisar existências isoladas, pois não são existências efetivas.



%%%%%%%%%%%%%%%%%%%%%%%%%%%%%%%%%%%%%%%%%%%%%%%%%%%%%%%%%%%%%%%%%%%%%%%%%%%%%%%%%%%%%%%%%%%%%%%%%%%%%%%%%%%%%%%%%%%%%%%%%%%%%%%%%%%%%%%%%%%%%%%%%%%%%%%%%%%
% Relação de Existência Mutável
\hspace{\baselineskip}

\textbf{Relação [ Interação Existencial ]} Algo existe se sua presença implica em alguma alteração reconhecível em confronto com a hipótese de inexistência. Vamos definir uma interação abstrata entre existências: Sejam $A$ e $B$ existências evolutivas, $A \overset{mod.}{\rightarrow} B$ significa que a presença de $A$ modifica a parte mutável de $B$ em comparação com a hipótese de que $A$ não age sobre $B$. \\

Dentro do contexto evolutivo, a relação caracteriza o conceito de existência efetiva, algo existe efetivamente para outro se ele pode alterar seu comportamento. É importante notar que este conceito é relativo, algo pode ser efetivo por transição, ou seja, altera alguém que altera diretamente a existência em questão. Devo também lembrar de que não estamos sob o domínio da lógica formal, assim todo raciocínio envolvendo existências evolutivas não é determinístico.

%%%%%%%%%%%%%%%%%%%%%%%%%%%%%%%%%%%%%%%%%%%%%%%%%%%%%%%%%%%%%%%%%%%%%%%%%%%%%%%%%%%%%%%%%%%%%%%%%%%%%%%%%%%%%%%%%%%%%%%%%%%%%%%%%%%%%%%%%%%%%%%%%%%%%%%%%%%
% Interação Existencial Destrutiva
\hspace{\baselineskip}

\textbf{Interação Existencial [ Destrutiva ]} Sejam duas existência $A$ e $B$ evolutivas, dizemos $A \overset{inibe}{\rightarrow} B$ caso a presença de $A$ prejudique a existência $B$, no sentido de diminuir sua estabilidade. Esta relação é uma interação existencial, que simboliza a competição na teoria evolutiva darwiniana.

%%%%%%%%%%%%%%%%%%%%%%%%%%%%%%%%%%%%%%%%%%%%%%%%%%%%%%%%%%%%%%%%%%%%%%%%%%%%%%%%%%%%%%%%%%%%%%%%%%%%%%%%%%%%%%%%%%%%%%%%%%%%%%%%%%%%%%%%%%%%%%%%%%%%%%%%%%%
% Interação Existencial Construtiva 
\hspace{\baselineskip}

\textbf{Interação Existencial [ Construtiva ]} Sejam duas existências $A$ e $B$, dizemos que $A \overset{ativa}{\rightarrow} B$ se $A$ contribui para existência de $B$ aumentando sua estabilidade.\\

Estas duas formas de interação permitem o começo de uma análise mais interessante, pois agora é possível considerar diferentes estratégias para aumento de estabilidade. Como se trata de uma interação entre existências a base trialética topológica parece ser adequada para modelar as diferentes estratégias. \\

Para refrescar a memória a base trialética topológica está associado a trialética humana. Naturalmente existe uma forte associação entre as estratégias definidas aqui e a trialética humana.

$$(\alpha, \gamma, \beta) \sim (\texttt{Emoção}, \texttt{Intelecto}, \texttt{Vontade}) $$

\textbf{Estratégia [ Alfa ]} A existência que busca evitar receber interações destrutivas e a busca receber interações construtivas de outras existências é mais estável. É chamado de estratégia alfa por considerar o recebimento de uma interação de outra existência. \\

Assim como a emoção está associado as sensações e percepções a estratégia alfa está associado a parte que recebe interferências de outras existências. O foco da estratégia alfa é maximizar as relações construtivas e portanto o coletivo passa ter importância. \\

\textbf{Estratégia [ Beta ]} A existência que busca inibir existências prejudiciais e ajudar existências que colaborem para sua estabilidade é mais estável. A estratégia é chamado de beta por considerar a ação da existência sobre outras existências. \\

O foco da estratégia beta é a capacidade de transformar a realidade, o poder que uma existência tem em moldar a realidade para seu próprio benefício está associado a ação da existência. \\

As estratégias alfa e beta, assim como a associação trialética topológica sugere, estão associados ao que se vê na realidade, quanto a estratégia gama, está associado ao interior da própria existência. Apesar da estratégia alfa e beta serem suficientes para descrever as estratégias que garantem a estabilidade entre existências evolutivas, a evolução propriamente dita somente pode ser compreendida pela estratégia gama, pois a evolução tem a ver com a retenção das informações que garantem externalizações de estratégias alfa e estratégia beta.\\

\textbf{Estratégia [ Gama ]} A existência que tem capacidade de reter informações que ajudem a externalizar estratégias alfa e beta e eliminem informações que prejudiquem a sua estabilidade é mais estável.\\

%%%%%%%%%%%%%%%%%%%%%%%%%%%%%%%%%%%%%%%%%%%%%%%%%%%%%%%%%%%%%%%%%%%%%%%%%%%%%%%%%%%%%%%%%%%%%%%%%%%%%%%%%%%%%%%%%%%%%%%%%%%%%%%%%%%%%%%%%%%%%%%%%%%%%%%%%%%
% Comentário das Estratégias
Estas estratégias permitem aplicar o conceito de evolução além da evolução biológica, que é a evolução da informação genética ao longo de gerações de seres vivos, por exemplo podemos usar como modelo para compreender a evolução comportamental de um indivíduo, com efeito, a informação retida em gama seriam as memórias, as boas memórias estão associados as boas estratégias alfa e beta.

%%%%%%%%%%%%%%%%%%%%%%%%%%%%%%%%%%%%%%%%%%%%%%%%%%%%%%%%%%%%%%%%%%%%%%%%%%%%%%%%%%%%%%%%%%%%%%%%%%%%%%%%%%%%%%%%%%%%%%%%%%%%%%%%%%%%%%%%%%%%%%%%%%%%%%%%%%%
% Dificuldade de Estabelecer quem é mais importante
\hspace{\baselineskip}

É difícil estabelecer qual estratégia é melhor para garantir a estabilidade, ou seja, qual a tendência evolutiva entre as três estratégias, pois como a trialética sugere as três estão associadas a partes importantes e distintas do conceito de existência. O caráter da base trialética é a harmonia e equilíbrio, o que ela nos fornece são direções distintas, como eixos de coordenadas na matemática.\\

Por hora vamos tentar refinar as definições de forma que seja possível construir e julgar cenários hipotéticos.

%%%%%%%%%%%%%%%%%%%%%%%%%%%%%%%%%%%%%%%%%%%%%%%%%%%%%%%%%%%%%%%%%%%%%%%%%%%%%%%%%%%%%%%%%%%%%%%%%%%%%%%%%%%%%%%%%%%%%%%%%%%%%%%%%%%%%%%%%%%%%%%%%%%%%%%%%%%
% Existência Tipo Alfa
\hspace{\baselineskip}

\textbf{Existência [ Alfa ]} É uma existência que prioriza a capacidade de escolher receber ou evitar boas interações de outras existências.

\begin{itemize}
	\item[ (i) ] \textbf{Existência Alfa [ Positiva ]} Prioriza a busca por receber boas interações.
	\item[ (ii) ] \textbf{Existência Alfa [ Negativa ]} Prioriza evitar receber más interações.
\end{itemize}

Sob o ponto de vista de alfa podemos definir atrator como uma entidade que fornece boas interações para alfas e repulsor como a entidade que prejudica, neste contexto fica evidente as tendências de que existências alfa são atraídas por seus respectivos atratores e repelidas por seus respectivos repulsores. Por conta dessa analogia faz sentido que as existências alfa tenham maior mobilidade, ser capaz de se mover é fundamental para a estratégia alfa. Faz sentido também que a alfa positiva seja mais móvel, pois busca interações, enquanto a alfa negativa somente se move quando encontra um repulsor. \\

\textbf{Existência [ Beta ]} É uma existência que prioriza a capacidade de estimular existências benéficas e inibir existências prejudiciais.

\begin{itemize}
	\item[ (i) ] \textbf{Existência Beta [ Positiva ]} É uma existência que prioriza estimular existências benéficas.
	\item[ (ii) ] \textbf{Existência Beta [ Negativa ]} É uma existência que prioriza inibir existências prejudiciais.
\end{itemize}

Em contrapartida às existências alfa, as existências beta por terem capacidade de ação sobre outras existências não tem como fundamental sua mobilidade, entretanto faz sentido que as existências beta positivas sejam mais móveis que as existências beta negativa. Por conta disso existências betas podem ser considerados como os atratores e repulsores de existências alfa. \\

\textbf{Tendência [ à Formação de Universos Existências Efetivos Centrados em Betas ]} Com base nas definições de existências alfa e beta podemos enxergar uma tendência das existências betas ficarem rodeadas por existências alfa compatíveis e se isolar de outras existências betas criando pequenos universos/sistemas existências efetivos.\\

Aqui podemos entender porque considerar existências alfa e beta já nos fornece recursos para entender interações dinâmicas (evolução), a verdade é que a informação está contida no posicionamento espacial de cada existência, então não seria necessário existências do tipo gama, pois a evolução da informação se dá pela cinética das entidades envolvidas. Existências do tipo gama permitem considerar a evolução de outras informações além de somente o posicionamento das existências, o que torna não somente mais interessante, mas essencial por exemplo para evolução biológica. \\

Com finalidade de aumentar a resolução, vamos chamar de evolução dialética como a evolução que considera apenas as estratégias alfa e beta, e evolução trialética a evolução que considera as três estratégias.\\

\textbf{Existência [ Gama ]} É uma existência que prioriza a retenção de boas informações e descartar informações prejudiciais.

\begin{itemize}
	\item[ (i) ] \textbf{Existência Gama [ Positiva ]} Prioriza a retenção de boas informações.
	\item[ (ii) ] \textbf{Existência Gama [ Negativa ]} Prioriza o descarte de boas informações.
\end{itemize}

Os processos que ocorrem dentro de uma existência gama, ou seja, a dinâmica da evolução das informações pode ser considerado como uma evolução dialética. Em outras palavras, podemos projetar a evolução da informação na base dialética alfa e beta.

$$\succ \gamma \bullet (\alpha, \beta) \text{?}$$

A realidade e seus efeitos diretos na fisiologia da existência são entidades do tipo beta, eles irão fornecer as informações para que a evolução do tipo gama aconteça, relativamente a estes temos as informação retidas na existência como entidades do tipo alfa, o que não impede que entre os alfa tenhamos interações alfa e beta.\\

A verdade é que a trialética humana (Emoção, Intelecto, Vontade) é uma projeção de $\gamma$ em $\alpha$ e $\beta$, pois tanto a emoção como a vontade residem na mente. Então esse procedimento não é só aceitável como talvez essencial para compreender a trialética existencial.

\subsection{Tendência a Existir e Réplicas}

\hspace{\baselineskip}

A tendência a existir externaliza o próprio conceito de estabilidade, tender a existir e tender a ser mais estável externalizam a mesma ideia. Entretanto não é natural enxergar a tendência de existir com apenas uma réplica da existência, pois ao enxergar um isolado, tanto faz se ele se torna mais estável ou mais instável, pois ele apenas depende de si mesmo, se este isolado tem processos auto-destrutivos então a tendência é que ele cesse de existir, se ele tem processos que garantem a estabilidade então a tendência é que ele se mantenha, é mais instrutivo considerar várias réplicas de uma mesma existência e considerar sua estabilidade em relação a um universo como um todo e não sua auto-estabilidade. Um exemplo para esta ideia é considerar o universo como uma existência, é estranho considerar a evolução do universo, tanto ele pode caminhar para estabilidade quanto para instabilidade dependendo apenas de seus componentes internos, neste sentido a evolução tem caráter dialético de considerar o indivíduo em relação a um todo e não o todo como um todo.

%%%%%%%%%%%%%%%%%%%%%%%%%%%%%%%%%%%%%%%%%%%%%%%%%%%%%%%%%%%%%%%%%%%%%%%%%%%%%%%%%%%%%%%%%%%%%%%%%%%%%%%%%%%%%%%%%%%%%%%%%%%%%%%%%%%%%%%%%%%%%%%%%%%%%%%%%%%
% Tendência a Existir
\hspace{\baselineskip}

\textbf{Tendência [ a Existir ]} Suponha uma existência e considere a sua estabilidade como sua interação efetiva em um universo, considere que hajam réplicas desta existência, as existências que tendem a existir no universo são aquelas que tem maior capacidade de interação com ela.\\

É preciso um pouco de cuidado ao lidar com esta definição, pois estabilidade da maneira como foi construída é relativa e portanto a própria definição de estabilidade de existência e portanto a tendência é relativa. Para ilustrar a ideia vamos considerar interações fortes como linhas fortes em uma rede de pontos e setas (pontos são existências, setas são interações), as linhas fracas representam interações inefetivas, imagine que cada vez apagamos as linhas fracas e apagamos os pontos isolados, logo as réplicas com linhas fracas tem a tendência neste contexto a serem eliminadas do modelo, o que não significa que elas não possam voltar a interagir em um contexto real, pois algo isolado não significa que não exista. Esse problema caracteriza sua aplicação em contextos locais, pois globalmente mesmo existências aparentemente inefetivas podem alterar o curso de evolução do sistema.\\

Essa tendência pode ser interpretada como um outro ponto de vista para tendência à Formação de Universos Existências Efetivos Centrados em Betas. O universo como um todo é como se fosse um grande atrator para existências alfas que buscam interagir, apesar da interação não ter caráter positivo ou negativo neste contexto.\\

Podemos aplicar este conceito para externalizar a ideia de relevância, quanto maior capacidade de interação, mais relevante uma existência, se quisermos abstrair o problema apenas aos elementos relevantes, é interessante considerar esta tendência.


% \textbf{Existências Alfa e Beta Compatíveis} Uma existência alfa é compatível com uma existência beta se a


%%%%%%%%%%%%%%%%%%%%%%%%%%%%%%%%%%%%%%%%%%%%%%%%%%%%%%%%%%%%%%%%%%%%%%%%%%%%%%%%%%%%%%%%%%%%%%%%%%%%%%%%%%%%%%%%%%%%%%%%%%%%%%%%%%%%%%%%%%%%%%%%%%%%%%%%%%%
% Trialética e Estratégias de Sobrevivência

%\textbf{Estabilidade [ Construtiva ]} Uma existência é mais estável se pode escolher interagir com relações construtivas.

%%%%%%%%%%%%%%%%%%%%%%%%%%%%%%%%%%%%%%%%%%%%%%%%%%%%%%%%%%%%%%%%%%%%%%%%%%%%%%%%%%%%%%%%%%%%%%%%%%%%%%%%%%%%%%%%%%%%%%%%%%%%%%%%%%%%%%%%%%%%%%%%%%%%%%%%%%%
% 
%\hspace{\baselineskip}

%\textbf{Estabilidade [ Construtiva Egoísta ]} É a estabilidade construtiva em que as relações escolhidas não são recíprocas, mas para benefício próprio.\\

%\textbf{Estabilidade [ Destrutiva ]} Uma existência é mais estável se pode escolher evitar relações destrutivas.\\

%Importante perceber que aqui escolha está sendo usado como analogia implícita, não necessariamente a existência escolhe no sentido próprio da palavra, pois é suposto que apenas existências conscientes possam ter tal habilidade.

%%%%%%%%%%%%%%%%%%%%%%%%%%%%%%%%%%%%%%%%%%%%%%%%%%%%%%%%%%%%%%%%%%%%%%%%%%%%%%%%%%%%%%%%%%%%%%%%%%%%%%%%%%%%%%%%%%%%%%%%%%%%%%%%%%%%%%%%%%%%%%%%%%%%%%%%%%%
% Formação de Novas Existências através de Simbiose
%\hspace{\baselineskip}

%\textbf{Relação [ Recíproca ]} Uma relação recíproca é uma relação recíproca de interação construtiva.



%%%%%%%%%%%%%%%%%%%%%%%%%%%%%%%%%%%%%%%%%%%%%%%%%%%%%%%%%%%%%%%%%%%%%%%%%%%%%%%%%%%%%%%%%%%%%%%%%%%%%%%%%%%%%%%%%%%%%%%%%%%%%%%%%%%%%%%%%%%%%%%%%%%%%%%%%%%
% Pesquisar sobre teoria dos jogos

%%%%%%%%%%%%%%%%%%%%%%%%%%%%%%%%%%%%%%%%%%%%%%%%%%%%%%%%%%%%%%%%%%%%%%%%%%%%%%%%%%%%%%%%%%%%%%%%%%%%%%%%%%%%%%%%%%%%%%%%%%%%%%%%%%%%%%%%%%%%%%%%%%%%%%%%%%%
% relação simbiótica provocando dependência

%%%%%%%%%%%%%%%%%%%%%%%%%%%%%%%%%%%%%%%%%%%%%%%%%%%%%%%%%%%%%%%%%%%%%%%%%%%%%%%%%%%%%%%%%%%%%%%%%%%%%%%%%%%%%%%%%%%%%%%%%%%%%%%%%%%%%%%%%%%%%%%%%%%%%%%%%%%
% Suavização e Parcial Transitividade


%%%%%%%%%%%%%%%%%%%%%%%%%%%%%%%%%%%%%%%%%%%%%%%%%%%%%%%%%%%%%%%%%%%%%%%%%%%%%%%%%%%%%%%%%%%%%%%%%%%%%%%%%%%%%%%%%%%%%%%%%%%%%%%%%%%%%%%%%%%%%%%%%%%%%%%%%%%
% Teoria de Evolução Trialética



%\subsection{Evolução da Informação e Evolução Gama}

%\hspace{\baselineskip}

%O problema principal desta subseção é enxergar a evolução gama como evolução de informações codificadas em uma existência.

%$$ \succ \gamma \text{ ?} $$

%%%%%%%%%%%%%%%%%%%%%%%%%%%%%%%%%%%%%%%%%%%%%%%%%%%%%%%%%%%%%%%%%%%%%%%%%%%%%%%%%%%%%%%%%%%%%%%%%%%%%%%%%%%%%%%%%%%%%%%%%%%%%%%%%%%%%%%%%%%%%%%%%%%%%%%%%%%
% FPF e Teoria de Evolução

\subsection{Evolução Trialética e Fundamento de Quatro Posições}

\hspace{\baselineskip}

A base trialética é um modelo que fornece uma interpretação simétrica e equilibrada dos componentes, para uma análise assimétrica podemos usar o fundamento de quatro posições, que essencialmente dialético divide em sujeito e objeto, sujeito mais ativo e objeto mais passivo. Naturalmente temos a relação:

$$ (\alpha, \beta) \sim (\texttt{Objeto}, \texttt{Sujeito}) $$

Pois uma existência tipo beta busca predominantemente interferir de forma ativa, enquanto que a existência alfa busca predominantemente interagir de forma passiva. A tendência de evolução é beta ficar rodeado por alfas (tendência à Formação de Universos Existências Efetivos Centrados em Betas), ambos são complementares neste sentido.\\

Quanto a existências do tipo gama, estas tem sua habilidade em reter informações e adquirir por meio destas informações qualidades alfa ou beta, como é algo interior sendo externalizado então existe um tempo para que o potencial gama possa ter efeito na realidade, existências com potencial gama no início são considerados objetos, mas depois de externalizados seu potencial podem vir a se tornar sujeitos dos outros tipos de existência. A analogia que combina neste contexto é a relação que ocorre dentro de uma família:

$$ (\texttt{Pai}, \texttt{Mãe}, \texttt{Filhos}) \sim (\beta, \alpha, \gamma) $$

No contexto do fundamento de quatro posições, a relação entre $\gamma$ e $(\alpha, \beta)$ é vertical enquanto que a relação entre $\alpha$ e $\beta$ é chamada de relação horizontal.\\

Foge do escopo deste documento tratar da evolução biológica, tal conteúdo será tratado em trabalhos posteriores. A finalidade desta subseção é pavimentar o caminho para análise da trialética humana sob o ponto de vista da consciência humana.







%%%%%%%%%%%%%%%%%%%%%%%%%%%%%%%%%%%%%%%%%%%%%%%%%%%%%%%%%%%%%%%%%%%%%%%%%%%%%%%%%%%%%%%%%%%%%%%%%%%%%%%%%%%%%%%%%%%%%%%%%%%%%%%%%%%%%%%%%%%%%%%%%%%%%%%%%%%
%%%%%%%%%%%%%%%%%%%%%%%%%%%%%%%%%%%%%%%%%%%%%%%%%%%%%%%%%%%%%%%%%%%%%%%%%%%%%%%%%%%%%%%%%%%%%%%%%%%%%%%%%%%%%%%%%%%%%%%%%%%%%%%%%%%%%%%%%%%%%%%%%%%%%%%%%%%
%%%%%%%%%%%%%%%%%%%%%%%%%%%%%%%%%%%%%%%%%%%%%%%%%%%%%%%%%%%%%%%%%%%%%%%%%%%%%%%%%%%%%%%%%%%%%%%%%%%%%%%%%%%%%%%%%%%%%%%%%%%%%%%%%%%%%%%%%%%%%%%%%%%%%%%%%%%


\section{A Consciência Humana}

\hspace{\baselineskip}

%%%%%%%%%%%%%%%%%%%%%%%%%%%%%%%%%%%%%%%%%%%%%%%%%%%%%%%%%%%%%%%%%%%%%%%%%%%%%%%%%%%%%%%%%%%%%%%%%%%%%%%%%%%%%%%%%%%%%%%%%%%%%%%%%%%%%%%%%%%%%%%%%%%%%%%%%%%
% Dificuldade inerente de explicar a consciência
A consciência é um fenômeno interessante e difícil de entender, pois a única ferramenta para sua compreensão é ela mesma. O próprio fenômeno em si é transcendental, a consciência é fruto de uma interação dinâmica que acontece no cérebro unindo a realidade e a memória, é o centro que confere a identidade da existência humana.



%%%%%%%%%%%%%%%%%%%%%%%%%%%%%%%%%%%%%%%%%%%%%%%%%%%%%%%%%%%%%%%%%%%%%%%%%%%%%%%%%%%%%%%%%%%%%%%%%%%%%%%%%%%%%%%%%%%%%%%%%%%%%%%%%%%%%%%%%%%%%%%%%%%%%%%%%%%
% Gancho com evolução existencial
\hspace{\baselineskip}

Com o auxílio do que foi construído na subseção de evolução existencial, a existência humana possui as três estratégias $(\alpha, \gamma, \beta)$ sendo a $\gamma$ mais próxima do que queremos e associado a evolução da informação dentro da mente.

%%%%%%%%%%%%%%%%%%%%%%%%%%%%%%%%%%%%%%%%%%%%%%%%%%%%%%%%%%%%%%%%%%%%%%%%%%%%%%%%%%%%%%%%%%%%%%%%%%%%%%%%%%%%%%%%%%%%%%%%%%%%%%%%%%%%%%%%%%%%%%%%%%%%%%%%%%%
% Vontade o que é vontade?

%%%%%%%%%%%%%%%%%%%%%%%%%%%%%%%%%%%%%%%%%%%%%%%%%%%%%%%%%%%%%%%%%%%%%%%%%%%%%%%%%%%%%%%%%%%%%%%%%%%%%%%%%%%%%%%%%%%%%%%%%%%%%%%%%%%%%%%%%%%%%%%%%%%%%%%%%%%
%%%%%%%%%%%%%%%%%%%%%%%%%%%%%%%%%%%%%%%%%%%%%%%%%%%%%%%%%%%%%%%%%%%%%%%%%%%%%%%%%%%%%%%%%%%%%%%%%%%%%%%%%%%%%%%%%%%%%%%%%%%%%%%%%%%%%%%%%%%%%%%%%%%%%%%%%%%

\subsection{Análise Trialética Existencial da Consciência}

%%%%%%%%%%%%%%%%%%%%%%%%%%%%%%%%%%%%%%%%%%%%%%%%%%%%%%%%%%%%%%%%%%%%%%%%%%%%%%%%%%%%%%%%%%%%%%%%%%%%%%%%%%%%%%%%%%%%%%%%%%%%%%%%%%%%%%%%%%%%%%%%%%%%%%%%%%%
% Melhor Modelo para Consciência é a Trialética Existencial
\hspace{\baselineskip}

Como o conceito de consciência está associado a existência humana, esta caracterizada por uma existência topológica interagindo dinamicamente com o seu entorno, então o melhor modelo dialético é a trialética existencial (Percepção, Processamento, Resposta).\\

Ao nos referir a consciência em si podemos usar a trialética existencial para estabelecer a relação conceitual abaixo:

$$ \texttt{Processamento} \sim \texttt{Informação} $$

$$ \texttt{Resposta} \sim \texttt{Ação} $$

$$ \texttt{(Percepção, Informação, Ação)} \sim (\alpha, \gamma, \beta) $$

\hrulefill

$$ \texttt{(Percepção, Informação, Ação)} \sim \texttt{(Realidade, Memória, Imaginação)} $$

Perceba que foi interessante modificar a trialética existencial para estabelecer a relação conceitual desejada, poderia ser feito de maneira direta, mas a fim de tornar natural a discussão, algumas vezes convém operar este raciocínio.\\

A imaginação é uma ação interessante, pois é uma ação que acontece exclusivamente na mente, é diferente de mexer um braço, movimentar os olhos, a imaginação é ser capaz de mexer com a organização das conexões entre os neurônios. A imaginação depende dos elementos adquiridos da realidade e preservados na memória, esta responsável pela capacidade de reter informações da realidade. Ao estarmos conscientes (no sentido de estarmos despertos) podemos dizer que a consciência (a nossa mente, identidade) navega entre estes três componentes de informação, a percepção da realidade, a memória e a imaginação, e temos como diferenciar com alguma certeza de qual fonte uma informação está associada.\\

Quando dizemos que uma existência é formada pela interação dinâmica entre existências finitas estamos definindo a mesma de maneira transcendental (associado ao tempo), é dito que essa existência é uma existência virtual, ou que ela é o centro virtual formado pela interação entre existências. A consciência pode ser definida desta maneira, como o centro que corresponde a interação entre as informações da realidade, as informações retidas na memória e a evolução da informação na imaginação.

$$ \texttt{(Realidade, Memória, Imaginação)} \sim \texttt{(Presente, Passado, Futuro)} $$

A capacidade de diferenciar estes três elementos parece estar associado aos conceitos de presente, passado e futuro. Portanto podemos dizer que presente, passado e futuro são conceitos a priori, pois estão atrelados a própria estrutura da consciência. É importante ressaltar que a análise dialética não fornece respostas conclusivas, assim como a modelagem, ela ajuda a compreensão e a proposição de perguntas.

%%%%%%%%%%%%%%%%%%%%%%%%%%%%%%%%%%%%%%%%%%%%%%%%%%%%%%%%%%%%%%%%%%%%%%%%%%%%%%%%%%%%%%%%%%%%%%%%%%%%%%%%%%%%%%%%%%%%%%%%%%%%%%%%%%%%%%%%%%%%%%%%%%%%%%%%%%%
% Inconsistência Topológica
\hspace{\baselineskip}

Existe um detalhe interessante na modelagem trialética da consciência: a memória e a imaginação estão mais próximas topologicamente entre si no interior do que a realidade que é algo exterior, enquanto que na trialética existencial a percepção e a ação estão próximas entre si no exterior e a informação está no interior. A associação só não é inconsistente, pois consideramos as bases trialéticas como associações de cada componente e não da preservação da estrutura topológica, poderíamos ser mais rigorosos na modelagem e tentar preservar a estrutura, entretanto teríamos que modificar a ordem dos componentes. Vamos representar a proximidade topológica por parênteses internos.

$$ \texttt{(Realidade , (Memória, Imaginação))} \sim \texttt{(Informação, (Percepção, Ação))} $$

Essa aparente inconsistência apenas reflete na natureza da modelagem conceitual dialética, quando fazemos uma associação, tentamos ajustar cada ideia contida nos conceitos de forma que faça sentido. Dessa forma, a associação dialética serve apenas como catalizador para diferentes ideias contidas nos conceitos e não uma nova informação (apesar de abstrações e atalhos são úteis, elas não são propriamente novas informações). Essa ideia é preciso ser assimilada para usar a análise dialética como ferramenta de raciocínio.\\

Observe que a inconsistência topológica não está presente em relação aos conceitos de presente, passado e futuro.

$$ \texttt{(Realidade, (Memória, Imaginação))} \sim \texttt{(Presente, (Passado, Futuro))} $$

%%%%%%%%%%%%%%%%%%%%%%%%%%%%%%%%%%%%%%%%%%%%%%%%%%%%%%%%%%%%%%%%%%%%%%%%%%%%%%%%%%%%%%%%%%%%%%%%%%%%%%%%%%%%%%%%%%%%%%%%%%%%%%%%%%%%%%%%%%%%%%%%%%%%%%%%%%%
% Insight inerente da interpretação topológica
Existe um insight que é inerente da interpretação topológica (dialética topológica ou existencial) em que a forma externa é a forma como podemos entender o caráter interno, assim na estrutura (Realidade, (Memória, Imaginação)) temos a (Memória, Imaginação) como forma externa à consciência para a compreensão da realidade.  \\

%%%%%%%%%%%%%%%%%%%%%%%%%%%%%%%%%%%%%%%%%%%%%%%%%%%%%%%%%%%%%%%%%%%%%%%%%%%%%%%%%%%%%%%%%%%%%%%%%%%%%%%%%%%%%%%%%%%%%%%%%%%%%%%%%%%%%%%%%%%%%%%%%%%%%%%%%%%
% Consciência e Perguntas Filosóficas

%%%%%%%%%%%%%%%%%%%%%%%%%%%%%%%%%%%%%%%%%%%%%%%%%%%%%%%%%%%%%%%%%%%%%%%%%%%%%%%%%%%%%%%%%%%%%%%%%%%%%%%%%%%%%%%%%%%%%%%%%%%%%%%%%%%%%%%%%%%%%%%%%%%%%%%%%%%
% Consciência e Vontade

%%%%%%%%%%%%%%%%%%%%%%%%%%%%%%%%%%%%%%%%%%%%%%%%%%%%%%%%%%%%%%%%%%%%%%%%%%%%%%%%%%%%%%%%%%%%%%%%%%%%%%%%%%%%%%%%%%%%%%%%%%%%%%%%%%%%%%%%%%%%%%%%%%%%%%%%%%%
% definição Preliminar de Consciência


Termino esta seção com uma definição preliminar de consciência, definimos consciência como o centro que dá identidade a relação entre as três entidades realidade, memória e imaginação. Podemos substituir realidade por sentidos ou sensações e portanto teríamos (Sensações, Memória, Imaginação) como a trialética que define a consciência.


%%%%%%%%%%%%%%%%%%%%%%%%%%%%%%%%%%%%%%%%%%%%%%%%%%%%%%%%%%%%%%%%%%%%%%%%%%%%%%%%%%%%%%%%%%%%%%%%%%%%%%%%%%%%%%%%%%%%%%%%%%%%%%%%%%%%%%%%%%%%%%%%%%%%%%%%%%%
%%%%%%%%%%%%%%%%%%%%%%%%%%%%%%%%%%%%%%%%%%%%%%%%%%%%%%%%%%%%%%%%%%%%%%%%%%%%%%%%%%%%%%%%%%%%%%%%%%%%%%%%%%%%%%%%%%%%%%%%%%%%%%%%%%%%%%%%%%%%%%%%%%%%%%%%%%%

%\subsection{Trialética da Consciência}

%\hspace{\baselineskip}

%%%%%%%%%%%%%%%%%%%%%%%%%%%%%%%%%%%%%%%%%%%%%%%%%%%%%%%%%%%%%%%%%%%%%%%%%%%%%%%%%%%%%%%%%%%%%%%%%%%%%%%%%%%%%%%%%%%%%%%%%%%%%%%%%%%%%%%%%%%%%%%%%%%%%%%%%%%
%%%%%%%%%%%%%%%%%%%%%%%%%%%%%%%%%%%%%%%%%%%%%%%%%%%%%%%%%%%%%%%%%%%%%%%%%%%%%%%%%%%%%%%%%%%%%%%%%%%%%%%%%%%%%%%%%%%%%%%%%%%%%%%%%%%%%%%%%%%%%%%%%%%%%%%%%%%

\subsection{Análise Trialética da Memória}

%%%%%%%%%%%%%%%%%%%%%%%%%%%%%%%%%%%%%%%%%%%%%%%%%%%%%%%%%%%%%%%%%%%%%%%%%%%%%%%%%%%%%%%%%%%%%%%%%%%%%%%%%%%%%%%%%%%%%%%%%%%%%%%%%%%%%%%%%%%%%%%%%%%%%%%%%%%
% Propriedades Induzidas da Memória com Interpretação de Ponto e Seta

\hspace{\baselineskip}

%\textbf{Questão [ Memória e Consciência ]} No paradigma trialético de estudo queremos dar uma atenção especial a memorização, queremos resolver o problema do excesso de informações. É interessante, antes de afirmar qualquer, coisa entender como funciona a nossa memória, como armazenamos conteúdo, como acessamos conteúdo. Quais os problemas, o que pode inibir o armazenamento de conteúdo, o que pode inibir o acesso do conteúdo?\\

%Vamos tentar resolver o problema com o método rudimentar da filosofia de projeções dialéticas. Rudimentar no sentido de que poderia ser feito com base na neurociência, mas que não pretendo faze-lo no momento. Mas antes vamos estabelecer uma notação.\\

%\textbf{Notação [ Projeção Dialética ]} 

%$$ \texttt{Objeto} \prec \texttt{Projeção} \bullet \texttt{Base Dialética} $$

%$$ \texttt{Objeto} \prec \texttt{Componentes} \bullet \texttt{Base Dialética} $$

%$$ \texttt{Objeto} \prec (a,b,c) \bullet (A,B,C) $$

%Com esta notação queremos dizer que o Objeto está sendo visto como uma projeção sob a luz de uma base dialética.

%$$ (a,b,c) \sim (f,g,h) $$ 

%$$ (a,b,c) \bullet (f,g,h) $$

%Representa a preservação de uma mesma estrutura entre as coleções (a,b,c) e (f,g,h).

%$$ \succ \texttt{Objeto} \bullet \texttt{Base Dialética} ? $$

%Significa a determinação do problema de encontrar os componentes/projeções do objeto sob a ótica da base dialética.

%$$ A := \texttt{Algo mais Longo} $$

%O simbolo $:=$ representa definição, no sentido usado acima representa a encapsulação de um conceito em uma representação mais simples.

%$$ A \prec B \bullet C \bullet D $$

%É importante notar que interno e externo é um conceito relativo, assim o terceiro elemento D se refere a quem está sendo usado como referência na comparação, a notação acima deve ser visto como A visto como B projeção sob a ótica de C base dialética relativo a D.  Outra forma de enxergar o conceito é interpretar C como um intermediário de comunicação entre B e D.

%\hspace{\baselineskip}

%De posse desta notação iremos realizar a análise dialética, a modelagem dialética, ou projeções dialéticas usando a representação dos conceitos simbolicamente, de forma que possamos nos libertar da linearidade natural do texto para algo com caráter mais espacial.

%%%%%%%%%%%%%%%%%%%%%%%%%%%%%%%%%%%%%%%%%%%%%%%%%%%%%%%%%%%%%%%%%%%%%%%%%%%%%%%%%%%%%%%%%%%%%%%%%%%%%%%%%%%%
% intro
A memória é o candidato natural para uma análise mais detalhada, pois dentre as entidades que pertencem ao interior da existência humana, a memória é mais fundamental, sendo que a imaginação depende das informações contidas na memória. A memória representa a retenção de informação, característica essencial para existências evolutivas.

%%%%%%%%%%%%%%%%%%%%%%%%%%%%%%%%%%%%%%%%%%%%%%%%%%%%%%%%%%%%%%%%%%%%%%%%%%%%%%%%%%%%%%%%%%%%%%%%%%%%%%%%%%%%
% Projeção Dialética da Memória
\hspace{\baselineskip}

\textbf{Projeção Dialética [ Memória e Bases Existenciais ]} 

$$ \succ \texttt{Memória} \bullet \texttt{Trialética Existencial} ? $$
$$ \texttt{Memória} \prec \texttt{(Recepção, Armazenamento, Acesso)} \bullet \texttt{Trialética Existencial} $$
\hrulefill

Podemos reaplicar a trialética existencial para a memória, de forma a enxergar a memória como um processo em que a consciência (o centro entre as entidades realidade, memória e imaginação reside) interage com a memória. Desta maneira temos os processos de recepção de informação da realidade ou da imaginação, armazenamento da informação recebida e acesso da informação armazenada.\\

$$ \succ \texttt{(Recepção, Armazenamento, Acesso)} \bullet \texttt{Dialética Existencial} ? $$
$$ \texttt{(Armazenamento, \{Recepção, Acesso\})} \sim \texttt{(Caráter Interno, Forma Externa)} \sim \texttt{Dialética Existencial} $$
$$ \texttt{(Qualidade de Armazenamento, \{Qualidade de Recepção, Qualidade de Acesso\})} \sim \texttt{Dialética Existencial} $$
\hrulefill

A qualidade de armazenamento pode ser verificada pela qualidade da recepção e pela qualidade de acesso.\\

$$ Q1:= \texttt{Qualidade de Recepção} $$
$$ Q2:= \texttt{Qualidade de Armazenamento} $$
$$ Q3:= \texttt{Qualidade de Acesso} $$
$$  \texttt{Memória} \prec (\{Q1,Q2\} , Q3) \bullet \texttt{Dialética Existencial} \bullet \texttt{Consciência} $$
\hrulefill

Para a nossa consciência podemos enxergar a qualidade da recepção e a qualidade de armazenamento de conteúdo através da qualidade do seu acesso. Queremos ter uma boa qualidade de recepção, ter uma boa qualidade de armazenamento, entretanto somente temos consciência através da qualidade de seu acesso. 

%%%%%%%%%%%%%%%%%%%%%%%%%%%%%%%%%%%%%%%%%%%%%%%%%%%%%%%%%%%%%%%%%%%%%%%%%%%%%%%%%%%%%%%%%%%%%%%%%%%%%%%%%%%%
% É preciso fazer uma discussão sobre o modelo da memória e suas implicações
%\hspace{\baselineskip}

%%%%%%%%%%%%%%%%%%%%%%%%%%%%%%%%%%%%%%%%%%%%%%%%%%%%%%%%%%%%%%%%%%%%%%%%%%%%%%%%%%%%%%%%%%%%%%%%%%%%%%%%%%%%
% Memória e consciência intelectual alfa
%\hspace{\baselineskip}

%%%%%%%%%%%%%%%%%%%%%%%%%%%%%%%%%%%%%%%%%%%%%%%%%%%%%%%%%%%%%%%%%%%%%%%%%%%%%%%%%%%%%%%%%%%%%%%%%%%%%%%%%%%%
%%%%%%%%%%%%%%%%%%%%%%%%%%%%%%%%%%%%%%%%%%%%%%%%%%%%%%%%%%%%%%%%%%%%%%%%%%%%%%%%%%%%%%%%%%%%%%%%%%%%%%%%%%%%

%%%%%%%%%%%%%%%%%%%%%%%%%%%%%%%%%%%%%%%%%%%%%%%%%%%%%%%%%%%%%%%%%%%%%%%%%%%%%%%%%%%%%%%%%%%%%%%%%%%%%%%%%%%%

%%%%%%%%%%%%%%%%%%%%%%%%%%%%%%%%%%%%%%%%%%%%%%%%%%%%%%%%%%%%%%%%%%%%%%%%%%%%%%%%%%%%%%%%%%%%%%%%%%%%%%%%%%%%
%%%%%%%%%%%%%%%%%%%%%%%%%%%%%%%%%%%%%%%%%%%%%%%%%%%%%%%%%%%%%%%%%%%%%%%%%%%%%%%%%%%%%%%%%%%%%%%%%%%%%%%%%%%%

\subsection{Redes e Sensação de Memória}

\hspace{\baselineskip}

Nesta subseção vamos tentar interpretar a memória pela sua representação natural em redes. Suponha que existam pontos que representem um determinado padrão de estímulo vindo da realidade, cada ponto assume valores 1 ou 0, para um determinado estímulo temos um determinado padrão de uns e zeros. Se adicionarmos apenas uma hipótese no modelo podemos entender três propriedades interessantes da sensação de memória.

%%%%%%%%%%%%%%%%%%%%%%%%%%%%%%%%%%%%%%%%%%%%%%%%%%%%%%%%%%%%%%%%%%%%%%%%%%%%%%%%%%%%%%%%%%%%%%%%%%%%%%%%%%%%
% Hipótese de Formação de Conexões por Coestímulo
\hspace{\baselineskip}

\textbf{Hipótese [ Formação de Conexões Positivas por Co-estímulo ]} Quando temos um evento externalizado por uma sensação vinda diretamente da realidade representado por um padrão entre zeros e uns, é formado conexões desordenadas conectando os pontos com valores uns. Estas conexões representam a memória deste evento. Tais conexões são chamadas de positivas, pois quando um ponto fica um ela estimula o outro ponto conectado a ficar um. \\

A primeira propriedade é a memória em si, chamamos de \textbf{Preservação de Informação por Conexões}, suponha que a sensação de um evento é obtida quando o padrão de uns é repetido, logo se a imaginação re-estimular algum ponto da fronteira das conexões que representem a memória deste evento e se tal estímulo promover um estímulo em cadeia baseado nas conexões formadas por co-estímulo então o padrão de uns é repetido e podemos interpretar isso como o processo de lembrança de um evento.\\

Tal propriedade é a essência da memória, a essência da preservação de informação. A segunda propriedade obtida é chamada de \textbf{Sensação do Todo pela Parte}, ela consiste na capacidade de sentir todas (ou quase todas) as informações de um determinado objeto ou evento a partir do estímulo real ou da imaginação de parte dos padrões de zeros e uns. Essa propriedade é derivada da formação de conexões por co-estímulo por conta da capacidade das conexões estimularem em cadeia os outros pontos que compõe a informação do objeto ou evento.\\

A terceira propriedade é interessante, chamada de \textbf{Conexões por Similares}, ela permite conectar eventos similares, por exemplo, objetos similares, suponha dois eventos com quase todos os pontos com valores iguais, pela propriedade anterior de sensação do todo pela parte então o estímulo de um evento irá estimular a sensação do estímulo do outro evento, o que promove uma conexão entre as duas informações e pode permitir modelar o processo de raciocínio como a evolução das sensações (externalizada pelos padrões de zeros e uns) através de sensações similares.\\

É interessante perceber que é natural neste modelo que o raciocínio se dê de maneira "contínua", pois transições contínuas são similares. \\

Perceba que podemos aplicar a base trialética existencial para tais propriedades da forma abaixo:

$$ \texttt{Sensação do Todo pela Parte} \sim \alpha $$

$$ \texttt{Preservação de Informação por Conexões} \sim \gamma $$

$$ \texttt{Conexões por Similares} \sim \beta $$

A propriedade de conexões por similares modela o raciocínio que é um tipo de ação, a sensação do todo pela parte está associado a recepção e percepção da informação, e a preservação de informação por conexões está associado a preservação da informação em si. Portanto tal associação faz sentido e é modelado por algo bem simples.\\

%%%%%%%%%%%%%%%%%%%%%%%%%%%%%%%%%%%%%%%%%%%%%%%%%%%%%%%%%%%%%%%%%%%%%%%%%%%%%%%%%%%%%%%%%%%%%%%%%%%%%%%%%%%%
% Estímulos negativos
Nesta hipótese estamos levando em conta apenas os estímulos positivos, entretanto apenas com estímulos positivos não há garantia de o sistema funcione, pois se todos os estímulos fossem positivos a tendência é que todos pontos fiquem com valor 1 e isso não pode ser capaz de estimular sensação alguma, pois não tem informação nenhuma. O papel dos estímulos negativos deve ser de discriminar informações distintas e também de colocar ordem no raciocínio, com essa habilidade pode ser que o sistema funcione e possamos entender como o raciocínio baseado em similares funciona.

%%%%%%%%%%%%%%%%%%%%%%%%%%%%%%%%%%%%%%%%%%%%%%%%%%%%%%%%%%%%%%%%%%%%%%%%%%%%%%%%%%%%%%%%%%%%%%%%%%%%%%%%%%%%
% hipótese de Conexões Negativas e informação
\hspace{\baselineskip}

\textbf{Hipótese [ Memória por Formação de Conexões Positivas e Negativas ]} Nesta hipótese, além de formação de conexões positivas por co-estímulo, também temos a formação de conexões negativas, de forma que os pontos de uns inibam os pontos de zero.\\

Tal hipótese é razoável, pois do ponto de vista da preservação de informação, tal padrão de formações de conexões preservam a informação original. Com a adição das conexões inibitórias temos como diferenciar informações, as informações inibem e competem entre si para estimularem sua própria sensação.



%%%%%%%%%%%%%%%%%%%%%%%%%%%%%%%%%%%%%%%%%%%%%%%%%%%%%%%%%%%%%%%%%%%%%%%%%%%%%%%%%%%%%%%%%%%%%%%%%%%%%%%%%%%%
%%%%%%%%%%%%%%%%%%%%%%%%%%%%%%%%%%%%%%%%%%%%%%%%%%%%%%%%%%%%%%%%%%%%%%%%%%%%%%%%%%%%%%%%%%%%%%%%%%%%%%%%%%%%

\subsection{Redes e Raciocínio da Imaginação}

\hspace{\baselineskip}

Nesta subseção iremos analisar as consequências que o modelo de memória por redes é capaz dizer sobre as propriedades do raciocínio. Para isso iremos definir duas duas representações em redes, a camada de sensação, que interpreta os pontos como padrões de estímulos de sensações, e a camada de informação, que interpreta os pontos como informações distintas.

%%%%%%%%%%%%%%%%%%%%%%%%%%%%%%%%%%%%%%%%%%%%%%%%%%%%%%%%%%%%%%%%%%%%%%%%%%%%%%%%%%%%%%%%%%%%%%%%%%%%%%%%%%%%
% Sensation Layer
\hspace{\baselineskip}

\textbf{[ Camada de Sensação ]} É a representação das diferentes sensações por pontos que podem ter valor um ou zero. As conexões são formadas de acordo com o que foi definido na subseção anterior.  \\

\textbf{[ Camada de Informação ]} É a representação das informações como pontos. As conexões são resultantes das conexões formadas na camada de sensação e possuem comportamento similar de estímulo negativo e positivo. \\

A camada de sensação é relativamente mais concreta que a camada de informação, a camada de informação é um modelo construído em cima das conexões formadas na camada de sensação, os pontos na camada de informação abstraem as conexões formadas na camada de sensação. Podemos fazer a seguinte associação nos componentes da base trialética.

$$ \texttt{Camada de Sensação} \sim \alpha $$
$$ \texttt{Camada de Informação} \sim \gamma $$

Como o raciocínio está associado a $\beta$ podemos sintetizar o objetivo desta subseção como:

$$ \succ \beta \bullet (\alpha, \gamma) \text{?} $$

O tipo de raciocínio induzido pela camada de sensação é chamado de \textbf{raciocínio por similaridade}. A camada de informação herda características da camada de sensação, entretanto a interpretação espacial das informações induz a ideia de proximidade, a proximidade entre as informações pode ser por similaridade, mas pode ser também pela proximidade física, e isto está associado a maneira como a informação é armazenada fisicamente. É de se esperar que não somente informações similares sejam estimuladas, mas também informações fisicamente próximas, assim defino o outro tipo de raciocínio como \textbf{raciocínio por proximidade}, sendo a proximidade física considerada.\\

Raciocínio por similaridade é mais livre, nos permite imaginar através de transformações contínuas de memórias, o raciocínio por proximidade é mais fiel as experiências reais (apesar de imaginações também se tornarem memórias).\\

%%%%%%%%%%%%%%%%%%%%%%%%%%%%%%%%%%%%%%%%%%%%%%%%%%%%%%%%%%%%%%%%%%%%%%%%%%%%%%%%%%%%%%%%%%%%%%%%%%%%%%%%%%%%
% Raciocínio como transição de memórias ~ realidade e verossimilhança
Apesar de chamarmos de raciocínio, ele não é propriamente um raciocínio, pois um raciocínio busca um objetivo, por exemplo preservar a verdade, chegar a uma conclusão, fazer uma previsão. Neste contexto raciocínio seria algo anterior a isso, uma transição de memórias não necessariamente com sentido, por exemplo em um sonho a transição de memórias parece seguir as mesmas regras, transições por similaridade e transições baseadas na proximidade física das informações no cérebro.\\

Este é um exemplo interessante, e o modelo também consegue explicar por que a transição de ideias no estado desperto é diferente do sonho, muito provavelmente a contínua exposição aos sentidos reais inibe memórias e transições que são estranhas a realidade (isso pelas conexões negativas que definem as informações) tornando as transições mais verossímeis à realidade. Durante o sonho, sem a exposição constante da realidade, a transições são mais livres e podem ser sem sentido.\\

Nada impede porém que se imagine coisas estranhas, entretanto temos a capacidade de sentir como uma sensação quando imaginamos coisas sem sentido e podemos assim fazer distinção entre o que é real e o que é ficção. Estas formas de transição das memórias estimuladas na camada de sensação modela o que chamamos de imaginação e isto conclui com alguma informação esta subseção.

%%%%%%%%%%%%%%%%%%%%%%%%%%%%%%%%%%%%%%%%%%%%%%%%%%%%%%%%%%%%%%%%%%%%%%%%%%%%%%%%%%%%%%%%%%%%%%%%%%%%%%%%%%%%
% Raciocínio ou Sensação
%Apesar de chamarmos de raciocínio, ambas induzem dois tipos de percepção, ou dois tipos de sensações, a sensação de proximidade temporal da aquisição da informação e a sensação de similaridade entre informações.

%%%%%%%%%%%%%%%%%%%%%%%%%%%%%%%%%%%%%%%%%%%%%%%%%%%%%%%%%%%%%%%%%%%%%%%%%%%%%%%%%%%%%%%%%%%%%%%%%%%%%%%%%%%%
%%%%%%%%%%%%%%%%%%%%%%%%%%%%%%%%%%%%%%%%%%%%%%%%%%%%%%%%%%%%%%%%%%%%%%%%%%%%%%%%%%%%%%%%%%%%%%%%%%%%%%%%%%%%

\subsection{Aspecto Prático da Interpretação em Redes}

\hspace{\baselineskip}

Entender através do modelo de redes como funciona o raciocínio na imaginação permite a compreensão da maneira como a atividade acadêmica evoluiu para otimizar o desenvolvimento do pensamento.

%%%%%%%%%%%%%%%%%%%%%%%%%%%%%%%%%%%%%%%%%%%%%%%%%%%%%%%%%%%%%%%%%%%%%%%%%%%%%%%%%%%%%%%%%%%%%%%%%%%%%%%%%%%%
% Externalização na Realidade
\hspace{\baselineskip}

\textbf{Paradigma [ Externalização ]} Assim como a realidade induz a verossimilhança inibindo ideias sem sentido podemos utilizar a própria realidade para otimizar um raciocínio abstrato colocando artificialmente os elementos que quisermos na realidade, seja em um quadro, em um papel ou mesmo através da verbalização das ideias. Ao fazer isso usamos a realidade para induzir a transição de memórias reduzindo o que sentimos como desconcentração.

%%%%%%%%%%%%%%%%%%%%%%%%%%%%%%%%%%%%%%%%%%%%%%%%%%%%%%%%%%%%%%%%%%%%%%%%%%%%%%%%%%%%%%%%%%%%%%%%%%%%%%%%%%%%
% Paradigma
\hspace{\baselineskip}

\textbf{Paradigma [ Nomeação ]} Nomear uma ideia abstrata permite a sua externalização, combinado com o paradigma de externalização permite que seja possível resolver eficientemente problemas abstratos. É o primeiro passo, nomear e conquistas.

%%%%%%%%%%%%%%%%%%%%%%%%%%%%%%%%%%%%%%%%%%%%%%%%%%%%%%%%%%%%%%%%%%%%%%%%%%%%%%%%%%%%%%%%%%%%%%%%%%%%%%%%%%%%
% Externalização ~ Sopa de Ideias
\hspace{\baselineskip}

\textbf{Externalização [ Sopa de Ideias ]} Uma sopa de ideias é uma estrutura externalizada pela escrita em que nomes, símbolos, desenhos e ideias são postas em proximidade física para induzir a formulação de elementos teóricos definidos, como perguntas, proposições, soluções. Ela possui pelo menos dois campos separados, o campo em que as ideias são livres e um campo em que caso se encontre alguma coisa útil seja escrito as conclusões.\\

A sopa de ideias permite pelo menos dois tipos de "raciocínios" (no sentido de externalizações para resolver problemas) importantes. O \textbf{raciocínio criativo} em que se busca criar coisas, como um texto, uma teoria, proposições ou resolver problemas ainda sem solução; o \textbf{raciocínio de resgate de memória} em que o objetivo é resgatar na memória uma informação complicada, por exemplo, uma fórmula, um teorema, neste tipo de problema se considera que o cérebro é melhor em reconhecer uma informação correta do que lembrar da mesma. O paradigma de sopa de ideias foi utilizado para escrever este documento.

%%%%%%%%%%%%%%%%%%%%%%%%%%%%%%%%%%%%%%%%%%%%%%%%%%%%%%%%%%%%%%%%%%%%%%%%%%%%%%%%%%%%%%%%%%%%%%%%%%%%%%%%%%%%
%%%%%%%%%%%%%%%%%%%%%%%%%%%%%%%%%%%%%%%%%%%%%%%%%%%%%%%%%%%%%%%%%%%%%%%%%%%%%%%%%%%%%%%%%%%%%%%%%%%%%%%%%%%%

\subsection{Sensações e o Elo Metafísico}

\hspace{\baselineskip}

Neste documento metafísica, dentre outras possíveis definições, se refere ao que não é objetivo, que não pode ser considerado real ou uma ferramenta para entender o real, mas que está mais ligado a espiritualidade e a hipóteses que não é possível demonstrar sua existência na realidade. Sensação define algo para entender a realidade, entretanto ele é um elo que induz naturalmente a ideias metafísicas. Nesta subseção irei externalizar esta naturalidade.\\

Sensação é algo que não pode ser explicado plenamente sem que o interlocutor seja capaz de sentir a mesma sensação, ou sensação similar. Sensação está associado com perceber o contraste entre o que é e o que não é aquela sensação. No modelo de redes, sensação é modelada como pontos "ativados" com uns e "desativados" com zeros. Da maneira como abordamos o conceito de sensação, ela parece estar atrelada a própria informação, talvez toda informação diferenciada corresponda a uma sensação, e da maneira como a informação é armazenada em nosso cérebro, então todo padrão de conexão parece corresponder a uma sensação.\\

Associar sensação com conexão no modelo de redes é suficiente para entender o porquê de ser o elo metafísico, o modelo abstrato de redes é uma abstração universal, talvez toda e qualquer coisa intelectível para o ser humano pode ser interpretada como conexões entre existências, portanto se é possível sentir através de conexões entre neurônios, por que não seria possível uma entidade superior sentir através de interações entre existências? Se nossa consciência é formada pela dinâmica da evolução de conexões entre neurônios, por que não seria possível que haja consciências formadas pela evolução das interações entre existências? Provavelmente o leitor já deve ter concluído que a existência de algo como Deus, ou menos ousado, como gaia (que seria aqui uma forma de consciência da natureza) são induções com origem nesta ideia.\\

Se sensações estão ligadas a informação, mas nem toda sensação é efetivamente sentida pelo ser humano, boa parte das transformações de informações de nosso corpo não são nem sentidas. O porquê de uma sensação ser ou não efetiva na consciência provavelmente tem relação com a sua localização virtual, como modelado anteriormente, a consciência está entre a percepção da realidade, a memória e a imaginação.

$$ \texttt{Ação} \sim \texttt{Objetivo} \sim \texttt{Unicidade} $$ 

$$ \texttt{Percepção} \sim \texttt{Multiplicidade} $$

É natural da consciência a unicidade, e esta é natural do objetivo e assim da ação, que na consciência cabe a imaginação. Portanto se tem alguém responsável pela seleção das percepções efetivamente sentidas deve estar associado a unicidade da ação.\\

%Para terminar esta subseção irei dois modelos trialéticos que correspondem a 

Como raciocínio é uma ação e ação está associado a escolha, a análise do raciocínio requer uma discussão acerca dos conceitos de vontade e desejo. Portanto postergo esta discussão.




%%%%%%%%%%%%%%%%%%%%%%%%%%%%%%%%%%%%%%%%%%%%%%%%%%%%%%%%%%%%%%%%%%%%%%%%%%%%%%%%%%%%%%%%%%%%%%%%%%%%%%%%%%%%
%%%%%%%%%%%%%%%%%%%%%%%%%%%%%%%%%%%%%%%%%%%%%%%%%%%%%%%%%%%%%%%%%%%%%%%%%%%%%%%%%%%%%%%%%%%%%%%%%%%%%%%%%%%%

\subsection{Desejo, Vontade e Escolha}

%%%%%%%%%%%%%%%%%%%%%%%%%%%%%%%%%%%%%%%%%%%%%%%%%%%%%%%%%%%%%%%%%%%%%%%%%%%%%%%%%%%%%%%%%%%%%%%%%%%%%%%%%%%%
% relação entre desejo vontade e escolha
\hspace{\baselineskip}

A relação entre desejo, vontade e escolha é simples quando se sabe o que cada conceito se refere, a vontade seria o nome que se dá a dinâmica da consciência em relação às escolhas dos desejos, desejos seriam opções enviadas a consciência e escolha seria uma ação natural da mesma. Tentar construir estes conceitos assumindo que não conhecemos é mais difícil, o que seria escolha? Realmente existe escolha em um universo determinístico? Escolha é um conceito a priori, assim como os desejos, é natural para a consciência, mas é suficientemente fundamental para ser difícil defini-las.

%%%%%%%%%%%%%%%%%%%%%%%%%%%%%%%%%%%%%%%%%%%%%%%%%%%%%%%%%%%%%%%%%%%%%%%%%%%%%%%%%%%%%%%%%%%%%%%%%%%%%%%%%%%%
% Escolha fundamental para a ação

%%%%%%%%%%%%%%%%%%%%%%%%%%%%%%%%%%%%%%%%%%%%%%%%%%%%%%%%%%%%%%%%%%%%%%%%%%%%%%%%%%%%%%%%%%%%%%%%%%%%%%%%%%%%
% Dialética Socrática o que é uma escolha
\hspace{\baselineskip}

\textbf{Dialética Socrática [ Escolha ]} ?` O que é uma escolha ?

\begin{itemize}
	\item[ (i) ] Ilusão ou não, escolha é resultado da evolução da informação dentro da consciência.
	\item[ (ii) ] Existe uma relação entre escolha e a transformação da multiplicidade em unicidade.
	\begin{itemize}
		\item[ (i) ] A consciência de um indivíduo é vista como algo único, então a escolha é especialmente importante para entender a consciência. 
		\item[ (ii) ] Cada escolha é vista como a escolha da consciência. A unicidade da consciência é externalizada pela unicidade da escolha.
	\end{itemize}
	\item[ (iii) ] Existe uma relação entre escolha e o aleatório, visto que o determinismo torna a escolha uma ilusão.
	\begin{itemize}
		\item[ (i) ] Se enxergamos as informações como entidades que podem se combinar e formar outras informações, temos o caos gerado pela combinatória, o mesmo caos gerado pela formação de novas existências. Talvez o caos seja responsável pela escolha, e o aleatório não seja algo intrínseco da existência, mas uma composição entre processos determinísticos com entrada aleatória. 
	\end{itemize}
	\item[ (iv) ] Escolha talvez seja a definição fundamental de ação.
	\begin{itemize}
		\item[ (ii) ] Uma ação é uma forma de informação que gera externalizações no corpo, entretanto a escolha é mais fundamental. Podemos reduzir a discussão entre o funcionamento da mente e do corpo para apenas a evolução da informação na mente.
	\end{itemize}	
	\item[ (v) ] Talvez a pergunta ?` Como surge a escolha ? seja mais simples
	\begin{itemize}
		\item[ (i) ] Seja a realidade, a imaginação, ou ao instinto, estes estimulam a memória (informação armazenada em conexões) à diferentes desejos (o que diferencia um desejo consciente de um inconsciente é que o desejo consciente está na imaginação, o inconsciente é fraco demais para ser resgatado de forma recorrente) e a escolha surge na transformação de diferentes desejos em apenas um desejo.
		%\item[ (ii) ] Como ocorre a transformação de diferentes desejos em um desejo. 
	\end{itemize}
	\item[ (vi) ] ?` Como surge a sensação de escolha ?
	\begin{itemize}
		\item[ (i) ] Sentir a capacidade de escolher é algo natural da consciência, talvez até uma forma de defini-la.
		\item[ (ii) ] ?` O que é necessário para poder sentir ?
		\begin{itemize}
			\item[ (i) ] Informação baseadas em conexões parece ser necessária, mas não é suficiente, visto que temos o subconsciente, são informações que parecem não serem fortes o suficiente.
			\item[ (ii) ] ?` O que seria força de consciência ?
			\item[ (iii) ] ?` É necessário para sentir que a informação esteja gravada na memória ?
			%\begin{itemize}
			%	\item[ (i) ] Uma sensação é forte
			%\end{itemize}
			\item[ (iv) ] Quem usa a base trialética para modelar conceitos pode se sentir tentado a dizer que para sentir seja necessário pelo menos dois requerimentos, cada um associado a um componente da base trialética que falta, visto que sentir já está associado ao componente $\alpha$. Neste raciocínio, para poder sentir algo, ele primeiro deve ser uma informação gravada na memória (componente $\gamma$), segundo é preciso que a informação tenha capacidade de estimular de forma recorrente a imaginação (componente $\beta$), ou seja, ela deve ser estável e forte o suficiente para poder ser resgatada pela imaginação (embora isso se configure uma escolha).
		\end{itemize}
		\item[ (iii) ] Talvez para poder sentir a capacidade de escolha seja necessário que ela se torne uma informação estável na memória e tal informação seja forte o suficiente para poder ser resgatada de maneira recorrente.
	\end{itemize}
	\item[ (vii) ] ?` Como a sensação de escolha se torna uma informação ?
	\begin{itemize}
		\item[ (i) ] A sensação de escolha é abstrata, resultado da informação de várias escolhas e de instintos. Talvez o contraste entre a informação que identifica um instinto e a informação que identifica a negação do instinto. Talvez a negação do instinto seja a formação da escolha e coincida com o desenvolvimento da consciência.
	\end{itemize}
\end{itemize}

Com base na dialética socrática acima vamos tentar entender o conceito de escolha.

%%%%%%%%%%%%%%%%%%%%%%%%%%%%%%%%%%%%%%%%%%%%%%%%%%%%%%%%%%%%%%%%%%%%%%%%%%%%%%%%%%%%%%%%%%%%%%%%%%%%%%%%%%%%
% Informação na Consciência
\hspace{\baselineskip}

\textbf{[ Informação na Consciência ]} Uma informação na consciência é um padrão de conexões que pode gerar um resultado no corpo ou na mente.

%%%%%%%%%%%%%%%%%%%%%%%%%%%%%%%%%%%%%%%%%%%%%%%%%%%%%%%%%%%%%%%%%%%%%%%%%%%%%%%%%%%%%%%%%%%%%%%%%%%%%%%%%%%%
% Modelo Trialético Baseado na Socrática
\hspace{\baselineskip}

\textbf{Modelo [ Trialético da Escolha ]} O modelo trialético da escolha tem dois componentes externalizados como informações da consciência, a primeira se constitui de múltiplas informações, que competem entre si e são possivelmente contraditórias, o segundo componente é o resultante das múltiplas informações e a dinâmica da consciência, se transformando em informações menos contraditórias e possivelmente em apenas uma única informação, o elo que que liga a multiplicidade para a unicidade é o componente fora do contexto de informações de consciência. 

\begin{itemize}
	\item[ (i) ] \textbf{[ Componente Gama ]} O componente gama representa a residência das múltiplas informações.
	\item[ (ii) ] \textbf{[ Componente Alfa ]} O componente alfa é a realidade atuando como uma direção para que as múltiplas informações se transformem em um. Sem a realidade as múltiplas informações se confundem e podem gerar sensações e experiências contraditórias, como em um sonho.
	\item[ (iii) ] \textbf{[ Componente Beta ]} É a transformação das múltiplas informações em poucas informações, estas podendo ter efeito na realidade através de ações ou nas sensações da mente. Os efeitos do componente beta da escolha pode ser visto novamente com a trialética:
	\begin{itemize}
		\item[ (i) ] \textbf{[ Alfa ]} O efeito das múltiplas informações se tornarem poucas pode provocar o efeito de sentir a informação.
		\item[ (ii) ] \textbf{[ Gama ]} A sensação naquele momento pode se tornar uma nova informação na memória, suficientemente forte para ser possível resgatar na imaginação.
		\item[ (iii) ] \textbf{[ Beta ]} O efeito das múltiplas informações se tornarem poucas pode resultar em ações efetivas na realidade.
		\item[ ] Existe uma ordem de força entre os diferentes efeitos, o mais fraco é a gama, que pode ficar na subconsciência, mas não necessariamente não existe na memória, o segundo é o alfa, que gera a sensação na consciência, e o mais forte dos efeitos é o efeito que gera uma decisão no corpo.
	\end{itemize}
\end{itemize}

%%%%%%%%%%%%%%%%%%%%%%%%%%%%%%%%%%%%%%%%%%%%%%%%%%%%%%%%%%%%%%%%%%%%%%%%%%%%%%%%%%%%%%%%%%%%%%%%%%%%%%%%%%%%
% 

%%%%%%%%%%%%%%%%%%%%%%%%%%%%%%%%%%%%%%%%%%%%%%%%%%%%%%%%%%%%%%%%%%%%%%%%%%%%%%%%%%%%%%%%%%%%%%%%%%%%%%%%%%%%
%%%%%%%%%%%%%%%%%%%%%%%%%%%%%%%%%%%%%%%%%%%%%%%%%%%%%%%%%%%%%%%%%%%%%%%%%%%%%%%%%%%%%%%%%%%%%%%%%%%%%%%%%%%%

\subsection{Depressão e Burnout}

\hspace{\baselineskip}

Este tópico é extremamente importante, a depressão poderá lhe acompanhar ao longo de sua vida, seja porque a sociedade fornece estímulos negativos constantes, seja porque sua vida não corresponde a suas expectativas (este método é um exemplo, o método trialético surgiu para lidar com o desejo de saber mais e com a limitação humana no aprendizado tradicional), seja por ignorância externalizada sob a forma de má alimentação, má cuidado com a saúde, seja por uma rotina de vida irregular, seja porque está cercado de pessoas ruins. Não bastasse isso o ser humano fica naturalmente entediado pela repetição, sua capacidade de abstração permite enxergar padrões repetitivos com maior facilidade, e a vida real exige repetição, estresse físico e psicológico. O que seria Depressão do ponto de vista dialético? \\

% importante para silenciar o mundo espiritual inimigo

%\textbf{Notação [ Relação ]} $ A \sim B $ se existe uma relação entre $A$ e $B$\\

%Vou agora fazer uma série de associações com a palavra depressão e outros conceitos, esse método é uma forma de brainstorm (caso não saiba o que é brainstorm, pesquise!). Brainstorm é basicamente criar um conjunto de estimuladores na esperança de que a mente organize os raciocínios associados.\\

\textbf{Análise Dialética [ Depressão ]}

$$ \texttt{Depressão} \sim \texttt{Consciência} $$
$$ \succ \texttt{Depressão} \bullet \texttt{Trialética ?} $$
$$ \texttt{(Depressão, Ausência de Vontade)} \sim \texttt{Dialética Existencial} \sim \texttt{Vontade} $$
$$ \texttt{(Depressão, Sentimentos Ruins)} \sim \texttt{Dialética Existencial} \sim \texttt{Emoção} $$
$$ \texttt{(Depressão, Pensamentos Ruins)} \sim \texttt{Dialética Existencial} \sim \texttt{Intelecto} $$

\hrulefill

Depressão age na consciência, mesmo que depressão tenha ligação física (no corpo, assim como a própria consciência), a depressão é sentida pela consciência e pode ser externalizada usando sua base natural (a trialética). Assim temos ausência de vontade e vontades ruins(como desejos autodestrutivos) associado a vontade (componente $\beta$ da trialética), temos sentimentos ruins, como tristeza, raiva, ansiedade associados a emoção (componente $\alpha$ da trialética) e temos pensamentos negativos, pensamentos que buscam diminuir sua autoestima, pensamentos que buscam relembrar eventos passados desagradáveis, pensamentos de inveja, estes associados ao intelecto (componente $\gamma$ da trialética).\\

Uma característica que talvez tenha passado desapercebido é que a base trialética é existencial entre si, cada componente contribui para existência da outra e juntas conferem estabilidade entre si, diferente da dialética de contraste, em que uma é posta em confronto com a outra. Assim na depressão expressa da forma trialética, a hipótese natural é que cada um dos componentes contribuem entre si e talvez sejam essenciais para definir o conceito representado. Ter ausência de vontade, ou sentimentos ruins ou pensamentos ruins de maneira isolada podem não ser suficientes para considerar que alguém tenha depressão, mas ter de forma recorrente os três elementos e os mesmos corroborando entre si para se sustentar na consciência é suficiente para dizer que alguém tenha depressão.

\hspace{\baselineskip}

É fatídico que a má saúde do corpo físico irá provocar um mal funcionamento da mente, assim a primeira suspeita que deve se ter é de que o corpo não está sendo privado de algum elemento essencial. Como se trata de filosofia e não medicina, vamos nos ater as causas que estão mais próximas da mente, tendo o cuidado de dizer que ambas as causas devem ser observadas, pois são existenciais entre si (corroboram entre si para coexistência). \\

\textbf{Análise Dialética [ Causas da Depressão ]}

$$ \texttt{Monotonicidade} \sim \gamma $$
$$ \texttt{Traumas} \sim \alpha $$
$$ ( \texttt{Esforço} \times \texttt{Resultado} ) \sim \beta $$
$$ \{ \texttt{Monotonicidade}, \texttt{Esforço} \times \texttt{Resultado} \} \sim \texttt{Falta de Entusiasmo} $$ 

\hrulefill

O monótono é desestimulante para o intelecto, por isso está associado a uma possível causa do tipo gama, traumas naturalmente afetam a parte alfa da consciência e o contraste entre o esforço e o resultado, neste caso, a sensação de muito esforço e pouco resultado desestimula a parte beta da consciência. Todos estes fatores parecem ter uma causa eficiente associada, é preferível que o cérebro armazene informações úteis, logo parece razoável que sua natureza desestimule repetições, muito esforço para pouco resultado indica ineficiência, logo parece razoável que sua natureza desestimule ações ineficientes, e evitar eventos traumáticos e experiências ruins apesar de mais complexo são também preferíveis, neste sentido a depressão agiria como um sinal natural para que algo seja mudado no seu comportamento.

$$ \texttt{Disciplina} \sim \{ \texttt{Monotonicidade} ,\texttt{Esforço}\times \texttt{Resultado} \} $$
$$ \texttt{Mente Forte} \sim \texttt{Traumas} $$

Estas palavras do senso comum parecem induzir uma forma de treinamento, ou um condicionamento da mente para poder suportar estas possíveis causas. Talvez uma das formas de condicionar é ter o conhecimento destes problemas, outra forma é tentar manter a continuidade de ações importantes.

\hspace{\baselineskip}



%Aparentemente há uma clara externalização da depressão por meio da vontade e da emoção, mas e quanto ao intelecto? O intelecto pode sofrer do mal da ignorância, do mal julgamento, da soberba, mas será que há depressão intelectual? Talvez a depressão intelectual esteja associado ao medo, ao chamado bloqueio mental, mas a ausência de vontade e as sensações ruins são as externalizações da depressão mais desgastantes, pois como não estão no intelecto, aparentemente não fazer sentido. Vamos considerar então primeiro a depressão externalizada através da ausência de vontade e das sensações ruins.\\

%\textbf{Questão [ Externalização da Depressão ]} Existe relação entre sensações ruins (emoção) e ausência de vontade? \\

%Essa pergunta foi de certa maneira respondida, sensações ruins podem levar a ausência de vontade, entretanto se as sensações ruins forem exclusivas do mundo externo então o ser humano seria incapaz de resolver problemas, o problema maior é quando as sensações ruins estão no interior, por causas desconhecidas.\\

%\textbf{Questão [ Ações sem Sentido ]} A depressão é externalizada na vontade apenas pela ausência de vontade?\\

%Podemos ter vontades ruins, por exemplo: suicídio, homicídio, automutilação, degradação do próprio corpo, isolamento. Estes são ações que externalizam a depressão, portanto uma definição mais precisa da externalização da depressão é a ausência de vontades consideradas boas pela própria consciência.\\

\textbf{Questão [ Estudo e Depressão ]} Vamos agora particularizar o problema, dentro do contexto de estudo, o que pode causar depressão (externalizando o conceito de depressão do ponto de vista trialético) ?\\

O que pode causar ausência de vontade e o que pode causar sensações ruins durante o estudo? Sabemos que o esforço sem resultado pode inibir a vontade, se o estudo for ineficiente e trabalhoso então irá inibir a vontade. Você precisa sentir que seu esforço está valendo a pena. Também sabemos que o esforço excessivo sem descanso pode esgotar o corpo, isso pode levar a sensações ruins, portanto é essencial tempos de descanso, tempos de lazer e distração.\\

%Existe um problema na sociedade moderna, por exemplo a ciência incremental, ciência incremental é a ausência de grandes descobertas, a ciência se desenvolve apenas com pequenas contribuições, entretanto o esforço para uma pequena contribuição é muito grande. Neste caso é preciso de uma mudança de perspectiva, a contribuição incremental de cada indivíduo da sociedade pode ser pequena se olhado individualmente, mas se olhado como um todo pode ser relevante, talvez o seu problema seja o desejo por trabalhar sozinho, talvez seja porque você não confie nas pessoas. \\

\textbf{Questão [ Asceticismo e Esforço ]} O excesso de esforço sem resultado pode causar depressão, mas uma vida de prazer pode levar também a depressão? A ausência de resultados, a ausência de conquistas pode levar a depressão?\\

Se uma pessoa gosta de sorvete e come todos os dias, todo tempo sorvete, do mesmo tipo, o do tipo preferido, é provável que ele deixe de gostar. A monotonicidade, a ausência de novidades pode provocar ausência de desejo, o saciamento do desejo também pode inibir o desejo. Mesmo uma vida rica e diversa irá se tornar monótona por que o cérebro irá atribuir similaridade ao comportamento, se a vida se estabilizar em algum padrão ela se tornará monótona, e ausência de padrões é algo muito difícil de acontecer, a realidade tem padrões, a realidade é limitada em novidades (ou a nossa realidade próxima, pois o universo é extenso e inacessível). 
Agora vem o pulo do gato, tenho como hipótese que exista um equilíbrio ótimo entre o Asceticismo e o Esforço para controlar a depressão. A pergunta que fica é onde se situa o equilíbrio, qual a proporção de lazer e qual a proporção de esforço é o ideal. Segundo pulo do gato, para mim não existe proporção ideal, ela é passível de mudança e deve ser controlado por meio de alguns sinais externos, se você se sente cansado, se você acha que está se esforçando muito então aumente o tempo de lazer, se você sente que você não sente mais prazer com as atividades que você gostava, reduza o tempo, reduza a frequência e aumente o esforço em sua vida. É preciso lembrar que você pode ganhar seu dia apenas com um segundo, então minimizar o lazer não é uma coisa ruim para o ser humano, algumas vezes pode ser bom, mas a ausência completa pode levar a depressão. Talvez a pessoa que consiga maximizar o esforço e minimizar o prazer seja uma pessoa mais feliz, pois irá realizar mais coisas sem entrar em depressão.\\

%Lembre-se um sorriso dura segundos na realidade, mas vale um dia inteiro. Um grande resultado vale um grande esforço. Não é errado aproveitar a vida. O ser humano busca se afirmar, busca existir para os outros, busca sentido em sua existência, não seja medíocre em sua vontade, mas seja ousado a sua maneira. Não consuma sua vida com sonhos irrealizáveis, crie um caminho realista e objetivo. Não externalize sua incapacidade pela desistência, mas externalize sua determinação com seu esforço. Não se esforce em algo que não funciona, busque soluções pela sua criatividade. Não use métodos adaptados para os outros, adapte o métodos dos outros a sua realidade. Não ignore sua intuição, mas também não abandone a objetividade. Questione, mas se permita explorar caminhos diferentes. Não se importe em fazer valer a pena sua vida se a falta de resultados estiver lhe deixando depressivo, ansioso ou excessivamente cansado, aproveite a vida em seus pequenos detalhes, em seus pequenos momentos e recomponha suas energias, mas não abandone o esforço. Se o mundo estiver lhe fazendo mal, se o mundo estiver tornando turvo sua intuição, dê um tempo para si mesmo, medite, dê um tempo para tudo que lhe faz mal, algumas vezes precisamos escrever em um papel em branco para resolver os problemas. Todos estes conselhos são conselhos de equilíbrio, o antagonismo entre os conselhos deve ser construtivo e não destrutivo, ele deve contribuir para sua existência e não dominar sua consciência. 

%\subsection{Consciência}

%\hspace{\baselineskip}

%\textbf{Questão [ Consciência ]} Supostamente entendemos que a memorização se constitui de três processos (recepção, armazenamento, acesso), sendo mais rigoroso, temos uma ideia fraca do que seria, temos uma heurística para podermos trabalhar. Mas como realmente funciona estes três processos?

%\begin{itemize}
%	\item[ (i) ] \textbf{Questão [ Acesso ]} Que entidade está acessando o conteúdo na memória?
%	\item[ (ii) ] \textbf{Questão [ Recepção e Armazenamento ]} Que entidade está controlando a recepção e o armazenamento?
%\end{itemize}

%Temos uma noção de que a mente não se constitui de algo abstrato, mas de algo concreto chamado cérebro e de algo abstrato que podemos chamar de consciência. O acesso estaria mais próximo da consciência, a entidade mais interior da existência humana, enquanto a recepção e o armazenamento estaria mais próximo do cérebro. Ou poderíamos afirmar que temos a consciência, uma existência virtual, abstrata que se comunica através de um mesmo processo de recepção e resposta com uma entidade concreta que armazena o conteúdo, o cérebro, as conexões entre os neurônios.

%$$ \texttt{Consciência} \bullet \texttt{Cérebro} \bullet \texttt{Corpo} \bullet \texttt{Realidade} $$

%Isto é uma organização topológica das entidade Consciência, Cérebro, Corpo e Realidade. Ela representa a topologia, a relação de proximidade entre as entidades. Para representações mais gerais podemos usar o grafo, que basicamente é uma coleção de entidades e uma coleção de relações entre as entidades.

%\hspace{\baselineskip}

%\textbf{Notação [ Interação ]}

%$$ \texttt{Objeto A}  \left( \begin{array}{c}
%\texttt{Recepção sob o Ponto de Vista de A} \\ \texttt{Ação sob o Ponto de Vista de A}
%\end{array} \right) \bullet \texttt{Objeto B} $$

%$$ \texttt{Objeto A}  \left( \begin{array}{c}
%\texttt{Recepção sob o Ponto de Vista de A} \\ \texttt{Ação sob o Ponto de Vista de A}
%\end{array} \right) \bullet \texttt{Base} \bullet \texttt{Objeto B} $$

%A notação segue a regra da mão direita o de cima é para esquerda e o de baixo é para direita no conteúdo dos parênteses, ele representa a interação do objeto A com B sob o ponto de vista de A.

%$$ \texttt{Objeto A} \bullet \left( \begin{array}{cc}
%\texttt{Ação de B} \\ \texttt{Recepção de B}
%\end{array} \right) \texttt{Objeto B}  $$

%A notação acima também segue a regra da mão direita, note que a relação se inverte.

%%%%%%%%%%%%%%%%%%%%%%%%%%%%%%%%%%%%%%%%%%%%%%%%%%%%%%%%%%%%%%%%%%%%%%%%%%%%%%%%%%%%%%%%%%%%%%%%%%%%%%%%%%%%
% Análise dialética da Consciência
%\hspace{\baselineskip}

%\textbf{Análise Dialética [ Consciência ]}

%$$ C := \texttt{Consciência}$$
%$$ R := \texttt{Realidade} $$
%$$ M := \texttt{Memória} $$
%$$ C \left(  \begin{array}{cc}
%\texttt{Percepção} \\ \texttt{Vontade}
%\end{array} \right)  \bullet \left( \begin{array}{cc}
%\texttt{Sentidos Captados} \\ \texttt{Alteração da Realidade}
%\end{array} \right) R $$
%$$ C \left(  \begin{array}{cc}
%\texttt{Imaginação} \\ \texttt{Acesso:Informações}
%\end{array} \right)  \bullet \left( \begin{array}{cc}
%\texttt{Envio:Informações} \\ \texttt{Memorização}
%\end{array} \right) M $$
%\hrulefill

%\textbf{Questão [ Intermediário ]} Existe uma pergunta topológica natural neste modelo, será que a memória atua conectando a realidade com nossa consciência como um intermediário exclusivo ou será que em nossa arquitetura cerebral temos duas vias de comunicação, uma via que se comunica diretamente com a realidade e uma via que se comunica com a memória e possibilita a imaginação?

%%%%%%%%%%%%%%%%%%%%%%%%%%%%%%%%%%%%%%%%%%%%%%%%%%%%%%%%%%%%%%%%%%%%%%%%%%%%%%%%%%%%%%%%%%%%%%%%%%%%%%%%%%%%

%\subsection{Raciocínio, Evolução e Fundamento de Quatro Posições}

%\hspace{\baselineskip}

%\textbf{Questão [ Raciocínio e Memorização ]} Existe raciocínio sem memorização?\\

%\textbf{Questão [ Substrato do Raciocínio ]} O raciocínio ocorre somente na imaginação ou ele ocorre também na realidade?\\

%%%%%%%%%%%%%%%%%%%%%%%%%%%%%%%%%%%%%%%%%%%%%%%%%%%%%%%%%%%%%%%%%%%%%%%%%%%%%%%%%%%%%%%%%%%%%%%%%%%%%%%%%%%%
% Evolução
%\textbf{Base Dialética [ Evolução ]} Temos a evolução como uma estrutura de três componentes (Existências, Transformações, Sentido), a existência é o objeto que evolui, as transformações são as possíveis mudanças nestes objetos, e o sentido é a direção da mudança. A transformação conecta uma existência com outra e possui uma direção. A evolução é um processo temporal, ela precisa de um substrato espacial, mas essencialmente ela é temporal. Algo somente pode evoluir se tiver estabilidade, a verdade é que algo somente existe se tiver alguma estabilidade de forma que seja possível estabelecer propriedades que externalizem sua existência para um ser consciente.

%$$ \texttt{(Informação, Raciocínio, Verdade)} \sim \texttt{Evolução} $$

%\textbf{Notação [ Existencial ]} (Caráter Interno/Natureza Interna)\{Forma Externa\} \\

%\textbf{Análise Dialética [ Consciência, Cérebro, Memória, Realidade, Raciocínio, Imaginação ]}

%$$ \mathcal{C} := \{ \texttt{Consciência, Cérebro, Memória, Realidade, Raciocínio, Imaginação} \} $$
%$$  \succ \mathcal{C} \bullet \{ \texttt{Fundamento de Quatro Posições} \} ? $$
%$$ FQP := \texttt{Fundamento de Quatro Posições} $$
%$$ \mathcal{C} \prec \texttt{(Natureza Interna, Corpo, Realidade, Memória)} \bullet FQP $$
%\hrulefill

%A natureza do ser humano impele o cérebro (componente do corpo) a interagir com a realidade através do corpo, o resultado da interação estável de dar e receber com a realidade cria um novo centro, uma nova existência chamada de memória, ela é externalizada pelas conexões que os neurônios fazem para associar diferentes locais da memória.

%$$ \texttt{(Natureza Interna, Memória, Realidade, Consciência da Realidade)} \sim FQP $$

%\hrulefill

%Uma vez existindo a memória pode interagir com as outras existências. A natureza interna estimula a relação entre a memória e a realidade, o estabelecimento de uma relação estável de dar e receber cria uma nova entidade, a consciência da realidade, a percepção da realidade. Supostamente todos os animais possuem esse grau de consciência. Podemos chamar a consciência da realidade de percepção.

%$$ \texttt{(Natureza Interna, Corpo, Memória, Imaginação)} \sim FQP $$

%\hrulefill 

%A relação estável entre a memória e o cérebro cria outra forma de consciência, a consciência de nossas experiências armazenadas na memória, que chamo de imaginação. Podemos agora definir a consciência de duas maneiras, a consciência é a resultante da interação dual entre (Imaginação, Percepção) que seria uma existência dialética ou podemos definir a consciência com uma existência trialética, em que as três entidades interagem dualmente entre si formando um equilíbrio entre três entidades o equilíbrio entre (Corpo, Realidade, Memória).

%Com a consciência temos dois substratos naturais para evolução da informação (raciocínio), o substrato da realidade, ou seja, o raciocínio é feito/acompanhado através da transformação dos objetos da realidade, e o substrato da imaginação em que a memória permite simular a realidade e o raciocínio se dá na memória.

%Agora o suprassumo da consciência:

%$$ \texttt{(Consciência, Raciocínio da Imaginação, Raciocínio na Realidade, Linguagem)} \sim FQP $$

%\hrulefill

%Ora temos agora um caminho natural para o desenvolvimento da linguagem, ela seria um intermediário resultante do raciocínio feito na imaginação e o raciocínio feito na realidade. Essa é uma leitura da linguagem de forma que seja utilizada para efetuar raciocínio lógico, mas a linguagem é bem mais que isso, pois está associado a interação entre seres humanos.\\

%É importante perceber que a ação de dar e receber requer tempo, pois requer um processo de interação até chegar a estabilidade, assim a consciência requer tempo até que fique estável, mas existe um processo que é o mais importante para o método trialético de estudo, a informação resgatada da memória requer tempo para ficar estável no substrato criado da consciência da imaginação, é importante associar o tempo como uma dificuldade inerente, associar a estabilidade como um processo de disputa entre ideias, o entendimento disso permitirá criar um método eficiente de estudo.\\

%\textbf{Questão [ Existência Trialética ou Dialética ]} A consciência humana é uma existência trialética (estabilidade de três elementos) ou é uma existência dialética (estabilidade de dois elementos)?\\

%\textbf{Análise Dialética [ Informação e Verdade ]}

%$$ \texttt{(Consciência, Informação(Consciência), Fenômeno(Realidade), Verdade)} \sim FQP $$

%\hrulefill

%A verdade é um atributo da informação, a verdade somente faz sentido com algo externo que possa comparar informações contraditórias, este ocupado pela realidade. A interação da consciência com a realidade permite a consciência do conceito de verdade.\\

%\textbf{Questão [ Sonho e Realidade ]} Qual a diferença entre as percepções do Sonho e as percepções da Realidade?\\

%A realidade tem regras mais rígidas em relação a verdade, em relação a direção dos fenômenos, o sonho não necessariamente segue as mesmas regras da realidade, algumas vezes contradiz, algumas vezes não tem sentido algum.\\

%\hspace{\baselineskip}

%\subsection{Escolha e Livre Arbítrio}

%\hspace{\baselineskip}

%\textbf{Questão [ Residência da Consciência ]} Onde reside a consciência? A consciência está mais próxima da realidade, está mais próxima da memória, está mais próxima do corpo?\\

%Ambas as formas de consciência (Percepção, Imaginação) tem uma entidade em comum, o raciocínio, o raciocínio parece ser elementar para consciência, podemos dizer que se a consciência está mais próxima de algum lugar, deve ser do raciocínio. É interessante notar que o raciocínio tem uma direção definida e que não percebemos raciocínios simultâneos, sempre vemos o raciocínio como único, até quando imaginamos processos simultâneos imaginamos eles como se fosse um raciocínio único. Essa unicidade será assumida como característico da consciência. Quem estimula a unicidade da consciência quando desperto é a realidade, quem estimula a consciência quando dormindo seria a memória da realidade, a perda de unicidade do raciocínio provavelmente (e isso é uma das muitas especulações feitas até agora) faz com que percamos a consciência, então temos a seguinte hipótese, quando dormimos perdemos a unicidade, entre o estado desperto e o sono temos uma breve etapa em que imaginamos coisas sem sentido e dormimos, essas coisas sem sentido são consequência da perda da unicidade, antes de acordar, durante o sono profundo temos o sonho que seria a recuperação da unicidade feita pelo cérebro, quando ficamos completamente conscientes durante o sono geralmente acordamos, ou seja, recuperamos a unicidade do raciocínio na consciência.

%Essa hipótese de unicidade do raciocínio é fundamental para o método trialético, pois ela configura uma disputa de ideias, ela permite entender o que é confusão, o que é desconcentração, o que é perda e dificuldade de raciocínio, pois tudo nada mais é do que a perda de unicidade do raciocínio. Portanto a perda de unicidade do raciocínio pode ser uma dificuldade inerente de qualquer arquitetura de consciência, até mesmo se criamos uma inteligência artificial, talvez essa disputa entre os raciocínios para definir o raciocínio dominante seja uma externalização do livre arbítrio, característico da consciência e que perturba e incentiva qualquer filósofo a estudar a mente humana.\\

%\textbf{Paradigma [ Escolha ]} A escolha é uma direção única e definida para realizar uma ação. É intrínseco da consciência a escolha, a liberdade, o livre arbítrio. A escolha, a liberdade talvez seja o real centro da consciência, escolhemos como raciocinar na imaginação e escolhemos como agir (que seria o raciocínio na realidade).\\

%\textbf{Questão [ Raciocínio e Ação ]} Dizemos que precisamos de uma direção única, mas o raciocínio e a ação são feitas simultaneamente ou uma vez fazemos uma e uma vez fazemos outra? Existe uma disputa entre o raciocínio e a ação?\\

%Não tenho segurança em responder a esta pergunta, o que posso dizer que quando nos concentramos em agir estamos de certa maneira raciocinando, e quando nos concentramos em imaginar é como se estivéssemos agindo em um simulador, parece haver uma disjunção mas parece também que estamos percebendo estes fenômenos ao mesmo tempo.\\

%\textbf{Análise Dialética [ Projeção da Consciência em Raciocínio ]}

%$$ \texttt{(Consciência, Direção do Raciocínio)} \sim \texttt{Dialética Existencial} $$

%\hrulefill

%A consciência é externalizada pela direção de raciocínio, ela está mais próxima da consciência.\\

%\textbf{Questão [ Perturbadores ]} O que pode estimular e alterar a direção do raciocínio?\\

%Temos os estímulos da realidade percebida através dos sentidos, temos a sensibilidade interna de nosso corpo correspondente a emoção, temos a memória que influencia na imaginação, ou podemos chamar de intelecto, temos a vontade que está associado com a realidade. Eis que surge naturalmente a trialética humana (Emoção, Intelecto, Vontade) como os pertubadores da consciência, eles direcionam os raciocínios dominantes na consciência e os raciocínios não dominantes (releitura da subconsciência). Emoção está associada as sensações do corpo, aos desejos sensoriais, intelecto está associado a memória, a busca por entender e modelar a realidade, a vontade está na escolha voluntária, nas ações. Emoção e Vontade estão mais próximos da realidade enquanto o intelecto está mais próximo da memória. Podemos dizer que a consciência é o equilíbrio trialético existencial da emoção, do intelecto e da vontade externalizado sob a forma do livre arbítrio. Também podemos dizer que é um equilíbrio dialético existencial entre a realidade (Emoção, Vontade) e a memória (Intelecto).





%%%%%%%%%%%%%%%%%%%%%%%%%%%%%%%%%%%%%%%%%%%%%%%%%%%%%%%%%%%%%%%%%%%%%%%%%%%%%%%%%%%%%%%%%%%%%%%%%%%%%%%%%%%%
%%%%%%%%%%%%%%%%%%%%%%%%%%%%%%%%%%%%%%%%%%%%%%%%%%%%%%%%%%%%%%%%%%%%%%%%%%%%%%%%%%%%%%%%%%%%%%%%%%%%%%%%%%%%
%%%%%%%%%%%%%%%%%%%%%%%%%%%%%%%%%%%%%%%%%%%%%%%%%%%%%%%%%%%%%%%%%%%%%%%%%%%%%%%%%%%%%%%%%%%%%%%%%%%%%%%%%%%%
%%%%%%%%%%%%%%%%%%%%%%%%%%%%%%%%%%%%%%%%%%%%%%%%%%%%%%%%%%%%%%%%%%%%%%%%%%%%%%%%%%%%%%%%%%%%%%%%%%%%%%%%%%%%
%%%%%%%%%%%%%%%%%%%%%%%%%%%%%%%%%%%%%%%%%%%%%%%%%%%%%%%%%%%%%%%%%%%%%%%%%%%%%%%%%%%%%%%%%%%%%%%%%%%%%%%%%%%%

%\section{Motivação e Propósito}

%\hspace{\baselineskip}

%Este paradigma foi feito (ou refeito por mim, eu não sei, reflita sobre isso) no século 21, mas o problema que motiva a criação de novos métodos de estudo acompanha o homem estudioso desde algum tempo, talvez desde que o homem começou a querer documentar as informações da natureza por meio da escrita e por meio de símbolos. 

%O estudioso já deve ter notado que o seu cérebro é limitado na capacidade de assimilar novas informações, e o seu desejo por entender é confrontado com a realidade de que seu cérebro não é o melhor método de armazenar conteúdos, seu cérebro erra, seu cérebro tem velocidade de assimilação limitada, seu cérebro tem velocidade de acesso limitada, seu cérebro não foi feito para lidar com excesso de informações. 

%Para você leitor que assim como eu recebeu a dádiva de ter uma memorização fotográfica de quatro pixels ou menos, mas um desejo ousado por entender talvez toda a natureza (no sentido de incluir também a cultura humana parte da natureza) então este método é para você, não é um método para pessoas medíocres, mas também não é um método para prodígios (nada impede que seja usado, ora pois).\\

%\textbf{Motivação [ Método Trialético de Estudo ]} Excesso de informação e limitação humana.\\

%Temos duas soluções para este problema, aceitar a limitação humana e usar a estratégia de dividir para conquistar, cada um se especializa como conseguir o seu cérebro, o que pode levar a possível dificuldade de comunicação entre áreas muito distintas, mas ainda é uma estratégia realista e eficiente. Outra solução é tentar lidar com o excesso de informações enfrentando o problema com métodos mais eficientes de estudo.\\

%\textbf{Propósito [ Método Trialético de Estudo ]} Assimilar e entender teorias com conteúdos difíceis de memorizar, seja por excesso de informações importantes, seja por similaridades entre as informações importantes que podem levar a frequência de erros e instabilidade, seja porque você quer saber com detalhes um conteúdo qualquer mesmo que isso te leve a uma crise existencial.\\

%É importante destacar que não é exclusividade de nenhuma área mais do intelecto humano o excesso de informações, muito se engana aquele que busca na matemática uma teoria completamente racional, porque irá ser confrontado com a necessidade de memorizar teoremas, como se não bastasse também entende-los e aplica-los de forma correta, isso vale para física, para engenharia, para química, para biologia, para história, para filosofia. 

%Também é importante ter a consciência de que vivemos um mundo competitivo, é necessário buscar se sobrepor em qualidade ao outro para estar em uma posição de destaque, estar nela sem ter devido conteúdo é arrogar para si um direito do qual não é digno, podemos nos questionar se existiria um sistema diferente, mas até mesmo para propor um novo sistema é necessário estudar com detalhes os outros sistemas, é necessário ter a argumentação necessária para poder criar um sistema confiável que possa substituir o sistema atual, tentar fazer sem estudo, sem esforço não é somente arrogância, é irresponsabilidade. Por hora viver fora da mediocridade é uma necessidade convergente.

%É seguro dizer que o desenvolvimento intelectual está intimamente ligado a escrita que foi um método mais eficiente de armazenar informações sobre o mundo, a eficiência que se traduz em estabilidade da informação (informações mais precisas) e a facilidade de acesso.







%%%%%%%%%%%%%%%%%%%%%%%%%%%%%%%%%%%%%%%%%%%%%%%%%%%%%%%%%%%%%%%%%%%%%%%%%%%%%%%%%%%%%%%%%%%%%%%%%%%%%%%%%%%%
%%%%%%%%%%%%%%%%%%%%%%%%%%%%%%%%%%%%%%%%%%%%%%%%%%%%%%%%%%%%%%%%%%%%%%%%%%%%%%%%%%%%%%%%%%%%%%%%%%%%%%%%%%%%
%%%%%%%%%%%%%%%%%%%%%%%%%%%%%%%%%%%%%%%%%%%%%%%%%%%%%%%%%%%%%%%%%%%%%%%%%%%%%%%%%%%%%%%%%%%%%%%%%%%%%%%%%%%%
%%%%%%%%%%%%%%%%%%%%%%%%%%%%%%%%%%%%%%%%%%%%%%%%%%%%%%%%%%%%%%%%%%%%%%%%%%%%%%%%%%%%%%%%%%%%%%%%%%%%%%%%%%%%
%%%%%%%%%%%%%%%%%%%%%%%%%%%%%%%%%%%%%%%%%%%%%%%%%%%%%%%%%%%%%%%%%%%%%%%%%%%%%%%%%%%%%%%%%%%%%%%%%%%%%%%%%%%%\\

\section{Método Trialético de Estudo}

\hspace{\baselineskip}

Esta seção tem um propósito mais prático com relação as outras subseções, não estaremos tentando modelar conceitos transcendentais. O objetivo desta seção é fornecer um método para aperfeiçoar a consciência pelo conhecimento diferente do método tradicional dialético para um método que utiliza a base trialética.

%\hspace{\baselineskip}

%\textbf{Questão [ Teoria ]} O que é uma teoria?\\

%Uma pergunta desse tipo para um filósofo dialético é o mesmo que induzir ao filósofo projetar o conceito "teoria" em alguma base dialética. Vamos nos contentar em dizer que uma teoria é uma trinca (Definições, Informações, Raciocínios), uma teoria está externalizada sob a forma de uma linguagem assim definições e informações devem estar externalizadas em forma de linguagem. Uma definição é composta de um nome e de uma coleção de propriedades, uma definição é sempre verdadeira, pois é uma criação, uma informação associa definições, elas podem ser verdadeiras ou falsas, propriedades são informações, uma definição é uma informação que associa um nome a propriedades, raciocínio conecta informações, um raciocínio pode ser visto também como informações e assim pode ser verdadeiro ou falso.\\

%\textbf{[ Interpretação de Teoria como Grafo ]} Um grafo é uma coleção de nodos e uma coleção relações entre os nodos, vamos considerar uma teoria como uma coleção de nodos de informações e que as relações são raciocínios ou paradigmas que conectam as informações. É importante perceber que o raciocínio também é uma informação, assim representar uma teoria por meio de grafos é arbitrário.\\

%\textbf{[ Grafo Existencial de Teoria ]} Aqui temos uma teoria como uma seleção de informações importantes, a função de uma teoria sob este ponto de vista é enfrentar problemas com tais informações. As relações que podem existir são, relação de cada teoria com um tópico ou categoria (por exemplo o nome da teoria, ou alguma subcategoria), o nome da categoria se relaciona com as informações presentes na teoria, esses nomes são estimuladores da consciência para lembrar das informações da teoria. O outro tipo de relação é a relação que existe entre as informações das teorias, suponha que força de existência é sua capacidade de influenciar o raciocínio principal de sua consciência (podemos chamar de força de consciência), tal relação é externalizada em quanto uma informação contribui para existência da outra informação. Existem tipos importantes de Grafos Existenciais para a memorização, a primeira é a linear, a segunda é a uniforme, e a terceira é a topologia de praticidade (ou de aplicação).

\subsection{Método Tradicional de Estudo}

\hspace{\baselineskip}

\textbf{[ Método Tradicional de Estudo ]} O método tradicional de estudo(ou método dialético de estudo) consiste em duas etapas bem definidas, a primeira é a exposição do conteúdo feita por um orador ou algo que mimetize um orador (por exemplo um texto) e a segunda etapa é a exercitação do conteúdo através da resolução de exercícios e problemas.\\

\textbf{Paradigma [ Trialética de Análise de Exposição ]} Suponha que você seja um ouvinte de uma palestra ou de um discurso e tem como objetivo questionar os conhecimentos e a lógica do expositor, este método simples consiste em fortalecer a consciência alfa durante a exposição permitindo que sua consciência possa explorar os conteúdos novos sem se perder durante o discurso.

\begin{itemize}
	\item[ (i) ] \textbf{[ Contexto ]} okayA primeira etapa consiste em vocalizar palavras chaves formando uma pilha mais ou menos ordenada de palavras, por exemplo se o título conter a palavra gato e o expositor falar sobre a relação entre o gato e o cachorro e entre o gato e o homem, temos as definições (gato,cachorro,homem), chamamos estas definições de contexto, a medida que outras informações são adicionadas, tentamos repetir vocalizando a última definição adicionada e depois pelo menos três outras definições dentro do contexto, por exemplo, se a próxima informação for os gatos não gostam de ficar dentro de casa, então vocalizamos (casa, homem,cachorro, gato).
	\item[ (ii) ] \textbf{[ Inventário ]} O inventário é uma tentativa de enumerar informações distintas sobre o que está sendo ouvido, é feito uma lista das informações utilizando o contexto como indutor da lembrança.
	\item[ (iii) ] \textbf{[ Dialética ]} A partir do contexto e do inventário de informações já é possível utilizar a sensibilidade a padrões ou contradições para fazer uma dialética socrática do assunto. Perguntas pertinentes devem ser feitas para o expositor.
\end{itemize}

\textbf{[ Método Trialético de Estudo ]} A primeira etapa é decomposta em duas, a memorização e a análise, enquanto a terceira etapa é a de exercitação, similar ao método tradicional.\\

O método tradicional de estudo está voltado ao diálogo entre o expositor e o ouvinte, a medida que o conhecimento começa a se tornar vasto e específico pelo desenvolvimento cultural, científico e tecnológico a consciência das informações apresentadas pelo expositor torna a compreensão do raciocínio aplicado ineficiente. Portanto a comunicação entre o expositor e o ouvinte torna-se inefetiva. Para contornar este problema o expositor deve reduzir o nível do seu conteúdo ou o ouvinte deve estudar previamente o assunto antes da interação. Em vez disso o método trialético de estudo busca usar o modelo natural da consciência (que é a trialética) para embasar o processo de estudo. \\

\textbf{Notação [ Pertubação ]} 

$$ A \overset{ativa}{\rightarrow} B $$

Significa que A contribui positivamente para existência e/ou estabilidade de B, B tende a existir com A.

$$ A \overset{inibe}{\rightarrow} B $$

Significa que A contribui negativamente para existência e/ou estabilidade de B, assim B tende a não existir com A.

\hspace{\baselineskip}

\textbf{Análise Dialética [ Excesso de Informações ]}

$$  (\texttt{Excesso de Informações} \overset{inibe}{\rightarrow} \texttt{Memorização} ) \overset{inibe}{\rightarrow} \texttt{Consciência das Informações}  $$

$$ ( (\overset{inibe}{\rightarrow} \texttt{Consciência das Informações}) \overset{inibe}{\rightarrow} \texttt{Análise} ) \overset{inibe}{\rightarrow} \texttt{Aplicação}  $$

\hrulefill

O excesso de informações inibe o aprendizado, pois cria um efeito em cadeia inibitório que resulta em uma capacidade de aplicar o conhecimento ineficiente, como a externalização do aprendizado é a capacidade de aplicação então podemos dizer que uma capacidade deficiente de aplicação é seguramente um indicador de aprendizado deficiente.\\

\textbf{Paradigma [ do Aprendizado Natural ]} O aprendizado natural se dá durante a exposição, o processo de memorização e análise é feito ao mesmo tempo, não há uma consulta prévia do conteúdo aprendido, ou seja, o conteúdo da exposição é novidade.\\

\textbf{Questão [ Problema do Aprendizado Natural ]} O processo de memorização e análise é feito ao mesmo tempo de maneira livre durante a etapa de exposição do método tradicional (aprendizado natural). Quais problemas podem surgir quando estamos em um contexto de excesso de informação?\\

No aprendizado natural a compreensão do conteúdo é externalizada na capacidade de interação entre o expositor e o aprendiz, ou seja, o aprendiz tem capacidade de formular questões, dúvidas e proposições (Quando não há um expositor, mais o conteúdo é exposto em um livro a externalização se dá na capacidade de aplicar a dialética socrática). Quando há um excesso de informações há inibição da capacidade analítica, ou seja, o aprendiz não consegue formular questões e proposições porque não tem consciência do conteúdo, pois a taxa de informações é maior que a capacidade de se fixar na memória e assim poder ser usado pela consciência (raciocínio).\\

Para resolver este problema usamos o paradigma de dividir para conquistar, separamos o aprendizado natural em memorização e análise. Vamos definir que memorização seria artificial no sentido de que é penoso e trabalhoso, e a análise e aplicação seria natural no sentido que é prazeroso e mais fácil.\\

\textbf{Paradigma [ Grafos e Decomposição do Aprendizado Natural ]} Supondo a representação de uma teoria em grafos (redes) e lembrando que categorizar em nodo(ponto) ou em relação(seta) uma informação é arbitrário, podemos usar um critério de escolha dos nodos e das relações, a definição dos nodos são as informações que serão memorizadas na etapa de memorização, a definição das relações são as informações secundárias resultantes da análise do conteúdo memorizado (dialética socrática).\\

\textbf{[ Grafo Existencial ]} Vamos agora abordar outra forma de ver uma teoria, suponha que os nodos sejam informações, pela digressão dialética da consciência nas seções anteriores uma informação para estar na consciência deve interagir e tornar-se estável, isso é resultado da relação entre a consciência e a memória. Também foi dito que a consciência é caracterizada por uma unicidade da direção do raciocínio, e que as informações tem a capacidade de perturbar a direção do raciocínio. Ora, temos que o raciocínio liga duas informações e que as informações tem a capacidade de pertubar a direção do raciocínio, uma conclusão possível é que uma informação contribui para existência de outra informação e isso é externalizado como raciocínio. Portanto estas relações de contribuições de existências entre as informações são as relações do grafo existencial. Chamamos de grafo existencial um grafo em que os nodos são existências e que as relações são as relações de contribuições (estabilidade e força de existência) entre as existências.\\

\textbf{[ Topologia Linear ]} É um grafo existencial em que as informações formam uma cadeia Linear de contribuições de existências.\\

\textbf{[ Topologia Uniforme ]} É um grafo existencial completo em que as informações estimula-se entre si com mesmo peso (mesma força de contribuição).\\

\textbf{[ Topologia de Aplicação ]} É um grafo existencial que tem pesos de relações diferentes para diferentes conjuntos de informações, tais relações são mais fortes quando uma informação e outra é estimulada ao mesmo tempo durante uma determinada aplicação. \\

\textbf{[ Topologia de Praticidade ]} É um grafo existencial que as relações existenciais são pela praticidade durante a aplicação dos conteúdos. É uma forma de topologia de aplicação particular para resolver problemas rapidamente.\\

\subsection{Topologia Linear, Topologia Uniforme e Topologia de Aplicação}

\hspace{\baselineskip}

\textbf{Questão [ Topologias de Existências ]} Supondo que uma teoria seja representada por um grafo existencial, qual topologia é ideal para armazenar uma teoria?\\

Uma das respostas mais naturais seria a topologia de aplicação ou a topologia de praticidade, pois o propósito de uma teoria é ser aplicado na realidade e queremos uma aplicação eficiente. Seria bom se a topologia de aplicação fosse tangível, a verdade é que não sabemos bem quando pode ser necessário um determinado padrão de raciocínio, pois não conhecemos todos os possíveis problemas da realidade, a topologia de praticidade pode ser herdada pela execução de exercícios. Suponha que o objetivo seja a capacidade de aplicar o conteúdo independente da eficiência, ora se não sabemos o que encontrar então o melhor a fazer é preservar um padrão uniforme, assim teríamos todas as possibilidades e poderíamos em tese enfrentar todos os problemas possíveis, a dificuldade inerente é que o esforço para memorização da topologia uniforme é maior que os outros, pois o grafo é completo, também que a contribuição uniforme pode gerar o efeito de confusão, ou seja, não há uma direção do raciocínio predominante e portanto há inibição da consciência externalizada pelo raciocínio. A topologia linear é a última topologia que queremos, primeiro pelo acesso, como o a estrutura é linear o acesso a seu conteúdo não é uniforme e o seu custo é proporcional a quantidade de informações (leia sobre memória ram em ciência da computação, leia também sobre estruturas de dados, listas encadeadas e arrays), como estamos lidando com o problema do excesso de informações isso é horrível, na prática a estrutura linear é como se não tivéssemos aprendido nada.\\

\textbf{Questão [ Topologia Linear ]} Qual topologia é preservada na leitura?\\

Um livro lido de ponta a ponta preserva uma topologia existencial linear (que é horrível), entretanto durante o sono essa topologia pode se alterar (durante o sono outras conexões são formadas), assim a leitura linear não necessariamente é ruim se você tiver uma boa noite de sono. Se você for um músico, sabe que a topologia linear é interessante no aprendizado de uma música, pois uma música é essencialmente linear (mas o aprendizado é diferente, no aprendizado de uma música muito difícil é necessário decompor a música em partes e depois juntar as peças).\\

\textbf{Questão [ Dispersão Temporal e Topologia Linear ]} Qual topologia é preservada se o aprendizado se dá durante longos períodos de tempo?\\

Sabemos que o tempo é essencialmente linear, então a topologia natural do tempo é a topologia linear, um conteúdo disperso em longo período de tempo irá preservar uma topologia linear, a real questão é se o sono pode reverter essa situação ou simplesmente precisamos de outra abordagem, mais artificial para modificar a topologia existencial. Um estudante de graduação deve estar ciente deste problema, pois se seu curso tiver muito conteúdo você chegará em períodos avançados da graduação com uma topologia linear e o mundo exigirá de você uma topologia uniforme ou de aplicação. \\

\textbf{Questão [ Tempo e Espaço ]} Vemos que o tempo tem caráter linear, em contraste, qual a topologia é natural para o espaço?\\

A consciência espacial requer a consciência de simultaneidade, poderíamos dizer que seria mais uniforme que o tempo, mas não seria totalmente uniforme porque nosso olho tem foco, não é um sensor multidirecional. Isto é importante, o espaço permite criar conexões mais uniformes, por causa da sua natureza simultânea, assim uma boa ferramenta para modificar e uniformizar a topologia do conhecimento é projetar o conhecimento em uma linguagem e coloca-lo no espaço (em um papel ou na tela do computador), isso permite que você junte caminhos de forma artificial e fortaleça com o uso da realidade e do sentido de visão.\\

\subsection{Esquecimento e Desconcentração}

%%%%%%%%%%%%%%%%%%%%%%%%%%%%%%%%%%%%%%%%%%%%%%%%%%%%%%%%%%%%%%%%%%%%%%%%%%%%%%%%%%%%%%%%%%%%%%%%%%%%%%%%%%%%
% 
\hspace{\baselineskip}

A análise do esquecimento é importante para criação de um método de estudo, pois não queremos um armazenamento instável de conteúdo, queremos ter a certeza de que a memorização e o entendimento podem ser acessados quando necessário e que são estáveis e precisos, algo que é difícil de obter com o cérebro humano.\\

\textbf{Questão [ Memória ]} Como funciona o acesso da memória pela consciência? Como funciona a memorização?\\

\textbf{Análise Dialética [ Memória ]}

$$ \texttt{Memória} \prec \texttt{(Conexões, Neurônios)} \bullet \texttt{Dialética Existencial} \bullet \texttt{Cérebro} $$

\hrulefill

Perguntar sobre como funciona a memorização é equivalente (externalizado) em como nosso cérebro forma conexões. Sabemos que memorizamos algo se podemos desencadear a informação completa através do estimulo de partes da informação. Sabemos que memorizamos um processo se podemos reviver o processo através do estado inicial do processo. Estas duas afirmações não estão aí por acaso, uma está associado ao espaço e outro está associado ao tempo. Portanto a externalização de uma informação temporal ou espacial se dá pela capacidade estímulos de parte da informação desencadear a informação completa.

O esquecimento por outro lado é quando mesmo que se estimule partes da informação o desencadeamento ou não acontece ou não é forte o suficiente para perturbar a consciência.

$$ C := \texttt{Consciencia} $$
$$ M := \texttt{Conceito de Memorização} $$
$$ E := \texttt{Conceito de Esquecimento} $$
$$ M.A \times M.I := \texttt{Conceito de Memória Ativa e Inativa} $$
$$ (C,M,E,M.A. \times M.I.) \sim FQP $$

\hrulefill

Existe uma parte da memória que está acessível de forma próxima a consciência enquanto outra parte da memória está inativa, distante da consciência. A parte ativa da memória "reconhece" as partes que compõe uma informação da memória, então quando estas partes (não necessariamente bem definidas) são estimuladas por alguma outra informação na consciência então a outra informação na memória ativa é acessada pela consciência.\\

\textbf{Questão [ Esquecimento ]} Quais as formas de aparente inativação de memórias (esquecimento)?\\

Por hipótese diria quando há uma informação mais forte, muito similar mas de direção contrária (diz coisas diferentes), uma informação compete com outra informação e uma delas é mais forte inibindo a outra. Podemos também ter uma informação muito distante, que requer uma sequência de ativações até chegar ao destino. Podemos ter uma informação temporal muito longa, que requer um acesso linear sequencial e por isso é sujeito a frequentes desvios e pertubações de outras ideias da mente.\\

\textbf{Análise Dialética [ Esquecimento ]}

$$ \texttt{Esquecimento} \prec (\texttt{Interação Destrutiva}, \texttt{Acesso Sequencial}) \bullet (\texttt{Espaço}, \texttt{Tempo}) $$

\hrulefill

É interessante fazer um paralelo com a ondulatória, em geral ondas diferentes mas similares (fases próximas) tem maior efeito destrutivo do que ondas menos similares (fases distantes). Quando temos informações similares, por conta da unicidade do raciocínio, elas irão competir pelos estímulos da consciência, por conta disso elas podem levar ao esquecimento, uma da outra. O acesso sequencial pode levar ao esquecimento se a distância entre o início e a informação final for muito longa, isso irá requerer tempo e a consciência estará sujeita a estímulos que desviem o caminho para lembrança.

\hspace{\baselineskip}

\textbf{[ Acesso da Memória por Similaridade ]} É o acesso da memória pelas partes de uma informação, ela é imediata desencadeada pela similaridade entre as informações\\

\textbf{[ Acesso da Memória por Temporalidade ]} É o acesso da memória por meio de uma sequência de eventos ordenados no tempo, o acesso é linear e consome tempo, não é imediato e por isso está sujeito a perturbações.\\

\textbf{Questão [ Sono e Temporalidade Discretizada ]} Existe algum mecanismo durante o sono que otimize os fenômenos temporais?\\

Se comparamos sonho com o estado desperto, vemos que o sonho é mais irregular, o acesso da memória parece ser mais susceptível a pertubações, enquanto o estado desperto parece ser mais regular, talvez por que o estímulo da realidade é regular, e essa regularidade serve como uma referência para o cérebro manter a unicidade do raciocínio. O que sugere que a constância da realidade pode ser uma fonte de concentração para nossa memória. Isso é importante, pois justifica o raciocínio simbólico efetuado pela matemática, justifica a escrita e o desenho utilizado em outras áreas. Todas essas ferramentas estão na realidade e parecem ajudar a manter a concentração durante um raciocínio complicado.\\


\subsection{Definição do Paradigma Trialético}

%%%%%%%%%%%%%%%%%%%%%%%%%%%%%%%%%%%%%%%%%%%%%%%%%%%%%%%%%%%%%%%%%%%%%%%%%%%%%%%%%%%%%%%%%%%%%%%%%%%%%%%%%%%%
% 
\hspace{\baselineskip}

\textbf{Paradigma [ Método Trialético de Estudo ]} Consiste em três etapas distintas de estudo, a primeira etapa que exige esforço maior é a da memorização, o objetivo desta etapa é criar uma acessibilidade mais próxima possível da topologia uniforme, esta etapa busca a consciência alfa da informação (ver o todo pela parte). A segunda etapa consiste na análise que essencialmente se dá pela leitura e pela dialética socrática, a topologia associada é chamada de topologia dedutiva, pois entrelaça as informações em relações de dedução, implicação, esta etapa está associado a consciência gama da informação (julgar o conteúdo de verdade). A terceira etapa é a etapa de aplicação do conhecimento, o objetivo é alcançar a topologia existencial prática, em que os caminhos mais fáceis são fortalecidos, esta etapa está associado a consciência beta da informação (capacidade de usar na prática o conhecimento). O método trialético de estudo ainda não está completamente definido, a seguir será detalhado com mais profundidade os paradigmas do método.

$$ \texttt{M.T.E.} \prec ( \texttt{Memorização}, \texttt{Análise} , \texttt{Aplicação} ) \bullet \texttt{Trialética}$$

\hspace{\baselineskip}

Quero que o método conserve a estrutura da trialética humana, na estrutura trialética humana a emoção busca beleza, harmonia, equilíbrio, o intelecto busca a verdade, a organização e a vontade busca a praticidade, a aplicabilidade. Quero que a memorização fique mais palatável a consciência atribuindo características da emoção, a memorização será como uma etapa artística, em que a organização, o equilíbrio, a harmonia e a beleza se tornem instrumentos que facilitem a memória.

$$ (\texttt{Memorização}, \texttt{Análise}, \texttt{Aplicação}) \sim (\texttt{Emoção}, \texttt{Intelecto}, \texttt{Vontade}) $$

\hspace{\baselineskip}

%%%%%%%%%%%%%%%%%%%%%%%%%%%%%%%%%%%%%%%%%%%%%%%%%%%%%%%%%%%%%%%%%%%%%%%%%%%%%%%%%%%%%%%%%%%%%%%%%%%%%%%%%%%%
% 

\textbf{Paradigma [ Ferramental ]} Uma teoria pode ser representada por um conjunto de pares (Nome, Informação), o nome deve estar na linguagem verbal (suponho que a linguagem verbal ou escrita está mais próximo da consciência humana, primeiro porque a fala é acessível, segundo porque aprendemos a linguagem mais cedo do que a representação abstrata da matemática), quanto a informação, ela pode estar em qualquer linguagem, desde que a informação esteja completa, de preferência o uso de abstrações, representações em fórmulas ou descrições mais curtas e simples. A representação de uma teoria em um conjunto de pares (Nome, Informação) é chamado de ferramental, com analogia clara a ferramenta. O conjunto dos nomes de uma teoria (podendo estar organizado em categorias de estimuladores) é chamado de inventário, o conjunto das informações é chamado de externalização do ferramental, ou externalização do inventário.

$$ \texttt{Teoria} \prec \{ \dots ( \texttt{Nome} , \texttt{Informação} ) \dots \} \bullet \texttt{Ferramental} \bullet \texttt{Dialética Existencial} $$

\hspace{\baselineskip}

\textbf{Paradigma [ Inventário ]} A memorização eficiente é externalizada para consciência humana pela sua acessibilidade, o método de inventário é uma forma de verificar a acessibilidade das informações, ele consiste em enumerar o máximo possível de conteúdo de um tópico. O tópico é o agente estimulador da consciência e a enumeração externaliza o acesso provocado pelo estimulador. O método do inventário pode ser feito sem externalizar o ferramental, neste procedimento não é necessário verificar se realmente o conteúdo foi memorizado, o que é avaliado aqui é uma certa acessibilidade com relação aos estimuladores, por isso digo que o inventário externaliza virtualmente a acessibilidade. \\

\textbf{Paradigma [ Externalização de Inventário ]} Depois de feito a externalização da acessibilidade de uma teoria por seu inventário, pode ser interessante ver se as informações são acessíveis corretamente, usando o inventário como estimulador o ferramental é externalizado para verificar o acesso concreto do conteúdo.\\

Aqui é importante dizer que a externalização de um inventario pode tomar muito tempo e tornar o aprendizado ineficiente, para aumentar a eficiência é interessante usar processos aleatórios para selecionar de maneira limitada quais itens do inventário serão externalizados.\\

\textbf{[ Taxa de Acesso ]} A taxa de acesso mede a eficiência do acesso, primeiro se estabelece um número que representa o ideal de informações acessadas (eu utilizo 100, pois considero que uma pessoa que é capaz de enumerar 100 informações de um tópico, mesmo que não domine o tópico, aparenta domina-lo), tentamos escrever ou falar o máximo de conteúdos distintos do tópico, contamos e dividimos por esse número ideal. A taxa de acesso irá externalizar a acessibilidade. Podemos usar uma versão alternativa da taxa de acesso estabelecendo um tempo limite (utilizo 2 horas) e considerando a taxa de acesso como enumeração dividido pelo tempo. Com esta variável será possível acompanhar a evolução do processo de memorização. A primeira versão da variável chamo de Resolução Espacial da Memorização, a segunda versão chamo de Eficiência Temporal da Memorização. \\

\textbf{[ Taxa de Acesso Virtual ]} Mede o acesso virtual feito pelo inventário. A taxa de acesso virtual é interessante na véspera de prova em que você não tem muito tempo para verificar todas as informações. \\

\textbf{[ Taxa de Acesso Concreto ]} Mede o acesso concreto feito na externalização do inventário. O acesso concreto é fundamental para garantir o aprendizado, mas é preciso ser realista, algumas vezes você não tem tempo para fazer tudo.\\

\textbf{Paradigma [ Acesso Randômico x Acesso Sequencial ]} Queremos preservar uma topologia uniforme das informações durante o processo de memorização, para isso é interessante que o acesso seja o mais caótico possível, randômico se você acreditar que o cérebro é realmente randômico e o mundo tem processos aleatórios. Portanto durante a externalização do acesso pelo paradigma de inventário tente não preservar ordem alguma, e claro, não preserve a ordem linear.\\

Em uma teoria com topologia uniforme, a etapa alfa do método trialético (a etapa de memorização) difere do método tradicional, pois é preferível o estudo randômico das informações do que o estudo linear feito por um expositor, a coleta pode ser via uma exposição tradicional, mas é interessante que a memorização seja estimulada aleatoriamente. \\

\textbf{Problema [ Estabilização da Taxa de Acesso ]} Pode ocorrer que a taxa de acesso pare de evoluir e se estabilize em um patamar inferior ao desejado, este é o problema da estabilização da taxa de acesso. Para resolver este problema podemos refinar o estimulador usando o paradigma de estimuladores.\\

\textbf{Paradigma [ Estimuladores em Árvore ]} Durante a externalização da memória pelo inventário, inicialmente temos um estimulador da memória, que é o tópico principal, o assunto ou o nome da teoria, caso haja uma estabilização da taxa de acesso, podemos utilizar o paradigma de dividir para conquistar e criar outros estimuladores, tais estimuladores podem ser subcategorias da teoria. Para conseguir um acesso de 100 conteúdos podemos dividir para conquistar e criar 10 subcategorias, cada uma responsável pelo acesso de 10, na soma conseguimos acessar todos os 100 com menor esforço, pois temos 1-10 de cada sub-estimulador mais 1-10 do tópico principal e das subcategorias. Tanto a escrita dos estimuladores quanto a escrita do inventário pode ser feita no papel, no computador ou em qualquer recurso disponível da realidade, pois isso irá aumentar a concentração herdando a estabilidade da realidade, de preferência escreva o inventário e os estimuladores em ambientes distintos.\\

\textbf{[ Externalização Trialética ]} Depender da memória para acessar o conteúdo pode ser instável, porque não está preso a realidade e está sujeito a inibição pelo excesso de informação. Por isso é interessante que o conteúdo a ser memorizado esteja organizado e externalizado em algo concreto (este algo concreto é este documento, a Externalização Trialética). O papel da externalização trialética é ser uma fonte confiável de consulta para poder memorizar. Como o objetivo é memorizar, deve ser sucinto e condensado como um formulário.\\

Com isso a etapa da memorização se subdivide em três novas etapas, a etapa de construção do documento externalização trialética, a etapa de estudo para memorização e a etapa de checagem de pela escrita de inventário e dos estimuladores.
$$ \texttt{Memorização Trialética de Teoria} \prec (\texttt{Seleção}, \texttt{Digestão}, \texttt{Invocação}) \bullet \texttt{Base Trialética} $$

Vamos agora definir a etapa interior, que agora está sendo chamada de digestão de conteúdo. A motivação do nome vem do processo de dividir para conquistar, temos uma teoria extensa, dividimos em partes, estudamos de forma randômica seu conteúdo para estabelecer uma conexão uniforme.

Da maneira como quero definir esta etapa, ela será a mais artificial de todas, pois o cérebro humano não é um bom gerador aleatório, e o conteúdo escrito é inerentemente sequencial, para que o acesso seja aleatório podemos utilizar softwares de memorização (exemplo ANKI, de memorização por repetição espaçada) ou construir pequenas rotinas para sugerir de maneira randômica o acesso a determinada informação, seja por repetição espaçada como no ANKI ou por algum algoritmo diferente de memorização (por exemplo podemos numerar as informações e usar um gerador aleatório no python para números inteiros).

Sugiro fortemente que se você tem alguma noção de programação e tenha tempo, crie seu próprio algoritmo, pois você poderá corrigir eventuais desvios de propósito. \\

\textbf{Paradigma [ Invocação ]} É o paradigma de ferramental e de árvore de estimuladores. Invocação Interna é apenas para o inventário e Invocação Externa é para o ferramental concreto. \\

\textbf{Paradigma [ Digestão ]} Utilizar um gerador aleatório (ou qualquer algoritmo) para estimular o acesso do conteúdo (o estimulador pode ser o nome da informação, pode ser uma parte da informação). Deve ser construído este gerador com base nas informações escritas na seleção (Externalização Trialética), por isso o documento deve ter alguma forma de organização para facilitar a consulta. O nome da implementação é chamado de digestor. \\

Durante um estudo de digestão é interessante utilizar a sopa de ideias, separando estimuladores secundários da informação final a ser acessada. O digestor deve fornecer aleatoriamente um estimulador principal (como o nome da informação), e a informação deve ser resgatada através da externalização em sopa de ideias. \\

Durante o processo de digestão algumas informações podem induzir a lembrança de outras, é interessante que tais informações sejam lembradas, pois a conexão entre elas pode facilitar a sua aplicação, mas não de forma a prejudicar o andamento do estudo, um bom paradigma é sempre tentar lembrar de outras informações caso a informação principal não consiga ser acessada imediatamente, ou quando uma informação é fácil demais lembrar posteriormente das informações que forem induzidas.\\

%\textbf{Paradigma [ Otimização da Taxa de Acesso pelo Digestor ]} O digestor pode funcionar como auxiliar para aumentar a taxa de acesso. Podemos usar o paradigma de estimuladores em árvore, mas o digestor é o que permite menos subcategorias e um acesso mais direto ao conteúdo.\\

\textbf{Paradigma [ Seleção ]} O conteúdo deve ser otimizado, trabalhado esteticamente para que fique confortável e mais fácil de memorizar durante esta etapa, deve selecionar o conteúdo de forma que seja de fácil consulta para a etapa de digestão.\\

Para resumir o método temos:\\

\textbf{Projeção Dialética [ Método de Estudo Trialético ]}

$$ \texttt{Estudo} \prec \texttt{(Memorização, Análise, Aplicação)} $$
$$ \texttt{Memorização} \prec \texttt{(Seleção, Digestão, Invocação)} $$
$$ \texttt{Análise} \prec \texttt{(Leitura, Dialética Socrática)} $$
$$ \texttt{Aplicação} \prec \texttt{(Fazer Exercícios, Criar Problemas)} $$
$$ \{ \texttt{Análise} \sim \texttt{Aplicação} \} \prec \texttt{Análise de Exemplos}$$

As etapas de análise e aplicação não serão detalhadas, pois são mais naturais contidos no método tradicional. Entretanto tanto a análise quanto a aplicação são externalizados no formato do texto. Por exemplo nas etapas de análise externalizamos o estudo escrevendo as demonstrações das informações como demonstrações dialética socráticas (é omitido as respostas, apenas as perguntas e construções auxiliares), na etapa da aplicação selecionamos exercícios que achamos importantes ou interessantes.


%%%%%%%%%%%%%%%%%%%%%%%%%%%%%%%%%%%%%%%%%%%%%%%%%%%%%%%%%%%%%%%%%%%%%%%%%%%%%%%%%%%%%%%%%%%%%%%%%%%%%%%%%%%%
%%%%%%%%%%%%%%%%%%%%%%%%%%%%%%%%%%%%%%%%%%%%%%%%%%%%%%%%%%%%%%%%%%%%%%%%%%%%%%%%%%%%%%%%%%%%%%%%%%%%%%%%%%%%\\
%%%%%%%%%%%%%%%%%%%%%%%%%%%%%%%%%%%%%%%%%%%%%%%%%%%%%%%%%%%%%%%%%%%%%%%%%%%%%%%%%%%%%%%%%%%%%%%%%%%%%%%%%%%%
%%%%%%%%%%%%%%%%%%%%%%%%%%%%%%%%%%%%%%%%%%%%%%%%%%%%%%%%%%%%%%%%%%%%%%%%%%%%%%%%%%%%%%%%%%%%%%%%%%%%%%%%%%%%
%%%%%%%%%%%%%%%%%%%%%%%%%%%%%%%%%%%%%%%%%%%%%%%%%%%%%%%%%%%%%%%%%%%%%%%%%%%%%%%%%%%%%%%%%%%%%%%%%%%%%%%%%%%%

\subsection{Estruturação da Externalização Trialética}

\hspace{\baselineskip}

Até aqui discutimos a argumentação filosófica dos paradigmas que formalizam este tipo de documento, agora vamos estabelecer uma série de notações e formalizações para a escrita do texto do tipo externalização trialética.

\hspace{\baselineskip}

\textbf{[ Definição ]} A definição é a primeiro elemento da organização de uma teoria, ela é sempre verdadeira, ela representa uma criação, um modelo para uma existência, geralmente composto de uma representação, um nome e um conjunto de propriedades. Como criação, ela é sempre verdadeira. No texto tudo que estiver em colchetes representa o nome de uma definição, e este é a definição de definição.\\

\textbf{[ Forma Externa ]} Uma forma externa está associado a externalização na consciência, a forma externa é uma reformulação de alguma entidade em outra mais fácil de memorizar, mais abstrata, ou mais organizada. \\

\textbf{[ Informações ]} Se temos uma teoria lógica axiomática, temos as informações assumidas como verdadeira de axiomas, as informações demonstradas pela lógica são teoremas (de maior importância), proposições (de importância menor) e Lema (utilizado apenas pontualmente em algum problema, ou algo muito simples de demonstrar), vamos simplificar e assumir que todas as informações são teoremas (o que implica em uma topologia existencial uniforme das informações). Paradigmas são informações não demonstradas formalmente, mas justificadas de alguma maneira, ou apenas consideradas verdade.

\hspace{\baselineskip}

Sempre que desejarmos estabelecer uma informação, usamos \textbf{Teorema [ Nome da Informação ]} com o teorema podendo ser substituído por paradigma, lema, proposição, informação etc. No caso da definição não escrevemos definição, apenas usamos \textbf{[ Nome da Definição ]}. Temos então o elemento de informação da externalização trialética como \textbf{Tipo de Elemento [ Nome do Elemento ]}, perceba o título do documento.

\hspace{\baselineskip}

\textbf{[ Notação ]} É uma definição simples, apenas diz como uma outra definição deve ser representada simbolicamente.

\hspace{\baselineskip}

\textbf{[ Estrutura ]} É uma lista da forma $(a,b,c)_{nome1}^{nome2}(A,B,C)$ que representa uma estrutura, um conjunto de definições que possuem uma relação entre si. Como convenção estabeleço que $(a,b,c)$ são elementos e $(A,B,C)$ são conjunto de elementos. O que caracteriza uma estrutura é sua relação interna.\\

\textbf{Paradigma [ Estrutural ]} Existem definições que sempre andam juntas, quando isso ocorre é interessante fazer uma exposição de uma vez utilizando a o elemento textual estrutura. \\

\textbf{Estrutura [ Teoremas ]} Um teorema pode ser dividido internamente como um conjunto de informações assumidas como verdade (Chamado de Hipóteses)  e um conjunto de informações que são deduzidas (Chamado de Kernel), eles formam a estrutura de um teorema $(\texttt{Hipótese}, \texttt{Kernel})_{\texttt{Teorema}}$, todo teorema pode ser dividido em hipótese e kernel, mas alguns teoremas podem ser representados apenas pelo kernel.\\

\textbf{Notação [ Teoremas ]} Neste texto um teorema é representado pela forma:

$$\Phi(\texttt{...Hipóteses...}) : \texttt{...Kernel...}$$

\textbf{[ Hipótese ]} Algumas vezes é interessante organizar hipóteses vista com frequência em diferentes teoremas, assim temos um tipo de classe de definição chamado hipótese.\\

\textbf{[ Kernel ]} Da mesma forma que as hipóteses, podemos ter kernels iguais em diferentes teoremas.\\

\textbf{[ Teorema ]} É uma informação que possui um nome, um corpo constituído de hipóteses, kernel e uma demonstração lógica (nestes tipos de documento irei optar pela dialética demonstrativa para representar uma demonstração, a dialética demonstrativa não demonstra, mas expõe as construções e as perguntas socráticas necessárias para resolver o teorema). Outros nomes para equivalentes de teorema são proposições e lemas, o teorema é considerado mais importante que as proposições, e as proposições são mais importantes que lemas. \\

\textbf{[ Exercício ]} São exercícios, informações que podem ser deduzidas dos teoremas, das propriedades de uma definição ou diretamente dos axiomas, é arbitrário definir o que é um exercício do que é um teorema/paradigma, mas um dos critérios objetivos possíveis é sua complexidade e aplicabilidade.\\

\textbf{[ Externalização ]} São exercícios que são aplicação direta de algum teorema/paradigma ou são deduzidos diretamente de propriedades de uma definição qualquer.\\

\textbf{[ Dialética Demonstrativa ]} É a representação de uma demonstração por meio de construções e perguntas, geralmente associado a um teorema. A ideia da dialética demonstrativa é expor o caminho da demonstração e deixar como exercício a escrita da demonstração propriamente dita.  \\

\textbf{[ Demonstração ]} É a demonstração de um teorema, a exposição passo a passo dos raciocínios dedutivos para chegar a conclusão de um teorema, mostrando sua validade.\\

\textbf{[ Externalização Trialética ]} É o tipo de documento que usa o formalismo definido nesta seção. O seu título também obedece ao formalismo: \textbf{Externalização Trialética [ Nome da Teoria ]}.  \\

\textbf{[ Motivação ]} Algumas vezes é interessante expor a motivação de uso de uma definição ou de um teorema.\\

\textbf{Paradigma [ Índice e Árvore Estimuladora ]} O índice e sua estrutura pode e deve ser a externalização de uma árvore de estimuladores, quanto mais organizado o conteúdo, mas fácil será para a etapa de digestão efetuada por geradores aleatórios e algoritmos de memorização.\\

\textbf{Questão [ Redundância ]} Pode ser que uma informação possa estar em uma categoria, quanto estar em outra, é vantajoso escolher uma ou é mais vantajoso colocar nas duas? \\

Deixo isso como arbitrário, algumas vezes se você quiser manter a teoria como uma teoria axiomática (que é estritamente linear quando externalizada em um documento) pode ser interessante que não inclua em muitas subcategorias uma mesma informação, algumas vezes se o objetivo é memorização apenas, ou que a teoria não seja uma teoria axiomática então é mais interessante incluir uma informação em diferentes categorias relacionadas, pois preserva um padrão de acessibilidade desejável.\\

O resultado de um estudo sob o método trialético bem sucedido é um documento do tipo externalização trialética, chamamos de externalização trialética, pois é a externalização da memorização (construção de fórmulas, estruturas, hipóteses, informações organizadas e categorizadas no índice), é a externalização da análise (nas demonstrações, nas motivações, nos paradigmas) e por fim é a externalização da aplicação (nos exercícios, nos exercícios de aplicação e em possíveis exercícios criados).

%%%%%%%%%%%%%%%%%%%%%%%%%%%%%%%%%%%%%%%%%%%%%%%%%%%%%%%%%%%%%%%%%%%%%%%%%%%%%%%%%%%%%%%%%%%%%%%%%%%%%%%%%%%%
%%%%%%%%%%%%%%%%%%%%%%%%%%%%%%%%%%%%%%%%%%%%%%%%%%%%%%%%%%%%%%%%%%%%%%%%%%%%%%%%%%%%%%%%%%%%%%%%%%%%%%%%%%%%\\
%%%%%%%%%%%%%%%%%%%%%%%%%%%%%%%%%%%%%%%%%%%%%%%%%%%%%%%%%%%%%%%%%%%%%%%%%%%%%%%%%%%%%%%%%%%%%%%%%%%%%%%%%%%%
%%%%%%%%%%%%%%%%%%%%%%%%%%%%%%%%%%%%%%%%%%%%%%%%%%%%%%%%%%%%%%%%%%%%%%%%%%%%%%%%%%%%%%%%%%%%%%%%%%%%%%%%%%%%
%%%%%%%%%%%%%%%%%%%%%%%%%%%%%%%%%%%%%%%%%%%%%%%%%%%%%%%%%%%%%%%%%%%%%%%%%%%%%%%%%%%%%%%%%%%%%%%%%%%%%%%%%%%%

%%%%%%%%%%%%%%%%%%%%%%%%%%%%%%%%%%%%%%%%%%%%%%%%%%%%%%%%%%%%%%%%%%%%%%%%%%%%%%%%%%%%%%%%%%%%%%%%%%%%%%%%%%%%%%%%%%%%%%%%%%%%%%%%%%%%%
% Paradigma de Topologia de Memória
%\hspace{\baselineskip}

\textbf{Paradigma [ Topologia de Memória ]} Topologia de um objeto pode se referir a como estão conectados seus componentes. Seja a memória composta de informações, queremos saber como estão conectados, mais especificamente, quais informações estão mais próximas entre si. Considero duas formas de estabelecer conexões e assim estabelecer proximidade topológica, a primeira é o co-estímulo (aprendizado é basicamente você co-estimular nome com significado), a segunda é a semelhança (supostamente semelhança não é um atributo do sujeito, mas dos objetos e assim não dependem muito da nossa vontade). Por paradigmas anteriores estabelecidos, uma demonstração intuitiva respeita a organização topológica da mente no estado anterior ao teorema, enquanto uma demonstração invocativa não. É natural, conforme discutido, desejarmos demonstrações intuitivas, entretanto para o caso dos teoremas desejamos (algumas vezes desobedecemos em nome do rigor da teorização axiomática) informações contra-intuitivas, pois caso fosse intuitivo não seria necessário expor seu conteúdo (algumas vezes é necessário como lema, caso sua aplicação seja recorrente e sua composição contra-intuitiva).

\hspace{\baselineskip}

\textbf{Paradigma [ Caminho Natural de Acesso da Memória ]} Em geral quando queremos usar um teorema(informação) é mais fácil lembrar da conclusão do teorema que das hipóteses. Considero que isso seja devido a maneira as informações estão organizadas em uma teoria, se uma teoria é organizada de forma que cada informação tenha várias hipóteses conectando uma ou um número menor de conclusões e as hipóteses são usadas várias vezes em mais de uma informação, então as informações desta teoria são mais fáceis de serem lembradas pela conclusão do que pela hipótese, pois a redundância das hipóteses irão induzir vários caminhos (estes vários caminhos irão provocar o efeito de confusão). Vamos representar uma teoria como descrito acima pela notação $(H>K)$ com $H$ representando a redundância das hipóteses e $K$ a redundância das conclusões, teorias do tipo $(H>K)$ são teorias com perfil demonstrativo, se você já tem posse da conclusão que você quer então você pode aplicar a teoria para demonstra-la, pois será mais fácil pela conclusão lembrar da informação. Para teorias $(H>K)$ é preferível que a conclusão venha antes das hipóteses nos teoremas. Por outro lado, teorias do tipo $(H<K)$, onde conjuntos limitados e organizados de hipóteses levam a várias conclusões são teorias que permitem a descoberta de informações, pois podem ser trabalhadas pelas hipóteses sem a necessidade de saber previamente do que se quer obter, portanto é preferível manter a ordem (hipótese, conclusões). Assim dependendo de como a teoria é organizada as informações devem estar organizadas de maneira que otimize a lembrança ou otimize a aplicação, por isso notações que alterem a ordem (hipótese, conclusão) para (conclusão, hipótese) ou notações que abstraiam o conteúdo das hipóteses em algo mais simples (como no caso da definição de estruturas) serão usadas com frequência neste documento.

\hspace{\baselineskip}

Uma teoria pode ser capaz de alterar a topologia entre informações, tornando informações que antes eram distantes em informações em que sua dedução é natural, análises rigorosas devem ser feitas antes de descartar qualquer abordagem. Podemos com alguma segurança dizer que a matemática busca alterar a topologia das informações de forma a encontrar novos caminhos para o futuro ou tornar os caminhos que antes eram trilhas em avenidas movimentadas.

%\hspace{\baselineskip}

%Existem riscos de se usar somente esse paradigma, o principal deles é achar ter demonstrado algo, é sempre bom manter um caderno separado com as demonstrações, comparar com demonstrações de livros tradicionais (Este formato não substitui os livros tradicionais, o objetivo é complementar) e também com colegas de estudo. 

%%%%%%%%%%%%%%%%%%%%%%%%%%%%%%%%%%%%%%%%%%%%%%%%%%%%%%%%%%%%%%%%%%%%%%%%%%%%%%%%%%%%%%%%%%%%%%%%%%%%%%%%%%%%%%%%%%%%%%%%%%%%%%%%%%%%%%%%%%%%%%%%%%%%%%%%%%%%%%%%%%%%%%%%%%
% Paradigma de Estruturação Projetado em Computação - Computador é um Experimento de Simulação da Mente humana

%%%%%%%%%%%%%%%%%%%%%%%%%%%%%%%%%%%%%%%%%%%%%%%%%%%%%%%%%%%%%%%%%%%%%%%%%%%%%%%%%%%%%%%%%%%%%%%%%%%%%%%%%%%%%%%%%%%%%%%%%%%%%%%%%%%%%%%%%%%%%%%%%%%%%%%%%%%%%%%%%%%%%%%%%%
% Consciência de Entidade Matemática - Externalização de Entidade Matemática - Projeção em Computação Orientação a Objeto

%%%%%%%%%%%%%%%%%%%%%%%%%%%%%%%%%%%%%%%%%%%%%%%%%%%%%%%%%%%%%%%%%%%%%%%%%%%%%%%%%%%%%%%%%%%%%%%%%%%%%%%%%%%%%%%%%%%%%%%%%%%%%%%%%%%%%%%%%%%%%%%%%%%%%%%%%%%%%%%%%%%%%%%%%%
% Paradigma de Exercitação Imediata e Consciência

%%%%%%%%%%%%%%%%%%%%%%%%%%%%%%%%%%%%%%%%%%%%%%%%%%%%%%%%%%%%%%%%%%%%%%%%%%%%%%%%%%%%%%%%%%%%%%%%%%%%%%%%%%%%%%%%%%%%%%%%%%%%%%%%%%%%%%%%%%%%%%%%%%%%%%%%%%%%%%%%%%%%%%%%%%
% Paradigma de Estudo Superficial - Declaração de Variáveis e Definição de Variáveis

%%%%%%%%%%%%%%%%%%%%%%%%%%%%%%%%%%%%%%%%%%%%%%%%%%%%%%%%%%%%%%%%%%%%%%%%%%%%%%%%%%%%%%%%%%%%%%%%%%%%%%%%%%%%%%%%%%%%%%%%%%%%%%%%%%%%%%%%%%%%%%%%%%%%%%%%%%%%%%%%%%%%%%%%%%
% Paradigma : Ideal Demonstrativo e Dialética Demonstrativa
%\hspace{\baselineskip}

%\textbf{Paradigma [ Ideal Demonstrativo e Dialética Demonstrativa ]} Suponha que todos os teoremas estivessem de acordo com o nosso ideal (teoremas contra-intuitivos e demonstrações intuitivas). Logo a dificuldade do teorema estaria na busca pelo caminho natural correto dentre diferentes caminhos naturais que poderiam levar a lugar nenhum, portanto seria o papel da dialética demonstrativa indicar o caminho correto e delegar ao leitor o papel de verificar com o rigor lógico cada etapa e obter a demonstração formal.


%%%%%%%%%%%%%%%%%%%%%%%%%%%%%%%%%%%%%%%%%%%%%%%%%%%%%%%%%%%%%%%%%%%%%%%%%%%%%%%%%%%%%%%%%%%%%%%%%%%%%%%%%%%%%%%%%%%%%%%%%%%%%%%%%%%%%%%%%%%%%%%%%%%%%%%%%%%%%%%%%%%%%%%%%%


%\newpage


%%%%%%%%%%%%%%%%%%%%%%%%%%%%%%%%%%%%%%%%%%%%%%%%%%%%%%%%%%%%%%%%%%%%%%%%%%%%%%%%%%%%%%%%%%%%%%%%%%%%%%%%%%%%%%%%%%%%%%%%%%%%%%%%%%%%%%%%%%%%%%%%%%%%%%%%%%%%%%%%%%%%%%%%%%
%%%%%%%%%%%%%%%%%%%%%%%%%%%%%%%%%%%%%%%%%%%%%%%%%%%%%%%%%%%%%%%%%%%%%%%%%%%%%%%%%%%%%%%%%%%%%%%%%%%%%%%%%%%%%%%%%%%%%%%%%%%%%%%%%%%%%%%%%%%%%%%%%%%%%%%%%%%%%%%%%%%%%%%%%%





\hspace{\baselineskip}

O método trialético de estudo consiste na consciência de três etapas distintas no processo de aprendizagem de uma teoria. As três etapas são a memorização, a análise e a aplicação, sendo que tais etapas são feitas preferencialmente nesta ordem. No estudo tradicional temos duas etapas distintas, a primeira é a exposição do conteúdo feita por meio da leitura ou por meio de um orador, nesta etapa não se exclui o diálogo, a interação é permitida, a segunda etapa é a etapa de aplicação, em que trabalhos, exercícios e aplicações são propostos. A diferença entre o estudo tradicional e o estudo trialético é que a exposição do conteúdo é decomposto em duas etapas, a memorização e a análise, enquanto a etapa de aplicação permanece inalterada.




%%%%%%%%%%%%%%%%%%%%%%%%%%%%%%%%%%%%%%%%%%%%%%%%%%%%%%%%%%%%%%%%%%%%%%%%%%%%%%%%%%%%%%%%%%%%%%%%%%%%%%%%%%%%
%%%%%%%%%%%%%%%%%%%%%%%%%%%%%%%%%%%%%%%%%%%%%%%%%%%%%%%%%%%%%%%%%%%%%%%%%%%%%%%%%%%%%%%%%%%%%%%%%%%%%%%%%%%%\\
%%%%%%%%%%%%%%%%%%%%%%%%%%%%%%%%%%%%%%%%%%%%%%%%%%%%%%%%%%%%%%%%%%%%%%%%%%%%%%%%%%%%%%%%%%%%%%%%%%%%%%%%%%%%
%%%%%%%%%%%%%%%%%%%%%%%%%%%%%%%%%%%%%%%%%%%%%%%%%%%%%%%%%%%%%%%%%%%%%%%%%%%%%%%%%%%%%%%%%%%%%%%%%%%%%%%%%%%%
%%%%%%%%%%%%%%%%%%%%%%%%%%%%%%%%%%%%%%%%%%%%%%%%%%%%%%%%%%%%%%%%%%%%%%%%%%%%%%%%%%%%%%%%%%%%%%%%%%%%%%%%%%%%
















































%%%%%%%%%%%%%%%%%%%%%%%%%%%%%%%%%%%%%%%%%%%%%%%%%%%%%%%%%%%%%%%%%%%%%%%%%%%%%%%%%%%%%%%%%%%%%%%%%%%%%%%%%%%%
%%%%%%%%%%%%%%%%%%%%%%%%%%%%%%%%%%%%%%%%%%%%%%%%%%%%%%%%%%%%%%%%%%%%%%%%%%%%%%%%%%%%%%%%%%%%%%%%%%%%%%%%%%%%
%%%%%%%%%%%%%%%%%%%%%%%%%%%%%%%%%%%%%%%%%%%%%%%%%%%%%%%%%%%%%%%%%%%%%%%%%%%%%%%%%%%%%%%%%%%%%%%%%%%%%%%%%%%%
%%%%%%%%%%%%%%%%%%%%%%%%%%%%%%%%%%%%%%%%%%%%%%%%%%%%%%%%%%%%%%%%%%%%%%%%%%%%%%%%%%%%%%%%%%%%%%%%%%%%%%%%%%%%
%%%%%%%%%%%%%%%%%%%%%%%%%%%%%%%%%%%%%%%%%%%%%%%%%%%%%%%%%%%%%%%%%%%%%%%%%%%%%%%%%%%%%%%%%%%%%%%%%%%%%%%%%%%%

\chapter{Teoria de Lógica Formal}

%%%%%%%%%%%%%%%%%%%%%%%%%%%%%%%%%%%%%%%%%%%%%%%%%%%%%%%%%%%%%%%%%%%%%%%%%%%%%%%%%%%%%%%%%%%%%%%%%%%%%%%%%%%%
% Introdução
\hspace{\baselineskip}

Este é o último capítulo da teoria principal, a teoria de lógica formal busca fazer uma releitura do raciocínio dedutivo e da lógica formal. Diferente da lógica matemática ela busca ser abrangente o suficiente para ser usado em qualquer atividade intelectual, por conta disso, ela difere do que é chamado de lógica formal pelos matemáticos. \\

Neste capítulo irei me concentrar muito mais na organização estrutural da atividade intelectual humana do que em regras de inferências, sendo estas apenas um dos componentes da atividade intelectual humana. Devo dar ênfase ao caráter humano desta teoria, é uma teoria adaptada para o uso de uma consciência similar a consciência humana, portanto tem como objetivo explorar o potencial intelectual humano como um indivíduo.\\

No fim do capítulo irei expor um pouco da lógica formal em forma de teoria axiomática utilizando a filosofia escrita neste livro, deve ser visto como uma introdução ao assunto e não uma investigação cuidadosa.

%Usando o que foi definido no capítulo anterior sobre a teoria de projeções dialéticas, a atividade intelectual humana possui uma decomposição natural na trialética.

%É uma teoria normativa, ela busca dizer o que deve ou não ser feito, como deve ser feito e como organizar o raciocínio.

%%%%%%%%%%%%%%%%%%%%%%%%%%%%%%%%%%%%%%%%%%%%%%%%%%%%%%%%%%%%%%%%%%%%%%%%%%%%%%%%%%%%%%%%%%%%%%%%%%%%%%%%%%%%
%%%%%%%%%%%%%%%%%%%%%%%%%%%%%%%%%%%%%%%%%%%%%%%%%%%%%%%%%%%%%%%%%%%%%%%%%%%%%%%%%%%%%%%%%%%%%%%%%%%%%%%%%%%%
%%%%%%%%%%%%%%%%%%%%%%%%%%%%%%%%%%%%%%%%%%%%%%%%%%%%%%%%%%%%%%%%%%%%%%%%%%%%%%%%%%%%%%%%%%%%%%%%%%%%%%%%%%%%

\section{Raciocínio Trialético}

%\hspace{\baselineskip}

$$ \texttt{Raciocínio} \prec ( \texttt{Memorização, Lembrança, Inferência} ) \bullet \texttt{Trialética} $$

A atividade intelectual humana sob o ponto de vista trialético pode ser visto como a interação entre a consciência (lugar do espírito, centro ou identidade do ser humano) e a memória, podendo fazer uso da realidade como substrato. Nesta trialética a memorização é o processo de armazenar uma informação na memória, a lembrança é o processo de resgate da informação e inferência é o uso efetivo das informações para fazer previsões, resolver problemas ou responder perguntas.\\

%%%%%%%%%%%%%%%%%%%%%%%%%%%%%%%%%%%%%%%%%%%%%%%%%%%%%%%%%%%%%%%%%%%%%%%%%%%%%%%%%%%%%%%%%%%%%%%%%%%%%%%%%%%%
% Raciocínio Vertical
\textbf{[ Raciocínio Vertical ]} É o par (Memorização, Lembrança) compondo os processos responsáveis pela interação entre a memória (lugar onde se armazena informações) e a parte responsável pelo seu uso efetivo (realidade ou imaginação), onde ocorre a inferência.\\

\textbf{[ Raciocínio Horizontal ]} É a inferência, ou o raciocínio propriamente dito.

%%%%%%%%%%%%%%%%%%%%%%%%%%%%%%%%%%%%%%%%%%%%%%%%%%%%%%%%%%%%%%%%%%%%%%%%%%%%%%%%%%%%%%%%%%%%%%%%%%%%%%%%%%%%
% O que seria armazenar uma informação na memória ?

%%%%%%%%%%%%%%%%%%%%%%%%%%%%%%%%%%%%%%%%%%%%%%%%%%%%%%%%%%%%%%%%%%%%%%%%%%%%%%%%%%%%%%%%%%%%%%%%%%%%%%%%%%%%
% Consciência Informacional
\hspace{\baselineskip}

\textbf{[ Consciência Informacional ]} É a definição de consciência associado a sensação de entendimento, a capacidade de julgamento do que é verdadeiro ou falso, a capacidade de abstração, é portanto a consciência associado a informação, ao uso do intelecto.

\begin{itemize}
	\item[ (i) ] \textbf{[ Consciência Alfa ]} É a capacidade de abstrair, de representar, de sentir toda informação apenas pelo estímulo de parte dela. %É com as conexões que é possível que parte de uma informação permita sentir a sensação como um todo, pois o estímulo transferido pelas conexões pode resgatar o estímulo original.
	\item[ (ii) ] \textbf{[ Consciência Gama ]} É a capacidade de julgar se uma informação é verdadeira ou falsa. 
	\item[ (iii) ] \textbf{[ Consciência Beta ]} É a capacidade de usar informações para resolver problemas ou de criar existências.
\end{itemize}

\textbf{Comentário [ Consciência Alfa ]} Sem a consciência alfa não poderíamos usar símbolos para efetuar inferências, não poderíamos representar objetos, não poderíamos criar lugares para memórias de existências. A consciência alfa é intimamente associada a arquitetura do cérebro humano baseado em conexões. Suponha que tenhamos sensação como um padrão de lâmpadas acesas e apagadas, e que as conexões entre as lâmpadas podem acender ou apagar outras lâmpadas adjacentes, as conexões representam a memória de alguma sensação, se algumas lâmpadas forem acesas, por conta da conexões memorizadas, outras lâmpadas são acesas e apagadas resgatando o padrão memorizado, portanto se um estímulo for suficientemente parecido com uma informação a sensação da informação como um todo pode ser resgatada, isto exemplifica a consciência alfa, o resgate do todo pela parte. A consciência alfa é treinada pelo estudo de memorização. \\

\textbf{Comentário [ Consciência Gama e Beta ]} Consciência gama é a capacidade de julgar se um raciocínio, uma inferência ou uma informação avulsa é correta. Ela depende da consciência alfa, o raciocínio abstrato depende da capacidade de sentir as informações e as propriedades dos símbolos abstratos usados para representar as informações brutas. Tanto a consciência gama e a beta correspondem a parte do raciocínio horizontal que é o raciocínio efetivo.



%%%%%%%%%%%%%%%%%%%%%%%%%%%%%%%%%%%%%%%%%%%%%%%%%%%%%%%%%%%%%%%%%%%%%%%%%%%%%%%%%%%%%%%%%%%%%%%%%%%%%%%%%%%%
%%%%%%%%%%%%%%%%%%%%%%%%%%%%%%%%%%%%%%%%%%%%%%%%%%%%%%%%%%%%%%%%%%%%%%%%%%%%%%%%%%%%%%%%%%%%%%%%%%%%%%%%%%%%
%%%%%%%%%%%%%%%%%%%%%%%%%%%%%%%%%%%%%%%%%%%%%%%%%%%%%%%%%%%%%%%%%%%%%%%%%%%%%%%%%%%%%%%%%%%%%%%%%%%%%%%%%%%%

\section{Representação Estrutural dos Raciocínios Trialéticos}

\hspace{\baselineskip}

Nesta subseção irei definir estruturas e problemas definidos associados ao raciocínio trialético. É recomendado que cada estrutura tenha uma notação diferente, dessa forma o cérebro condiciona suas competências de maneira especializada para cada tipo de problema através da notação. Para quem tem conhecimento de programação esta subseção será bem similar a definição de uma interface gráfica para representar um raciocínio.

%%%%%%%%%%%%%%%%%%%%%%%%%%%%%%%%%%%%%%%%%%%%%%%%%%%%%%%%%%%%%%%%%%%%%%%%%%%%%%%%%%%%%%%%%%%%%%%%%%%%%%%%%%%%
% Estrutura Recursiva
\hspace{\baselineskip}

\textbf{[ Estrutura Recursiva ]} É o componente elementar para representar um raciocínio.

\begin{itemize}
	\item[ (i) ] \textbf{[ Abertura de Problema ]} Deve existir uma notação para abrir um tipo de problema de forma que este tenha um título, ou uma descrição do objetivo ou algo que indique o início da estrutura recursiva.
	\item[ (ii) ] \textbf{[ Parte Interna ]} Nesta parte é possível começar outras estruturas recursivas. Deve ser evidente na notação que esta parte pertence a abertura do problema e a notação deve ser suficientemente definida para que não haja confusão entre subestruturas dentro da parte interna.
	\item[ (iii) ] \textbf{[ Conclusão ]} Nesta parte devemos obter alguma conclusão com relação ao problema.
\end{itemize}

A notação em si fica a cargo do usuário, sendo interessante que cada tipo de estrutura tenha uma notação diferente.

%%%%%%%%%%%%%%%%%%%%%%%%%%%%%%%%%%%%%%%%%%%%%%%%%%%%%%%%%%%%%%%%%%%%%%%%%%%%%%%%%%%%%%%%%%%%%%%%%%%%%%%%%%%%
% Resgate de Informação
\hspace{\baselineskip}

\textbf{Estrutura Recursiva [ Resgate de Informação ]} Representa o problema de resgatar uma informação da memória a partir de um nome, de uma representação abstrata. A abertura do problema deve ter o nome da informação. A parte interna é livre como um rascunho, podendo ser composta de estimuladores (elementos que podem ajudar na lembrança) ou de subestruturas. A conclusão deve conter uma formulação da informação resgatada pela memória, com comentários ou sinais para indicar certeza, erro, quase certo e etc.\\

\textbf{Estrutura Recursiva [ Definição ]} É o processo inverso de resgate de informação, busca criar abstrações, sendo que a abstração pode ser o nome da informação ou pode estar contida na conclusão como algum símbolo ou notação. A parte interna é livre como um rascunho.\\

\textbf{Estrutura Recursiva [ Verificação (Gama) ]} Busca verificar se uma informação é verdadeira, a parte interna além de subestruturas possuem raciocínios e a conclusão diz se a informação foi verificada ou não. O nome, uma formulação ou uma descrição da informação ou do problema deve ser colocado na abertura do problema. \\

\textbf{Estrutura Recursiva [ Criação Objetiva (Beta) ]} A abertura do problema deve conter um objetivo associado a criação de algo. O conteúdo interno é livre e a conclusão pode conter a criação ou indicar que a criação foi realizada com sucesso.\\

\textbf{Estrutura Recursiva [ Indução / Criação (Alfa) ]} A abertura do problema pode ser qualquer coisa assim como o conteúdo interno, a conclusão pode ser formulações, criações e etc. Uma estrutura indutiva não tem um objetivo bem definido, a ideia é que o indutor seja capaz de gerar algo que não é conhecido ainda.\\

As duas primeiras estruturas recursivas estão associados ao raciocínio vertical (interação com a memória). As três últimas podem ser consideradas uma decomposição trialética do processo de raciocínio horizontal. É recomendado que para cada uma das cinco formas de raciocínio exista uma notação diferente.\\

\textbf{[ Problema ]} Com base nas cinco estruturas recursivas definidas, associamos para cada estrutura um tipo de problema. Problema de Resgate de Informação, Problema de Definição, Problema de Verificação (ou Demonstração), Problema Objetivo ( Beta : com propósito suficientemente definido), e Problema Criativo ( Alfa : com Propósito indefinido, desconhecido ).

%%%%%%%%%%%%%%%%%%%%%%%%%%%%%%%%%%%%%%%%%%%%%%%%%%%%%%%%%%%%%%%%%%%%%%%%%%%%%%%%%%%%%%%%%%%%%%%%%%%%%%%%%%%%
%%%%%%%%%%%%%%%%%%%%%%%%%%%%%%%%%%%%%%%%%%%%%%%%%%%%%%%%%%%%%%%%%%%%%%%%%%%%%%%%%%%%%%%%%%%%%%%%%%%%%%%%%%%%
%%%%%%%%%%%%%%%%%%%%%%%%%%%%%%%%%%%%%%%%%%%%%%%%%%%%%%%%%%%%%%%%%%%%%%%%%%%%%%%%%%%%%%%%%%%%%%%%%%%%%%%%%%%%

\subsection{Formas de Denotar as Estruturas de Raciocínio}

\hspace{\baselineskip}

Podem existir infinitas formas de denotar as diferentes estruturas de raciocínio, entretanto existem algumas formas que são interessantes para atender a propósitos definidos.

%%%%%%%%%%%%%%%%%%%%%%%%%%%%%%%%%%%%%%%%%%%%%%%%%%%%%%%%%%%%%%%%%%%%%%%%%%%%%%%%%%%%%%%%%%%%%%%%%%%%%%%%%%%%
% Indicar componentes e subsestruturas da parte interna de uma estrutura
\hspace{\baselineskip}

\textbf{[ Hierarquia de Estruturas por Indentação ]} Para indicar que uma estrutura ou componente de raciocínio pertence a parte interna de uma estrutura anterior podemos usar o que em computação é chamado de indentação, basicamente você deve colocar uma margem à esquerda do conteúdo para indicar que tudo que tiver à direita desta margem pertence a abertura de problema imediatamente acima que não tem a margem (ou tem uma margem menor). É possível denotar níveis hierárquicos de estruturas usando quantidades definidas de espaçamento das margens.\\

\textbf{Comentário :} Esta forma de representar hierarquias de estruturas é interessante para representar estruturas em editores de texto em um computador. Entretanto se o substrato de raciocínio for um caderno pode não ser muito conveniente usar indentação através de espaçamento de margens, neste substrato podemos substituir o espaçamento por linhas verticais tracejadas, estes são mais eficientes de indicar os níveis hierárquicos, permite que estruturas extensas possam continuar nas páginas seguintes sem perder seu significado.

%%%%%%%%%%%%%%%%%%%%%%%%%%%%%%%%%%%%%%%%%%%%%%%%%%%%%%%%%%%%%%%%%%%%%%%%%%%%%%%%%%%%%%%%%%%%%%%%%%%%%%%%%%%%
% Indicar componentes e subsestruturas da parte interna de uma estrutura atraves de simbolo com repeticao
\hspace{\baselineskip}

\textbf{[ Hierarquia de Estruturas por Repetição de Símbolos em Abertura de Problema ]} A estratégia desta forma de notação é colocar alguma marca na abertura do problema que indique a hierarquia. Por exemplo, poderíamos utilizar $\alpha :$ para indicar uma abertura de problema, para indicar subproblemas poderíamos usar $\alpha::$ e $\alpha :::$ sucessivamente.

%%%%%%%%%%%%%%%%%%%%%%%%%%%%%%%%%%%%%%%%%%%%%%%%%%%%%%%%%%%%%%%%%%%%%%%%%%%%%%%%%%%%%%%%%%%%%%%%%%%%%%%%%%%%
% Thread
\hspace{\baselineskip}

\textbf{[ Thread de Raciocínio ]} Um raciocínio pode ser interpretado como uma transformação de informação através do tempo. Frequentemente começamos um raciocínio que se prova infrutífero e precisamos recomeçar o raciocínio, é interessante indicar através de notações este processo. O propósito do thread de raciocínio é indicar explicitamente que uma sequência de informações pertencem a um mesmo raciocínio e indicar se tal raciocínio teve alguma conclusão ou se mostrou infrutífero, ou sua continuidade foi postergada.

\begin{itemize}
	\item[ (i) ] Para indicar que um raciocínio pode continuar futuramente podemos usar o símbolo de elipse (...), os três pontos.
	\item[ (ii) ] Podemos representar um thread de raciocínio por uma linha vertical, para indicar que houve alguma conclusão basta colocar um traço horinzontal separando a conclusão das outras informações. Para indicar o início de um thread podemos fazer um ponto marcado (uma bolinha preenchida). Para não confundir threads de raciocínio com marcas de indentação podemos colocar os threads à direita e as indentações à esquerda do conteúdo.
	\item[ (iii) ] É interessante que haja uma notação para indicar que um thread de raciocínio foi descartado como infrutífero, podemos usar a rasura, ou indicar através de comentários que o thread aparenta não dar resultados.
\end{itemize}

\textbf{Comentário :} Os threads de raciocínio podem ser hierarquizados, entretanto um thread de raciocínio é mais informal, apenas auxilia na interpretação de um possível leitor da mesma maneira que um comentário. Portanto não é interessante que seja feito hierarquias muito rigorosas de forma a dificultar o raciocínio e sua compreensão.

%%%%%%%%%%%%%%%%%%%%%%%%%%%%%%%%%%%%%%%%%%%%%%%%%%%%%%%%%%%%%%%%%%%%%%%%%%%%%%%%%%%%%%%%%%%%%%%%%%%%%%%%%%%%
% Thread Justificado
\hspace{\baselineskip}

\textbf{[ Thread Justificado ]} Para abstrair raciocínios podemos indicar outras informações ou definições que validem tal raciocínio, tal processo é chamado de justificação. Em geral as informações e definições são abstraídas por nomes, por isso as justificações fazem referência a tais nomes. Um thread justificado é especialmente importante para a inferência, para o processo de raciocínio usado na lógica matemática.


%%%%%%%%%%%%%%%%%%%%%%%%%%%%%%%%%%%%%%%%%%%%%%%%%%%%%%%%%%%%%%%%%%%%%%%%%%%%%%%%%%%%%%%%%%%%%%%%%%%%%%%%%%%%
\hspace{\baselineskip}

\textbf{[ Interface de Raciocínio ]} Uma vez definido as notações para os componentes de raciocínio estrutural recursivo, chamamos tal notação de interface de raciocínio. Ela medeia a resolução de um problema externalizado em um substrato real que representa algum raciocínio. É vantajoso jogarmos nossas ideias na realidade, pois podemos captar boas propriedades da realidade para ajudar a resolver problemas, por exemplo, a realidade é mais estável, portanto podemos guardar informações sem o medo de perde-las, podemos verificar, podemos comparar diferentes objetos e formular perguntas mais elaboradas e refinadas para resolver um problema e também podemos nos concentrar melhor, pois a realidade funciona como um estímulo permanente na consciência para resolução de um problema.

%%%%%%%%%%%%%%%%%%%%%%%%%%%%%%%%%%%%%%%%%%%%%%%%%%%%%%%%%%%%%%%%%%%%%%%%%%%%%%%%%%%%%%%%%%%%%%%%%%%%%%%%%%%%
%%%%%%%%%%%%%%%%%%%%%%%%%%%%%%%%%%%%%%%%%%%%%%%%%%%%%%%%%%%%%%%%%%%%%%%%%%%%%%%%%%%%%%%%%%%%%%%%%%%%%%%%%%%%
%%%%%%%%%%%%%%%%%%%%%%%%%%%%%%%%%%%%%%%%%%%%%%%%%%%%%%%%%%%%%%%%%%%%%%%%%%%%%%%%%%%%%%%%%%%%%%%%%%%%%%%%%%%%

\subsection{Definição de Estudo por Interface de Raciocínio}

\hspace{\baselineskip}

A interface de raciocínio representa toda atividade intelectual, assim o estudo deve ser olhado e externalizado nesta interface da mesma maneira que o uso de informações. Quando estamos resolvendo um problema usamos nomes de informações para acessar informações reais, portanto abstraímos os conteúdos reais em nomes ou símbolos que permitam o seu resgate, assim uma teoria sobre um determinado assunto é representado como um conjunto de nomes de informações, como nomes são abstrações não temos consciência alfa imediata, portanto dar nome a algo não implica que possamos resgatar a informação através deste nome, é necessário um processo chamado estudo para criar conexões suficientes que liguem o nome às propriedades desta informação.\\

\textbf{[ Trialética de Estudo ]} 

\begin{itemize}
	\item[ (i) ] \textbf{[ Declaração de Nomes ]} Criar um conjunto de nomes de informações de uma teoria através de uma leitura superficial.
	\item[ (ii) ] \textbf{[ Definição de Nomes ]} Treinar para cada nome de informação o resgate do conteúdo preciso e efetivar as abstrações feitas na declaração.
	\item[ (iii) ] \textbf{[ Dialética do Estudo Tradicional ]}
	\begin{itemize}
		\item[ (i) ] \textbf{[ Análise ]} Compreender, julgar o sentido, julgar o conteúdo de verdade em problemas, em raciocínios, ou no texto original da teoria.
		\item[ (ii) ] \textbf{[ Aplicação ]} Utilizar as informações para resolver problemas, fazer exercícios, refazer raciocínios ou criar outros elementos.
	\end{itemize}
\end{itemize}

\textbf{Comentário [ Programação ]} Para quem tem afinidade com programação, o processo foi inspirado no processo análogo de declaração de variáveis e funções e definição de variáveis e funções definidos pela ciência da computação.\\

\textbf{Comentário [ Memorização ]} Efetivamente existe três estudos distintos além do processo de declaração de nomes, existe a etapa de memorização feita na definição de nomes (componente alfa da trialética), a etapa de análise dos raciocínios (componente gama da trialética) e a etapa de uso da teoria para resolver problemas (componente beta da trialética). Para muitos a internet eliminou a necessidade do estudo de memorização, de fato muitas informações não precisam mais serem definidas, apenas declaradas, entretanto a consciência do uso de uma determinada informação requer uma consciência alfa do conteúdo declarado, sem isso o computador se torna um baú cheio de papéis velhos e cheios de traças com informações mantidas na ignorância, é necessário algum estudo para que a informação seja utilizada, quanto maior a consciência alfa de uma informação, maior a capacidade de reconhecer aplicações, maior a capacidade analítica e a capacidade de formular problemas.\\

\textbf{Comentário [ Algoritmos de Memorização ]} A verdade que a computação não elimina a necessidade de memorização, ela permite um estudo de memorização mais eficiente através de algoritmos de estudo. É possível construir um algoritmo bem simples usando gerador aleatório para sortear nomes de informações de um inventário de nomes que represente as informações de teoria, gosto de chamar tais programas de Brainwashers, fazendo referência a lavagem cerebral visto em filmes de ficção em que uma pessoa vítima de lavagem cerebral é exposta contra a vontade a estímulos capazes de refazer memórias, apesar do nome, os brainwashers tem propósito mais modesto, estabelecer a consciência alfa entre o símbolo e a existência representada. Recomendo o uso de tais algoritmos para qualquer campo intelectual de nível superior. Exemplos de softwares de memorização disponíveis (ANKI, Memrise, Mnemosyne, SuperMemo), entretanto recomendo que seu interesse em programação o leve a criar um programa próprio para tal finalidade.

%%%%%%%%%%%%%%%%%%%%%%%%%%%%%%%%%%%%%%%%%%%%%%%%%%%%%%%%%%%%%%%%%%%%%%%%%%%%%%%%%%%%%%%%%%%%%%%%%%%%%%%%%%%%
%%%%%%%%%%%%%%%%%%%%%%%%%%%%%%%%%%%%%%%%%%%%%%%%%%%%%%%%%%%%%%%%%%%%%%%%%%%%%%%%%%%%%%%%%%%%%%%%%%%%%%%%%%%%
%%%%%%%%%%%%%%%%%%%%%%%%%%%%%%%%%%%%%%%%%%%%%%%%%%%%%%%%%%%%%%%%%%%%%%%%%%%%%%%%%%%%%%%%%%%%%%%%%%%%%%%%%%%%

%\section{Descrição Estrutural de Teoria}

%\subsection{Estruturação de Razões Alfa - Lembrança}

%\subsection{Estruturação de Razões Beta - Modelagem}

%\subsection{Estruturação de Razões Gama - Inferência}

%%%%%%%%%%%%%%%%%%%%%%%%%%%%%%%%%%%%%%%%%%%%%%%%%%%%%%%%%%%%%%%%%%%%%%%%%%%%%%%%%%%%%%%%%%%%%%%%%%%%%%%%%%%%
%%%%%%%%%%%%%%%%%%%%%%%%%%%%%%%%%%%%%%%%%%%%%%%%%%%%%%%%%%%%%%%%%%%%%%%%%%%%%%%%%%%%%%%%%%%%%%%%%%%%%%%%%%%%
%%%%%%%%%%%%%%%%%%%%%%%%%%%%%%%%%%%%%%%%%%%%%%%%%%%%%%%%%%%%%%%%%%%%%%%%%%%%%%%%%%%%%%%%%%%%%%%%%%%%%%%%%%%%
%%%%%%%%%%%%%%%%%%%%%%%%%%%%%%%%%%%%%%%%%%%%%%%%%%%%%%%%%%%%%%%%%%%%%%%%%%%%%%%%%%%%%%%%%%%%%%%%%%%%%%%%%%%%
%%%%%%%%%%%%%%%%%%%%%%%%%%%%%%%%%%%%%%%%%%%%%%%%%%%%%%%%%%%%%%%%%%%%%%%%%%%%%%%%%%%%%%%%%%%%%%%%%%%%%%%%%%%%
%%%%%%%%%%%%%%%%%%%%%%%%%%%%%%%%%%%%%%%%%%%%%%%%%%%%%%%%%%%%%%%%%%%%%%%%%%%%%%%%%%%%%%%%%%%%%%%%%%%%%%%%%%%%

\section{Tipos de Teorias}

\hspace{\baselineskip}

Nesta subseção irei descrever os tipos de teoria com base na interface de raciocínio. A teoria mais simples é a teoria de nomes, ela busca apenas abstrair informações complicadas em nomes. Uma teoria de nomes não requer nenhuma conexão com a realidade, é um a teoria que sua ênfase está na consciência alfa da informação, é extremamente importante para teorias do tipo alfa técnicas de memorização. Como exemplo de teoria do tipo beta temos o que chamo de teoria homeostática, ela é uma teoria baseada no conceito de tendência, de ativação e inibição de eventos. Uma teoria homeostática fornece uma previsão curta, imprecisa da realidade, entretanto sob ponto de vista prático ela permite a tomada de decisões, uma teoria homeostática é essencialmente uma teoria em lógica informal, a teoria homeostática define uma interface para criação de estratégias. Como teorias do tipo gama em que a verdade e o raciocínio é enfatizado temos as teorias inferente, dedutiva e axiomática, sendo o terceiro tipo a teoria mais completa dentre as três.

%%%%%%%%%%%%%%%%%%%%%%%%%%%%%%%%%%%%%%%%%%%%%%%%%%%%%%%%%%%%%%%%%%%%%%%%%%%%%%%%%%%%%%%%%%%%%%%%%%%%%%%%%%%%
%%%%%%%%%%%%%%%%%%%%%%%%%%%%%%%%%%%%%%%%%%%%%%%%%%%%%%%%%%%%%%%%%%%%%%%%%%%%%%%%%%%%%%%%%%%%%%%%%%%%%%%%%%%%
%%%%%%%%%%%%%%%%%%%%%%%%%%%%%%%%%%%%%%%%%%%%%%%%%%%%%%%%%%%%%%%%%%%%%%%%%%%%%%%%%%%%%%%%%%%%%%%%%%%%%%%%%%%%

\subsection{Teoria de Nomes}

\hspace{\baselineskip}

Apesar de ser o tipo mais simples é o mais elementar, no sentido de que toda teoria construída é também uma teoria de nomes. Uma boa teoria de nomes deve permitir uma memorização rápida de seu conteúdo e uma boa estabilidade na memória.

%%%%%%%%%%%%%%%%%%%%%%%%%%%%%%%%%%%%%%%%%%%%%%%%%%%%%%%%%%%%%%%%%%%%%%%%%%%%%%%%%%%%%%%%%%%%%%%%%%%%%%%%%%%%
% Dialética de Consciência Alfa de Informação
\hspace{\baselineskip}

\textbf{Dialética [ de Consciência Alfa de Informação ]}

\begin{itemize}
	\item[ (i) ] \textbf{[ Consciência Alfa Vertical ]} É sentir a informação através do seu nome, é sentir a informação através da sua representação.
	\item[ (ii) ] \textbf{[ Consciência Alfa Horizontal ]} Suponha uma teoria de nomes, os nomes funcionam como forma externa das informações, de forma que tais formas externas possam interagir, a consciência alfa horizontal de informação é sentir através das interações entre tais formas externas o conteúdo real de informação.
\end{itemize}

\textbf{Comentário :} A consciência alfa horizontal é especialmente importante para teorias reais (teorias homeostáticas e teorias dedutivas). Uma teoria exclusivamente de nomes provavelmente não existe, pois o propósito de uma teoria reside em representar interações horizontais entre existências, o que pode existir é são teorias de nomes auxiliares dentro de uma teoria dedutiva ou homeostática.

%%%%%%%%%%%%%%%%%%%%%%%%%%%%%%%%%%%%%%%%%%%%%%%%%%%%%%%%%%%%%%%%%%%%%%%%%%%%%%%%%%%%%%%%%%%%%%%%%%%%%%%%%%%%
%%%%%%%%%%%%%%%%%%%%%%%%%%%%%%%%%%%%%%%%%%%%%%%%%%%%%%%%%%%%%%%%%%%%%%%%%%%%%%%%%%%%%%%%%%%%%%%%%%%%%%%%%%%%
%%%%%%%%%%%%%%%%%%%%%%%%%%%%%%%%%%%%%%%%%%%%%%%%%%%%%%%%%%%%%%%%%%%%%%%%%%%%%%%%%%%%%%%%%%%%%%%%%%%%%%%%%%%%

\subsection{Teoria Inferente}

\hspace{\baselineskip}

Uma teoria de nomes é uma coleção de informações das definições dos nomes, uma teoria de nomes é uma teoria estática, você não estende a lista das informações a não ser que se defina novos nomes. Uma teoria inferente cria um mecanismo para inserção de outras informações. Além de nomes (definições) uma teoria inferente tem as inferências, são mecanismos que permitem adicionar outras informações a teoria.

%%%%%%%%%%%%%%%%%%%%%%%%%%%%%%%%%%%%%%%%%%%%%%%%%%%%%%%%%%%%%%%%%%%%%%%%%%%%%%%%%%%%%%%%%%%%%%%%%%%%%%%%%%%
% Inferência
\hspace{\baselineskip}

\textbf{[ Inferência ]} Uma inferência é uma sequência de fórmulas. Uma inferência tem duas partes, as hipóteses e a conclusão. Uma inferência específica permite sua reconstrução completa (hipóteses mais conclusão) a partir das hipóteses. Em computação isso significa que existe um algoritmo bem definido para obter a conclusão a partir das hipóteses. \\

\textbf{Comentário :} Por hora uma inferência não está dizendo nada sobre raciocínios verdadeiros ou falsos, uma inferência é usada para tal finalidade, entretanto sua definição é anterior e não tem nenhuma responsabilidade em preservar raciocínios verdadeiros, o objetivo é definir uma sequência de fórmulas. Inferência deve ser vista inicialmente como uma tentativa de programar um computador para que efetue raciocínios. 

%%%%%%%%%%%%%%%%%%%%%%%%%%%%%%%%%%%%%%%%%%%%%%%%%%%%%%%%%%%%%%%%%%%%%%%%%%%%%%%%%%%%%%%%%%%%%%%%%%%%%%%%%%%%
%%%%%%%%%%%%%%%%%%%%%%%%%%%%%%%%%%%%%%%%%%%%%%%%%%%%%%%%%%%%%%%%%%%%%%%%%%%%%%%%%%%%%%%%%%%%%%%%%%%%%%%%%%%%
%%%%%%%%%%%%%%%%%%%%%%%%%%%%%%%%%%%%%%%%%%%%%%%%%%%%%%%%%%%%%%%%%%%%%%%%%%%%%%%%%%%%%%%%%%%%%%%%%%%%%%%%%%%%
\hspace{\baselineskip}

\textbf{Estrutura [ Razão, Inferência e Consciência ]}

\begin{itemize}
	\item[ (i) ] \textbf{[ Razão ]} É uma sequência de informações que pode representar todas as informações verdadeiras de uma teoria. Componente Alfa.
	\item[ (ii) ] \textbf{[ Inferência ]} Inferência são mecanismos que permite uma consciência estender a razão adicionando novas informações. É interessante que a inferência mantenha a propriedade de verdade da razão. Componente Gama.
	\item[ (iii) ] \textbf{[ Consciência ]} É aquele que interpreta a razão e utiliza a inferência para obter novas informações. Componente Beta.
\end{itemize}

\textbf{Comentário :} Inferência é considerado uma informação, mas uma informação não é uma inferência, ela pode possuir regras de inferência associadas. Inferência faz a interface entre a razão e a consciência. Para a consciência estender a razão é necessário habilidade de reconhecer na razão os componentes das hipóteses e verificar se existem inferências que podem ser aplicadas.

%%%%%%%%%%%%%%%%%%%%%%%%%%%%%%%%%%%%%%%%%%%%%%%%%%%%%%%%%%%%%%%%%%%%%%%%%%%%%%%%%%%%%%%%%%%%%%%%%%%%%%%%%%%%
%%%%%%%%%%%%%%%%%%%%%%%%%%%%%%%%%%%%%%%%%%%%%%%%%%%%%%%%%%%%%%%%%%%%%%%%%%%%%%%%%%%%%%%%%%%%%%%%%%%%%%%%%%%%
%%%%%%%%%%%%%%%%%%%%%%%%%%%%%%%%%%%%%%%%%%%%%%%%%%%%%%%%%%%%%%%%%%%%%%%%%%%%%%%%%%%%%%%%%%%%%%%%%%%%%%%%%%%%

\subsection{Teoria Homeostática}

\hspace{\baselineskip}

A teoria homeostática é inspirada no conceito fisiológico (biológico) de homeostase, comportamento presente em sistemas biológicos que busca preservar a condição anterior a uma determinada pertubação, tende a preservar o equilíbrio do sistema através de respostas contrárias. A homeostase externaliza a estabilidade, componente fundamental de uma existência.

\hspace{\baselineskip}

\textbf{Dialética [ Teoria Homeostática ]} 

\begin{itemize}
	\item[ (i) ] \textbf{[ Componente Vertical ]} É dado o nome de \textbf{Evento} a classe de informações encapsuladas de forma vertical (consciência alfa horizontal, abstração).
	\item[ (ii) ] \textbf{[ Componente Horizontal ]} São relações entre eventos.
	\begin{itemize}
		\item[ (i) ] \textbf{[ Ativação ]} É dito que um evento ativa outra existência objeto se ele acentua o objeto, se ele torna tal objeto mais estável corroborando para sua existência dentro de um contexto.
		\item[ (ii) ] \textbf{[ Inibição ]} É dito que um evento inibe outra existência objeto se ele reduz a estabilidade de tal objeto, se ele prejudica tal objeto, se ele destrói ou impede que tal objeto exista dentro de um contexto.
	\end{itemize}
\end{itemize}

\textbf{Comentário :} Uma teoria homeostática possui além de nomes (componentes verticais) que são chamados de definições, ele possui o que podemos chamar de informações (componentes horizontais). É possível modelar processos e eventos argumentando com base nas informações de sistemas modelados por uma teoria homeostática. A biologia molecular, bioquímica, fisiologia são exemplos refinados de teorias homeostáticas, a lei de oferta e procura é um exemplo simples da teoria econômica que também se enquadra em uma teoria homeostática.

%%%%%%%%%%%%%%%%%%%%%%%%%%%%%%%%%%%%%%%%%%%%%%%%%%%%%%%%%%%%%%%%%%%%%%%%%%%%%%%%%%%%%%%%%%%%%%%%%%%%%%%%%%%%
% Raciocínio Preditivo de uma teoria Homeostática
\hspace{\baselineskip}

\textbf{Raciocínio Preditivo  [ Fundamental da Teoria Homeostática ]} 

\begin{itemize}
	\item[ (i) ] Tenha modelado um sistema que você quer prever situações com base em eventos e informações de ativação e inibição.
	\item[ (ii) ] Suponha uma situação real que é vista inicialmente em equilíbrio. 
	\item[ (iii) ] O raciocínio preditivo da teoria homeostática consiste em imaginar o que acontece através de pequenas pertubações neste sistema, tais pertubações podem ser ativações ou inibições de eventos singulares.
\end{itemize}

\textbf{Comentário :} Uma teoria homeostática pode ser usada como prevenção de eventos indesejados. Sua finalidade consiste em de alguma maneira criar estratégias que reforcem a própria existência do sujeito (do que usa a teoria). Naturalmente a teoria homeostática pode também ser usada para ativar eventos desejados, eventos estes que reforçam a própria existência do sujeito.

%%%%%%%%%%%%%%%%%%%%%%%%%%%%%%%%%%%%%%%%%%%%%%%%%%%%%%%%%%%%%%%%%%%%%%%%%%%%%%%%%%%%%%%%%%%%%%%%%%%%%%%%%%%%
% raciocínio Preditivo para Prevenção de Eventos Indesejados
\hspace{\baselineskip}

\textbf{[ Eventos Próximos ]} São eventos que o sujeito tem maior controle (pode diretamente, ou mais facilmente interferir (ativar ou inibir) ). Faz contraste com Eventos Distantes que são eventos que o sujeito tem menor controle, ou controle indireto.\\

\textbf{[ Ações ]} Dependendo do contexto, eventos próximos ao sujeito também são chamados de ações. \\

\textbf{[ Problemas ]} Eventos distantes e indesejáveis ao sujeito são chamados de problemas.\\

\textbf{[ Soluções ]} Uma rede que conecte uma ação a um problema é chamado de solução.\\

\textbf{Raciocínio Preditivo [ Prevenção de Eventos Indesejados ]} O raciocínio de prevenção de eventos indesejados dentro de uma teoria homeostática é identificar eventos indesejados, estabelecer uma teoria de ativação e inibição através de eventos próximos e criar estratégias (ativação ou inibição de eventos próximos) que inibam os eventos indesejados distantes.\\

\textbf{Comentário :} Esta é uma aplicação prática de uma teoria homeostática, não requer determinismo, pois sob o ponto de vista preventivo, basta a redução das chances de um evento indesejado acontecer.

%%%%%%%%%%%%%%%%%%%%%%%%%%%%%%%%%%%%%%%%%%%%%%%%%%%%%%%%%%%%%%%%%%%%%%%%%%%%%%%%%%%%%%%%%%%%%%%%%%%%%%%%%%%%
% Raciocínio Preditivo para Identificação de Problemas Existenciais
\hspace{\baselineskip}

\textbf{[ Problema Existencial ]} Um problema (evento indesejado) é dito existencial se ela forma uma rede de ativação com outro problema, ou seja, pelo menos dois problemas distintos corroboram entre si para estabilidade mútua.\\

\textbf{Comentário :} Um problema existencial é mais estável que um problema isolado, portanto a estratégia de prevenção deve ser diferente em relação a um problema isolado. Sob o ponto de vista teórico, identificar problemas existenciais é interessante como etapa anterior a tentativa de criar soluções.\\

\textbf{Raciocínio [ Identificação de Problemas Existenciais ]} O problema existencial induz a necessidade de identificar tais problemas, o nome dado a este processo é raciocínio de identificação de problemas existenciais. A identificação de problemas existenciais é um componente de teorização, portanto tem um aspecto intelectual, diferente do aspecto prático do raciocínio de prevenção de eventos indesejados.

%%%%%%%%%%%%%%%%%%%%%%%%%%%%%%%%%%%%%%%%%%%%%%%%%%%%%%%%%%%%%%%%%%%%%%%%%%%%%%%%%%%%%%%%%%%%%%%%%%%%%%%%%%%%
% Raciocínio Preditivo [ Ativação de Eventos Desejados ]
\hspace{\baselineskip}

\textbf{Raciocínio Preditivo [ Ativação de Eventos Desejados ]} Analogamente temos a estratégia de ativação de eventos desejados, em que são criados soluções que conectem eventos desejados distantes a ações do sujeito.

%%%%%%%%%%%%%%%%%%%%%%%%%%%%%%%%%%%%%%%%%%%%%%%%%%%%%%%%%%%%%%%%%%%%%%%%%%%%%%%%%%%%%%%%%%%%%%%%%%%%%%%%%%%%
% Indeterminismo por Caminhos
\hspace{\baselineskip}

\textbf{[ Indeterminismo por Caminhos Contraditórios ]} Suponha $S$ sujeito, $B$ evento distante desejado, $A$ uma ação tal que $A \rightarrow B$ e $A \rightarrow E$ e $E \overset{inibe}{\rightarrow} B$. Temos duas soluções entre $A$ e $B$, a primeira $A \rightarrow B$ implica que ativar $A$ é desejável, a segunda $A \rightarrow E \overset{inibe}{\rightarrow} B$ implica que inibir $A$ é desejável. Não podemos dizer nada sobre a ação $A$, pois não sabemos qual dos caminhos é mais relevante.\\

\textbf{Comentário :} Quanto mais refinado um modelo, maiores as chances de criar teorias com soluções indeterminadas. A biologia soluciona este problema através da experimentação, basta verificar através do método científico qual caminho é mais relevante e qual pode ser desprezado em determinado contexto. As ciências biológicas sob este ponto de vista parecem ser inerentemente científicas (há necessidade de experimentação), um argumento que corrobora com isso é que o conteúdo teórico da biologia consiste em redes relevantes de uma teoria homeostática, entretanto a biologia é mutável por conta da evolução, sempre pode surgir alguma variante genética que modifique quais caminhos são relevantes ou não, induzindo a necessidade de experimentação para verificar inconsistências.\\

\textbf{Comentário :} Tal inerente empirismo é algo considerado senso comum na biologia, entretanto não é preservado na teoria econômica, que também possui componentes homeostáticos. A ciência econômica é peculiar, pois o método científico não consegue ser aplicado (não temos como construir sistemas isolados, em geral o que é analisado são séries históricas que passam longo de um experimento científico propriamente dito). Apesar de não poder ser feito experimentos ao rigor do método científico, é possível obter respostas mais fracas usando a estatística através da análise de séries históricas para discriminar entre os diferentes caminhos o caminho mais relevante.

%%%%%%%%%%%%%%%%%%%%%%%%%%%%%%%%%%%%%%%%%%%%%%%%%%%%%%%%%%%%%%%%%%%%%%%%%%%%%%%%%%%%%%%%%%%%%%%%%%%%%%%%%%%%
% Raciocínio Transitivo
\hspace{\baselineskip}

\textbf{Raciocínio [ Transitividade de Informações ]} Se temos $A \overset{(1)}{\rightarrow} B$ e $B \overset{(2)}{\rightarrow} C$, podemos obter uma informação derivada $A \overset{(1)*(2)}{\rightarrow} C$ tal que se (1) e (2) forem ativações então (1)*(2) será também uma ativação, se (1) e (2) forem inibições então (1)*(2) será uma ativação, se apenas um deles for inibição e o outro ativação temos (1)*(2) como inibição. Esse processo de gerar informações derivadas é chamado de raciocínio transitivo.

%%%%%%%%%%%%%%%%%%%%%%%%%%%%%%%%%%%%%%%%%%%%%%%%%%%%%%%%%%%%%%%%%%%%%%%%%%%%%%%%%%%%%%%%%%%%%%%%%%%%%%%%%%%%
% Justificação de Raciocínio Transitivo com base na Consciência Alfa de uma teoria
\hspace{\baselineskip}

\textbf{Comentário [ Raciocínio Transitivo ]} Suponha que tenhamos memorizado os componentes de uma teoria individualmente \{ $A \rightarrow B$ , $B \rightarrow C$ \}, portanto temos consciência alfa de cada componente individualmente, o que significa dizer que sabemos a relação entre \{$A, B$\} e sabemos entre \{$B, C$\}, entretanto não sabemos sobre a relação entre \{$A,C$\}, a não ser que efetuemos o raciocínio transitivo $A \rightarrow C$, portanto não é imediato concluir nada de $A$ e $C$, o que faz a inferência do raciocínio transitivo importante para reduzir a quantidade de informações a serem memorizadas necessárias para resolver um problema.

%%%%%%%%%%%%%%%%%%%%%%%%%%%%%%%%%%%%%%%%%%%%%%%%%%%%%%%%%%%%%%%%%%%%%%%%%%%%%%%%%%%%%%%%%%%%%%%%%%%%%%%%%%%%
% Informação Direta
\hspace{\baselineskip}

\textbf{[ Informação Direta ]} É uma informação entre eventos que são próximos temporalmente ou espacialmente de forma que interferências tem menor chance de acontecer. Um caminho formado por caminhos diretos não é uma informação direta, pois há maior chance de ocorrer interferências.\\

\textbf{Comentário :} A ideia da informação direta é que quanto mais etapas de informações diretas entre uma ação e um evento distante, mais distantes temporalmente ou espacialmente a ação vai estar do evento, o que aumenta as chances de interferências e torna a solução menos relevante. Podemos considerar uma probabilidade de não ocorrer interferência associado a uma informação direta e considerar a probabilidade de eventos compostos como a multiplicação das probabilidades de cada evento direto.

\begin{itemize}
	\item[ (i) ] \textbf{[ Informação Direta com $0.9$ de probabilidade ]} $ (1)0.9 ; (2)0.81 ; (3)0.729; (4)0.6561; (5)0.59049 ; (6)0.531441 ; (7)\sim0.47 $. Portanto um caminho com 7 informações diretas não é mais confiável.
	\item[ (ii) ] \textbf{[ Informação Direta com $0.8$ de probabilidade ]} $ (1)0.8 ; (2)0.64; (3)0.512; (4)0.4056 $. Portanto um caminho com 4 informações diretas não é mais confiável.
	\item[ (iii) ] \textbf{[ Informação Direta com $0.7$ de probabilidade ]} $ (1)0.7; (2)0.49 $. Portanto não é possível obter novas informações por raciocínio transitivo, pois apenas com duas etapas a informação não é mais confiável.
\end{itemize}

%%%%%%%%%%%%%%%%%%%%%%%%%%%%%%%%%%%%%%%%%%%%%%%%%%%%%%%%%%%%%%%%%%%%%%%%%%%%%%%%%%%%%%%%%%%%%%%%%%%%%%%%%%%%
% Teoria Homeostática Topológica
%\hspace{\baselineskip}

\textbf{Teoria Homeostática [ Topológica ]} É uma teoria homeostática em que temos informações diretas.\\

\textbf{Teoria Homeostática Topológica [ Transitiva ]} É uma teoria homeostática topológica em que podemos gerar novas informações por transitividade (informação direta tem probabilidade maior que 0.7).\\

\textbf{Comentário :} Teorias homeostáticas topológicas transitivas são as teorias homeostáticas desejadas, perde-se o propósito da teoria caso contrário, portanto quando me referir a teorias homeostáticas topológicas estarei implicitamente considerando que sejam transitivas. Entretanto não deve ser dado muito rigor ao cálculo de probabilidade, pois neste documento serve apenas como uma ilustração da ideia.\\

\textbf{Comentário [ Lógica Formal ]} A lógica formal, a lógica usada pela matemática para construir teorias axiomáticas pode ser vista como uma teoria homeostática transitiva em que o raciocínio transitivo tem probabilidade 1 (100 porcento de chance, sempre irá ocorrer), portanto podemos efetuar raciocínios transitivos quantos for necessário e seu resultado também será certo de acontecer. Por conta disso a lógica usada na matemática é uma forma de raciocínio bem particular e não é aplicável para todos os campos do conhecimento humano, a teoria homeostática generaliza a lógica usada na matemática.

%%%%%%%%%%%%%%%%%%%%%%%%%%%%%%%%%%%%%%%%%%%%%%%%%%%%%%%%%%%%%%%%%%%%%%%%%%%%%%%%%%%%%%%%%%%%%%%%%%%%%%%%%%%%
% Definição de Estratégia
\hspace{\baselineskip}

\textbf{[ Estratégia ]} É uma teoria homeostática descrita como uma lista de pares (Evento, Ação), o evento é alguma ocorrência observável no presente ou no passado ou uma previsão futura, a ação é a ação sugerida ou inibida pela ação. Uma estratégia tem um objetivo definido, afim de realizar tal objetivo a estratégia cria uma teoria homeostática em que para alcançar o objetivo alguém deve ao reconhecer um determinado evento efetuar uma determinada ação.\\

\textbf{Estratégia [ Lemas ]} É uma estratégia com objetivo de eliminar a perda de tempo em tomada de decisões fazendo um registro e definição de ações para problemas recorrentes de decisão.\\

\textbf{Comentário :} Estes tipos de teorias homeostáticas são formas de criar paradigmas, fazendo analogia com computação, é a maneira de programar a nós mesmos. A aplicação de estratégias permite agilidade na tomada de decisões, é uma forma genérica e informal de teorização com natureza normativa. Estes tipos de teorias são úteis apenas quando há sobrecarga de informações.


%%%%%%%%%%%%%%%%%%%%%%%%%%%%%%%%%%%%%%%%%%%%%%%%%%%%%%%%%%%%%%%%%%%%%%%%%%%%%%%%%%%%%%%%%%%%%%%%%%%%%%%%%%%%
%%%%%%%%%%%%%%%%%%%%%%%%%%%%%%%%%%%%%%%%%%%%%%%%%%%%%%%%%%%%%%%%%%%%%%%%%%%%%%%%%%%%%%%%%%%%%%%%%%%%%%%%%%%%
%%%%%%%%%%%%%%%%%%%%%%%%%%%%%%%%%%%%%%%%%%%%%%%%%%%%%%%%%%%%%%%%%%%%%%%%%%%%%%%%%%%%%%%%%%%%%%%%%%%%%%%%%%%%

\subsection{Teoria Dedutiva}

\hspace{\baselineskip}

Na teoria homeostática temos o raciocínio transitivo como mecanismo de inferência, tal mecanismo tem a vantagem de permitir reduzir a quantidade de informações a serem memorizadas para resolver um problema. Na teoria dedutiva teremos outro mecanismo com a mesma propriedade, entretanto diferente do raciocínio transitivo que resulta em informações prováveis, o raciocínio dedutivo terá conclusões certas.\\

\textbf{[ Informações Universais ]} Uma informação universal possui componentes que fazem referências a vários objetos, por exemplo (um gato é um animal) possui dois componentes com caráter universal \{ animal, gato \}, o primeiro mais genérico e o segundo genérico, porém particular em relação ao primeiro.\\

\textbf{Comentário :} A teoria dedutiva tem seu raciocínio baseado na relação que existe entre o particular e o universal, assim uma teoria com informações universais se verdadeira implica necessariamente que as formas particulares sejam também verdadeiras.\\

\textbf{[ Informação Particular ]} Faz dialética contrapositiva com informação universal, uma informação universal diz sobre várias informações particulares.\\

\textbf{Comentário :} A linguagem verbal humana permite a construção de teoria dedutivas de maneira natural, uma palavra em geral faz referência a uma classe de objetos ou entidades, quando queremos dizer sobre algo particular a gramática cria regras de sintaxe para podermos expressar esta vontade.




%%%%%%%%%%%%%%%%%%%%%%%%%%%%%%%%%%%%%%%%%%%%%%%%%%%%%%%%%%%%%%%%%%%%%%%%%%%%%%%%%%%%%%%%%%%%%%%%%%%%%%%%%%%%
% Teoria Dedutiva
\hspace{\baselineskip}

\textbf{Teoria [ Dedutiva ]} É uma teoria cuja razão é composta por informações universais. A dedução é o processo de adicionar a razão informação que são particulares à informação universal contida na razão da teoria. % É uma teoria cuja razão é composta por informações universais, estas podendo ser representadas por fórmulas fechadas. A dedução é o processo de adicionar a razão através da inferência de dedução a partir de uma informação universal contida na razão da teoria.


%%%%%%%%%%%%%%%%%%%%%%%%%%%%%%%%%%%%%%%%%%%%%%%%%%%%%%%%%%%%%%%%%%%%%%%%%%%%%%%%%%%%%%%%%%%%%%%%%%%%%%%%%%%%
%%%%%%%%%%%%%%%%%%%%%%%%%%%%%%%%%%%%%%%%%%%%%%%%%%%%%%%%%%%%%%%%%%%%%%%%%%%%%%%%%%%%%%%%%%%%%%%%%%%%%%%%%%%%

\subsection{Teoria Axiomática}

\hspace{\baselineskip}

Uma teoria axiomática é uma teoria dedutiva que projetada na estrutura razão, inferência e consciência, é visto pela consciência de um indivíduo como uma razão com um conjunto de informações iniciais que são chamadas axiomas, nestes axiomas temos regras de inferência assim como informações consideradas verdadeiras. A lógica é a teoria axiomática principal, todas as outras teorias axiomáticas herdam seus axiomas.\\

Dentro de uma teoria axiomática temos como enxergar diferentes razões. Temos a razão exposta como texto, temos a razão contendo apenas os axiomas e regras de inferência, temos a razão intangível que é o conjunto de todas as possíveis razões efetuando infinitas inferências e temos também a razão em que elementos de outras teorias, ou elementos de um problema específico são colocadas juntas da teoria com a finalidade de obter uma aplicação da teoria para um determinado problema.\\

\textbf{Exposição [ de Teoria ]} Uma teoria axiomática é exposta como texto com o objetivo de que o leitor construa uma razão em sua consciência iniciando pelos axiomas e regras de inferência e que cada nova informação exposta pelo texto o expositor deve demonstrar que tal nova informação pode ser obtida através dos axiomas e regras de inferência em uma sequência finita de etapas de inferência.

\begin{itemize}
	\item[ (i) ] \textbf{[ Proposição ]} É uma informação que o expositor deseja adicionar a razão e que é obtida pelos axiomas e regras de inferência através de finitas etapas de inferência. Existem outros nomes associados ao propósito ou importância (Teorema, Lema) que possuem mesmo significado que proposição.
	\item[ (ii) ] \textbf{[ Demonstração ]} É onde o expositor deseja demonstrar através de uma sequência finita de etapas de inferência a proposição.
\end{itemize}

\textbf{Comentário :} A medida que uma teoria se torna muito complicada a compreensão de uma demonstração pode se tornar laboriosa, dentro da filosofia deste livro existe o que se chama de dialética demonstrativa, ela busca simplificar a compreensão da demonstração criando uma linguagem para explicitar as estratégias utilizadas na demonstração. Como já deve ter ficado evidente, a base trialética é considerada como uma base mais eficiente que uma base dialética de dois componentes, assim na exposição trialética de uma teoria, além da proposição e da demonstração deve ser exposto a dialética da demonstração e também a construção, organização e classificação de dialéticas comuns encontradas nas demonstrações da teoria.

%%%%%%%%%%%%%%%%%%%%%%%%%%%%%%%%%%%%%%%%%%%%%%%%%%%%%%%%%%%%%%%%%%%%%%%%%%%%%%%%%%%%%%%%%%%%%%%%%%%%%%%%%%%%
% Exposição Trialética
\hspace{\baselineskip}

\textbf{Exposição Trialética [ de Teoria ]} A exposição trialética de uma teoria axiomática além da exposição de definições, axiomas, regras de inferência, proposições e demonstrações, ela deve lidar com a consciência alfa humana através da construção e exposição de dialéticas demonstrativas e de resolução de problemas.

%%%%%%%%%%%%%%%%%%%%%%%%%%%%%%%%%%%%%%%%%%%%%%%%%%%%%%%%%%%%%%%%%%%%%%%%%%%%%%%%%%%%%%%%%%%%%%%%%%%%%%%%%%%%
%%%%%%%%%%%%%%%%%%%%%%%%%%%%%%%%%%%%%%%%%%%%%%%%%%%%%%%%%%%%%%%%%%%%%%%%%%%%%%%%%%%%%%%%%%%%%%%%%%%%%%%%%%%%
%%%%%%%%%%%%%%%%%%%%%%%%%%%%%%%%%%%%%%%%%%%%%%%%%%%%%%%%%%%%%%%%%%%%%%%%%%%%%%%%%%%%%%%%%%%%%%%%%%%%%%%%%%%%
%%%%%%%%%%%%%%%%%%%%%%%%%%%%%%%%%%%%%%%%%%%%%%%%%%%%%%%%%%%%%%%%%%%%%%%%%%%%%%%%%%%%%%%%%%%%%%%%%%%%%%%%%%%%
%%%%%%%%%%%%%%%%%%%%%%%%%%%%%%%%%%%%%%%%%%%%%%%%%%%%%%%%%%%%%%%%%%%%%%%%%%%%%%%%%%%%%%%%%%%%%%%%%%%%%%%%%%%%
%%%%%%%%%%%%%%%%%%%%%%%%%%%%%%%%%%%%%%%%%%%%%%%%%%%%%%%%%%%%%%%%%%%%%%%%%%%%%%%%%%%%%%%%%%%%%%%%%%%%%%%%%%%%

\section{Teoria Inferente [ Teoria de Substituição ]}

\hspace{\baselineskip}

A teoria inferente de substituição é uma teoria com uma proposta bem simples, definir a notação utilizada para escrever informações (expressões, variáveis, operadores, relações) e definir as chamadas inferências de substituição que permite escrever informações equivalentes com base na substituição de nomes de variáveis ou funções.

%%%%%%%%%%%%%%%%%%%%%%%%%%%%%%%%%%%%%%%%%%%%%%%%%%%%%%%%%%%%%%%%%%%%%%%%%%%%%%%%%%%%%%%%%%%%%%%%%%%%%%%%%%%%
% Substantivo
\hspace{\baselineskip}

\textbf{[ Variável ]} É uma palavra ou símbolo para representar uma existência dentro de uma expressão.\\

\textbf{Inferência [ Reconhecimento de Variável ]} Sempre podemos reconhecer que uma variável está presente em uma expressão. Da mesma forma saber que não está presente.

%\textbf{[ Lista ]} É uma sequência finita de variáveis.

%%%%%%%%%%%%%%%%%%%%%%%%%%%%%%%%%%%%%%%%%%%%%%%%%%%%%%%%%%%%%%%%%%%%%%%%%%%%%%%%%%%%%%%%%%%%%%%%%%%%%%%%%%%%
% Liberdade de Criar Nomes
\hspace{\baselineskip}

\textbf{[ Liberdade ]} É a capacidade de estender a razão sem mecanismos de inferência.\\

\textbf{Liberdade [ de Definição ]} Podemos sempre criar definições $a := b$ para qualquer $b$ expressão. Tal que $a$ não deve ter sido utilizado anteriormente na razão e ele deve representar um nome e não uma expressão geral.\\

\textbf{Liberdade [ de Repetição ]} Qualquer elemento da razão pode ser repetido e adicionado novamente a razão.\\

\textbf{[ Operador ]} Um operador conecta variáveis, suponha a notação $f(\phi,\psi)$, $f$ é um operador que conecta as variáveis $\phi$ e $\psi$. No contexto de operador $\phi$ e $\psi$ são chamados de argumentos. \\

\textbf{Comentário :} $f$ é também uma variável, pois é o nome simbólico do operador. Podemos definir $y := f(\phi, \psi)$ como uma variável, entretanto $y$ é diferente de $f$, pois $y$ é o resultado da operação $f$ com as variáveis $\phi$ e $\psi$.\\

\textbf{Liberdade [ de Definição de Operador ]} Sempre podemos definir operadores a partir de expressões.



%%%%%%%%%%%%%%%%%%%%%%%%%%%%%%%%%%%%%%%%%%%%%%%%%%%%%%%%%%%%%%%%%%%%%%%%%%%%%%%%%%%%%%%%%%%%%%%%%%%%%%%%%%%%
% Definição de Operador Lista
\hspace{\baselineskip}

\textbf{Operador [ Lista ]} O operador lista é um operador que 'lista' um 'conjunto' finito de variáveis. Sua notação é feita com parênteses e vírgulas. Por exemplo $(x)$ é uma lista com uma variável, $(x,y) $ é uma lista com duas variáveis.\\

\textbf{Comentário :} Assumimos que podemos reconhecer lista de variáveis, isso requer que tenhamos uma inferência para reconhecer uma lista de variáveis.\\

\textbf{Comentário :} Uma lista de variáveis pode ser chamada como uma variável e um operador digamos $f(x,y,z)$ com $X:=(x,y,z) $ pode ser expresso como $fX$.\\

\textbf{Notação :} $ X^0 \overset{part}{\rightarrow} \texttt{LV}^* $ representa dizer que $X$ é uma lista de variáveis.\\

\textbf{Inferência [ Reconhecimento de Lista de Variáveis ]} Sempre podemos reconhecer que uma lista de variáveis tem as variáveis presentes em uma expressão, em particular, podemos reconhecer que uma lista contem todas as variáveis de uma expressão e também que ela tem apenas variáveis daquela expressão.\\

\textbf{Inferência [ Retirada de Parênteses em Lista ]} Sempre que um operador tiver como algum de seus argumentos uma lista de variáveis ele equivale a recompor a lista de variáveis sem considerar parênteses internos. Por exemplo $ f((x,y),z) = f(x,y,z) $, ou $ ((x,y),z) = (x,y,z) $.\\

\textbf{Comentário :} A retirada de parênteses somente pode ser feita em listas.

%%%%%%%%%%%%%%%%%%%%%%%%%%%%%%%%%%%%%%%%%%%%%%%%%%%%%%%%%%%%%%%%%%%%%%%%%%%%%%%%%%%%%%%%%%%%%%%%%%%%%%%%%%%%
\hspace{\baselineskip}

\textbf{[ Relação ]} É uma expressão ou símbolo que conecta dois ou mais substantivos. Sendo ela mesma um substantivo especial.  \\

\textbf{Comentário :} Uma relação é um operador por definição, entretanto faz sentido no contexto de lógica definir operadores e relações como classes separadas, relação é uma informação, portanto pode ter valores de verdadeiro ou falso, uma expressão resultante de um operador não necessariamente é uma informação.\\

\textbf{Notação [ Funcional ]} $f(\phi,\psi)$ é um exemplo de relação representado em notação funcional, $f$ é o nome da relação, $f$ está conectando as existências $\phi$ e $\psi$.\\

\textbf{Comentário :} É esperado que a entidade que faz uso da teoria seja capaz de reconhecer relações através de uma notação funcional.\\

\textbf{[ Fórmula ]} Usamos variáveis e operadores para construir fórmulas.\\

\textbf{Comentário :} $f(x,y)$ é uma fórmula, $x $ é uma fórmula, $g(f(x),h(x,y) ) $ é uma fórmula.\\

\textbf{[ Expressão vs Fórmulas ]} Definimos \textbf{[ Sintaxe ]} as regras que permitem identificar fórmulas válidas de expressões quaisquer que podem não ter significado algum. Fórmulas obedecem regras de sintaxe enquanto expressões podem ou não obedecer.\\

%\textbf{Comentário :} A lógica de primeira ordem ou lógica por quantificadores usada na matemática é construída de maneira similar. Entretanto irei construir (reconstruir) uma lógica intermediária entre a linguagem comum e a linguagem matemática.\\

%%%%%%%%%%%%%%%%%%%%%%%%%%%%%%%%%%%%%%%%%%%%%%%%%%%%%%%%%%%%%%%%%%%%%%%%%%%%%%%%%%%%%%%%%%%%%%%%%%%%%%%%%%%

\textbf{Relação [ Igualdade ]} $a=b$ ou $(=)(a,b)$ significa que os símbolos $a$ e $b$ fazem referência a um mesmo objeto, são representantes de uma mesma coisa na realidade ou em um raciocínio.\\

\textbf{Comentário :} A relação de igualdade pode ser representado pela notação funcional $ a = b $ pode ser escrito como $(=)(a,b)$ por exemplo. Definir igualdade dessa forma permite construir expressões de uma maneira consistente.\\

\textbf{Comentário :} A relação de igualdade depende do contexto onde se insere as variáveis, por exemplo, não podemos dizer que bananas são iguais a maçãs, portanto $x$ bananas não podem ser iguais a $y$ maçãs. Um exemplo presente no texto é a igualdade de variáveis e a igualdade de operadores, operadores são símbolos que estão em contexto diferente das variáveis que o operador opera.\\

\textbf{Inferência [ de Igualdade ]} Representa a propriedade de comutação dos símbolos ao redor do símbolo de igualdade.
$$a=b$$
\hrule
$$b=a$$

\textbf{Comentário :} Somente aplicável quando $a$ e $b$ estão em contextos permitidos pela igualdade $=$.\\

\textbf{Inferência [ de Substituição ]} Representa o conceito de que se uma coisa é igual a outra podemos substituir livremente aonde uma for achada pela outra.
$$ a = b $$
$$ x^0 \overset{part}{\rightarrow} \texttt{LV}^* $$
$$  \phi(a,x) $$
\hrule
$$ \phi(b,x) $$

\textbf{Comentário :} Deve estar claro a notação que representa fórmulas abertas, relações e variáveis. A notação funcional ( símbolo seguido de parênteses ) representa fórmulas abertas, enquanto símbolos avulsos podem representar relações ou existências. Também deve ser claro que temos a liberdade de criar nomes (representar existências por símbolos) tanto para fórmulas (notação funcional) quanto para existências.\\

\textbf{Comentário :} A substituição precisa ser completa, todas as ocorrências de $a$ de uma fórmula $\phi$ devem ser substituídas pelo símbolo $b$ na inferência de substituição.\\



\hspace{\baselineskip}

\textbf{Inferência [ de Definição ]}
$$ a:=b $$
\hrule
$$ a=b $$

\textbf{Comentário :} A inferência de substituição e de definição somente pode ser aplicado se respeitar o contexto em que a igualdade se insere. Por exemplo $ f(x) = x*x $ é uma fórmula definida com o operador $*$ e $f(y)=y*y$ é a mesma fórmula apenas utilizando uma outra representação para variável. $ (f(x) = x*x) \overset{op}{=} (f(y) = y*y)$ com $\overset{op}{=}$ igualdade de operador não significa que temos $f(x) \overset{var}{=} f(y)$ com $\overset{var}{=}$ igualdade de variável, portanto não podemos substituir $f(x)$ por $f(y)$.\\

\textbf{Notação :} $ \overset{op}{=} $ é a igualdade no contexto de operadores e $\overset{var}{=}$ é a igualdade no contexto de variáveis. Sempre que um conceito é inserido e é possível nomeá-los, convém criar uma notação de igualdade apropriada para separar os contextos de substituição. A omissão da identificação somente pode ser feita caso não haja dúvidas do seu significado.

%\textbf{Informação [ Relação ]} Relação representada por $\sim$ conecta dois substantivos, por exemplo $A \sim B$, e representa a informação mais genérica possível.



%%%%%%%%%%%%%%%%%%%%%%%%%%%%%%%%%%%%%%%%%%%%%%%%%%%%%%%%%%%%%%%%%%%%%%%%%%%%%%%%%%%%%%%%%%%%%%%%%%%%%%%%%%%%
%%%%%%%%%%%%%%%%%%%%%%%%%%%%%%%%%%%%%%%%%%%%%%%%%%%%%%%%%%%%%%%%%%%%%%%%%%%%%%%%%%%%%%%%%%%%%%%%%%%%%%%%%%%%
%%%%%%%%%%%%%%%%%%%%%%%%%%%%%%%%%%%%%%%%%%%%%%%%%%%%%%%%%%%%%%%%%%%%%%%%%%%%%%%%%%%%%%%%%%%%%%%%%%%%%%%%%%%%
%%%%%%%%%%%%%%%%%%%%%%%%%%%%%%%%%%%%%%%%%%%%%%%%%%%%%%%%%%%%%%%%%%%%%%%%%%%%%%%%%%%%%%%%%%%%%%%%%%%%%%%%%%%%
%%%%%%%%%%%%%%%%%%%%%%%%%%%%%%%%%%%%%%%%%%%%%%%%%%%%%%%%%%%%%%%%%%%%%%%%%%%%%%%%%%%%%%%%%%%%%%%%%%%%%%%%%%%%
%%%%%%%%%%%%%%%%%%%%%%%%%%%%%%%%%%%%%%%%%%%%%%%%%%%%%%%%%%%%%%%%%%%%%%%%%%%%%%%%%%%%%%%%%%%%%%%%%%%%%%%%%%%%

\section{Teoria Inferente [ Teoria de Fórmulas Universais ]}

\hspace{\baselineskip}

A teoria de fórmulas universais tem como objetivo permitir a escrita simbólica de informações universais para teorias dedutivas, uma teoria de fórmula universal herda o conteúdo da teoria de substituição, assim como as demais teorias expostas posteriormente.\\

\textbf{Fórmula [ Aberta ]} É uma fórmula em que não é feito discriminação entre os componentes particulares e os componentes universais. Os componentes que não estão discriminados são chamados de abertos.\\

\textbf{Notação [ Fórmula Aberta ]} Quando representamos uma fórmula aberta usamos a notação de função para mostrar quais símbolos estão abertos, por exemplo $ \phi(a) $ significa que o símbolo $a$ está aberto, se temos a fórmula representada como $\phi(a,b)$ então $a$ e $b$ estão abertos. \\

\textbf{Fórmula [ Fechada ]} É uma fórmula em que cada componente tem discriminado a sua universalidade e a sua particularidade.\\

\textbf{Comentário :} Fórmulas fechadas são de fato informações, enquanto fórmulas abertas, substantivos avulsos ou relações avulsas são expressões incompletas.\\

\textbf{[ Quantificador ]} Quantificador é uma expressão especial usada para fechar componentes de fórmulas.\\

\textbf{Comentário :} Quantificador é um operador ternário, por exemplo: $\forall x \in X: \phi(x)$ sendo $\forall x \in X$ o quantificador da fórmula $\phi(x)$ que diz que $\phi(x)$ vale para qualquer $x$ pertencente a $X$, neste exemplo o quantificador está conectando $x$, $X$ e $\phi$.\\

\textbf{Quantificador [ Universal ]} $\forall$ é o símbolo de quantificação universal.\\

\textbf{Quantificador [ Dialética de Representação e Realidade ]}

\begin{itemize}
	\item[ (i) ] \textbf{[ Quando se refere a Algo da Realidade ]} Quando usamos $A^{*}$ significa que fazemos referência algum conceito ou existência da realidade ou que está bem definido, ou que assumimos como bem definido.
	\item[ (ii) ] \textbf{[ Quando se refere a Representação ]} $A^{0}$ significa que fazemos referência a $A$ como símbolo, como representação.
\end{itemize}

\textbf{Notação [ Quantificador ]} Para denotar um quantificador usamos parênteses duplos ((expressão quantificadora)). Essa notação é uma interface de raciocínio que denota o início de uma estrutura recursiva. \\

\textbf{Notação [ Definição ]} Para definir um conceito, algo da realidade, ou simplesmente assumir que algo está bem definido usamos $A^{*}$\{ Expressão em linguagem usual que define a entidade em questão \} ou (( $A^{*}$\{Expressão ...\} )) com a notação de quantificador.

%%%%%%%%%%%%%%%%%%%%%%%%%%%%%%%%%%%%%%%%%%%%%%%%%%%%%%%%%%%%%%%%%%%%%%%%%%%%%%%%%%%%%%%%%%%%%%%%%%%%%%%%%%%%
% 
\hspace{\baselineskip}

\textbf{Relações [ Quantificadoras ]} Relações usadas dentro de expressões quantificadoras para definir quantificadores especiais.

\begin{itemize}
	\item[ (i) ] \textbf{[ Elemento Genérico ]} (( $a^0 \overset{gen.}{\rightarrow} A^{*}$ )) é o quantificador que significa que o símbolo $a$ representa alguém qualquer do conceito $A$. % Pode ser usado como indutor externo
	\item[ (ii) ] \textbf{[ Elemento Particular ]} (( $a^0 \overset{part.}{\rightarrow} A^{*}$ )) é o quantificador que significa que o símbolo $a$ é um representante particular do conceito $A$. % Talvez possa ser usado como indutor externo
	%	\item[ (iii) ] \textbf{[ Elemento Particular ]} % Diferente do Genérico não induz inferência dedutiva
\end{itemize}

\textbf{Comentário :} Poderíamos usar $\forall a \in A$ com o mesmo significado que $a^0 \overset{gen.}{\rightarrow} A^{*}$, entretanto $\in$ é um símbolo da teoria de conjuntos, que é algo com propriedades definidas pela teoria axiomática de conjunto, para nos atermos a generalidade usamos uma notação que está entre o formalismo e a dubiedade da lógica informal. Como a inferência usada é similar podemos chamar o quantificador de elemento genérico como quantificador universal, entretanto a rigor os conceitos são diferentes.\\

\textbf{Comentário :} Quantificadores também são fórmulas, portanto a notação de quantificador é necessária para discriminar se ele age como quantificador ou apenas como uma fórmula.

%%%%%%%%%%%%%%%%%%%%%%%%%%%%%%%%%%%%%%%%%%%%%%%%%%%%%%%%%%%%%%%%%%%%%%%%%%%%%%%%%%%%%%%%%%%%%%%%%%%%%%%%%%%
% Inferência
\hspace{\baselineskip}

%\textbf{[ Inferência ]} Uma inferência é uma sequência de fórmulas. Uma inferência tem duas partes, as hipóteses e a conclusão. Uma inferência específica permite sua reconstrução completa (hipóteses mais conclusão) a partir das hipóteses. Em computação isso significa que existe um algoritmo bem definido para obter a conclusão a partir das hipóteses. \\

%\textbf{Comentário :} Por hora uma inferência não está dizendo nada sobre raciocínios verdadeiros ou falsos, uma inferência é usada para tal finalidade, entretanto sua definição é anterior e não tem nenhuma responsabilidade em preservar raciocínios verdadeiros, o objetivo é definir uma sequência de fórmulas. Inferência deve ser vista inicialmente como uma tentativa de programar um computador para que efetue raciocínios. \\

\textbf{Comentário :} Uma inferência quando aceita em uma teoria é considerada uma informação e podemos representar como fórmulas. Por exemplo $\phi \implies \psi$ representa uma inferência em que $\phi$ são as hipóteses e $\psi$ são as conclusões. Para representar várias hipóteses podemos utilizar o símbolo $\&$, por exemplo $A \& B \implies C$ significa que as duas hipóteses $A$ e $B$ levam a conclusão $C$.

%%%%%%%%%%%%%%%%%%%%%%%%%%%%%%%%%%%%%%%%%%%%%%%%%%%%%%%%%%%%%%%%%%%%%%%%%%%%%%%%%%%%%%%%%%%%%%%%%%%%%%%%%%%
% Inferência
\hspace{\baselineskip}

\textbf{Inferência [ Desacoplamento ]} 

$$ A \& B $$
\hrule
$$ A $$

\textbf{Inferência [ Desacoplamento ]} 

$$ A \& B $$
\hrule
$$ B $$

%%%%%%%%%%%%%%%%%%%%%%%%%%%%%%%%%%%%%%%%%%%%%%%%%%%%%%%%%%%%%%%%%%%%%%%%%%%%%%%%%%%%%%%%%%%%%%%%%%%%%%%%%%%
% Inferência de Extensão
%\hspace{\baselineskip}

%\textbf{Inferência [ de Extensão ]} Seja uma sequência de fórmulas



%  Inferências são mecânicas, logo não precisam ser verificadas 

%%%%%%%%%%%%%%%%%%%%%%%%%%%%%%%%%%%%%%%%%%%%%%%%%%%%%%%%%%%%%%%%%%%%%%%%%%%%%%%%%%%%%%%%%%%%%%%%%%%%%%%%%%%
% Inferência Universal
\hspace{\baselineskip}

\textbf{Inferência [ de Dedução ]} A inferência universal ou dedutiva é uma sequência de duas fórmulas fechadas, uma universal e outra particular como mostrado abaixo:

$$ ((a^0 \overset{gen.}{\rightarrow} A^*)) : \phi(a,x) $$
$$ b^0 \overset{part.}{\rightarrow} A^* $$
\hrule
$$  \phi(b,x) $$

\textbf{Comentário :} A inferência dedutiva representa a ideia de que se uma informação universal é verdadeira então todas as particulares são verdadeiras da mesma forma. A inferência dedutiva possui os conceitos de substituição e de igualdade das outras duas inferências, o conceito de substituição é evidente na transformação de $\phi(a)$ por $\phi(b)$ , a igualdade reside na propriedade comum $\phi$ que todos os $A$ preservam.

%%%%%%%%%%%%%%%%%%%%%%%%%%%%%%%%%%%%%%%%%%%%%%%%%%%%%%%%%%%%%%%%%%%%%%%%%%%%%%%%%%%%%%%%%%%%%%%%%%%%%%%%%%%%
% Inferência 
\hspace{\baselineskip}

\textbf{Inferência [ de Dedução (2) ]} 

$$ x^0 \overset{gen.}{\rightarrow} A^* \implies \phi(x,a) $$
\hrule
$$ ((x^0 \overset{gen.}{\rightarrow} A^*)) \phi(x,a) $$

%%%%%%%%%%%%%%%%%%%%%%%%%%%%%%%%%%%%%%%%%%%%%%%%%%%%%%%%%%%%%%%%%%%%%%%%%%%%%%%%%%%%%%%%%%%%%%%%%%%%%%%%%%%%
% Quantificador Universal
%\hspace{\baselineskip}

%\textbf{Quantificador [ Universal ]} (( $S^{*}$ ))

%%%%%%%%%%%%%%%%%%%%%%%%%%%%%%%%%%%%%%%%%%%%%%%%%%%%%%%%%%%%%%%%%%%%%%%%%%%%%%%%%%%%%%%%%%%%%%%%%%%%%%%%%%%
% Razoes
\hspace{\baselineskip}

\textbf{[ Razão ]} É uma sequência de fórmulas. Uma inferência pode agir sobre uma razão para estender a sequência. As hipóteses da inferência devem estar contidas anteriormente na razão para que a conclusão possa ser adicionada na sequência.

%%%%%%%%%%%%%%%%%%%%%%%%%%%%%%%%%%%%%%%%%%%%%%%%%%%%%%%%%%%%%%%%%%%%%%%%%%%%%%%%%%%%%%%%%%%%%%%%%%%%%%%%%%%%
% Inferência Fundamental
\hspace{\baselineskip}

\textbf{Inferência [ Fundamental ]} 
$$ \phi \implies \psi $$
$$ \phi $$
\hrule
$$ \psi $$
\textbf{Comentário :} É chamado de fundamental, pois ele representa a própria ideia de inferência, que se $\phi$ implica em $\psi$ então se temos $\phi$ podemos concluir $\psi$.

%%%%%%%%%%%%%%%%%%%%%%%%%%%%%%%%%%%%%%%%%%%%%%%%%%%%%%%%%%%%%%%%%%%%%%%%%%%%%%%%%%%%%%%%%%%%%%%%%%%%%%%%%%%%
% Transitividade de Inferência
\hspace{\baselineskip}

\textbf{Inferência [ Transitividade ]}

$$ A \implies B $$
$$ B \implies C $$
\hrule
$$ A \implies C $$

\textbf{Notação [ Sequência de Inferências Transitivas ]} Representamos uma sequência de inferências transitivas por $ A \{\implies\} B$, significando que existem uma sequência de inferências conectando $A$ para $B$.\\

\textbf{Inferência [ Reconhecimento de Definição ]} Uma definição deve ser dotada de um método para discriminar se algo pertence ou não aquela definição. Se este método $\phi$ é satisfeito então podemos concluir que ele pertence a uma definição.

$$ \phi(x) $$
\hrule
$$ x^0 \overset{part.}{\rightarrow} X^* $$

\textbf{Comentário :} Em geral este método $\phi$ foge do escopo dos símbolos da teoria exposta, requer que a entidade consciente tenha suficiente capacidade de compreender o que o expositor quer dizer. Portanto convém utilizar linguagem verbal quando for demonstrar que algo pertence a uma definição. Por exemplo, reconhecimento de inferências $\{\implies\}$ ou $\implies$ ou $A \& B$ são reconhecimentos de definição que precisam sair do escopo dos símbolos da teoria de inferência. Para constatar isso basta tentar responder a pergunta, por que eu não posso descrever a inferência [ de reconhecimento de inferência ] através das notações de inferência? Basta observar que o reconhecimento requer a entidade consciente enxergando a inferência como ela é e isso não pode ser descrito por inferências, apenas a linguagem tem a flexibilidade necessária para comunicar ao leitor do processo. \\

%\textbf{Inferência [ Reconhecimento de Operadores ]}

%Suponha um operador $f$, sempre podemos reconhecer em expressões a definição do operador

\textbf{Inferência [ Reconhecimento de Inferência Fundamental ]}
\begin{itemize}
	\item[ ] Se de $A$ conseguimos uma sequência de hipóteses com elementos anteriores da razão ou encadeadas por inferências da teoria tal que tenha em sua última linha tenhamos $B$
\end{itemize}
\hrule

$$ A \{ \implies \} B $$

\textbf{Comentário :} Podemos usar $\implies$ em vez de $\{ \implies \}$ quando não houver dúvidas, perde-se pouco de seu significado entretanto $\{ \implies \}$ tem um sentido mais refinado.\\

\textbf{Inferência [ Reconhecimento de Conjunção de Hipóteses ]}
\begin{itemize}
	\item[ ] Suponha que a razão esteja projetada como linhas representando informações horizontalmente e verticalmente a sequência de raciocínios. Se na razão temos em qualquer ordem $A$ em uma linha e $B$ em outra linha temos.
\end{itemize}
\hrule

$$ A \& B $$

%%%%%%%%%%%%%%%%%%%%%%%%%%%%%%%%%%%%%%%%%%%%%%%%%%%%%%%%%%%%%%%%%%%%%%%%%%%%%%%%%%%%%%%%%%%%%%%%%%%%%%%%%%%%
%%%%%%%%%%%%%%%%%%%%%%%%%%%%%%%%%%%%%%%%%%%%%%%%%%%%%%%%%%%%%%%%%%%%%%%%%%%%%%%%%%%%%%%%%%%%%%%%%%%%%%%%%%%%
%%%%%%%%%%%%%%%%%%%%%%%%%%%%%%%%%%%%%%%%%%%%%%%%%%%%%%%%%%%%%%%%%%%%%%%%%%%%%%%%%%%%%%%%%%%%%%%%%%%%%%%%%%%%
%%%%%%%%%%%%%%%%%%%%%%%%%%%%%%%%%%%%%%%%%%%%%%%%%%%%%%%%%%%%%%%%%%%%%%%%%%%%%%%%%%%%%%%%%%%%%%%%%%%%%%%%%%%%
%%%%%%%%%%%%%%%%%%%%%%%%%%%%%%%%%%%%%%%%%%%%%%%%%%%%%%%%%%%%%%%%%%%%%%%%%%%%%%%%%%%%%%%%%%%%%%%%%%%%%%%%%%%%
%%%%%%%%%%%%%%%%%%%%%%%%%%%%%%%%%%%%%%%%%%%%%%%%%%%%%%%%%%%%%%%%%%%%%%%%%%%%%%%%%%%%%%%%%%%%%%%%%%%%%%%%%%%%

\section{Teoria Axiomática [ Lógica ]}

\hspace{\baselineskip}

Nesta seção vamos efetivamente estar construindo uma teoria axiomática para expor a lógica, entretanto para expor uma teoria axiomática é necessário a lógica, o que faz com que uma exposição desse tipo seja extremamente desconfortável para uma mente cuidadosa. Para tentar amenizar este desconforto é necessário enxergar o processo de inferência como algo diferente da lógica ou do raciocínio, tente pensar que a cada proposição e cada demonstração estamos fazendo um processo mecânico em que aumentamos uma lista (razão) de informações, mas que eventualmente tais informações possam ser interpretadas como elas realmente são e não meramente um processo mecânico.\\

Este é a última parte da meta-teoria principal, com ela busco exemplificar a construção de teorias axiomáticas sob o método trialético. A meta-teoria principal tem como objetivo ser a base para construção e exposição de teorias, portanto não busco adentrar muito nos detalhes da lógica matemática, ela deve expor os conceitos de forma ágil e suficiente para que novas teorias possam ser construídas e analisadas.

\subsection{Lógica Proposicional}

%\newpage

%%%%%%%%%%%%%%%%%%%%%%%%%%%%%%%%%%%%%%%%%%%%%%%%%%%%%%%%%%%%%%%%%%%%%%%%%%%%%%%%%%%%%%%%%%%%%%%%%%%%%%%%%%%%
% Dialética Socrática
\hspace{\baselineskip}

\textbf{Dialética Socrática [ Lógica ]}

\begin{itemize}
	\item[ ] ?` O que é lógica ? : Lógica é um campo de estudo.
	\begin{itemize}
		\item[ ] ?` O que caracteriza um campo de estudo ? : Um campo de estudo tem um problema, um objeto a ser estudado.
		\item[ ] ?` Qual o objeto de estudo da lógica ?
		\begin{itemize}
			\item[ (i) ] Modelar o conceito de informação. Neste sentido modelar é tentar aproximar através da linguagem algo intangível que neste caso é a informação.
			\item[ (ii) ] Modelar o conceito de raciocínio. Intrínseco em ambos os conceitos é o conceito de verdade, é um atributo desejável da informação a verdade, é um atributo desejável do raciocínio a verdade. A informação é uma interpretação estática, o raciocínio é uma interpretação dinâmica.
		\end{itemize}
	\end{itemize}
	
\end{itemize}

\textbf{Comentário :} O raciocínio lógico projetado tradicionalmente no substrato papel possui duas direções de leitura, a primeira horizontal, da esquerda para direita é onde a informação é externalizada, a segunda vertical, de cima para baixo é onde o raciocínio é projetado. É possível externalizar a lógica (informação,raciocínio) utilizando outras notações diferentes da utilizada pela lógica matemática, por exemplo em um texto uma informação pode ser uma frase, terminada em ponto e um raciocínio pode estar organizado em parágrafos.\\

%%%%%%%%%%%%%%%%%%%%%%%%%%%%%%%%%%%%%%%%%%%%%%%%%%%%%%%%%%%%%%%%%%%%%%%%%%%%%%%%%%%%%%%%%%%%%%%%%%%%%%%%%%%%
% Teoria Axiomática
% \hspace{\baselineskip}

\textbf{Axioma [ de Axioma ]} Qualquer axioma declarado neste texto como axioma válido da teoria lógica é considerado como adicionado à razão da lógica.\\

\textbf{Axioma [ de Inferência ]} Qualquer inferência declarada neste texto como inferência válida da teoria lógica é considerado como adicionado à razão da lógica.\\

\textbf{Axioma [ de Proposição ]} Qualquer proposição devidamente demonstrada neste texto por uma sequência finita de inferências e de elementos anteriores da razão é considerado como adicionado à razão da lógica.\\

\textbf{Comentário [ Pré Axiomas Trialéticos ]} Podemos interpretar uma teoria axiomática com a trialética (Axioma, Proposição, Inferência). Axioma são as informações iniciais da razão da teoria, Proposição são informações que desejamos adicionar à razão da teoria, e Inferência é a interface entre o leitor consciente e o raciocínio utilizado para adicionar a proposição à razão. Uma demonstração de uma proposição é construída através de inferências. \\

\textbf{Axioma [ de Herança ]} Façamos que a teoria axiomática de lógica herde todas as inferências de teorias axiomáticas, incluindo as inferências expostas na seção de teoria dedutiva como axiomas em sua razão e as teorias inferentes de substituição e de fórmulas universais.\\




\textbf{[ Informação ]} É qualquer coisa que possui atributo de verdade. Uma informação pode ser verdadeira ou pode ser falsa. Uma informação pode ser representada por símbolos, por nomes. Representamos uma informação qualquer pelo quantificador $ (( x^0 \overset{gen.}{\rightarrow} \mathcal{I}^* )) $ de forma que $x$ é o símbolo dessa informação qualquer. Quando uma informação $x$ é verdadeira representamos como $x=1$, quando uma informação é falsa representamos como $x=0$. \\

\textbf{Comentário :} Uma inferência representa um raciocínio olhando verticalmente (de cima para baixo), uma inferência requer em cada linha (componente horizontal da esquerda para direita) de informações, portanto é necessário para escrever uma inferência que seus componentes horizontais pertençam às informações $\mathcal{I}$.\\

\textbf{Notação :} $\overset{inf}{=}$ é o símbolo de igualdade utilizado no contexto de informações que podem assumir valores de verdadeiro e falso.\\



\textbf{Axioma [ Lei do Meio Excluído ]} Uma informação se for verdadeira somente pode ser verdadeira, uma informação se for falsa somente pode ser falsa. Uma informação não pode ser ao mesmo tempo verdadeira ou falsa.\\

\textbf{Comentário :} Representamos $x \in \{0,1\}$ para dizer que $x$ pode assumir os valores de verdade ou falsidade respeitando o meio excluído. É importante notar que existe uma inferência para reconhecer $x \in \{0,1 \}$, é interessante para tal inferência utilizar a linguagem verbal.\\

\textbf{Notação :} $\texttt{LME} := \{ \texttt{Axioma Lei do Meio Excluido} \} $\\

\textbf{Axioma [ Falsidade ]} Não podemos sob qualquer tipo de inferência adicionar na razão da lógica informações que são falsas. Qualquer inferência, axioma ou proposição que adicione informações falsas não podem ser consideradas como inferências válidas para a teoria lógica.\\

\textbf{Notação :} $\texttt{AF} := \{ \texttt{Axioma de Falsidade} \} $




%%%%%%%%%%%%%%%%%%%%%%%%%%%%%%%%%%%%%%%%%%%%%%%%%%%%%%%%%%%%%%%%%%%%%%%%%%%%%%%%%%%%%%%%%%%%%%%%%%%%%%%%%%%%
% Inferência de Identificação de Informação
\hspace{\baselineskip}

\textbf{Inferência [ Identificação de Informação ]} 

$$ x \in \{0,1\} $$
\hrule
$$ x \overset{part.}{\rightarrow} \mathcal{I}^* $$

\textbf{Inferência [ de Demonstração de Informação Universal ]}

$$ x^0 \overset{gen.}{\rightarrow} \mathcal{X}^* \{\implies\} \phi(x)  $$
\hrule
$$ ((x^0 \overset{gen.} {\rightarrow} \mathcal{X}^* )) \phi(x) $$

\textbf{Inferência [ de Demonstração de Informação Universal ]}

$$ x^0 \overset{gen.}{\rightarrow} \mathcal{X}^* \& y^0 \overset{gen.}{\rightarrow} \mathcal{Y}^* \{\implies\} \phi(x,y)  $$
\hrule
$$ ((x^0 \overset{gen.} {\rightarrow} \mathcal{X}^* )) ((y^0 \overset{gen.} {\rightarrow} \mathcal{Y}^* )) \phi(x,y) $$

%%%%%%%%%%%%%%%%%%%%%%%%%%%%%%%%%%%%%%%%%%%%%%%%%%%%%%%%%%%%%%%%%%%%%%%%%%%%%%%%%%%%%%%%%%%%%%%%%%%%%%%%%%%%
% Operador de Conjunção
\hspace{\baselineskip}

\textbf{Operador [ Conjunção ]} $ (( x^0 \overset{gen.}{\rightarrow} \mathcal{I}^* )) (( y^0 \overset{gen.}{\rightarrow} \mathcal{I}^* )) $

\begin{itemize}
	\item[ (i) ] $x \wedge y = 0$ se $x=0$ e $y=0$	
	\item[ (ii) ] $x \wedge y = 0$ se $x=0$ e $y=1$
	\item[ (iii) ] $x \wedge y = 0$ se $x=1$ e $y=0$
	\item[ (iv) ] $ x \wedge y = 1$ se $x=1$ e $y=1$
\end{itemize}

% \textbf{Notação :} $ ((x \in \mathcal{I})) := (( x^0 \overset{gen.}{\rightarrow} \mathcal{I}^* )) $

%%%%%%%%%%%%%%%%%%%%%%%%%%%%%%%%%%%%%%%%%%%%%%%%%%%%%%%%%%%%%%%%%%%%%%%%%%%%%%%%%%%%%%%%%%%%%%%%%%%%%%%%%%%%
% Proposição Exemplo Demonstrar que Operador Conjunção é uma informação
\hspace{\baselineskip}

\textbf{Proposição [ Operador Conjunção ]} O operador conjunção é uma informação.

$$ (( x^0 \overset{gen.}{\rightarrow} \mathcal{I}^* )) (( y^0 \overset{gen.}{\rightarrow} \mathcal{I}^* )) x \wedge y \overset{part.}{\rightarrow} \mathcal{I}^* $$

\textbf{Dialética :}

\begin{itemize}
	\item[ ] !` Inferência de Demonstração de Informação Universal !
	\begin{itemize}
		\item[ ] ?` $ x^0 \overset{gen.}{\rightarrow} \mathcal{I}^* \& y^0 \overset{gen.}{\rightarrow} \mathcal{I}^* \{ \implies \} x\wedge y \in \{0,1\} $ ?
		\begin{itemize}
			\item[ ] !` Inferência de Reconhecimento de $ \dots \in \{ 0,1 \}$ !
			\item[ ] !` Inferência de Reconhecimento de $ \dots \{\implies \} \dots $ !
		\end{itemize}
		\item[ ] !` Inferência de Identificação de Informação !
		\item[ ] !` Inferência de Transitividade !
		\item[ ] !` Inferência de Demonstração de Informação Universal !
	\end{itemize}
\end{itemize}

\textbf{Comentário :} (?`, ?, !`, !) são as notações de interface de raciocínio recursivo, (?) representa uma pergunta ou subpergunta, (!) representa o uso de uma inferência para completar um objetivo.\\

\textbf{Demonstração :} Se $x$ elemento genérico de informação e $y$ elemento genérico de informação então $x$ e $y$ assumem valores de verdadeiro e falso. Se $x$ e $y$ podem assumir valores de verdadeiro e falso então temos que se $x=0$ e $y=0$ então $x\wedge y = 0$ se $x=1$ e $y=1$ então $x \wedge y =1$ se $x$ e $y$ possuem valores distintos então $x \wedge y=0$. Usando transitividade concluímos que se $x$ elemento genérico e $y$ elemento genérico de informação então $x\wedge y$ assume valores de verdadeiro e falso. Identificamos que $x\wedge y$ pode assumir valores de verdadeiro e falso e portanto $x \wedge y \in \{0,1\} $. Usando a transitividade concluímos que $x^0 \overset{gen.}{\rightarrow} \mathcal{I}^* \& y^0 \overset{gen.}{\rightarrow} \mathcal{I}^* \{ \implies \} x\wedge y \in \{0,1\}$. Como $x\wedge y \in \{0,1\} \implies x \wedge y \overset{part}{\rightarrow} \mathcal{I}^*$ e $x^0 \overset{gen.}{\rightarrow} \mathcal{I}^* \& y^0 \overset{gen.}{\rightarrow} \mathcal{I}^* \{ \implies \} x\wedge y \in \{0,1\}$ podemos concluir usando a transitividade que $x^0 \overset{gen.}{\rightarrow} \mathcal{I}^* \& y^0 \overset{gen.}{\rightarrow} \mathcal{I}^* \{ \implies \} x\wedge y \overset{part}{\rightarrow} \mathcal{I}^*$. Adicionamos a razão a proposição utilizando a inferência de demonstração de informação universal.\\

\textbf{Comentário :} É recomendado que o leitor tente demonstrar usando algum substrato conveniente de raciocínio apenas fazendo uso da dialética antes de tentar entender a demonstração de uma proposição. Considere as proposições como exercícios e a demonstração como a resolução.

%%%%%%%%%%%%%%%%%%%%%%%%%%%%%%%%%%%%%%%%%%%%%%%%%%%%%%%%%%%%%%%%%%%%%%%%%%%%%%%%%%%%%%%%%%%%%%%%%%%%%%%%%%%%
% Definição de Símbolo de Recíproco
\hspace{\baselineskip}

\textbf{[ Recíproco ]} Definimos o símbolo de recíproco pelas inferências:

\textbf{Inferência :}

$$ A \implies B $$
$$ B \implies A $$
\hrule
$$ A \iff B $$

\textbf{Inferência :}

$$ A \iff B $$
\hrule
$$ (A \implies B) \&  (B \implies A )$$

%%%%%%%%%%%%%%%%%%%%%%%%%%%%%%%%%%%%%%%%%%%%%%%%%%%%%%%%%%%%%%%%%%%%%%%%%%%%%%%%%%%%%%%%%%%%%%%%%%%%%%%%%%%%
% Axioma 
\textbf{Axioma [ Notação Funcional de Informação ]} $ (( \phi^0 \overset{gen.}{\rightarrow} \mathcal{I}^* )) : \phi = 1 \iff \phi $

%%%%%%%%%%%%%%%%%%%%%%%%%%%%%%%%%%%%%%%%%%%%%%%%%%%%%%%%%%%%%%%%%%%%%%%%%%%%%%%%%%%%%%%%%%%%%%%%%%%%%%%%%%%%
% Proposição [ Notação Funcional de Conjunção ]
\hspace{\baselineskip}

\textbf{Proposição [ Notação Funcional de Conjunção ]} O símbolo $\&$ implica no símbolo $\wedge$.

$$ ((x^0 \overset{gen.}{\rightarrow} \mathcal{I}^*)) ((y^0 \overset{gen.}{\rightarrow} \mathcal{I}^*)) : x \& y \implies x \wedge y  $$

\textbf{Dialética :} 

\begin{itemize}
	\item[ ] Basta usar as inferências de desacoplamento, os axiomas de notação funcional de informação e a definição de operador conjunção.
\end{itemize}

\textbf{Demonstração :}

$$ (x^0 \overset{gen.}{\rightarrow} \mathcal{I}^*) \& (y^0 \overset{gen.}{\rightarrow} \mathcal{I}^*)  $$

?` $ x \& y \implies x \wedge y  $  ?

$$ x \& y $$
$$ x $$
$$ y $$
$$ x = 1 $$
$$ y = 1 $$
$$ x \wedge y = 1 $$
$$ x \wedge y $$
$$ x\& y \{ \implies \} x \wedge y $$
$$ \blacksquare $$

$$ ((x^0 \overset{gen.}{\rightarrow} \mathcal{I}^*)) ((y^0 \overset{gen.}{\rightarrow} \mathcal{I}^*)) : x \& y \implies x \wedge y  $$

$$ \blacksquare $$

\textbf{Comentário :} A compreensão de uma demonstração requer a consciência alfa da razão exposta que pode ser externalizada através da evidenciação das inferências utilizadas para cada etapa da conclusão. Portanto é um exercício interessante para verificar a compreensão da demonstração de um teorema fazer anotações evidenciando cada inferência utilizada.\\

\textbf{Comentário [ Notação ]} $ \blacksquare $ significa o fechamento de uma razão de demonstração, quando uma demonstração chegou a uma conclusão.\\

\textbf{Comentário :} Parece trivial que o recíproco $ x \wedge y \implies x \& y $ seja verdadeiro, entretanto essa conclusão não pode ser feita facilmente pois $x \wedge y$ não tem mecanismos de desacoplamento como $x \& y$.\\

%%%%%%%%%%%%%%%%%%%%%%%%%%%%%%%%%%%%%%%%%%%%%%%%%%%%%%%%%%%%%%%%%%%%%%%%%%%%%%%%%%%%%%%%%%%%%%%%%%%%%%%%%%%%
% Dialética da Informação Indiferente
\textbf{Dialética Socrática [ Problema da Informação Indiferente ]} 

\begin{itemize}
	\item[ ] ?` Entre $A \implies B=1$ e $A \implies B=0 $ somente uma dessas informações está certa e necessariamente pelo menos uma deve estar certa ? 
	\begin{itemize}
		\item[ (i) ]  Essa conclusão parece natural, mas não se encaixa na realidade se considerarmos o tempo futuro e fenômenos aleatórios. Suponha que se aceite fenômenos aleatórios, e que $B$ seja um desses fenômenos aleatórios que está em um evento futuro ao evento $A$, portanto $A$ não tem como ter informação necessária para concluir se $B=1$ ou $B=0$.
		\item[ (ii) ] Suponha que se aceite que aliens estão presentes em galáxias muito distantes de tal forma que seja impossível eles se comunicarem conosco. Suponha as informações (Aliens existem) e sua negação (Aliens não existem). Como assumimos que aliens não podem se comunicar conosco, tanto faz se uma informação ou outra for verdadeira, não podemos dizer nada a respeito da proposição (aliens existem), pois tal proposição não interfere em nenhuma conclusão de nosso sistema no momento anterior a essa proposição. Essas informações vamos chama-las de informações indiferentes. Muitos definem metafísica como sendo conhecimentos que são indiferentes para realidade, são informações que não interferem na realidade, por exemplo, quantos anjos cabem em uma ponta de agulha?
		\item[ ] ?` Existe relação entre uma determinada existência ser existencial a um sistema e informações indiferentes ?
		\begin{itemize}
			\item[ (i) ] Se uma existência não se comunica através de informações com um sistema, então ela não tem capacidade de interferir nela, qualquer particularidade dessa existência não tem menor capacidade de modificar o sistema.
			\item[ ] ?` Em um universo em que é possível criar ou gerar novas existências, todas as existências apenas geradas no futuro são não existenciais ao sistema presente. Um sistema com capacidade geração de existências pode deterministicamente predizer sobre a formação de uma existência, seria teoricamente (contrário de efetivamente) possível dizer sobre suas propriedades ?
		\end{itemize}
	\end{itemize}
\end{itemize}

%%%%%%%%%%%%%%%%%%%%%%%%%%%%%%%%%%%%%%%%%%%%%%%%%%%%%%%%%%%%%%%%%%%%%%%%%%%%%%%%%%%%%%%%%%%%%%%%%%%%%%%%%%%%
% Negação
\hspace{\baselineskip}

\textbf{Operador [ Negação ]} $((x^0 \overset{gen.}{\rightarrow} \mathcal{I}^* )) :$ $\neg x = 0$ se $x = 1$ e $\neg x = 1$ se $x=0$.\\

\textbf{Comentário :} Negação é uma informação caso $x$ seja uma informação. Não há nenhum ganho em demonstrar isso.\\

\textbf{Comentário :} O conceito de negação é fundamental para a lógica, pois a capacidade de identificar uma existência requer a capacidade de identificar o que não é aquela existência.

%%%%%%%%%%%%%%%%%%%%%%%%%%%%%%%%%%%%%%%%%%%%%%%%%%%%%%%%%%%%%%%%%%%%%%%%%%%%%%%%%%%%%%%%%%%%%%%%%%%%%%%%%%%%
% Proposição [ Negação de Negação ]
\hspace{\baselineskip}

\textbf{Proposição [ Negação de Negação ]} A negação da negação de uma informação equivale a própria informação dentro de uma razão.

$$ (( x^0 \overset{gen.}{\rightarrow} \mathcal{I}^*  )) : ( \neg \neg x \iff x )$$

\textbf{Dialética :}

\begin{itemize}
	\item[ ] $\texttt{INFS}:=$ \{ Inferência de substituição, de Reconhecimento de Inferência Fundamental, Definição de Operador Negação, Axioma de Notação Funcional de Informação., Liberdade de Definição \}
	\item[ ]  ?` $ \neg \neg x \implies x $ ? $\bullet$ $\texttt{INFS}$ 
	\item[ ] ?` $ x \implies \neg \neg x $ ? $\bullet$ $\texttt{INFS}$
\end{itemize}

\textbf{Comentário [ Notação ]} $\bullet$ indica um conjunto de razões e inferências utilizadas especificamente para uma conclusão ou demonstração. A dificuldade de compreender ou de demonstrar através da dialética uma demonstração indica que o leitor não possui consciência alfa das razões indicadas pelo símbolo.

%%%%%%%%%%%%%%%%%%%%%%%%%%%%%%%%%%%%%%%%%%%%%%%%%%%%%%%%%%%%%%%%%%%%%%%%%%%%%%%%%%%%%%%%%%%%%%%%%%%%%%%%%%%%
% Definição de Conjunção
\hspace{\baselineskip}

\textbf{Operador [ Disjunção ]} $x\vee y := ( \neg ( (\neg x) \wedge (\neg y) ) )$\\

%\begin{itemize}
	%\item[ ] Se $x=0$ e $y=0$ então $x \vee y =0$
	%\item[ ] Se $x=0$ e $y=1$ então $x \vee y = 1$
	%\item[ ] Se $x=1$ e $y=0$ então $x \vee y = 1$
	%\item[ ] Se $x=1$ e $y=1$ então $x \vee y = 1$
%\end{itemize}

\textbf{Comentário :} Como exercício demonstre (crie uma sequência de inferências) que $x \vee y$ é uma informação.\\

\textbf{Comentário [ Universo Existencial ]} O conceito de negação está associado ao conceito de identidade, algo é uma existência $X$ se podemos distinguir daquilo que não é uma existência $X$, o conceito de conjunção permite dizer sobre a interação entre as existências, algo é verdadeiro se outro é verdadeiro. É interessante observar que os operadores de negação e conjunção permitem a construção de novos operadores.\\

\textbf{Comentário [ Notação de Negação ]} A partir de agora o operador unário de negação tem precedência sobre qualquer operador, isso simplifica a notação para definir a disjunção $x \vee y = \neg( \neg x \wedge \neg y ) $.

%%%%%%%%%%%%%%%%%%%%%%%%%%%%%%%%%%%%%%%%%%%%%%%%%%%%%%%%%%%%%%%%%%%%%%%%%%%%%%%%%%%%%%%%%%%%%%%%%%%%%%%%%%%%
% Operador de Implicação Booleana
\hspace{\baselineskip}

\textbf{Operador [ Inferência Booleana ]} $ x \rightarrow y := ( (\neg x) \vee y ) $\\

\textbf{Comentário :} A notação funcional permite saber a precedência das operações, entretanto na notação utilizada tradicionalmente na lógica é preciso indicar a precedência das informações manualmente, por isso colocamos os parêntes para indicar que tal operação enclausurada por parênteses deve ser feito primeiro.\\

\textbf{Comentário :} A inferência booleana é um pouco estranha se comparar com o que representa, pois $x\rightarrow y$ possui valores quando $x=0$, e tais valores são verdadeiros para qualquer valor de $y$, isso não parece ser o comportamento da inferência $\implies$, pois durante a construção de uma teoria faz sentido que a pergunta $x \implies y$ seja posterior a constatação de que $x$ pertence a razão ($x$ é verdadeiro). \\

Vamos explorar um pouco o conceito de inferência assumindo que talvez a inferência booleana não represente a inferência $\implies$ usada na interface de construção de teorias. Essa discussão vai tornar um pouco mais interessante a teoria.

%%%%%%%%%%%%%%%%%%%%%%%%%%%%%%%%%%%%%%%%%%%%%%%%%%%%%%%%%%%%%%%%%%%%%%%%%%%%%%%%%%%%%%%%%%%%%%%%%%%%%%%%%%%%
% Inferência 1
\hspace{\baselineskip}

\textbf{Operador [ Inferência A ]} Vamos definir $x \overset{A}{\rightarrow} y$ como $ (1 \overset{A}{\rightarrow} 1 )=1 $, $ (1 \overset{A}{\rightarrow} 0 )=0 $ e $ (x \overset{A}{\rightarrow} y) = ( \neg y \overset{A}{\rightarrow} \neg x ) $\\

\textbf{Comentário :} Se $x$ implica em $y$ então se não temos $y$ como verdade devemos ter que $x$ não é verdade, caso contrário teríamos $y$ e $\neg y$ ao mesmo tempo verdadeiros contradizendo o meio excluído. A propriedade $ (x \overset{A}{\rightarrow} y) = ( \neg y \overset{A}{\rightarrow} \neg x ) $ descreve essa ideia fundamental na inferência e portanto é interessante testar se a definição de inferência satisfaz essa propriedade.\\

\textbf{Notação [ Substituição de Valores ]} Vamos supor um operador $*$ a notação $[ x * y ](a,b)$ significa que $x=a$ e $y=b$ e portanto $[x*y](a,b) = (a*b)$. Suponha também que $ [ f(x,y,z) ](a,b,c) =  f(a,b,c)$ respeitando a ordem em que cada variável aparece na expressão $f(x,y,z)$.\\

\textbf{Proposição [ Inferência Nula ]} 

$$ [ x \overset{A}{\rightarrow} y ](0,0) = 1 $$

\textbf{Dialética :} Basta utilizar as inferências de definição de variáveis, substituição e transitividade.\\

\textbf{Demonstração :}

$$ x \overset{A}{\rightarrow} y = \neg y \overset{A}{\rightarrow} \neg x $$

$$ [ x\overset{A}{\rightarrow} y ](x=0,y=0) = [ \neg y \overset{A}{\rightarrow} \neg x ](y=0,x=0) $$

$$ q:=\neg y $$

$$ r:= \neg x $$

$$ [ \neg y \overset{A}{\rightarrow} \neg x ](y=0,x=0) = [ q \overset{A}{\rightarrow} r ](q=1,r=1) = (1 \overset{A}{\rightarrow} 1) = 1 $$

$$ [ x\overset{A}{\rightarrow} y ](x=0,y=0) = 1 $$

$$\blacksquare $$

\textbf{Comentário :} Isso mostra que se quisermos que o operador de inferência satisfaça $(x \implies y ) \iff (\neg y \implies \neg x)$ o operador deve ser verdadeiro mesmo que $x=0$ e $y=0$, o que inicialmente pode parecer estranho, pois se $x=0$ e $y=0$ a pergunta $x \implies y$ não parece fazer sentido. É importante notar que $\neg x = 1$ e $\neg y = 1$ equivale a $x=0$ e $y=0$ e faz sentido para eles a pergunta $\neg y \implies \neg x$ que equivale a $x \implies y$, o que explica de alguma maneira a necessidade destes valores.

%%%%%%%%%%%%%%%%%%%%%%%%%%%%%%%%%%%%%%%%%%%%%%%%%%%%%%%%%%%%%%%%%%%%%%%%%%%%%%%%%%%%%%%%%%%%%%%%%%%%%%%%%%%%
% Proposição
\hspace{\baselineskip}

\textbf{Proposição :}

$$ [ x \overset{A}{\rightarrow} y ](0,1) = [ \neg y \overset{A}{\rightarrow} \neg x ](1,0) $$

\textbf{Dialética :} Basta usar as propriedades definidas da inferência $A$ e inferências já utilizadas. Como não é uma tarefa difícil não irei demonstrar.\\

\textbf{Comentário :} A propriedade $(x \implies y ) \iff (\neg y \implies \neg x)$ permite determinar o valor de $ 0 \overset{A}{\rightarrow} 0 $, entretanto ficou faltando determinar o valor de $ 0 \overset{A}{\rightarrow} 1 $. Na proposição a seguir temos que tanto faz seu valor para essa propriedade não conseguimos determinar seu valor.\\

\textbf{Proposição [ Raciocínio Inverso de Inferência ]}

$$\Phi( ( 0 \overset{A}{\rightarrow} 1 ) := 1 ) : ((x^0 \overset{gen}{\rightarrow} \mathcal{I}^*)) ((y^0 \overset{gen}{\rightarrow} \mathcal{I}^*)) : \neg y \overset{A}{\rightarrow} \neg x = x \overset{A}{\rightarrow} y $$

$$\Phi( ( 0 \overset{A}{\rightarrow} 1 ) := 0 ) : ((x^0 \overset{gen}{\rightarrow} \mathcal{I}^*)) ((y^0 \overset{gen}{\rightarrow} \mathcal{I}^*)) : \neg y \overset{A}{\rightarrow} \neg x = x \overset{A}{\rightarrow} y $$

\textbf{Dialética :} Basta verificar para $x=0$ e $y=1$, pois a definição de $( 0 \overset{A}{\rightarrow} 1 )$ diz respeito a esses valores.\\

\textbf{Comentário :} $\Phi$ é uma notação para dizer sobre as informações que são assumidas como verdadeiras. $\Phi$ é um tipo de quantificador dentro da teoria de fórmulas universais, entretanto tais informações não podem ser negadas como no operador universal, por exemplo $\forall x :\phi(x)$ tem sua negação $\exists x : \neg \phi(x)$ com $\exists$ significando que existe um $x$ tal que $\phi(x)$ é falso, $\Phi( x : \phi(x) )$ não pode ser negado.\\

\textbf{Relação [ de Diferença ]} $ x \neq y := \neg (x = y) $\\

\textbf{[ Tautologia ]} Uma tautologia é uma fórmula lógica em que as variáveis são informações tal que para qualquer valor de verdadeiro e falso da informação a fórmula é verdadeira. Por exemplo $ x = \neg \neg x $ é uma tautologia, pois se $x=0$ então temos $(0=0)=1$ e se $x=1$ temos $ (1=1) = 1 $.\\

\textbf{Comentário :} A inferência de reconhecimento de definição de tautologia será chamada de inferência de reconhecimento de tautologia. Essencialmente para reconhecer que uma fórmula é uma tautologia precisamos testar todas as possibilidades de valores de verdadeiro e falso das variáveis.\\

\textbf{Comentário :} Todas as tautologias são verdadeiras, pois são adicionadas a razão da teoria lógica através do Axioma [ Notação Funcional de Informação ]. Note que isso vale tanto para informações particulares quanto para a fórmula universal da tautologia.\\

\textbf{Comentário :} $ ( \phi(\cdot) \overset{part}{\rightarrow} \texttt{Tautologia}^* ) \iff (( x^0 \overset{gen}{\rightarrow} \mathcal{I}^* )) : \phi(x) $ é verdadeiro por causa da definição de tautologia, assim como $ ( \phi(\cdot, \cdot) \overset{part}{\rightarrow} \texttt{Tautologia}^* ) \iff (( x^0 \overset{gen}{\rightarrow} \mathcal{I}^* )) (( y^0 \overset{gen}{\rightarrow} \mathcal{I}^* )) : \phi(x,y) $ e análogos.\\

\textbf{Comentário :} Uma tautologia tem a interessante propriedade de poder gerar novas tautologias substituindo as variáveis informações por outras fórmulas.

%%%%%%%%%%%%%%%%%%%%%%%%%%%%%%%%%%%%%%%%%%%%%%%%%%%%%%%%%%%%%%%%%%%%%%%%%%%%%%%%%%%%%%%%%%%%%%%%%%%%%%%%%%%%
% Proposição de Assimetria de Inferência
\hspace{\baselineskip}

\textbf{Proposição [ Assimetria de Inferência ]}

%$$ \Phi( (x \overset{A}{\rightarrow} y) \neq  (y \overset{A}{\rightarrow} x ) ): [ x \overset{A}{\rightarrow} y ](0,1) = 1 $$

$$ \Phi( [ x \overset{A}{\rightarrow} y ](0,1) = 0 ) : ( x \overset{A}{\rightarrow} y ) = ( y \overset{A}{\rightarrow} x ) $$

\textbf{Dialética :} 

\begin{itemize}
	\item[ ] ?` $ ( x \overset{A}{\rightarrow} y ) = ( y \overset{A}{\rightarrow} x ) \overset{part}{\rightarrow} \texttt{Tautologia}^* $  ? $\rightarrow$
	\begin{itemize}
		\item[ ] ?` $ \Phi( x=y ) : ( x \overset{A}{\rightarrow} y ) = ( y \overset{A}{\rightarrow} x )  $ ?
		\item[ ] ?` $ \Phi( x=0, y=1 ) : ( x \overset{A}{\rightarrow} y ) = ( y \overset{A}{\rightarrow} x ) $   ?
		\item[ ] ?` $ \Phi( x=1, y=0 ) : ( x \overset{A}{\rightarrow} y ) = ( y \overset{A}{\rightarrow} x ) $  ?
		\item[ ] !` Reconhecimento de Tautologia !
	\end{itemize}
\end{itemize}

\textbf{Comentário :} Sempre que a proposição for suficientemente fácil de verificar através da dialética será omitido a demonstração. A dialética é ao mesmo tempo um exercício como uma maneira mais elegante de expor a demonstração de um teorema. \\

\textbf{Comentário :} É interessante que a inferência $\implies$ seja uma informação assimétrica, a equivalência $\iff$ pode ser simétrica, mas quando dizemos que $x\implies y$ não podemos dizer automaticamente que $y \implies x$. Suponha que $x$ seja um cachorro, ser um cachorro implica que $x$ é um animal, entretanto se $x$ for um animal isso não implica necessariamente que $x$ seja um cachorro.\\

\textbf{Comentário :} As proposições feitas para inferência A mostram que se quisermos que a inferência definida como operador de informação seja um bom modelo para o conceito de inferência ela deve ser justamente igual a inferência booleana. Se quisermos que o operador inferência satisfaça as propriedades de assimetria e raciocínio inverso de inferência, além das propriedades que caraterizam a inferência, então ela deve ser a inferência booleana.

%%%%%%%%%%%%%%%%%%%%%%%%%%%%%%%%%%%%%%%%%%%%%%%%%%%%%%%%%%%%%%%%%%%%%%%%%%%%%%%%%%%%%%%%%%%%%%%%%%%%%%%%%%%%
% Negação de Quantificador Universal
\hspace{\baselineskip}

\textbf{Quantificador [ Universal de Informação ]} $ (\forall x \in \mathcal{I} : \phi(x)) $ significa que para toda informação $x$, $\phi(x)$ é verdadeiro.\\

\textbf{Quantificador [ Existencial de Informação ]} $ (\exists x \in \mathcal{I} : \phi(x) ) := (\forall x \in \mathcal{I}  : \neg \phi(x) )$ que é a negação do quantificador universal.\\

\textbf{Quantificador [ Neutro de Informação ]} $ \Phi( \psi ) : \phi $ é chamado de quantificador neutro, qualquer coisa dentro de $\Phi(\dots)$ é assumido como verdadeiro sem possibilidade de questionamento. São chamados de hipóteses do problema e em geral são informações já demonstradas ou definições.\\

\textbf{Comentário :} O quantificador neutro funciona como um elemento de interface de raciocínio. É similar ao quantificador universal entretanto não pode ser negado, por exemplo se demonstrarmos que $ \Phi( \dots ) : \forall(\dots) : \exists(\dots) : \psi(\dots) $ é falso, então temos que $ \Phi(\dots) : \exists(\dots) : \forall(\dots) : \neg \psi(\dots) $ é verdadeiro.\\

\textbf{Comentário :} Para demonstrar problemas com quantificador neutro basta utilizar as inferências que são utilizadas no quantificador universal, seja  ?` $ \Phi(A) : B $ ? então basta demonstrar ?` $ A \implies B$  ? da mesma maneira que as demonstrações de $\forall A : B$ seriam feitas.\\

\textbf{Propriedade [ Troca de Posição do Quantificador Neutro ]} O quantificador neutro pode mudar livremente de posição, ele representa as informações que são estabelecidas a priori da construção da informação. Portanto $( \Phi (\psi) : \forall x : \phi(x) )=( \forall x : \Phi(\psi) : \phi(x) )=( \forall x : \phi(x) : \Phi(\psi) )$. O mesmo não é sempre verdade para os outros quantificadores.

%%%%%%%%%%%%%%%%%%%%%%%%%%%%%%%%%%%%%%%%%%%%%%%%%%%%%%%%%%%%%%%%%%%%%%%%%%%%%%%%%%%%%%%%%%%%%%%%%%%%%%%%%%%%
% Axioma de Inferência 
\hspace{\baselineskip}

\textbf{Axioma [ Inferência como Operador ]} $\implies$ é um operador binário de informações. \\

\textbf{Axioma [ Raciocínio Inverso de Inferência ]} $ \forall x, y \in \mathcal{I} : (x \implies y )= (\neg y \implies \neg x) $ \\

\textbf{Axioma [ Assimetria de Inferência ]} $ \exists x,y \in \mathcal{I} : (x \implies y) \ne (y \implies x) $

%%%%%%%%%%%%%%%%%%%%%%%%%%%%%%%%%%%%%%%%%%%%%%%%%%%%%%%%%%%%%%%%%%%%%%%%%%%%%%%%%%%%%%%%%%%%%%%%%%%%%%%%%%%%
% Teorema
\hspace{\baselineskip}

\textbf{Teorema [ Equivalência de Inferência e Inferência Booleana ]} 

$$ \forall x,y \in \mathcal{I} : (x \implies y) = (x \rightarrow y) $$

\textbf{Dialética :} !` Reconhecimento de Tautologia  !

\begin{itemize}
	\item[ ] ?` $ [ x \implies y ](1,1) = 1 $ ? $\rightarrow$
	\begin{itemize}
		\item[ ] ?` $ (x=1 \& y = 1) \implies (x \implies y) $ ?
		\begin{itemize}
			\item[ ] $y=1$
			\item[ ] ?` $ x \implies y $  ? $\rightarrow$ !` Reconhecimento de Inf. Fund. !
		\end{itemize}
	\end{itemize}
	\item[ ] ?` $ [x \implies y](1,0) = 0 $  ? $\rightarrow$
	\begin{itemize}
		\item[ ] $\texttt{AF}$:= (Axioma de Falsidade)
		\item[ ] ?` $ [x \implies y](1,0) \implies \neg \texttt{AF} $  ? $\rightarrow$ 
		\begin{itemize}
			\item[ ] ?` $1$ ? $\rightarrow$ !` Tautologia e Notação Fundamental de Informação !
			\item[ ] ?` $ ( 1 \& (1 \implies 0) ) \implies 0 $  ?
			\item[ ] !` Por inferência fundamental a inclusão de informações falsas nega AF !
		\end{itemize}
		\item[ ] !` Axioma de Raciocínio Inverso !
	\end{itemize}
	\item[ ] ?` $ [x \implies y](0,0) = 1 $ ? $\rightarrow$ !` Axioma de Raciocínio Inverso de Inferência !
	\item[ ] ?` $ [ x \implies y] (0,1) = 1 $  ? $\rightarrow$ !` Axioma de Assimetria de Inferência !
	\item[ ] !` Reconhecimento de Tautologia  !
\end{itemize}

\textbf{Comentário :} Quando construir uma teoria axiomática, é interessante levar em conta o propósito fundamental da filosofia que é criar um caminho natural entre as ideias, portanto devemos sair do que nos é natural para tentar entender conceitos mais complicados. Com este teorema mostramos que a definição $\rightarrow$ pouco natural de inferência pode ser tangenciado através de axiomas que são mais naturais de serem aceitos sobre o conceito de inferência.\\

\textbf{Comentário :} Verificar tautologias através de lógica booleana (interpretar a lógica por operadores e informações por 0 e 1) é uma tarefa efetiva, basta substituir por zeros e uns para cada situação e verificar o resultado. Conectar a inferência e o operador $\rightarrow$ definido pela lógica booleana faz com que a demonstração das proposições da lógica se torne a verificação de alguma tautologia.\\



\textbf{Proposição [ Conjunção de Inferência e Conjunção Booleana ]}

$$ \forall x, y \in \mathcal{I} : ( (x \wedge y) \implies (x \& y) )$$

\textbf{Dialética :} 

\begin{itemize}
	\item[ ] ?` $ \forall x,y : (x \wedge y) \implies x $  ? $\rightarrow$ !` Verificar se é uma tautologia !
	\item[ ] ?` $ \forall x,y : (x \wedge y) \implies y $ ? $\rightarrow$ !` Verificar se é uma tautologia !
	\item[ ] $ x,y $ informações particulares.
	\item[ ] ?` $ (x \wedge y) \implies (x \& y) $ ? $\rightarrow$ !` Reconhecimento de $\&$ !
	\item[ ] !` Reconhecimento de $\implies$ !
\end{itemize}

\hspace{\baselineskip}

\textbf{Igualdade [ Operador Booleano ]} Estabelecemos que a igualdade operacional de operadores booleanos é $ ( f \overset{op}{=} g ) \iff ( f(\cdot) = g(\cdot) \overset{part}{\rightarrow} \texttt{Tautologia}^* ) $ considerando $\cdot$ como uma ou mais variáveis. Dessa forma $ \implies \overset{op}{=} \rightarrow $ e $ \& \overset{op}{=} \wedge $ e portanto elas podem ser substituídas nas fórmulas usando inferência de substituição.

%%%%%%%%%%%%%%%%%%%%%%%%%%%%%%%%%%%%%%%%%%%%%%%%%%%%%%%%%%%%%%%%%%%%%%%%%%%%%%%%%%%%%%%%%%%%%%%%%%%%%%%%%%%%
% Dialética Socrática de Construção dos Operadores Lógicos Booleanos
\hspace{\baselineskip}

\textbf{Dialética Socrática [ Reconstrução dos Operadores Lógicos Booleanos ]}

\begin{itemize}
	\item[ (1) ] É necessário o operador de negação. Pelo menos faz sentido ser necessário por conta da conexão entre o conceito de identidade de existência, reconhecimento e etc.
	\item[ (2) ] ?` É possível reconstruir os operadores booleanos a partir da definição de operador de negação e da definição de um dos operadores binários ?
	\begin{itemize}
		\item[ (1) ] É interessante que seja possível, pois a teoria lógica começa a se conectar com a teoria existencial exposta no capítulo de projeções dialéticas, basta algo que modele o conceito de reconhecimento de existência (identidade) e basta algo que modelo o conceito de interação (no que foi feito anteriormente foi utilizado a conjunção) que é necessária para a propriedade de existir para que um universo possa ser construído (na analogia o universo seria a lógica).
		\item[ (2) ] Já sabemos que podemos construir os operadores por negação e conjunção.
		\item[ (3) ] ?` É possível construir os operadores a partir da negação e disjunção ?
		\begin{itemize}
			\item[ (1) ] Podemos definir conjunção por $x \wedge y := \neg( \neg x \vee \neg y )$, portanto se partirmos de negação e disjunção temos a conjunção. Com a negação e conjunção podemos definir os demais operadores, portanto a resposta é afirmativa.
		\end{itemize}
		\item[ (4) ] ?` É possível construir os operadores a partir da negação e inferência ?
		\begin{itemize}
			\item[ (1) ] De maneira análoga temos a disjunção capaz de ser definida pela inferência $ x \vee y := (\neg x \rightarrow y) $, portanto a resposta é afirmativa.
		\end{itemize}
	\end{itemize}
	\item[ (3) ] \textbf{[ Exposição de Teoria Lógica Projetada no Tempo ]} É uma teoria lógica que reconstrói os operadores usando negação e inferência.
	\item[ (4) ] \textbf{[ Exposição de Teoria Lógica Projetada no Espaço ]} É uma teoria lógica que reconstrói os operadores usando negação e conjunção (espaço positivo) ou negação e disjunção (espaço negativo).
\end{itemize}

\textbf{Comentário [ Base Espaço e Tempo ]} A interpretação da exposição da lógica no tempo e espaço é interessante, pois a intuição do tempo e espaço sempre estão presentes na maneira de pensar humana, a constatação da possibilidade de projetar a exposição da lógica nesta base dialética reforça a tese de que nossa consciência da realidade é dependente da consciência do espaço e tempo.\\

\textbf{Comentário [ Base Trialética ]} Podemos constatar algumas características da base trialética, uma das características é que existem dois operadores parecidos (conjunção e disjunção), o que faz parecer que estão na mesma topologia, e ambos foram associados ao espaço, e temos um operador (inferência) que está associado ao tempo. O operador associado a interpretação temporal podemos associar a gama, enquanto alfa e beta são um dos operadores de interpretação espacial, apesar de eu não me atrever a dizer que um é alfa e outro é gama, pois é o mesmo que querer caracterizar cargas negativas e positivas, sendo que o que se sabe é que são apenas diferentes.


%%%%%%%%%%%%%%%%%%%%%%%%%%%%%%%%%%%%%%%%%%%%%%%%%%%%%%%%%%%%%%%%%%%%%%%%%%%%%%%%%%%%%%%%%%%%%%%%%%%%%%%%%%%%
% Inferência de Reconhecimento de Negação de Falsidade
\hspace{\baselineskip}

\textbf{Inferência [ Reconhecimento de Negação de Falsidade ]}
$$ 0 $$
\hrule
$$ \neg \texttt{AF}$$

\textbf{Comentário :} Essa inferência pode ser considerada como uma inferência de reconhecimento do axioma de falsidade, em que não podemos adicionar a razão informações falsas. Explicitar essa inferência torna a argumentação de demonstrações por raciocínio inverso mais fáceis.\\

\textbf{Comentário :} Essa inferência implica que em toda teoria que herda a razão da lógica $1$ está presente na razão. Basta usar o axioma de raciocínio inverso de inferência.

%%%%%%%%%%%%%%%%%%%%%%%%%%%%%%%%%%%%%%%%%%%%%%%%%%%%%%%%%%%%%%%%%%%%%%%%%%%%%%%%%%%%%%%%%%%%%%%%%%%%%%%%%%%%
% Inferência de Reconhecimento de 

%%%%%%%%%%%%%%%%%%%%%%%%%%%%%%%%%%%%%%%%%%%%%%%%%%%%%%%%%%%%%%%%%%%%%%%%%%%%%%%%%%%%%%%%%%%%%%%%%%%%%%%%%%%%
% Proposição [ Teoria sem Informação ]
\hspace{\baselineskip}

\textbf{Proposição [ Teoria sem Informação ]}

$$ \forall x \in \mathcal{I} : (0 \implies x) $$

\textbf{Dialética :} Basta verificar que $ (0 \implies x) $ é uma tautologia.\\

\textbf{Comentário :} A ideia desta proposição é que se em algum momento pudermos adicionar uma informação falsa a razão de uma teoria, é possível concluir que todas as informações são verdadeiras. Para uma teoria isso significa que além de violar o meio excluído ela não consegue distinguir o certo do errado para nenhuma informação e portanto é uma teoria sem propósito. É importante deixar claro que isso não significa dizer que o axioma de falsidade é necessário, já foi assumido o axioma de falsidade na demonstração da proposição acima e portanto a questão de ser verdadeiro ou falso não é possível ser feita para o axioma de falsidade. O que podemos dizer aqui é, se for construído uma teoria em que em algum momento concluímos que uma informação falsa é verdadeira então devemos abandonar a teoria como uma teoria axiomática, pois neste exato momento a teoria não tem capacidade de responder a mais nada. Isso nos leva a pergunta, será que apenas por sorte o ser humano não encontrou uma conclusão deste tipo (informação falsa como verdadeira) entretanto existe alguma informação que pela lógica e por regras de inferência essa conclusão desastrosa possa ser feita de maneira válida? No momento que provarem que é possível tal conclusão a lógica precisa ser refeita assim como todas as teorias dela dependentes precisam ser novamente demonstradas.

%%%%%%%%%%%%%%%%%%%%%%%%%%%%%%%%%%%%%%%%%%%%%%%%%%%%%%%%%%%%%%%%%%%%%%%%%%%%%%%%%%%%%%%%%%%%%%%%%%%%%%%%%%%%
% Fórmulas Booleanas
\hspace{\baselineskip}

\textbf{[ Fórmulas Booleanas ]} Fórmulas booleanas são fórmulas construídas por informações que podem assumir verdadeiro e falso e operadores de negação e um operador relacional ($\wedge, \vee, \rightarrow$). Fórmulas universais fogem do escopo de fórmulas booleanas a não ser que seja utilizado $\rightarrow$ como operador para descrever e as variáveis serem exclusivamente informações.\\

\textbf{[ Tautologia Booleana ]} Uma fórmula booleana que é uma tautologia é chamada de tautologia booleana. 

%%%%%%%%%%%%%%%%%%%%%%%%%%%%%%%%%%%%%%%%%%%%%%%%%%%%%%%%%%%%%%%%%%%%%%%%%%%%%%%%%%%%%%%%%%%%%%%%%%%%%%%%%%%%
% Externalização Booleana de Axioma de Falsidade e Axioma de Lei do Meio Excluído
\hspace{\baselineskip}

\textbf{Representação [ Axioma de Falsidade ]} $ \texttt{AF} =^{*} ( \neg (1 \rightarrow 0) ) $\\

\textbf{Representação [ Axioma de Lei do Meio Excluído ]} $ \texttt{LME} =^* \forall x \in \mathcal{I} : (\neg (\neg x \wedge x)) $\\

\textbf{Comentário :} O símbolo $=^{*}$ representa a ideia de que negar essa relação de igualdade é o mesmo que negar a capacidade de reconhecer alguma definição presente na fórmula, nestes casos a capacidade de reconhecer os axiomas. Na prática ela é uma relação de igualdade, entretanto para o leitor estamos dizendo que essa relação é evidente por definição e a incapacidade de compreender significa que o leitor não possui a consciência informacional alfa de algum conteúdo presente na expressão. \\

\textbf{Comentário :} Qualquer teoria axiomática a cerca de teorias fundamentais para o raciocínio, como a lógica ou teoria de conjuntos, requer inicialmente a descrição verbal de seu conteúdo, pois não há símbolos e notações definidos, entretanto uma vez construído notações para escrever fórmulas tais axiomas podem ser escritos utilizando a notação que foi construída mas a demonstração dessas fórmulas torna-se normalmente sem sentido ou talvez impossível e requer acreditar na capacidade da consciência humana de compreender.\\

\textbf{Comentário :} É interessante observar que $ \texttt{LME} = ( \forall x \in \mathcal{I} : (x \rightarrow x) ) $, basta aplicar identidades ou constatar a tautologia, porém $x \rightarrow x$ parece ter menos informação que $ \neg (\neg x \wedge x) $, $x \rightarrow x $ foi introduzido como uma inferência pela necessidade de conseguir construir a razão de uma teoria, enquanto $ \neg ( \neg x \wedge x ) $ caracteriza uma propriedade fundamental da informação.\\

\textbf{Comentário :} Alguém poderia pensar que seria melhor começar a lógica definindo os operadores booleanos e depois atribuindo significado a cada um dos operadores, isso poderia ser feito, mas o operador $\rightarrow$ não é um operador trivial de ser compreendido enquanto $\implies$ é algo mais intuitivo.

%%%%%%%%%%%%%%%%%%%%%%%%%%%%%%%%%%%%%%%%%%%%%%%%%%%%%%%%%%%%%%%%%%%%%%%%%%%%%%%%%%%%%%%%%%%%%%%%%%%%%%%%%%%%
% Inferência de Absurdo
\hspace{\baselineskip}

\textbf{Inferência [ Reconhecimento de Absurdo ]}

$$ x $$
$$ \neg x $$
\hrule
$$ \neg \texttt{LME} $$

\textbf{Comentário :} É possível demonstrar que $\texttt{LME}$ equivale a $\texttt{AF}$ ? É possível considerar um deles como sendo derivado do outro? Estrategicamente eles não são equivalentes, pois a inferência de reconhecimento de absurdo é mais interessante para tirar conclusões por estar em função de uma informação qualquer, enquanto o $\texttt{AF}$ necessita operar para obter $0$. Ambos porém permitem o raciocínio por absurdo que é demonstrar algo a partir da argumentação de que se assumir sua negação então temos uma contradição no raciocínio.

%%%%%%%%%%%%%%%%%%%%%%%%%%%%%%%%%%%%%%%%%%%%%%%%%%%%%%%%%%%%%%%%%%%%%%%%%%%%%%%%%%%%%%%%%%%%%%%%%%%%%%%%%%%%
% Dialética Demonstrativa
\hspace{\baselineskip}

\textbf{Dialética Demonstrativa [ Demonstração por Absurdo via LME ]} A dialética demonstrativa é uma dialética socrática para descrever a estratégia de demonstração por absurdo, ela utiliza reconhecimento de raciocínio inverso e reconhecimento de absurdo. \\

?` $\Phi( x^0 \overset{part}{\rightarrow} \mathcal{I}^* ) : x $  ? $\rightarrow$

\begin{itemize}
	\item[ ] ?` $ \neg x \implies \neg \texttt{LME} $ ? $\rightarrow$
	\begin{itemize}
		\item[ (1) ] !` Razão que inicia com $\neg x$ mas não pode conter $x$ isolado !
		\item[ (2) ] !` Inferência de Reconhecimento de Absurdo de alguma outra informação assumida como verdade diferente de $x$ !
	\end{itemize}
	\item[ ] !` Inferência de Raciocínio Inverso !
\end{itemize}

\textbf{Comentário :} É importante perceber a importância que (1) tem em não poder conter $x$ isolado. Se você obteve $x$ isolado você demonstrou diretamente e basta utilizar a inferência de reconhecimento de inferência $\implies$, muito provavelmente se você obteve $x$ no raciocínio você utilizou alguma informação derivada do próprio $x$, porém não é interessante utilizar informações e conclusões de algo que se quer demonstrar, pois este algo pode estar errado.\\

\textbf{Comentário :} Em (2) a cláusula de que a informação tenha que ser diferente de $x$ é importante pois não podemos esperar obter $x$ e $\neg x$ em um raciocínio válido, então devemos tentar obter $\neg x \implies \neg y$ com $y$ uma informação que antes do raciocínio era verdadeira.\\

\textbf{Comentário :} Estrategicamente, demonstrar por absurdo permite que possamos utilizar uma nova fórmula ($\neg x$) no raciocínio, diferente do raciocínio direto que precisa partir de fórmulas das hipóteses, no raciocínio por absurdo temos tanto as fórmulas das hipóteses como essa nova fórmula ($\neg x$) que eventualmente vai tornar a razão inválida. Ter mais fórmulas permite explorar mais opções de raciocínio, o que torna mais interessante para resolver um problema. \\

\textbf{Comentário :} Outro argumento estratégico é a noção intuitiva de que destruir algo é mais fácil do que construir, no raciocínio por absurdo queremos 'destruir' a razão achando uma contradição e portanto parece mais fácil do que o raciocínio direto que necessita construir a razão. \\

\textbf{Dialética Demonstrativa [ Demonstração por Absurdo via AF ]} A dialética demonstrativa é uma dialética socrática para descrever a estratégia de demonstração por absurdo, ela utiliza reconhecimento de raciocínio inverso e reconhecimento de negação de axioma de falsidade. \\

?` $\Phi( x^0 \overset{part}{\rightarrow} \mathcal{I}^* ) : x $  ? $\rightarrow$

\begin{itemize}
	\item[ ] ?` $ \neg x \implies \neg \texttt{AF} $ ? $\rightarrow$
	\begin{itemize}
		\item[ (1) ] !` Razão que inicia com $\neg x$ mas não pode conter $x$ isolado !
		\item[ (2) ] !` Inferência de Reconhecimento de Negação de Axioma de Falsidade de alguma outra informação assumida como verdade diferente de $x$ !
	\end{itemize}
	\item[ ] !` Inferência de Raciocínio Inverso !
\end{itemize}


\textbf{Comentário :} Recomendo que se for a intenção do leitor ser um teórico, memorizar dialéticas demonstrativas até obter consciência alfa. Uma vez com consciência alfa será permitido que se faça demonstrações de maneira mais ágil.\\

Para revisar e ter um set de estratégias de demonstrações irei definir dialéticas demonstrativas para demonstrações genéricas da lógica.

%%%%%%%%%%%%%%%%%%%%%%%%%%%%%%%%%%%%%%%%%%%%%%%%%%%%%%%%%%%%%%%%%%%%%%%%%%%%%%%%%%%%%%%%%%%%%%%%%%%%%%%%%%%%
% Dialética Demonstrativa
\hspace{\baselineskip}

\textbf{Dialética Demonstrativa [ Raciocínio Direto ]} \\

?` $ \Phi(\dots) : \phi \implies \psi $ ? $\rightarrow$

\begin{itemize}
	\item[ (1) ] !` Começar uma razão com $\phi$ herdando razões de $\Phi(\dots)$ e terminar em $\psi$ !
	\item[ (2) ] !` i($\implies$) !
\end{itemize}

\textbf{Notação :} i($\implies$) denota a inferência de reconhecimento de $\{ \implies \}$.

%%%%%%%%%%%%%%%%%%%%%%%%%%%%%%%%%%%%%%%%%%%%%%%%%%%%%%%%%%%%%%%%%%%%%%%%%%%%%%%%%%%%%%%%%%%%%%%%%%%%%%%%%%%%
% Dialética Demonstrativa
\hspace{\baselineskip}

\textbf{Dialética Demonstrativa [ Demonstração por Absurdo ]}\\

?` $ \Phi(\dots) : \phi $ ? $\rightarrow$

\begin{itemize}
	\item[ (1) ] ?` $ \neg \phi \implies \emptyset $  ?
	\item[ (2) ] !` ira !
\end{itemize}

\textbf{Notação :} $\emptyset$ representa tanto $\neg \texttt{AF}$ como $\neg \texttt{LME}$ e ira representa a inferência associada para reconhecer essas ocorrências.

%%%%%%%%%%%%%%%%%%%%%%%%%%%%%%%%%%%%%%%%%%%%%%%%%%%%%%%%%%%%%%%%%%%%%%%%%%%%%%%%%%%%%%%%%%%%%%%%%%%%%%%%%%%%
% Dialética Demonstrativa
\hspace{\baselineskip}

\textbf{Dialética Demonstrativa [ Informação Universal ]} \\

?` $ \Phi(\dots) : (( x^0 \overset{gen}{\rightarrow} X^* )) : \phi(x) $ ? $\rightarrow$

\begin{itemize}
	\item[ (1) ] ?` $ ( x^0 \overset{gen}{\rightarrow} X^* ) \implies \phi(x) $ ?
	\item[ (2) ] !` idiu !
\end{itemize}

\textbf{Notação :} idiu denota inferência de demonstração de informação universal.

%%%%%%%%%%%%%%%%%%%%%%%%%%%%%%%%%%%%%%%%%%%%%%%%%%%%%%%%%%%%%%%%%%%%%%%%%%%%%%%%%%%%%%%%%%%%%%%%%%%%%%%%%%%%
% Dialética Demonstrativa
\hspace{\baselineskip}

\textbf{Dialética Demonstrativa [ Raciocínio Inverso ]}\\

?` $ \Phi( \dots ) : \phi \implies \psi $ ? $\rightarrow$

\begin{itemize}
	\item[ (1) ] ?` $ \neg \psi \implies \neg \phi $  ? 
	\item[ (2) ] !` iri !
\end{itemize}

\textbf{Notação :} iri é a inferência de raciocínio inverso do símbolo ($\implies$).

%%%%%%%%%%%%%%%%%%%%%%%%%%%%%%%%%%%%%%%%%%%%%%%%%%%%%%%%%%%%%%%%%%%%%%%%%%%%%%%%%%%%%%%%%%%%%%%%%%%%%%%%%%%%
% Dialética Demonstrativa de Implicação por absurdo
\hspace{\baselineskip}

\textbf{Dialética Demonstrativa [ Raciocínio Direto e Absurdo ]}\\

?` $ \Phi(\dots) : \phi \implies \psi $  ? $\rightarrow$ 

\begin{itemize}
	\item[ (1) ] ?` $ (\phi \wedge \neg \psi) \implies \emptyset  $ ?
	\item[ (2) ] !` ira !
\end{itemize}

\textbf{Comentário :} Apesar de estranho dizer, você deve ser capaz de demonstrar que a dialética demonstrativa de raciocínio direto é essencialmente o problema de raciocínio direto resolvido pela dialética de demonstração por absurdo. Se comparar com o raciocínio inverso podemos notar que temos duas informações que podemos utilizar $\phi$ e $\neg \psi$ para tentar destruir a razão, enquanto no raciocínio inverso o temos apenas $\neg \psi$ sem ganhos de fórmulas comparado ao raciocínio direto.\\

\textbf{Comentário :} É interessante para o leitor de uma demonstração explicitar a dialética utilizada, é como dizer a estratégia utilizada para demonstrar algo, isso facilita a compreensão da demonstração, compreensão esta que é externalizada pela evidenciação das inferências utilizadas e agora também pelas dialéticas demonstrativas. Escrever uma teoria axiomática em forma de exposição trialética é interessante para o leitor por levar em conta a sua consciência alfa.\\

\textbf{Comentário :} Na interface de dialética demonstrativa, dizemos a dialética demonstrativa utilizada usando a notação ?`$ \phi $ ? $\rightarrow$ \{ Nome da Dialética Demonstrativa a ser utilizada \}. Isso torna a demonstração mais clara, pois prepara o leitor a reconhecer o raciocínio desenvolvido.\\

\textbf{Comentário :} Também na interface de dialética demonstrativa podemos colocar a tentativa de demonstrar algo utilizando uma dialética específica como um problema. Por exemplo : !` ( ?` $\phi$ ? $\rightarrow$ \{ Dialética Demonstrativa [ X ] \} ) ! é o problema de demonstrar $\phi$ utilizando a Dialética X. Essa notação facilita o desenvolvimento de teorias, pois podemos criar um mapa dos caminhos explorados para tentar responder a uma pergunta. Desenvolver teorias axiomáticas pode ser dito como um dos objetivos principais da lógica formal.

\subsection{Classes e Conjuntos}

\hspace{\baselineskip}

%%%%%%%%%%%%%%%%%%%%%%%%%%%%%%%%%%%%%%%%%%%%%%%%%%%%%%%%%%%%%%%%%%%%%%%%%%%%%%%%%%%%%%%%%%%%%%%%%%%%%%%%%%%%
% 
A partir de agora irei começar a tentar definir as existências que são chamadas de classes e conjuntos. A matemática costuma tratar a lógica e a teoria de conjuntos de maneira separada, pois ambos possuem conteúdos muito vastos, entretanto tanto um quanto outro são essenciais para o raciocínio. Neste texto não busco aprofundar nem em um e nem em outro assunto, mas abordar-los de uma maneira que seja possível ter ferramentas para construir teorias axiomáticas sem perder muito tempo. Nesta abordagem não irei utilizar os axiomas de Zermello-Frankel por uma questão estética, considero que alguns dos axiomas não são muito elegantes, irei aproveitar o máximo do que foi feito anteriormente e se necessário introduzir axiomas espelhando nos axiomas tradicionais da teoria axiomática de conjuntos. Por conta disso minha teoria pode ter falhas que talvez eu não tenha disposição para investigar, estarei atento na medida do possível, cabe a outro uma investigação mais profunda, recomendo que investigue a teoria de conjuntos tradicional.

%%%%%%%%%%%%%%%%%%%%%%%%%%%%%%%%%%%%%%%%%%%%%%%%%%%%%%%%%%%%%%%%%%%%%%%%%%%%%%%%%%%%%%%%%%%%%%%%%%%%%%%%%%%%
% Proposição [ Tautologia de Igualdade de Informação ]
\hspace{\baselineskip}

\textbf{Proposição [ Tautologia de Igualdade de Informação ]}

$$ (\phi \iff \psi) = ( \phi = \psi ) \texttt{   } \overset{part}{\rightarrow} \texttt{Tautologia}^* $$

\textbf{Dialética :} !` Demonstração de Tautologia !\\

\textbf{Comentário :} Apesar de simples sua demonstração é importante chamar atenção a essa tautologia, ela essencialmente estabelece que $( \overset{inf}{=} ) \overset{op}{=} ( \iff )$ o que faz com que a ocorrência de $\iff$ em uma razão permita a inferência de substituição de informações.



%%%%%%%%%%%%%%%%%%%%%%%%%%%%%%%%%%%%%%%%%%%%%%%%%%%%%%%%%%%%%%%%%%%%%%%%%%%%%%%%%%%%%%%%%%%%%%%%%%%%%%%%%%%%
% Definição Formal de Relação de Pertence
\hspace{\baselineskip}

\textbf{Operador [ Pertence ]} $  (x \in X) := (\in) (x,X) $\\

\textbf{Propriedade [ Assimetria de Pertence ]} $\exists x,y : (x \in y) \neq (y \in x)$\\

\textbf{Propriedade [ Relação Pertence ]} $ (x \in X) \overset{part}{\rightarrow} \mathcal{I}^* $\\

\textbf{Propriedade [ Contexto Trialético de Igualdade ]} O operador pertence ele possui três contextos distintos de igualdade e portanto três contexto de inferência de substituição, na fórmula $ x \in A$ o primeiro contexto é o contexto da variável $x$, ou contexto de elementos com a igualdade $\overset{el}{=}$, o segundo contexto $\in$ é o contexto de operador, entretanto não é um operador booleano, pois um operador booleano requer que suas variáveis sejam do contexto de informações, utilizamos a mesma notação $ \overset{op}{=} $ na esperança de que não encontremos nenhum problema no caminho, o terceiro contexto é o contexto de classes (conjunto é considerado uma classe e portanto se insere neste contexto.) $ A $ é uma classe e a igualdade é denotada por $\overset{cl}{=}$.\\

\textbf{Comentário :} Estamos definindo inicialmente um operador formal, no sentido de que ele não possui o significado que queremos dar ainda. Assumimos como parte da definição suas propriedades. Esse operador é uma informação entretanto suas variáveis podem não ser informações no sentido de que pode não fazer sentido as variáveis terem valores 0 ou 1 ou serem verdadeiros ou falsos, portanto não é um operador booleano.\\

\textbf{Comentário :} A relação pertence é intrínseco da definição de conjunto, $x \in X$ ela associa o elemento $x$ com o conjunto $X$ dizendo que $x$ pertence a $X$, sabe-se lá o que isso quer dizer, essa relação conceitual é similar a relação que uma informação tem com os conceitos de verdadeiro e falso, tanto que usamos a notação $x \in \{0,1\}$ no início da seção. Da teoria de conjuntos sabemos que não necessariamente a parte direita precisa ser um conjunto, mas vou deixar os detalhes para depois.\\

\textbf{Notação [ Conjunto ]} $ A^0 \overset{part}{\rightarrow} \mathcal{S}^* $ é a notação para dizer que o símbolo $A$ representa uma existência do tipo conjunto. Podemos utilizar a notação $ \forall A \in \mathcal{S}^* $ ou $A \in \mathcal{S}^*$, entretanto é importante ressaltar que $\mathcal{S}^*$ não é considerado um conjunto, pois há inconsistências em definir o conjunto de todos os conjuntos na teoria de conjuntos tradicional.\\

\textbf{Notação :} Para facilitar a notação façamos $\in$ ter maior prioridade que $\implies$.\\

\textbf{Notação :} $ (x \notin A) := \neg( x \in A) $

%%%%%%%%%%%%%%%%%%%%%%%%%%%%%%%%%%%%%%%%%%%%%%%%%%%%%%%%%%%%%%%%%%%%%%%%%%%%%%%%%%%%%%%%%%%%%%%%%%%%%%%%%%%%
% Todo conjunto tem um identificador de elementos
\hspace{\baselineskip}

\textbf{Axioma [ Identificador de Elementos ]} $ \forall A  \in \mathcal{S}^* : \exists \phi_{efc} \in \mathcal{I}^* : A:=\{ x \in A \iff \phi(x) \} $\\

\textbf{Comentário [ Paradoxo de Russel ]} A notação $ A := \{ \psi \} $ para um conjunto significa que $ A $ está definido pela informação $\psi$. Essa fórmula $\phi$ é chamado de identificador de conjunto, não necessita ser única, entretanto ter um identificador não implica que a entidade seja um conjunto, temos o famoso paradoxo de Russel utilizando a construção $ X:= \{ a \in X \iff a \notin a \} $, o identificador $\phi(x) = (x \notin x)$ leva a um paradoxo, um identificador deve testar elementos de algum conjunto controlado, da maneira que o identificador foi escrito podemos testar o próprio conjunto, suponha que $ X \notin X $, então $X$ pertence a $X$, o que significa $X \in X$, mas $ X \in X $ contradiz $X \notin X$, suponha agora que $X \in X$, se $X \in X$ então ele satisfaz o identificador $\phi(X)$, portanto $X \notin X$, novamente uma contradição. Mais tarde veremos que é necessário para um identificador ele mesmo ser um conjunto (funções são conjuntos em teoria de conjuntos) no sentido de que ele deve ser uma função que chamamos de indicadora, que leva um conjunto a um valor de verdadeiro ou falso.\\

\textbf{Comentário :} Mais um exemplo de que ter um identificador não significa ser um conjunto, podemos utilizar uma identidade $ A:= \{ x \in A \iff x \in A \} $, essa é uma formulação para um conjunto inacessível, não podemos efetivamente testar nenhum elemento de $A$. $ \phi_{efc} $ é a notação para dizer que $\phi$ é um identificador efetivo, ou seja, é possível testar efetivamente se um elemento pertence ou não ao conjunto, no exemplo temos que $\phi(x) = (x \in A)$ não é um identificador efetivo. \\

\textbf{Comentário :} Chamamos um identificador de informal quando a definição do identificador não pode ser feita por símbolos da lógica ou de elementos matemáticos definidos e precisa ser descrito utilizando linguagem verbal. Muitas vezes é necessário definir identificadores informais por conta da natureza do problema requerer competências da consciência do leitor, assim como visto em definições de inferências.\\

\textbf{Notação [ Estrutura de Conjunto ]} $ (A,\phi)_{set} $ significa dizer que $A$ é um conjunto e $\phi$ é um identificador efetivo para $A$.

%%%%%%%%%%%%%%%%%%%%%%%%%%%%%%%%%%%%%%%%%%%%%%%%%%%%%%%%%%%%%%%%%%%%%%%%%%%%%%%%%%%%%%%%%%%%%%%%%%%%%%%%%%%%
% Operador de Continência
\hspace{\baselineskip}

\textbf{Operador [ Continência ]} $ \forall A \in \mathcal{S}^* : \forall B \in \mathcal{S}^* : (A \subseteq B) \iff ( x \in A \implies x \in B ) $\\

\textbf{[ Subclasse ou Subconjunto ]} Em uma relação $A \subseteq B$, temos que $A$ é chamado de subconjunto/subclasse de $B$ e $B$ superconjunto/superclasse de $A$.

%%%%%%%%%%%%%%%%%%%%%%%%%%%%%%%%%%%%%%%%%%%%%%%%%%%%%%%%%%%%%%%%%%%%%%%%%%%%%%%%%%%%%%%%%%%%%%%%%%%%%%%%%%%%
% Inferência de Reconhecimento de Definição
\hspace{\baselineskip}

\textbf{Notação [ Definição ]} $ \phi \overset{def}{\iff} \psi $ significa que $\overset{def}{\iff}$ é a inferência de reconhecimento de definição. Para facilitar a leitura, esse operador deve ter a menor prioridade que qualquer outro símbolo operador.

%%%%%%%%%%%%%%%%%%%%%%%%%%%%%%%%%%%%%%%%%%%%%%%%%%%%%%%%%%%%%%%%%%%%%%%%%%%%%%%%%%%%%%%%%%%%%%%%%%%%%%%%%%%%
% Introdução a Classes
\hspace{\baselineskip}

Para introduzir o conceito de classes vamos explorar um pouco os problemas que existem na teoria de conjuntos, todos correlatos ao paradoxo de Russel. Em ultima instância vamos contornar esse problema definindo classes e posteriormente definindo classes especiais chamadas de conjuntos.

%%%%%%%%%%%%%%%%%%%%%%%%%%%%%%%%%%%%%%%%%%%%%%%%%%%%%%%%%%%%%%%%%%%%%%%%%%%%%%%%%%%%%%%%%%%%%%%%%%%%%%%%%%%%
% Axioma de Separação
\hspace{\baselineskip}

\textbf{[ Axioma Esquema de Especificação ]} 

$$\forall A \in \mathcal{S}^* : \forall \phi \in \mathcal{I}^* : B:= \{ x \in B \iff (x \in A \wedge \phi(x) ) \} \implies B \in \mathcal{S}^* $$

\textbf{Comentário :} Esse axioma não é assumido como um axioma válido na nossa teoria lógica, portanto denoto ele como um nome [ Axioma Esquema de Especificação ], esse nome não implica que o axioma seja verdadeiro para nossa razão. Esse axioma é um dos axiomas de Zermello-Frankel da teoria de conjuntos.\\

\textbf{Comentário :} Era esperado que o axioma fosse da forma $ \forall A \in \mathcal{S}^* : \forall \phi \in \mathcal{I}^* : B:= \{ x \in B \iff \phi(x) \} \implies B \in \mathcal{S}^* $, entretanto conjuntos como $ R:=\{ x \in R \iff x \notin R \} $ ou $ R:=\{ x \in R \iff x \notin x \} $ do paradoxo de Russel fazem com que a liberdade de definir conjuntos permitam conjuntos que gerem contradição lógica. Vamos explorar um dos problemas definindo primeiro o conjunto vazio e depois o conjunto complementar ao vazio que vou chamar de conjunto de todas as coisas.\\

\textbf{Notação :} $\texttt{AEE} :=$ Axioma Esquema de Especificação $^*$\\

\textbf{Classe [ Vazio ]} $ \emptyset := \{  x \in \emptyset \iff 0 \} $ \\

\textbf{Comentário :} Na teoria de conjuntos usando o axioma esquema de especificação a classe vazio é um conjunto. Vou assumir que seja uma classe mesmo não definindo o conceito de classe ainda, pois acredito que isso não será um grande problema.\\

\textbf{Comentário :} O identificador da classe é uma função constante $ \phi(x) = 0 $ seja qual for $x$ a ser testado.\\

\textbf{Classe [ de Todas as Coisas ]} $ \Xi := \{ x \in \Xi \iff 1 \} $\\

\textbf{Comentário :} Parece ser razoável definir algo do tipo, se assumirmos o vazio como um conjunto por que a classe de todas as coisas não seria um conjunto? Vamos mostrar que não pode ser pois pelo axioma esquema de especificação temos o paradoxo de Russel.\\

\textbf{Proposição [ Conjunto de Todas as Coisas ]}

$$ \Phi( \texttt{AEE} ) : \Xi \notin \mathcal{S}^* $$

\textbf{Dialética :} $\rightarrow$ !` Dialética de Demonstração por Absurdo !\\

?` $ \Xi \in S^* \implies \emptyset $  ? $\rightarrow$

\begin{itemize}
	\item[ ] $ R := \{ x \in R \iff ( x \in \Xi \wedge x \notin x ) \} $
	\item[ ] ?` $ \texttt{AEE} \implies R \in \mathcal{S}^* $ ?
	\item[ (i) ] Na teoria de conjuntos quando temos um conjunto podemos a inferência de identificação no caso $ x \in R \iff ( x \in \Xi \wedge x \notin x ) $ é adicionado na razão. Esse detalhe é importantíssimo, no presente momento vamos assumir que podemos utilizar essa inferência para todo conjunto.
	\item[ ] ?` $ R \in R \implies R \notin R $ ?
	\item[ ] ?` $ R \notin R \implies R \in R $ ?
	\item[ ] !` Inferência de Reconhecimento de Absurdo !
\end{itemize}

\textbf{Comentário :} A verdade é que na teoria deste texto, essa dialética não leva a uma demonstração válida por conta de $(i)$. Justamente $(i)$ será o mecanismo para contornar o paradoxo. Convém por hora demonstrar como exercício a partir da dialética assumindo $(i)$.\\

\textbf{Comentário :} Na teoria de conjuntos de Zermello-Frankel, $\Xi$ não pode ser um conjunto, pois leva ao paradoxo de Russel.\\

\textbf{[ Detalhe Escondido da Teoria de Conjuntos ]} Vamos chamar $(i)$ na dialética da proposição do conjunto de todas as coisas como $\texttt{DETC}$, esse nome vai se tornar claro quando definirmos classes.

%%%%%%%%%%%%%%%%%%%%%%%%%%%%%%%%%%%%%%%%%%%%%%%%%%%%%%%%%%%%%%%%%%%%%%%%%%%%%%%%%%%%%%%%%%%%%%%%%%%%%%%%%%%%
% 
\hspace{\baselineskip}

\textbf{Dialética Socrática :} Suponha $ R:= \{ x \in R \iff ( x \in \{R\} \wedge x \notin x ) \} $, esse conjunto pode ser construído utilizando $\texttt{AEE}$ ?

\begin{itemize}
	\item[ (1) ] Singletons são conjuntos de apenas um elemento entretanto talvez singletons formados por conjuntos desconhecidos não possam ser. $\{ R \}$
	\item[ (2) ] Se assumirmos $\{R\}$ um conjunto mesmo que $R$ esteja indefinido, então $R$ por $\texttt{AEE} $ é um conjunto também. Esse conjunto $R$ leva a alguma contradição ? : Vamos mostrar que leva na proposição a seguir considerando $ I_1 $ como essa possibilidade de escrever singletons com elementos indefinidos e desconhecidos.
\end{itemize}

%%%%%%%%%%%%%%%%%%%%%%%%%%%%%%%%%%%%%%%%%%%%%%%%%%%%%%%%%%%%%%%%%%%%%%%%%%%%%%%%%%%%%%%%%%%%%%%%%%%%%%%%%%%%
% Proposição [ Paradoxo de Russel por Singleton ] 
% \hspace{\baselineskip}

\textbf{Proposição [ Paradoxo de Russel por Singleton ]}

$$ \Phi( \texttt{AEE} \wedge \texttt{DETC} \wedge I_1 ) : R := \{ x \in R \iff ( x \in \{R\} \wedge x \notin x ) \} \implies \emptyset $$

\textbf{Dialética :} $\rightarrow$ 

\begin{itemize}
	\item[ ] ?` $R \in \mathcal{S}^*$ ?
	\item[ ] ?` $ R \in R \implies R \notin R $ ? $\rightarrow$
	\begin{itemize}
		\item[ ] ?` $ ( (A \wedge B) \wedge A ) \implies B $ ? $\rightarrow$ !` Tautologia !
	\end{itemize}
	\item[ ] ?` $ R \notin R \implies R \in R $ ?
	\item[ ] !` $ (\iff) \overset{op}{=} (=) $ !
	\item[ ] !` Reconhecimento de Absurdo !
\end{itemize}

\textbf{Comentário :} Esse paradoxo parece gerar um problema bem maior, parece que deve existir uma restrição na criação de singletons, não pode ser feito conjuntos com elementos desconhecidos dentro da teoria de conjuntos? Se existe possibilidade de contornar esse problema usando a teoria de conjuntos não tive a curiosidade de investigar, deixo isso como um trabalho posterior.

%%%%%%%%%%%%%%%%%%%%%%%%%%%%%%%%%%%%%%%%%%%%%%%%%%%%%%%%%%%%%%%%%%%%%%%%%%%%%%%%%%%%%%%%%%%%%%%%%%%%%%%%%%%%
% Coleção
\hspace{\baselineskip}

\textbf{[ Classe ]} Classe é uma informação definição da forma $ A := ( (x \in A) \iff \phi(x) )$ tal que $\phi$ é uma informação, $A$ é o nome da classe. Neste exemplo utilizamos a notação $A:=( (x \in A) \iff \phi(x) )$ para definir a classe.

\begin{itemize}
	\item[ (i) ] \textbf{[ Nome ]} $A$ é tanto o nome da informação quanto o nome da classe.
	\item[ (ii) ] \textbf{[ Efetividade ]} A efetividade é a informação $ (x \in A \iff \phi(x)) $ que é a fórmula do nome $A$. Se a efetividade da classe é aceita na razão da teoria dizemos que a classe é efetiva no sentido de que a efetividade é verdadeira ou assumido como verdadeira.
	\item[ (iii) ] \textbf{[ Indicadora ]} $ \phi(x) $ é a fórmula indicadora, ela somente pode ser utilizada para testar elementos se a classe for efetiva, em outras palavras se $ \phi(x) = 1 $ então utilizando inferência em $ (x \in A \iff \phi(x)) $ podemos concluir que $ x \in A $, se $\phi(x) = 0$ então usando inferência $ (x \notin A \iff \neg \phi(x)) $ ( Informacionalmente igual a efetividade ) temos $ x \notin A $.
\end{itemize}

\textbf{Comentário :} Observe que $A:=$ define uma informação, isso resolve o problema que tinhamos na teoria de conjuntos, se assumirmos que uma classe é uma informação então ela pode ser uma informação falsa, podemos escrever $A:= \dots$ sem problema algum na razão da teoria usando a liberdade de definição, entretanto não podemos colocar $A$ sozinho na razão da teoria, pois ele pode ser falso. Se $A=0$ então não podemos testar elementos da classe, para podermos testar se existem elementos $A$ precisa ser verdadeiro, isso não é automático, ou devemos demonstrar que ele é verdadeiro ou devemos adicionar $A$ como axioma.\\

\textbf{Comentário :} É importante lembrar que na fórmula $ x \in A $ temos $A$ dentro do contexto de classe com igualdade $\overset{cl}{=}$, portanto não podemos substituir por inferência de igualdade a igualdade do contexto de informação $\overset{inf}{=}$. Com efeito $ x \in ( x \in A \iff \phi(x) ) $ é uma expressão errada, pois viola a compartimentação dos contextos de igualdade. \\% É possível mostrar que se violarmos a compartimentalização temos um absurdo.

\textbf{Comentário :} Para $X,Y$ classes se $ X \overset{inf}{=} Y $ então se uma classe for efetiva a outra também é efetiva, isso é diferente de dizer que $ X \overset{cl}{=} Y $ que diz que ele tem os mesmos elementos. Essa compartimentalização das igualdades é extremamente importante para teoria. \\

\textbf{Comentário :} É razoável pensar que é possível criar teorias inconsistentes, assim definir coisas inconsistentes também, portanto não é surpresa que na teoria de conjuntos haja paradoxos, a razão é que na teoria de conjuntos a criação de um conjunto coincide com a sua efetivação na razão da teoria, na teoria aqui desenvolvida a criação de uma classe não implica na sua efetivação.\\

\textbf{Notação :} $A \in \mathcal{C}^*$ é a maneira de representar que $A$ é uma classe.

%%%%%%%%%%%%%%%%%%%%%%%%%%%%%%%%%%%%%%%%%%%%%%%%%%%%%%%%%%%%%%%%%%%%%%%%%%%%%%%%%%%%%%%%%%%%%%%%%%%%%%%%%%%%
% Operador de Continência
\hspace{\baselineskip}

\textbf{Operador [ Continência ]} $ \forall A \in \mathcal{C}^* : \forall B \in \mathcal{C}^* : (A \subseteq B) \iff ( x \in A \implies x \in B ) $\\

\textbf{Notação [ Estrutura de Classe ]} $ (A,\phi)_{class} $ significa dizer que $A$ é um conjunto e $\phi$ é um identificador efetivo para $A$.\\

\textbf{Comentário :} É interessante definir algo que represente coleção como informação, diferente da teoria de conjuntos que define conjunto como algo a parte do que é informação a teoria de classes aqui desenvolvida constrói classe com base na informação identificando ela mesma como uma informação.\\

\textbf{Igualdade [ de Classe ]} $( A \overset{cl}{=} B) \iff ( (A \subseteq B) \wedge (B \subseteq A) ) $

%%%%%%%%%%%%%%%%%%%%%%%%%%%%%%%%%%%%%%%%%%%%%%%%%%%%%%%%%%%%%%%%%%%%%%%%%%%%%%%%%%%%%%%%%%%%%%%%%%%%%%%%%%%%
% Classe verdadeira
\hspace{\baselineskip}

\textbf{Classe [ Verdade ]} $ 1 := ( x \in 1 \iff x \in 1 ) $\\

\textbf{Comentário :} A classe verdade é uma classe que precisa ser efetiva para que a relação de pertence faça sentido. É importante observar que ela efetiva no sentido de que a efetividade da classe precisa ser verdadeira, não que é possível testar elementos nela. Observe que não faz sentido testar elementos da classe verdade. \\

\textbf{Classe [ Mentira ]} $ 0 := ( x \in 0 \iff x \notin 0 ) $\\

\textbf{Comentário :} A classe mentira é uma classe que não pode ser efetiva sob o mesmo argumento.\\

\textbf{Comentário :} Vamos utilizar ?` $ \phi \implies 0 $ ? em vez de ?` $ \phi \implies \emptyset $ ? para dizer que queremos demonstrar por absurdo. Essa alteração na notação faz com que a mecânica do raciocínio se torne consistente, pois se aceitarmos a efetividade da classe vazio então $\phi \implies \emptyset$ sempre é verdadeiro.

%\textbf{Comentário :} Apesar dessas classes existirem ambas são classes inefetivas, a classe mentira não pode testar elementos, pois ela é falsa e $ x \in 0 \iff x \notin 0 $ não pode ser adicionado a razão da teoria, enquanto a classe 

\hspace{\baselineskip}

\textbf{[ Axioma da Tautologia Não Booleana ]} Seja $ \forall x : \phi(x) $ então chamamos $\phi(x)$ uma tautologia não booleana. O axioma da tautologia não booleana aceita como verdade que $\phi(x)=1$ se $\forall x : \phi(x)$.\\

\textbf{Notação :} $\texttt{ATNB} :=$ Axioma da Tautologia não booleana$^*$

%%%%%%%%%%%%%%%%%%%%%%%%%%%%%%%%%%%%%%%%%%%%%%%%%%%%%%%%%%%%%%%%%%%%%%%%%%%%%%%%%%%%%%%%%%%%%%%%%%%%%%%%%%%%
% Proposição Igualdade de Classes por Efetividade
\hspace{\baselineskip}

\textbf{Proposição [ Igualdade de Classes por Efetividade ]}

$$ \Phi( A \wedge ( A:= ( x \in A \iff \phi(x)) ) \wedge B \wedge ( B:=( x \in B \iff \phi(x) )  ) : A \overset{cl}{=} B $$

\textbf{Dialética :} Essa questão é simples, é importante notar que $\overset{cl}{=}$ somente pode ser feito porque $A$ e $B$ estão inclusos na razão pelo quantificador neutro.

%%%%%%%%%%%%%%%%%%%%%%%%%%%%%%%%%%%%%%%%%%%%%%%%%%%%%%%%%%%%%%%%%%%%%%%%%%%%%%%%%%%%%%%%%%%%%%%%%%%%%%%%%%%%
% Classe é uma Classe
\hspace{\baselineskip}

\textbf{Proposição [ Classe é uma Classe ]} $ \mathcal{C}^* \in \mathcal{C}^* $\\

\textbf{Demonstração :} $ \mathcal{C}^* := ( X \in \mathcal{C}^* \overset{def}{\iff}  ( \exists \phi \in \mathcal{I}^* : X:=( x \in X \iff \phi(x) ) ) ) $\\

\textbf{Comentário :} A classe não é somente uma classe como também é efetiva, pois em toda definição está implícito sua efetividade.

%%%%%%%%%%%%%%%%%%%%%%%%%%%%%%%%%%%%%%%%%%%%%%%%%%%%%%%%%%%%%%%%%%%%%%%%%%%%%%%%%%%%%%%%%%%%%%%%%%%%%%%%%%%%
% Axioma de Definição
\hspace{\baselineskip}

\textbf{Axioma [ de Definição ]} Toda definição é uma classe, deve ser separado o nome da definição da sua efetividade. Suponha uma definição de um conceito $X^*$ então ela é $ X^* := ( x \in X^* \iff \phi^*(x) ) $ com $\phi^*$ uma regra de reconhecimento que pode ser formal (construída por operadores e variáveis) ou informal (que requer habilidades da consciência humana para reconhecer), como o informal. Toda definição anterior a esse axioma é considerado efetivo entretanto toda definição posterior somente é efetivo caso explicitamente considerado e caso não forme uma contradição óbvia. \\

\textbf{Comentário :} Esse axioma reinterpreta a definição como definição de classe, é interessante notar que a inferência de reconhecimento de uma definição corresponde a efetividade da definição.\\

\textbf{Comentário :} É importante diferenciar o processo de nomear algo (liberdade de definição de nomes) que sempre é possível fazer do processo de definir e aceitar a inferência de reconhecimento na teoria. O processo de nomear continua livre mas o processo de definição requer que a definição funcione.\\

\textbf{Comentário :} Quando criamos uma definição, pode não ser evidente que a definição seja inconsistente, também pode ser que não possamos demonstrar que a efetividade da definição seja verdadeira pela própria natureza da definição, suponha que definamos um tipo de alien, o alien é um alien de varginha se e somente se ele for verde, olhos grandes e e tenha antenas, nesta definição não podemos com as informações da realidade provar que a efetividade é verdadeira ou falsa, e a princípio ela não leva a uma contradição óbvia, portanto para prosseguir com uma teoria sobre os aliens de varginha é conveniente axiomatizar sua efetividade, se em algum momento for provado que a efetividade é contraditória devemos refazer toda teoria que faz menção a elementos dessa definição.\\

\textbf{Comentário :} No mesmo exemplo alienígena, a efetividade dessa definição parece ser nem verdadeira nem falsa, se considerarmos que não haja confronto de nomes então a efetividade da definição de alien de varginha é uma informação indiferente a realidade, tanto faz se adicionarmos ela ou sua negação, como a realidade é inexistencial nem os elementos da realidade interferem na efetividade nem a efetividade interfere no que foi feito até então sobre a realidade. A efetividade de uma definição reabre a discussão sobre as informações indiferentes. \\

\textbf{Comentário :} Existe um problema prático com relação a informações indiferentes, é interessante considerar a indiferença com relação a tudo que poderia ser feito (inclui todas as razões válidas da teoria), entretanto essa qualidade pode requerer uma habilidade humanamente impossível, talvez a indiferença com relação ao que já foi feito (que é finito) é uma conceituação razoável da qualidade de indiferença.\\

\textbf{Proposição [ Obtenção de Classe Vazio ]}
$$ \Phi( \emptyset \wedge A ) : \Phi( A := ( x \in A \iff \phi(x) )  \wedge B:= (x \in B \iff (x \in A \wedge x \notin B)  )  : B \implies B \overset{cl}{=} \emptyset $$

\textbf{Dialética :} $\rightarrow$

\begin{itemize}
	\item[ ] ?` $ x \in B \implies x \in \emptyset $ ? $ \rightarrow $ !` Raciocínio Direto !
	\begin{itemize}
		\item[ ] ?` $ x \in B \implies  0 $  ?
	\end{itemize}
	\item[ ] ?` $ x \in \emptyset \implies x \in B $ ?	$\rightarrow$ !` Dialética de Raciocínio Direto e Absurdo !
\end{itemize}

%%%%%%%%%%%%%%%%%%%%%%%%%%%%%%%%%%%%%%%%%%%%%%%%%%%%%%%%%%%%%%%%%%%%%%%%%%%%%%%%%%%%%%%%%%%%%%%%%%%%%%%%%%%%
% Discussão sobre a teoria homeostática de informação indiferente
%\hspace{\baselineskip}
\hspace{\baselineskip}

\textbf{Dialética Socrática [ Informação Indiferente e Interpretação de Tempo ]}

\begin{itemize}
	\item[ (i) ] Quando construímos uma teoria qualquer de nomes, estamos projetando o conhecimento em um modelo que precisa ser construído, como construção é uma adição gradual de elementos e uso de elementos anteriores para construção de novas coisas a construção de uma teoria se comporta como se estivéssemos reconstruindo a realidade (ou qualquer coisa que estamos estudando).
	\item[ (ii) ] Na intuição que temos da realidade existem informações que consideramos indiferentes, por exemplo, informações que são aleatórias sobre eventos futuros são informações indiferentes (não interfere no passado, caso contrário não são aleatórias), informações entre um sistema fechado e outro são indiferentes entre si, pois não se comunicam (essa é a ideia da indiferença por propriedade de existência). Podemos tentar tangenciar o conceito de indiferença que parece estar presente no processo de definir utilizando a intuição que temos sobre o tempo na realidade.
	\item[ ] ?` O que seria uma informação indiferente na intuição de tempo da realidade ? $\rightarrow$
	\begin{itemize}
		\item[ (1) ] O tempo divide a realidade de forma trialética, temos (passado, presente, futuro). Uma informação indiferente no futuro pode ser uma informação que é genuinamente aleatória. Uma informação indiferente no passado ou no presente pode ser devido ao modelo de nosso sistema não estar conectado por interações com os elementos que a informação possui, ou seja, estão em sistemas que são fechados entre si ou não existenciais.
		\item[ ] ?` Toda informação indiferente futura é por conta de um processo aleatório ? $\rightarrow$
		\begin{itemize}
			\item[ (1) ] Se considerarmos o modelo que temos da realidade, pode ser que o modelo que temos da realidade não consiga determinar a verdade ou mentira de uma informação, entretanto isso não se deve a alguma falha da realidade, mas alguma incapacidade do próprio modelo em responder questões. Dentro deste modelo essa informação pode ser considerada indiferente, mas para o conceito de realidade talvez essa informação não seja indiferente.
			\item[ (2) ] Suponha que os aliens vão chegar a terra no ano de 2050, todas as informações envolvendo aliens para nossa realidade são indiferentes até o momento em que eles chegarem.
		\end{itemize}
		\item[ ] ?` Toda informação indiferente está no futuro ? $\rightarrow$
		\begin{itemize}
			\item[ (1) ] Essa pergunta parece ter sido respondida, entretanto se considerarmos que não temos consciência simultânea da realidade, se eu estou em uma sala fechada sem janelas e quero saber se tem alguém fora dela eu preciso sair desta sala ou perguntar se tem alguém de forma que alguém de fora possa ouvir, toda ação está no futuro então mesmo responder a uma questão de simultaneidade pode ser interpretada como um processo que somente pode ser respondido no futuro. A verdade é que toda informação para ser verificada necessita de uma ação que está no futuro.
		\end{itemize}
	\end{itemize}
	\item[ ] ?` O que seria uma informação indiferente para razão de uma teoria ? $\rightarrow$
	\begin{itemize}
		\item[ (1) ] Podemos definir informalmente uma informação indiferente como uma informação $\phi$ que ela mesma não gera nenhuma contradição com os elementos anteriores da razão e que $\neg \phi$ não gera nenhuma contradição da mesma forma. Se uma teoria é muito extensa ou o raciocínio para provar a contradição precisar ser muito elaborado, ou mesmo for necessário encontrar um contra exemplo específico, pode ser impraticável demonstrar que uma informação é indiferente, portanto faz sentido em vez de construir uma teoria axiomática lógica construir uma teoria homeostática que auxilie na identificação dos sinais que podem responder o problema.
	\end{itemize}
\end{itemize}

%%%%%%%%%%%%%%%%%%%%%%%%%%%%%%%%%%%%%%%%%%%%%%%%%%%%%%%%%%%%%%%%%%%%%%%%%%%%%%%%%%%%%%%%%%%%%%%%%%%%%%%%%%%%
% Inferência Homeostática
\hspace{\baselineskip}

\textbf{[ Inferência Homeostática ]} Uma inferência homeostática não é uma inferência determinística, a hipótese não implica na conclusão, uma inferência homeostática tem como propósito dizer uma informação provável e não uma informação certa, em vez de escrever $A \implies B$ será escrito $A \overset{ativa}{\rightarrow} B$ para representar informações homeostáticas em uma razão de teoria.\\

\textbf{Comentário :} Raciocínios utilizando inferências homeostáticas somente podem gerar outras inferências homeostáticas pelo raciocínio transitivo de inferência. É interessante que Inferências homeostáticas sejam justificadas, pois pode haver inferências e raciocínios de inferências homeostáticas contraditórias entre si.

%%%%%%%%%%%%%%%%%%%%%%%%%%%%%%%%%%%%%%%%%%%%%%%%%%%%%%%%%%%%%%%%%%%%%%%%%%%%%%%%%%%%%%%%%%%%%%%%%%%%%%%%%%%%
% 
\hspace{\baselineskip}

\textbf{Inferência Homeostática [ Identificador Orobólico ]}\\

Seja uma definição de classe $A := ( x \in A \iff \phi(x) )$, $\phi$ é orobólico se ele faz referência ao próprio $A$ ou ele contém alguma relação que relaciona algo com ele mesmo, como exemplo $x \in x$ é uma relação orobólica, $\phi(x) = (x \notin A )$ é um identificador orobólico.\\
\hrule 

\hspace{\baselineskip}

Se temos um identificador orobólico é bem provável que a classe seja falsa no contexto de informação e a demonstração da efetividade da classe provavelmente é imediata.

%%%%%%%%%%%%%%%%%%%%%%%%%%%%%%%%%%%%%%%%%%%%%%%%%%%%%%%%%%%%%%%%%%%%%%%%%%%%%%%%%%%%%%%%%%%%%%%%%%%%%%%%%%%%
% 
\hspace{\baselineskip}

\textbf{Inferência Homeostática [ Análise Existencial ]}\\

Suponha que queiramos demonstrar que uma definição de classe é indiferente, quais elementos da razão da teoria devem ser investigados.\\
\hrule 

\hspace{\baselineskip}

Todos as informações que possuem lista de variáveis livres que são referenciados na definição de classes são os elementos da razão da teoria que são interessantes serem investigados.

%%%%%%%%%%%%%%%%%%%%%%%%%%%%%%%%%%%%%%%%%%%%%%%%%%%%%%%%%%%%%%%%%%%%%%%%%%%%%%%%%%%%%%%%%%%%%%%%%%%%%%%%%%%%
% 
\hspace{\baselineskip}

\textbf{Inferência Homeostática [ Identificador Não Orobólico ]}\\

Se temos um identificador não orobólico para uma classe.\\
\hrule 

\hspace{\baselineskip}

Provavelmente essa classe tem sua efetividade indiferente.

%%%%%%%%%%%%%%%%%%%%%%%%%%%%%%%%%%%%%%%%%%%%%%%%%%%%%%%%%%%%%%%%%%%%%%%%%%%%%%%%%%%%%%%%%%%%%%%%%%%%%%%%%%%%
% 
\hspace{\baselineskip}

\textbf{Inferência Homeostática [ Axioma Esquema de Especificação ]}\\

Se temos um identificador que parte de uma classe já bem estabelecida e uma conjunção com outra informação. \\
\hrule 

\hspace{\baselineskip}

Provavelmente essa classe tem sua efetividade indiferente e caso haja alguma contradição ela deve ser igual em classe a classe vazio.

%%%%%%%%%%%%%%%%%%%%%%%%%%%%%%%%%%%%%%%%%%%%%%%%%%%%%%%%%%%%%%%%%%%%%%%%%%%%%%%%%%%%%%%%%%%%%%%%%%%%%%%%%%%%
% 
\hspace{\baselineskip}

\textbf{Inferência Homeostática [ Identificação Posterior de Definição Falsa ]}\\

Se por ironia do destino identificarmos que uma definição estabelecida anteriormente é uma definição falsa. \\
\hrule 

\hspace{\baselineskip}

Provavelmente a teoria é segura antes de sua definição, todas as conclusões posteriores podem estar erradas e devem ser verificadas.

%%%%%%%%%%%%%%%%%%%%%%%%%%%%%%%%%%%%%%%%%%%%%%%%%%%%%%%%%%%%%%%%%%%%%%%%%%%%%%%%%%%%%%%%%%%%%%%%%%%%%%%%%%%%
% 
\hspace{\baselineskip}

\textbf{Comentário :} Com essa mini teoria homeostática é possível com alguma segurança desenvolver teorias axiomáticas mesmo que a resposta sobre se uma informação é indiferente ou não, ou um axioma pode ou não ser adicionado, não for possível responder.

\hspace{\baselineskip}

\textbf{Inferência Homeostática [ Orobólico e Transcendental ]}\\

Se uma definição tem como objetivo modelar um conceito transcendental. \\
\hrule 

\hspace{\baselineskip}

Provavelmente é necessário incluir um identificador orobólico\\

\textbf{Comentário :} O exemplo mais clássico é a definição de números naturais que modela o infinito enumerável e pode ser definido como $ \mathbb{N} := ( (x\in \mathbb{N} ) \iff (x = 0 \wedge \exists y \in \mathbb{N} : x = y+1) ) $, o conceito de infinito é transcendental segundo investigações informais anteriores à esta seção.

%%%%%%%%%%%%%%%%%%%%%%%%%%%%%%%%%%%%%%%%%%%%%%%%%%%%%%%%%%%%%%%%%%%%%%%%%%%%%%%%%%%%%%%%%%%%%%%%%%%%%%%%%%%%
% Classe Regular
\hspace{\baselineskip}

\textbf{Classe [ Regular ]} $ A \in \mathcal{CR}^* \overset{def}{\iff} \exists \mathcal{U} \in \mathcal{C}^* \exists \phi \in \mathcal{I}^* : A := ( x \in A \iff \Phi( x \in \mathcal{U} ) : \phi(x) ) $\\

\textbf{Comentário :} A classe regular possui uma classe $\mathcal{U}$ onde podemos fazer o teste de pertencimento usando a indicadora. Observe que $ (x \notin A ) = \neg (x \in A) =  \neg ( \Phi(x \in \mathcal{U}) : \phi(x) ) =  ( \Phi( x \in \mathcal{U}) : \neg \phi(x) )$.\\

\textbf{Comentário :} O propósito da definição de classe regular é interpretar corretamente a negação da relação pertence. Para uma classe geral, $x \notin A$ não é bem definido, se for aceito a efetividade da classe $\Xi$ então coisas desconhecidas podem fazer $x \notin A$ ser verdadeiro, não queremos testar que aliens não pertencem a um conjunto $A$, assim queremos um universo que possamos controlar o que podemos testar como elementos.\\

\textbf{Comentário :} Na teoria de conjuntos o problema da negação é resolvida pelo axioma esquema de especificação, na definição de classe regular escrevemos algo parecido utilizando as notações de fórmulas universais. Vale a pena comparar e perceber o sentido da criação do quantificador neutro.\\

\textbf{Comentário :} Uma informação é bem definida se sua negação é bem definida, sem a negação como podemos garantir o meio excluído? Portanto uma classe regular não apenas define $ x \notin A$ como define a própria relação $x \in A$, é portanto essencial para o desenvolvimento de uma teoria lógica utilizar classes regulares, pois classes gerais podem não definir $x \in A$ como uma informação. Podemos nos atrever a dizer que a teoria de classes somente é uma teoria axiomática se for utilizado classes regulares.\\

\textbf{Comentário :} As classes $\emptyset$ e $ \Xi $ dentro da teoria de classes regulares tomam forma como $\emptyset := ( x \in \emptyset \iff \Phi(x\in \mathcal{U}) : 0 )$ e $ \Xi := ( x \in \Xi \iff \Phi( x \in \mathcal{U} ) : 1 )$, temos que $ \Xi \overset{cl}{=} \mathcal{U} $ se considerarmos a efetividade de ambos.

%%%%%%%%%%%%%%%%%%%%%%%%%%%%%%%%%%%%%%%%%%%%%%%%%%%%%%%%%%%%%%%%%%%%%%%%%%%%%%%%%%%%%%%%%%%%%%%%%%%%%%%%%%%%
% Proposição [ Preservação de Regularidade por Especificação ]
\hspace{\baselineskip}

\textbf{Proposição [ Preservação de Regularidade por Especificação ]}

$$ \Phi( A \wedge B \wedge A \in \mathcal{CR}^* \wedge B:=( x \in B \iff ( x \in A \wedge \phi(x) ) ) ) : B \in \mathcal{CR}^*  $$

\textbf{Dialética :} !` Inferência de Substituição em Informações !\\

\textbf{Comentário :} Observe que o conjunto universo de $B$ é o mesmo que o universo de $A$.\\

\textbf{Comentário :} Existem várias formas de regularizar uma classe, por exemplo poderíamos regularizar $B$ como $ B':= ( x \in B' \iff \Phi( x \in A ) : \phi(x) ) $, isso daria uma relação diferente, a definição de $B'$ não seria a mesma que $B$ entretanto eles teriam os mesmos elementos caso fossem ambos efetivos. Portanto temos um exemplo em que a igualdade informacional difere da igualdade de classes: $B' \overset{inf}{\neq} B$ mas $ B' \overset{cl}{=} B $.\\

\textbf{Comentário :} Regularidade não implica em efetividade, basta observar o contra exemplo $ A:=(x \in A \iff \Phi(x \in \mathcal{U}): x \notin A ) $, $A\overset{inf}{=}0$ é uma tautologia e portanto essa classe apesar de regular não é efetiva.













%Com base no que foi investigado aqui o processo de definir algo consistente está associado a identificação de informação indiferente

%%%%%%%%%%%%%%%%%%%%%%%%%%%%%%%%%%%%%%%%%%%%%%%%%%%%%%%%%%%%%%%%%%%%%%%%%%%%%%%%%%%%%%%%%%%%%%%%%%%%%%%%%%%%
% Coleção Formal de Conjuntos
\hspace{\baselineskip}

\textbf{[ Família Formal ]} $ (f(x), Y)_{ff} \overset{def}{\iff} (\forall x : f(x) \overset{part}{\rightarrow} Y )$\\

\textbf{Notação :} Para simplificar e identificar mais facilmente em uma família usamos a notação $ f_x $ com mesmo significado que $f(x)$

%%%%%%%%%%%%%%%%%%%%%%%%%%%%%%%%%%%%%%%%%%%%%%%%%%%%%%%%%%%%%%%%%%%%%%%%%%%%%%%%%%%%%%%%%%%%%%%%%%%%%%%%%%%%
% Uniões Formais
\hspace{\baselineskip}

\textbf{Classe [ Uniões Formais ]} $ \Phi( (X_{\xi} , \mathcal{C}^*)_{ff} \wedge A \in \mathcal{C}^* ) : $

\begin{itemize}
	\item[ (i) ] \textbf{[ União Formal por Operador Pertence ]} 
	\item[ ]$ \left( \bigcup\limits_{ \xi \in A } X_{\xi}  , \phi(a) :=( \exists \xi \in A : a \in X_{\xi} ) \right)_{class}  $
	\item[ (ii) ] \textbf{[ União Formal por Operador Continência ]} 
	\item[ ] $ \left( \bigcup\limits_{ B \subseteq A } X_B  , \phi(a) := ( \exists B \subseteq A : a \in  X_B ) \right)_{class}  $ 
\end{itemize}

\textbf{Comentário :} É uma tarefa simples provar que a regularidade é preservada na união caso a família formal seja regular e efetiva para um mesmo conjunto universo. Não podemos entretanto afirmar que a efetividade seja passada para união, convém porém que isso não é algo muito difícil de assumir, pois a união é essencialmente juntar todas as informações de efetividade em uma efetividade que considera todas as outras.

%%%%%%%%%%%%%%%%%%%%%%%%%%%%%%%%%%%%%%%%%%%%%%%%%%%%%%%%%%%%%%%%%%%%%%%%%%%%%%%%%%%%%%%%%%%%%%%%%%%%%%%%%%%%
% Inferência de Quantificador Nulo
\hspace{\baselineskip}

\textbf{Inferência [ Quantificador Neutro ]} $ ( \Phi(\psi) : \phi ) \overset{1}{\iff} ( \forall (\psi) : \phi ) $\\

\textbf{Notação :} $ \overset{1}{\iff} $ significa que a inferência somente pode ser utilizada quando se considera ambos os lados verdadeiros, precisamos desta inferência para o quantificador nulo por conta do seu comportamento na negação. Essa inferência é diferente da igualdade de informações, $ (\overset{1}{\iff} ) \ne (\overset{inf}{=}) $, isso porque a igualdade informacional deve considerar a negação. Também não devemos utilizar a identidade $(A \implies B) = (\neg B \implies \neg A)$ que também é falsa se considerar todas as possibilidades para o quantificador nulo.\\

\textbf{Comentário :} Essa inferência é assumida como verdade, como uma propriedade inerente do quantificador neutro. Isso irá facilitar os raciocínios utilizando o quantificador neutro. É importante perceber que o quantificador neutro é topologicamente distinto dos quantificadores existencial e universal, essa afirmação é consistente com a trialética em que um dos componentes não está sob as mesmas regras do que os outros dois, no caso do quantificador neutro, ele não sofre influência da negação enquanto os quantificadores universal e existencial são complementares entre si.\\

\textbf{Comentário [ Quantificador Neutro Escondido ]} Quando utilizamos $ \forall x \in S : \phi(x) $ estamos utilizando um quantificador neutro escondido, a informação equivale a $ \Phi( x \in S ) : \forall x : \phi(x) $. Problemas podem surgir na demonstração por absurdo, por exemplo, suponha que $ (\forall x \in S : \phi(x) ) \implies 0 $ seja demonstrado, uma pessoa poderia concluir que $ \exists x \in S : \neg \phi(x) $ usando o raciocínio inverso ($1 \implies \dots$), o que ele deveria concluir na verdade seria $ x \notin S \vee ( \exists x \in S : \neg \phi(x) ) $, pois a falsidade de  $ x \in S \implies \forall x : \phi(x) $ é $ x \notin S \vee ( \exists x \in S : \neg \phi(x) ) $. Por outro lado se em $ \Phi(x \in S) : \forall x : \phi(x) $ fosse demonstrado que $ ( \forall x : \phi(x) ) \iff 0 $, ou seja, que $ ( \forall x : \phi(x) ) = 0 $ e portanto não apenas para um caso particular, mas para todos os casos de $x \in S$ tivéssemos $\neg \phi(x)$ então teríamos demonstrado que $ \Phi(x \in S) : \exists x : \neg \phi(x) $, entretanto também teríamos demonstrado que $ \Phi(x \in S) : \forall x : \neg \phi(x) $ que é mais forte. Recomendo que seja evitado o máximo possível conclusões baseado em raciocínios inversos quando houver quantificadores neutros escondidos.

%%%%%%%%%%%%%%%%%%%%%%%%%%%%%%%%%%%%%%%%%%%%%%%%%%%%%%%%%%%%%%%%%%%%%%%%%%%%%%%%%%%%%%%%%%%%%%%%%%%%%%%%%%%%
% Conjunto Complementar
\hspace{\baselineskip}

\textbf{Classe [ Complementar ]} $\Phi( (A,\phi)_{class} ) : A^c := ( x \in A^c \iff x \notin A ) $ \\

\textbf{Comentário :} Convém dizer que a classe complementar somente faz sentido para classes regulares.\\

\textbf{Classe [ Interseção Formal ]} $ \bigcap\limits_{\xi \in H} X_{\xi} := ( \bigcup\limits_{ \xi \in H } X_{\xi}^c  )^c  $ e $ \bigcap\limits_{\xi \subseteq H} X_{\xi} := ( \bigcup\limits_{ \xi \subseteq H } X_{\xi}^c  )^c $\\

\textbf{Comentário :} É uma tarefa simples demonstrar que $ \bigcap\limits_{\xi \in H} X_{\xi} := ( x \in \bigcap\limits_{\xi \in H} X_{\xi} \iff \forall \xi \in H : x \in X_{\xi} ) $.

%%%%%%%%%%%%%%%%%%%%%%%%%%%%%%%%%%%%%%%%%%%%%%%%%%%%%%%%%%%%%%%%%%%%%%%%%%%%%%%%%%%%%%%%%%%%%%%%%%%%%%%%%%%%
% Classe Singleton
\hspace{\baselineskip}

\textbf{Classe [ Singleton ]} $ \{x\} := ( a \in \{x\} \overset{1}{\iff} a = x ) $\\

\textbf{Comentário :} Poderíamos definir o singleton como uma classe regular $ \{x\} := ( a \in \{x\} \iff \Phi(x \in \mathcal{U}) : a = x ) $, convém porém conhecer ou estabelecer a classe universo antes de definir o singleton como classe regular, caso contrário a negação da relação não fará sentido. É esperado que o leitor tenha capacidade para definir singleton como classe regular.

%%%%%%%%%%%%%%%%%%%%%%%%%%%%%%%%%%%%%%%%%%%%%%%%%%%%%%%%%%%%%%%%%%%%%%%%%%%%%%%%%%%%%%%%%%%%%%%%%%%%%%%%%%%%
% Classe Par
\hspace{\baselineskip}

\textbf{Classe [ Par ]} $ \{x,y\} := ( a \in \{ x,y \} \overset{1}{\iff}  (a = x \vee a = y) ) $\\

\textbf{Comentário :} Toda classe possui múltiplas definições se considerar as diferentes possibilidades da classe universo, o vazio também possui múltiplas definições e o leitor não pode considerar tais definições como idênticas entre si, pois considerando uma classe como informação as informações podem não são equivalentes.\\

\textbf{Comentário :} Na teoria de conjuntos de Zermello-Frankel a existência de pares é um axioma.\\

\textbf{Comentário :} É importante perceber que na notação $ \{x\} $ ou $ \{x,y\} $, $x$ e $y$ estão no contexto de elementos e a inferência de substituição deve ser feita somente com a igualdade $ \overset{el}{=} $ do contexto de elementos, não podemos substituir apenas com $ \overset{cl}{=} $. Suponha que $ x $ ou $y$ sejam conjuntos, como definir sua igualdade no contexto de elementos?\\

\textbf{Igualdade [ no Contexto de Elementos para Classes ]} 

$$ \Phi( A \in \mathcal{C}^* \wedge B \in \mathcal{C}^* ) : A \overset{el}{=} B \iff ( A \overset{class}{=} B \wedge A \overset{inf}{=} B ) $$

\textbf{Comentário :} O conceito de igualdade possui uma íntima relação com o conceito de substituição, algo é considerado igual ou equivalente se podemos substituir e ao substituir não provocamos nenhuma alteração no sistema. Existem conceitos de igualdade mais fortes, por exemplo, igualdade de expressões considerando caracteres idênticos. No contexto de classes, para interpretar classe como um elemento introduzimos esta definição, não basta a igualdade de elementos, é preciso que sob o ponto de vista de informação eles também sejam iguais.

%%%%%%%%%%%%%%%%%%%%%%%%%%%%%%%%%%%%%%%%%%%%%%%%%%%%%%%%%%%%%%%%%%%%%%%%%%%%%%%%%%%%%%%%%%%%%%%%%%%%%%%%%%%%
% Classe Par Ordenado
\hspace{\baselineskip}

\textbf{Classe [ Par Ordenado ]} $ (x,y) := ( a \in (x,y) \iff ( a \in \{ \{x\} , \{ x,y \}  \} ) ) $\\

\textbf{Comentário :} O objetivo da classe par ordenado é modelar uma lista, já tínhamos definido o conceito de listas de variáveis, a classe par ordenado preserva importantes propriedades presentes no conceito de lista de variáveis.\\

\textbf{Comentário :} A classe par ordenado permite a construção de várias classes que são muito importantes para matemática.

%%%%%%%%%%%%%%%%%%%%%%%%%%%%%%%%%%%%%%%%%%%%%%%%%%%%%%%%%%%%%%%%%%%%%%%%%%%%%%%%%%%%%%%%%%%%%%%%%%%%%%%%%%%%
% Proposição [ Assimetria de Par Ordenado ]
\hspace{\baselineskip}

\textbf{Proposição [ Assimetria de Par Ordenado ]} $$ \Phi( (x,y) \wedge (y,x) \wedge x \overset{el}{\ne} y ) : ( x ,y ) \overset{class}{=} (y,x) $$

\textbf{Dialética :} $\rightarrow$ !` Absurdo !

\begin{itemize}
	\item[ ] ?` $ \Phi( x \overset{el}{\ne} y \wedge \{x\} \wedge \{x,y\} ) : \{x\} \overset{class}{\ne} \{x,y\} $  ?
	\begin{itemize}
		\item[ ] ?` $ \{x\} \overset{class}{=} \{x,y\} \implies 0 $ ? $\rightarrow$ ?` $ y \in \{x\} \implies 0 $ ?
	\end{itemize}
	\item[ ] ?` $ (x,y) \overset{class}{=} (y,x) \implies 0 $ ? $\rightarrow$ !` Use $\{x\}$ !
\end{itemize}

\textbf{Comentário [ Igualdade Ponto a Ponto ]} Seria muito tedioso demonstrar ?` $ (x,y) \overset{cl}{=} (a,b) \iff (x = a \wedge y = b) $ ?, basta saber que é verdade que a definição de par ordenado preserva esta propriedade.

%%%%%%%%%%%%%%%%%%%%%%%%%%%%%%%%%%%%%%%%%%%%%%%%%%%%%%%%%%%%%%%%%%%%%%%%%%%%%%%%%%%%%%%%%%%%%%%%%%%%%%%%%%%%
% Produto Cartesiano
\hspace{\baselineskip}

\textbf{Classe [ Produto Cartesiano ]} $ X \times Y := ( (x,y) \in X \times Y \iff ( x\in X \wedge y \in Y ) ) $\\

\textbf{Comentário :} Podemos definir $ X \times Y \times Z $ tripla ordenada definindo $ (x,y,z) $ como um conjunto similar ao par ordenado que preserva a assimetria, de forma que permutações gerem triplas diferentes. Vou me abster de definir e explorar suas propriedades visto que a intuição sobre listas ordenadas finitas são requeridas para compreensão de textos e também porque seria muito trabalhoso explorar conclusões óbvias sobre listas ordenadas. Basta saber que alguns livros definem uma tripla ordenada como $ (x,y,z) := ( (x,y),z ) $, ele preserva a propriedade de assimetria e também da propriedade de igualdade ponto a ponto.\\

\textbf{Classe [ Relação ]} $ R \in \texttt{Rel.}^*(X,Y) \overset{def}{\iff} R \subseteq X \times Y $\\

\textbf{Classe [ de Funções ]} $ Y^X := ( f \in Y^X \iff ( \forall x \in X: \exists ! y \in Y : (x,y) \in f ) ) $\\

\textbf{Comentário :} $f$ é chamado de função de $X$ para $Y$, normalmente uma função é representada pela notação $ f : X \rightarrow Y $. O quantificador $\exists !$ significa que existe e é único. $X$ é chamado de domínio de $f$ e $Y$ é chamado de contradomínio. Como em $(x,y) \in f$ temos $y$ único representamos $y = f(x)$ o valor da função em $x$, desenvolvendo a notação temos que $f(X)$ denota todos os valores $f(x)$ para $x \in X$ chamado de imagem.  \\

\textbf{Comentário :} Faz sentido que a classe de funções seja uma classe regular com o conjunto universo associado ao conjunto universo de $ X \times Y $. A verdade é que definir uma função dessa forma torna ela regular de várias formas, ela pode ser regular por conta de ser uma função dentro de $ X \times Y $ ou ela pode ser regular por herdar a regularidade de $X$ e de $Y$. Em uma função sempre sabemos onde procurar seus elementos. \\

\textbf{Comentário :} Seja $ A:= ( x \in A \iff \phi(x) ) $ uma definição de classe, podemos definir $ \phi $ como uma função tal que $ \phi := ( \Phi(x \in X \wedge y \in \{0,1\}) : (x,y) \in \phi \iff \psi(x,y) ) $, como $\phi$ é uma classe regular tanto $x$ quanto seu valor $\phi(x)$ estão dentro de alguma outra classe e portanto podemos considerar $A$ como uma classe regular se $\phi$ for uma função.\\

\textbf{Comentário :} Alguém poderia definir uma função como $ f := ( \Phi( x \in X \wedge y \in Y ) : (x,y) \in f \iff \phi_f(x,y) ) $ entretanto essa construção nem sempre é uma função, é preciso que todo $x \in X$ tenha um valor $y$ associado $y =f(x)$ e se existir tal valor ele deve ser único. A seguir a dialética demonstrativa para dizer que uma relação é uma função.\\

\textbf{Dialética Demonstrativa [ Boa Definição de Função ]}\\

?` $ \Phi( f :=( \Phi(x \in X \wedge y \in Y) : (x,y) \in f \iff \phi_f(x,y)  ) ) : f \in Y^X $ ? $\rightarrow$

\begin{itemize}
	\item[ (1) ] \textbf{[ Teste de Domínio ]} ?` $ \forall x \in X :\exists y \in Y: \phi_f(x,y) $ ?
	\item[ (2) ] \textbf{[ Unicidade de Imagem ]} ?` $ ( y = f(x) \wedge z = f(x) )\implies y=z $ ?
\end{itemize}

%%%%%%%%%%%%%%%%%%%%%%%%%%%%%%%%%%%%%%%%%%%%%%%%%%%%%%%%%%%%%%%%%%%%%%%%%%%%%%%%%%%%%%%%%%%%%%%%%%%%%%%%%%%%
% Função Injetiva
%\hspace{\baselineskip}

\textbf{Função [ Injetiva ]} $ \Phi^*( f \in Y^X ) : f_{\rightarrow} X \iff ( f(x) = f(y) \implies x = y ) $\\

\textbf{Comentário :} $\Phi^*$ significa que estamos omitindo algumas hipóteses que são razoavelmente evidentes. Neste caso que $x$ e $y$ pertencem a $X$ ou que $X$ e $Y$ são classes. Omitir informações evidentes permite a escrita de fórmulas mais elegantes, entretanto é preciso separar o que é evidente por conta da consciência alfa do que é evidente para um leigo.\\

\textbf{Função [ Sobrejetiva ]} $ \Phi^*( f \in Y^X ) : f_{\leftarrow} X \iff f(X) = Y $\\

\textbf{Comentário :} É fácil obter uma função sobrejetiva a partir de uma função qualquer, basta mudar o contradomínio pela imagem da função anterior.\\

\textbf{Função [ Bijetiva ]} $ \Phi^*( f \in Y^X ) : f_{\leftrightarrow}X \iff ( f_{\rightarrow} \wedge f_{\leftarrow}  )X $\\

\textbf{Comentário :} Assim como é fácil obter uma função sobrejetiva a partir de uma função qualquer, é fácil obter uma função bijetiva a partir de uma função injetiva.\\

\textbf{Comentário :} Funções bijetivas serão especialmente importantes para comparar infinitos com infinitos, segundo a teoria de conjuntos existem infinitos que são 'maiores' que outros, e infinitos que são equivalentes, os infinitos equivalentes possuem bijeções entre si enquanto infinitos 'diferentes' não podem ter tais funções.\\

\textbf{Função [ Composta ]} $ \Phi( f \in Y^X \wedge g \in Z^Y ) : g \circ f := ( (x,z) \in (g \circ f) \iff z = g( f(x) ) ) $\\

\textbf{Comentário :} É fácil ver que a composta de duas injetivas continua injetiva e a composta de duas bijetivas continua bijetiva. Em contextos menos formais vamos assumir como propriedade elementar das funções visto que a demonstração é tediosa.

\subsection{O Infinito e Algoritmos}

\hspace{\baselineskip}

Existem pelo menos dois conjuntos (classes) de infinito que são fundamentais para a compreensão da realidade, a classe dos números naturais e a classe dos números reais, estudar profundamente tais conjuntos é papel da Teoria dos Números e da Análise Real que são respectivamente os campos de estudo da matemática para lidar com tais conjuntos. Sob o ponto de vista da lógica vamos nos ater a compreender o conceito de infinito que é um conceito transcendental utilizando estes dois conjuntos e também criar novas classes utilizando recursos obtidos pela aceitação destes conjuntos.

%%%%%%%%%%%%%%%%%%%%%%%%%%%%%%%%%%%%%%%%%%%%%%%%%%%%%%%%%%%%%%%%%%%%%%%%%%%%%%%%%%%%%%%%%%%%%%%%%%%%%%%%%%%%
% O conjunto dos caracteres repetidos
\hspace{\baselineskip}

\textbf{[ Caracteres Repetidos ]} $R(\texttt{x})$ é a coleção de nomes formados por caracteres $\texttt{x}$ sem espaços, ele é uma classe que possui uma indicadora informal, sabemos reconhecer que um nome contém apenas os mesmos caracteres repetidos como $\texttt{xxx}$ ou $\texttt{xxxxx}$, também sabemos como reconhecer nomes que não estão nesta coleção como $\texttt{xyx}$ ou $\texttt{asdrtvx}$.\\

\textbf{Comentário :} Se utilizarmos esses caracteres para nomear quantidade de coisas é fácil perceber que ele representa os números naturais na base 1 (os números naturais da base 10 possui dez caracteres de 0 a 9, a base 1 possui apenas um caractere).\\

\textbf{Notação [ Fórmulas Formais ]} Vamos designar $ \texttt{FF}^* $ como a classe de operadores formais, um operador formal não é exatamente uma função, ele é definido através de alguma fórmula, por exemplo $ f(x) = x*y $ que permite a substituição $f(y) = y*y$ ou $f(z) = z*y$.\\

\textbf{Classe [ Caracteres Repetidos ]} $ \Phi( ( S(\xi) = \xi \texttt{x} ) \overset{part}{\rightarrow} \texttt{FF}^* ) : $

$$R(\texttt{x}) := ( n \in R(\texttt{x}) \iff n = \texttt{x} \vee \exists m \in R(\texttt{x}) : n = S(m) ) )$$

\textbf{Comentário :} Essa é uma das definições de números naturais mais próxima da ideia de representar números por símbolos, entretanto perdemos a capacidade de comparar números grandes muito rapidamente, é bem provável que acima de 20 caracteres tenhamos dificuldade de saber diferenciar nomes próximos. A função $S$ é a função sucessor que no caso dos números naturais na base decimal é definido como $S(x) = x + 1$.\\

\textbf{Comentário :} Existem várias propriedades que sabemos reconhecer que podem ser complicados de demonstrar formalmente em uma teoria axiomática, por exemplo, a função sucessora gera sempre um novo nome, diferente dos nomes anteriores, a função sucessora nunca repete saídas para diferentes entradas (chamamos isso de função injetiva), a classe de caracteres repetidos é infinita, podemos perceber da intuição de infinito como processo sem fim, a função sucessora estabelece uma ordem, se pegarmos qualquer subclasse (classe contida) existe um elemento de menor ordem.

%%%%%%%%%%%%%%%%%%%%%%%%%%%%%%%%%%%%%%%%%%%%%%%%%%%%%%%%%%%%%%%%%%%%%%%%%%%%%%%%%%%%%%%%%%%%%%%%%%%%%%%%%%%%
% Classe dos Naturais
\hspace{\baselineskip}

\textbf{Naturais [ Base Decimal ]} $\Phi^*$

$$ \Phi( S(0) = 1 , S(1) = 2, S(2) = 3, S(3) = 4, S(4) = 5, S(5) = 6, S(6) = 7, S(7) = 8, S(8) = 9  ) : $$

$$ \Phi( 0 \in \mathbb{N} \wedge x \in \{0,1,2,3,4,5,6,7,8,9\} \wedge \xi \in \mathbb{N} ) : $$

$$ \Phi( S(\xi x) = \xi S(x) [ x \ne 9 ] + S(\xi) 0 [ x = 9 ] ) : $$

$$ \mathbb{N} := ( y \in \mathbb{N} \iff  (y = 0 \vee \exists z \in \mathbb{N} : y = S(z) )  ) $$

\textbf{Comentário :} Na função formal sucessor $ S(\xi x) = \xi S(x) [ x \ne 9 ] + S(\xi) 0 [ x = 9 ] $ significa que se $x\ne 9$ então $ S(\xi x) = \xi S(x) $, se $x =9$ então $S(\xi x) = S(\xi) 0$ com $0$ um algarismo.\\

\textbf{Comentário [ Infinitos Nomes ]} Aceitar a existência dos números naturais não é difícil, existem conceitos que não estão presentes na realidade mas estão presentes na consciência humana, como a vontade e o infinito, a vontade é algo que não consegue ser interpretado através dos fenômenos e o infinito é a percepção de que existem processos que sempre podem continuar mais um pouco se não houver nenhum evento que os impeça. Uma vez aceito os números naturais conseguimos ter uma lista de infinitos nomes apenas criando uma sequência de caracteres, por exemplo existem infinitos nomes da forma $(X_1, X_2, X_3, \dots )$.

%%%%%%%%%%%%%%%%%%%%%%%%%%%%%%%%%%%%%%%%%%%%%%%%%%%%%%%%%%%%%%%%%%%%%%%%%%%%%%%%%%%%%%%%%%%%%%%%%%%%%%%%%%%%
% Cardinalidade
\hspace{\baselineskip}

\textbf{Relação [ Cardinalidade ]} $ \Phi( A \in \mathcal{C}^* \wedge B \in \mathcal{C}^* ) : A \overset{card}{=} B \iff ( \exists f: f_{\leftrightarrow} \in B^A )  $\\

\textbf{Comentário :} Podemos facilmente concluir que a cardinalidade é transitiva, por conta da preservação de bijeções em funções compostas. A composta de duas bijeções é também uma bijeção, portanto se $A \overset{card}{=} B$ e $ B \overset{card}{=}C $ isso implica que $ A \overset{card}{=} C $.\\

\textbf{Comentário :} A cardinalidade de conjuntos finitos é quantidade de elementos.\\

\textbf{Relação [ Ordem de Cardinalidade ]} $ A \overset{card}{>} B \iff ( \exists f_{\rightarrow} A^B \wedge A \overset{card}{\ne} B ) $ \\

\textbf{Comentário :} Uma das coisas mais interessantes da teoria de conjuntos infinitos é demonstrar que o conjunto das partes de um infinito é um infinito de cardinalidade maior e que os reais possuem cardinalidade maior que os números naturais.

%%%%%%%%%%%%%%%%%%%%%%%%%%%%%%%%%%%%%%%%%%%%%%%%%%%%%%%%%%%%%%%%%%%%%%%%%%%%%%%%%%%%%%%%%%%%%%%%%%%%%%%%%%%%
% Função Contagem
\hspace{\baselineskip}

\textbf{Função [ Contagem ]} A função contagem é uma função que para cada elemento de $X$ classe qualquer associamos um único elemento de $\mathbb{N}$ tal que seja sobrejetivo.

$$ \Phi( f \in \mathbb{N}^X ) : f \in \texttt{Cont}( \mathbb{N}^X ) \iff f_{\leftrightarrow} \mathbb{N}^X  $$

%%%%%%%%%%%%%%%%%%%%%%%%%%%%%%%%%%%%%%%%%%%%%%%%%%%%%%%%%%%%%%%%%%%%%%%%%%%%%%%%%%%%%%%%%%%%%%%%%%%%%%%%%%%%
% 
\textbf{Comentário :} Estabelecer uma contagem é um processo intuitivo para um conjunto, com a contagem podemos comparar conjuntos infinitos com o conjunto dos naturais. Chamamos de conjunto enumerável ou contável todo o conjunto que tem mesma cardinalidade que os naturais. Em geral até mesmo os matemáticos não exigem demonstrações mais rigorosas quando o processo de contagem dos elementos é evidente.\\

\textbf{Comentário [ Exemplos de Contáveis ]} Perceba que os números pares são contáveis apenas criando a função contagem apropriada, os números ímpares são contáveis, se retirarmos conjuntos finitos de números dos naturais os naturais ainda continuam contáveis, os números inteiros (que possuem inteiros negativos) são contáveis, os produtos cartesianos finitos $ \mathbb{N} \times \mathbb{N} $ ou $ Q \times Q $ para um $Q$ contável continua contável, qualquer subconjunto infinito de números é contável, se representarmos os racionais como pares ordenados de inteiros percebemos que ele é quase o produto cartesiano de inteiros, o que faz com que os racionais também sejam contáveis. Toda sequência que pode ser escrita como um processo de etapas através de um algoritmo com finitas instruções, se o processo for infinito então a sequência é contável.\\

\textbf{Comentário :} Funções $X^{\mathbb{N}}$ são chamados de sequências em $X$, assim uma contagem é uma sequência infinita e bijetiva.\\

\textbf{Classe [ Enumerável ]} $ \Phi(A \in \mathcal{C}^*) : A_{enum} \iff \exists f \in \texttt{Cont}(\mathbb{N}^A) $

%%%%%%%%%%%%%%%%%%%%%%%%%%%%%%%%%%%%%%%%%%%%%%%%%%%%%%%%%%%%%%%%%%%%%%%%%%%%%%%%%%%%%%%%%%%%%%%%%%%%%%%%%%%%
% Proposição [ União Contável de Contáveis ]
\hspace{\baselineskip}

\textbf{Proposição [ União Contável de Contáveis ]}

$$ \Phi^*( (\xi_n)_{enum} ) : (\bigcup\limits_{n \in \mathbb{N}} \xi_n )_{enum} $$

\textbf{Dialética :} $\rightarrow$ !` Construção de Contagem !

\begin{itemize}
	\item[ (1) ] !` Organize mentalmente cada classe $\xi_n$ como colunas, os elementos dispostos em cada linha e em cada coluna temos um $\xi_n$ ordenadamente !
	\item[ (2) ] !` Comece a contagem no primeiro elemento de $\xi_0$ !
	\item[ (3) ] !` Se formar um quadrado avance a contagem na coluna seguinte !
	\item[ (4) ] !` Faça a contagem completando um quadrado !
	\item[ (5) ] !` Sempre que formar um quadrado de elementos nomeados vá para o item (3) !
\end{itemize}

\textbf{Comentário :} Essa dialética somente faz sentido se assumirmos que todos os elementos da união são contados em algum momento e que possamos reconhecer a contagem como uma bijeção. Como isso é um tópico de análise ou de teoria de conjuntos vamos assumir que isso seja suficiente para demonstrar a proposição, entretanto o leitor deve convir que o argumento de contagem é um argumento informal (requer que o leitor tenha habilidade de reconhecer como contagem).

%%%%%%%%%%%%%%%%%%%%%%%%%%%%%%%%%%%%%%%%%%%%%%%%%%%%%%%%%%%%%%%%%%%%%%%%%%%%%%%%%%%%%%%%%%%%%%%%%%%%%%%%%%%%
% Algoritmo
\hspace{\baselineskip}

\textbf{[ Algoritmo ]} Algoritmo é uma fórmula que representa a construção de maneira contável de uma sequência. Essa fórmula necessariamente precisa de instruções finitas.\\

\textbf{Comentário :} Um algoritmo é uma maneira de expressar uma sequência contável de ações, assim como na lógica proposicional e na definição de classes e conjuntos construímos uma notação para representar estas entidades, vamos construir uma notação apropriada para representar diferentes tipos de algoritmos espelhando os algoritmos da computação.\\

\textbf{Comentário :} Algoritmos são definidos matematicamente pela teoria de Máquina de Turing, entretanto não vamos explorar muito seu conteúdo por ser muito extenso para o objetivo deste texto.

%%%%%%%%%%%%%%%%%%%%%%%%%%%%%%%%%%%%%%%%%%%%%%%%%%%%%%%%%%%%%%%%%%%%%%%%%%%%%%%%%%%%%%%%%%%%%%%%%%%%%%%%%%%%
% Ação
\hspace{\baselineskip}

\textbf{[ Ação ]} Uma ação, assim como a informação em uma proposição, é o elemento mais simples de um algoritmo. Construímos algoritmos complicados por ações, uma ação isolada é um algoritmo assim como um algoritmo é uma ação.\\

\textbf{Comentário :} Na construção de fórmulas algorítmicas a ação equivale a informação na lógica proposicional. Todas os elementos nas fórmulas algorítmicas são ações.\\

\textbf{Notação :} $ \alpha \in \texttt{Action}^* $ é a notação para representar $\alpha$ como uma ação.\\

\textbf{Ação [ Sequência Simples ]} Uma sequência simples é uma sequência finita de ações, como exemplo temos $ \alpha_1 \alpha_2 \alpha_3$ com as ações individuais $\alpha_i \in \texttt{Action}^*$, neste exemplo as ações são executadas da esquerda para direita.\\

\textbf{Comentário :} Na programação os algoritmos são vistos através de outra sintaxe, é interessante que o leitor pesquise e compare as diferenças. Como o propósito do texto é criar uma teoria lógica vamos tentar aproximar a linguagem dos algoritmos com a linguagem utilizada na matemática.\\

\textbf{Ação [ Condicional ]} Uma ação condicional é uma ação que somente é feita se determinada condição for satisfeita, essa condição é uma informação a respeito dos valores de variáveis do sistema.\\

\textbf{Notação :} $ (\phi : \alpha ) $ denota a ação $\alpha$ somente é executada se $\phi$ for verdadeiro. $ (\phi : \alpha ; \beta) $ denota que a ação $\alpha$ é executada somente se $\phi$ for verdadeiro, caso $\phi$ seja falso então $\beta$ é executado. $ (\phi : \alpha ; \psi : \beta ) := ( \phi : \alpha ; (\psi : \beta) ) $ e para $ ( \phi_1 : \alpha_1 ; \phi_2 : \alpha_2 ; \phi_3 : \alpha_3 ) := ( \phi_1 : \alpha_1 ; ( \phi_2 : \alpha_2 ; \phi_3 : \alpha_3 ) ) $, com isso temos um padrão para definir listas condicionais maiores.\\

\textbf{Comentário :} Esta estrutura condicional é uma maneira de representar as estruturas if else else if da computação utilizando uma notação para representar algoritmos em uma linha.\\

\textbf{Ação [ Retorno ]} Uma ação de retorno faz com que seja repetido uma sequência de ações, seja $\alpha_1 \alpha_2 \alpha_3$ uma sequência de ações e $(\phi \rightarrow \alpha_1)$ então se temos uma ação da forma $ \alpha_1 \alpha_2 \alpha_3 (\phi \rightarrow \alpha_1) $, se $\phi$ for verdadeiro a sequência de ações será como se tivesse reiniciado em $\alpha_1$.\\

\textbf{Comentário :} A ação de retorno equivale ao $\texttt{go to}$ presentes em linguagens de programação mais antigas como fortran e C. Através dessa ação é possível construir o que são chamados de loops ou laços, que são repetições de sequência de ações.

%%%%%%%%%%%%%%%%%%%%%%%%%%%%%%%%%%%%%%%%%%%%%%%%%%%%%%%%%%%%%%%%%%%%%%%%%%%%%%%%%%%%%%%%%%%%%%%%%%%%%%%%%%%%
% Repetição Infinita
\hspace{\baselineskip}

\textbf{Operador [ Repetição Infinita ]} $ \Phi^* : R^{ \otimes} \alpha := ( \alpha ( 1 \rightarrow \alpha) ) $\\

\textbf{Comentário :} Essa fórmula representa a repetição da ação $\alpha$ de maneira infinita.\\

\textbf{Notação :} $ R^{\otimes \otimes} R^{\otimes} \alpha := ( R^{\otimes} ( R^{\otimes} \alpha ) ) $ representa um loop dentro de loop.\\

\textbf{Comentário :} A notação de operadores de repetição segue a lógica de estruturas e subestruturas, por isso temos uma notação para denotar a hierarquia de estruturas, $R^{\otimes \otimes}$ tem maior hierarquia que $ R^{\otimes} $, temos em $R^{\otimes} \alpha$ que $\alpha$ está dentro da estrutura do operador $R^{\otimes}$ e podemos definir estruturas dentro de estruturas.\\

\textbf{Comentário :} Quando queremos realizar uma ação antes das estruturas devemos colocar o operador de hierarquia ou usar parênteses por exemplo $ \beta^{\otimes} R^{\otimes} (\alpha)  = \beta ( R^{\otimes} \alpha )  $

%%%%%%%%%%%%%%%%%%%%%%%%%%%%%%%%%%%%%%%%%%%%%%%%%%%%%%%%%%%%%%%%%%%%%%%%%%%%%%%%%%%%%%%%%%%%%%%%%%%%%%%%%%%%
% Operador Until
\hspace{\baselineskip}

\textbf{Operador [ Until ]} $ \Phi^* : \mathcal{U}^{ \phi \otimes } \alpha := ( \alpha ( \neg \phi \rightarrow \alpha ) ) $\\

\textbf{Comentário :} O operador until executa $\alpha$ pelo menos uma vez, caso $\phi$ seja verdadeiro então ele sai do laço. Esse operador não é um operador muito comum em programação, em geral se utiliza o laço de repetição que será definido a seguir.

\hspace{\baselineskip}

\textbf{Operador [ While ]} $ \Phi^* :  \mathcal{W}^{\phi \otimes} \alpha := ( ( \phi : \alpha (1 \rightarrow \texttt{INI}) ; (1 \rightarrow \texttt{OUT} ) ) ) $\\

\textbf{Comentário :} $\texttt{INI}$ representa o início da sequência dentro da estrutura do loop while e $\texttt{OUT}$ significa a saída da estrutura.\\

\textbf{Comentário :} Com estes operadores temos quase tudo pronto para representar algoritmos, assumimos que possamos definir funções utilizando tais fórmulas da mesma maneira que definimos funções formais para informações ou para outras entidades vistas anteriormente. Para demonstrar mais formalmente a existência de contagens de conjuntos é efetivo mostrar a construção da sequência utilizando fórmulas algorítmicas.

%%%%%%%%%%%%%%%%%%%%%%%%%%%%%%%%%%%%%%%%%%%%%%%%%%%%%%%%%%%%%%%%%%%%%%%%%%%%%%%%%%%%%%%%%%%%%%%%%%%%%%%%%%%%
% Proposição Contagem das Partes de Naturais
\hspace{\baselineskip}

\textbf{Proposição [ Contagem das Partes de Naturais ]} $$ \{0,1\}^{\mathbb{N}} \overset{card}{>} \mathbb{N} $$

\textbf{Dialética :} $\rightarrow$ !` Absurdo !

\begin{itemize}
	\item[ ] ?` $  \exists f \in \texttt{Cout}(  \mathbb{N} \rightarrow \{0,1\}^{\mathbb{N}} ) \implies 0 $ ? $\rightarrow$ !` Argumento Diagonal de Cantor !
	\begin{itemize}
		\item[ ] !` Faça a sequência $F_n \in \{0,1\}^{\mathbb{N}}$ representar cada elemento !
		\item[ ] $ G:= \{ G(m) = \neg F_m(m) \} $
		\item[ ] !` Argumente que $G$ não é contado de $ \{0,1\}^{ \mathbb{N} } $ !
	\end{itemize}
	\item[ ] ?` $ \exists f_{\rightarrow} : \mathbb{N} \rightarrow \{0,1\}^{\mathbb{N}} $ ? $\rightarrow$ $\texttt{ALFA}$
\end{itemize}

\textbf{Comentário :} Argumento Diagonal de Cantor é uma técnica de demonstrar que um conjunto tem cardinalidade superior através de uma ideia bem intuitiva. Disponha cada $F_n$ como linhas de uma matriz, construa $G$ pegando a diagonal e alterando seus valores, $G$ é um elemento válido de $\{0,1\}^{\mathbb{N}}$ mas ele não é representado em nenhuma forma de contagem, portanto não é possível fazer uma contagem bijetiva de $ \{0,1\}^{ \mathbb{N} } $.\\

\textbf{Comentário :} Existe uma bijeção entre $\{0,1\}^{ \mathbb{N} }$ e os subconjuntos de $\mathbb{N}$, basta observar que cada elemento de $ \{0,1\}^{\mathbb{N}} $ é como se fosse uma indicadora de um subconjunto de $\mathbb{N}$. Isso permite concluir que o conjunto dos conjuntos de $\mathbb{N}$ tem cardinalidade superior a $\mathbb{N}$. Chamamos um conjunto infinito com cardinalidade superior aos naturais como incontável. \\

\textbf{Comentário :} Um conjunto contável não pode ser gerado por um algoritmo, pois isso requer que algoritmo seja executado infinitamente, entretanto existe um algoritmo em que todo elemento de um conjunto contável pode ser alcançado em um número finito de etapas. Um conjunto incontável é um conjunto em que essa propriedade algorítmica não é valida, existe elementos desse conjunto que não podem ser alcançados com um número finito de etapas para um algoritmo qualquer. Isso faz com que do ponto de vista de algoritmos geradores, temos pelo menos três ordens de efetividade (neste caso efetivo significa que o algoritmo consegue alcançar seu objetivo em um número finito de etapas), um conjunto finito tem o próprio conjunto como efetivo, um conjunto infinito contável tem como efetivo um elemento qualquer do conjunto, enquanto um incontável apenas uma parcela enumerável de elementos são efetivos computacionalmente.\\

\textbf{Comentário :} É fácil concluir que a cardinalidade de $ \{ 0,1,2,3,4,5,6,7,8,9 \}^{ \mathbb{N} } $ é pelo menos a cardinalidade $ \{0,1\}^{ \mathbb{N} } $, ou senão maior, se considerarmos os números reais como uma sequencia de algarismos mais um natural para indicar a posição do decimal podemos facilmente ver que os números reais são não enumeráveis. Deixarei os detalhes mais interessantes para teoria de conjuntos e análise ficando estes conjuntos como exemplos de infinitos.

%%%%%%%%%%%%%%%%%%%%%%%%%%%%%%%%%%%%%%%%%%%%%%%%%%%%%%%%%%%%%%%%%%%%%%%%%%%%%%%%%%%%%%%%%%%%%%%%%%%%%%%%%%%%
% Conjunto das Partes
\hspace{\baselineskip}

\textbf{[ Classe das Partes ]} $ \Phi( A \in \mathcal{C}^* ) : \mathcal{P}(A) := \bigcup\limits_{ B \subseteq A } \{B\} $\\

\textbf{Notação :} $2^A := \mathcal{P}(A)$

%%%%%%%%%%%%%%%%%%%%%%%%%%%%%%%%%%%%%%%%%%%%%%%%%%%%%%%%%%%%%%%%%%%%%%%%%%%%%%%%%%%%%%%%%%%%%%%%%%%%%%%%%%%%
% Ordens de Infinito
\hspace{\baselineskip}

\textbf{Proposição [ Ordens Arbitrariamente Maiores de Infinito ]} $$ \Phi( A_{\infty} \in \mathcal{C}^* ) : 2^A \overset{card}{>} A $$

\textbf{Dialética :} $\rightarrow$ !` Argumento Diagonal de Cantor !

\begin{itemize}
	\item[ ] ?` $ \{0,1\}^A \overset{card}{=} A \implies 0 $ ? $\rightarrow$
	\begin{itemize}
		\item[ ] ?` $ \exists F_{\bullet} \in ( \{0,1\}^A )^A : \{0,1\}^A \overset{class}{=} \bigcup\limits_{\alpha \in A} \{ F_{\alpha} \}  $ ?
		\item[ ] $ G := \{ \forall x \in A : G(x) = \neg F_x(x) \} $
		\item[ ] !` Argumento Diagonal de Cantor !
	\end{itemize}
\end{itemize}

\textbf{Comentário :} Essa proposição mostra que não somente existe infinitos maiores que o natural, existem infinitos arbitrariamente maiores, bastando aceita a existência da classe das partes.\\

\textbf{Classe [ Produtório Cartesiano ]} $$ \prod\limits_{\xi \in A} X_{\xi} := ( (x_{\alpha} )_{ \alpha \in A } \in \prod\limits_{\xi \in A} X_{\xi} \iff ( \forall \alpha \in A : x_{\alpha} \in X_{\alpha}   ) )  $$\\

\textbf{Comentário :} $ (x_{\alpha})_{\alpha \in A} $ é uma generalização do conceito de listas, entretanto não é possível dizer que são ordenados, podemos perceber que essa entidade é como uma função de $ A $ para a união dos $X_{\alpha}$'s.\\

\textbf{Comentário :} Na teoria de conjuntos essa construção é responsável por definir um importante axioma, o axioma de escolha, essencialmente ele diz que o produtório deste tipo para $A$ diferente de vazio e cada $X_{\xi}$ diferente de vazio tem pelo menos um elemento, este chamado de função escolha, porque escolhe um elemento de cada $X_{\xi}$.\\

\textbf{Comentário :} Podemos perceber uma conexão entre a notação de conjuntos e a álgebra de números observando que $ \prod\limits_{\alpha \in A} X \overset{class}{=} X^A  $.

%%%%%%%%%%%%%%%%%%%%%%%%%%%%%%%%%%%%%%%%%%%%%%%%%%%%%%%%%%%%%%%%%%%%%%%%%%%%%%%%%%%%%%%%%%%%%%%%%%%%%%%%%%%%
%%%%%%%%%%%%%%%%%%%%%%%%%%%%%%%%%%%%%%%%%%%%%%%%%%%%%%%%%%%%%%%%%%%%%%%%%%%%%%%%%%%%%%%%%%%%%%%%%%%%%%%%%%%%
%%%%%%%%%%%%%%%%%%%%%%%%%%%%%%%%%%%%%%%%%%%%%%%%%%%%%%%%%%%%%%%%%%%%%%%%%%%%%%%%%%%%%%%%%%%%%%%%%%%%%%%%%%%%

\chapter{Considerações Finais}

\hspace{\baselineskip}

A meta-teoria principal tem como propósito ser uma ferramenta pessoal de estudo e teorização (e como objetivo mais nobre uma ferramenta que pode ser utilizada por outros seres conscientes), sendo abrangente pelo menos do ponto de vista teórico. É interessante enumerar algumas boas propriedades que eu acho que uma teoria deve ter antes de começar a se aventurar na teorização, estas considerações são sugestões e não devem ser seguidas a risca, apenas recomendadas.

\begin{enumerate}
	\item O estudo deve seguir o caminho $(\alpha, \beta, \gamma)$ da trialética, isso significa que primeiro devemos memorizar conteúdos já estabelecidos pela humanidade, depois devemos resolver problemas com tais conteúdos, somente depois pode ser interessante entender o processo da descoberta.
	\item Uma teoria deve considerar os três aspectos da consciência informacional, ela deve ser do ponto de vista alfa, simples o suficiente para ser eficientemente memorizável, deve sob o ponto de vista beta ter aplicações e exercitar para resolução de problemas e deve sob o ponto de vista gama ser esclarecido e justificado.
	\item Sob o ponto de vista gama devemos partir do natural para o transcendental, devemos começar o aceitável para justificar o misterioso. Sob o ponto de vista beta devemos dar ênfase ao importante e não demandar tempo nos conteúdos auxiliares. Sob o ponto de vista alfa devemos partir do simples e com ele construir estruturas mais complexas.
	\item Devemos considerar que o ser humano erra, e portanto sempre que possível as afirmações devem ser verificadas. Por conta dessa natureza a arrogância não fica muito distante da ignorância.
	\item Devemos considerar o heurístico, o empírico e o lógico como recursos complementares da ciência (no sentido abrangente) humana, como uma decomposição trialética da ciência.
	\item Nós seres humanos temos intuições e muitas vezes consideramos certeza utilizando a intuição que varia conforme a experiência intelectual de cada ser humano, devemos substituir o óbvio pela conotação da consciência alfa, algo somente é evidente se o ouvinte tiver a mesma consciência alfa de quem faz a proposição.
	\item Sob o ponto de vista gama, teorias axiomáticas utilizando fórmulas (essencialmente lógica) são as teorias mais satisfatórias em justificar suas proposições, sob o ponto de vista beta, teorias homeostáticas são mais práticas em se teorizar e em possibilitar a criação de estratégias, sob o ponto de vista alfa, a dialética permite que o leitor possa atingir um nível de consciência maior sobre um assunto. Esses três pontos de vistas são três direções distintas de evolução de uma teoria, mas tais direções não são disjuntas entre si, eles são complementares como toda base trialética.
	\item Se queremos demonstrar a verdade utilizamos a filosofia gama, se queremos resolver problemas utilizamos a filosofia beta, se queremos aumentar o nível de consciência utilizamos a filosofia alfa. Essas são três escolas filosóficas complementares entre si, é sugestivo dizer que sejam seus três anjos da guarda da ciência, seus guias para o caminho do conhecimento.
	\item Uma teoria que investe muito em uma escola filosófica da trialética pode estar deficiente na qualidade das outras. Um texto analítico que busca investigar minuciosamente a verdade pode ser difícil de aplicar e difícil de compreender, um texto empírico e prático pode ser misterioso e difícil de compreender, um texto filosófico pode carecer de justificações aceitáveis e pode fugir da realidade.
	\item O conhecimento é um investimento que ultrapassa gerações, está além da vida individual e deve ser visto desta forma. Aquele que se apega a sua vida própria e desconsidera a vida dos outros dificilmente irá investir em criar teorias.
	\item Nossa vida é um caminho para conhecermos nosso próprio espírito, não conseguimos ver nosso espírito, para confirmar a qualidade de nosso espírito é interessante projeta-lo externalizando na trialética (disciplina, sinceridade, bondade) $\sim$ ($\beta$, $\gamma$, $\alpha$). Disciplina e Sinceridade associados ao indivíduo e bondade para existir. Disciplina para construir, Sinceridade para entender e Bondade para nos tornarmos humanos.
	\item Uma pessoa que arroga para si o direito de liderar os outros deve ser uma pessoa que se comunica com as três escolas filosóficas da trialética, liderar em um sentido abrangente poderia ser equivalente a orientar e influenciar a vida de outras pessoas. Uma pessoa que apenas valoriza um ou dois dos aspectos tem o dever se colocar em uma posição que não prejudique as outras pessoas.
\end{enumerate}

Agradeço a quem chegou até aqui, espero que de alguma forma tenha contribuído para seu intelecto e sua humanidade. Este não é o fim, mas o começo de uma jornada intelectual.

%%%%%%%%%%%%%%%%%%%%%%%%%%%%%%%%%%%%%%%%%%%%%%%%%%%%%%%%%%%%%%%%%%%%%%%%%%%%%%%%%%%%%%%%%%%%%%%%%%%%%%%%%%%%
%%%%%%%%%%%%%%%%%%%%%%%%%%%%%%%%%%%%%%%%%%%%%%%%%%%%%%%%%%%%%%%%%%%%%%%%%%%%%%%%%%%%%%%%%%%%%%%%%%%%%%%%%%%%
%%%%%%%%%%%%%%%%%%%%%%%%%%%%%%%%%%%%%%%%%%%%%%%%%%%%%%%%%%%%%%%%%%%%%%%%%%%%%%%%%%%%%%%%%%%%%%%%%%%%%%%%%%%%






\end{document}